% !TeX program = pdflatex
% Tom III: Analiza stochastyczna; rozdzialy 1--16, audyt rozdzialow 13--14.
% Przygotowano jako samodzielny dokument; zasady i symbole jak w tomie II.
\documentclass[11pt,a4paper,oneside]{report}
\usepackage[T1]{fontenc}
\usepackage[utf8]{inputenc}
\usepackage[polish]{babel}
\usepackage{lmodern}
\usepackage{amsmath,amssymb,amsthm,mathtools}
\usepackage{microtype}
\usepackage{geometry}
\usepackage{enumitem}
\usepackage{xcolor}
\usepackage{tikz}
\usetikzlibrary{arrows.meta,positioning,calc,shapes.geometric,decorations.pathmorphing}
\usepackage[unicode]{hyperref}
% Odsylacze miedzytomowe; zbudowac wszystkie .aux i przechowywac PDF-y obok siebie.
\usepackage{xr-hyper}
\externaldocument[I-]{rachunek_prawdopodobienstwa_tom_I}[rachunek_prawdopodobienstwa_tom_I.pdf]
\externaldocument[II-]{rachunek_prawdopodobienstwa_tom_II}[rachunek_prawdopodobienstwa_tom_II.pdf]
\externaldocument[IV-]{rachunek_prawdopodobienstwa_tom_IV}[rachunek_prawdopodobienstwa_tom_IV.pdf]

\geometry{margin=2.65cm}
\hypersetup{colorlinks=true,linkcolor=blue!45!black,urlcolor=blue!45!black}
\setlength{\parindent}{0pt}
\setlength{\parskip}{4pt}
\setlength{\emergencystretch}{2em}
\setlist[enumerate]{topsep=3pt,itemsep=2pt}
\setlist[itemize]{topsep=3pt,itemsep=2pt}
\numberwithin{equation}{chapter}
\theoremstyle{plain}
\newtheorem{twierdzenie}{Twierdzenie}[chapter]
\newtheorem{lemat}[twierdzenie]{Lemat}
\newtheorem{wniosek}[twierdzenie]{Wniosek}
\newtheorem{stwierdzenie}[twierdzenie]{Stwierdzenie}
\theoremstyle{definition}
\newtheorem{definicja}{Definicja}[chapter]
\newtheorem{przyklad}{Przykład}[chapter]
\newtheorem{zadanie}{Zadanie rozwiązane}[chapter]
\theoremstyle{remark}
\newtheorem{uwaga}{Uwaga}[chapter]
\renewcommand{\proofname}{Dowód}
\newcommand{\R}{\mathbb{R}}
\newcommand{\N}{\mathbb{N}}
\newcommand{\Z}{\mathbb{Z}}
\newcommand{\Q}{\mathbb{Q}}
\newcommand{\Prob}{\mathbb{P}}
\newcommand{\E}{\mathbb{E}}
\newcommand{\Var}{\operatorname{Var}}
\newcommand{\Cov}{\operatorname{Cov}}
\newcommand{\Law}{\operatorname{Law}}
\newcommand{\calF}{\mathcal{F}}
\newcommand{\calE}{\mathcal{E}}
\newcommand{\calH}{\mathcal{H}}
\newcommand{\calN}{\mathcal{N}}
\newcommand{\calB}{\mathcal{B}}
\newcommand{\calP}{\mathcal{P}}
\newcommand{\calG}{\mathcal{G}}
\newcommand{\calD}{\mathcal{D}}
\newcommand{\calS}{\mathcal{S}}
\newcommand{\calC}{\mathcal{C}}
\newcommand{\ind}[1]{\mathbf{1}_{#1}}
\newcommand{\dd}{\mathrm{d}}
\newcommand{\norm}[1]{\lVert #1\rVert}
\newcommand{\abs}[1]{\lvert #1\rvert}
\newcommand{\iid}{\mathrel{\overset{\mathrm{iid}}{\sim}}}
\newcommand{\Normal}{\mathcal{N}}
\DeclareMathOperator{\dist}{dist}
\DeclareMathOperator{\supp}{supp}
\DeclareMathOperator{\spanop}{span}
\begin{document}
\hypersetup{pageanchor=false}
\begin{titlepage}
\centering\vspace*{3.4cm}
{\Large Podstawy teorii prawdopodobieństwa i statystyki\par}
\vspace{1.2cm}
{\Huge\bfseries Tom III\\[0.25cm] Analiza stochastyczna\par}
\vspace{0.9cm}
{\LARGE Całki Itô, procesy Lévy'ego i całki po trajektoriach\par}
\vspace{0.7cm}
{\Large Rozdziały 1--16\par}
\vfill
{\large Konstrukcje, dowody, intuicje i zadania rozwiązane\par}
\vspace{1.1cm}
\end{titlepage}
\hypersetup{pageanchor=true}
\tableofcontents
\chapter*{Wprowadzenie do tomu III}
\addcontentsline{toc}{chapter}{Wprowadzenie do tomu III}
W tomie II poznaliśmy proces Wienera, martyngały w czasie ciągłym,
czasy zatrzymania oraz wariację kwadratową. Teraz zmieniamy punkt
widzenia: zamiast wyłącznie badać procesy, chcemy względem nich
\emph{całkować}. Nasuwa się wzór $\int_0^T H_t\,\dd W_t$, ale
sam zapis nie określa jeszcze znaczenia całki.

Przeszkoda jest fundamentalna. Prawie każda trajektoria procesu Wienera
jest ciągła, ale ma na każdym nietrywialnym przedziale nieskończoną
wariację całkowitą. Klasyczna teoria Riemanna--Stieltjesa nie daje więc
od razu całki dla wszystkich interesujących nas funkcji podcałkowych.
Rozwiązaniem jest wykorzystanie informacji probabilistycznej:
przyrosty procesu Wienera są niezależne, mają średnią zero i znaną
wariancję. Dzięki temu sumy całkowe możemy kontrolować w przestrzeni
$L^2$, nawet gdy nie mają sensownej interpretacji w klasycznym
rachunku całkowym wzdłuż każdej pojedynczej trajektorii.

W rozdziale pierwszym budujemy \emph{całkę Wienera} dla
\emph{deterministycznych} funkcji z $L^2([0,T])$. Jest ona szczególnym
przypadkiem całki Itô, lecz jej prostsza konstrukcja pozwala zobaczyć
istotę izometrii i związku z procesami gaussowskimi. W rozdziale
drugim dopuszczamy losową funkcję podcałkową zależną od przeszłości
procesu. Konstrukcję przeprowadzimy najpierw dla funkcji prostych
przewidywalnych, a następnie przez domknięcie w $L^2$ i lokalizację.

Od początku odróżniamy trzy obiekty: zwykłą całkę względem czasu
$\int_0^T H_t\,\dd t$, całkę stochastyczną
$\int_0^T H_t\,\dd W_t$ oraz całkę \emph{po przestrzeni trajektorii}
$\int F(w)\,\dd\mu_W(w)$. Ostatni zapis jest całką względem miary
Wienera, nie zaś synonimem całki Itô. Związki z całką Feynmana,
wzorem Feynmana--Kaca i obrotem Wicka omówimy w dalszej części tomu.

\paragraph{Założenia stałe.}
Niech $T>0$, $(\Omega,\calF,\Prob)$ będzie przestrzenią probabilistyczną,
$(\calF_t)_{0\le t\le T}$ filtracją spełniającą warunki zwykłe,
a $W=(W_t)_{0\le t\le T}$ rzeczywistym standardowym procesem Wienera
\emph{względem tej filtracji}: $W_0=0$, ma ciągłe trajektorie,
jest adaptowany oraz $W_t-W_s\sim\Normal(0,t-s)$ jest niezależny
od $\calF_s$ dla $0\le s<t\le T$. Ostatni warunek jest mocniejszy
niż samo stwierdzenie, że $W$ jest procesem Wienera jako rodzina
zmiennych losowych; będziemy z niego korzystać przy całkowaniu
losowych funkcji podcałkowych.

Przez $L^p(\Omega)$ rozumiemy $L^p(\Omega,\calF,\Prob)$,
a przez $L^p([0,T])$ przestrzeń względem miary Lebesgue'a.
Wartość całki stochastycznej w danej chwili jest określona z dokładnością
do równości prawie na pewno, a cały proces całkowy będziemy wybierać
w wersji o ciągłych trajektoriach.

\chapter{Całka Wienera dla deterministycznych funkcji}
\label{ch:wiener-integral}
\paragraph{Odsyłacze międzytomowe.} Własności procesu Wienera przypomina tom II, rozdział~\ref{II-ch:wiener}, a całkowanie Monte Carlo i redukcję wariancji omawia tom IV, rozdział~\ref{IV-ch:monte-carlo}.

\section{Dlaczego nie zaczynamy od Riemanna--Stieltjesa?}
Dla gładkiej funkcji $g$ klasyczne wyrażenie
$\int_0^T f(t)\,\dd g(t)$ można rozumieć jako
$\int_0^T f(t)g'(t)\,\dd t$. Jeżeli $g$ ma jedynie
skończoną wariację, całka Riemanna--Stieltjesa nadal istnieje
dla każdej ciągłej funkcji $f$. Tymczasem dla prawie każdej
trajektorii $w(t)=W_t(\omega)$ żadne z tych założeń nie zachodzi.

Nie oznacza to, że nie da się obliczyć żadnej całki względem $w$:
na przykład przy wystarczającej regularności funkcji $f$ można
zastosować całkowanie przez części. Problem polega na tym,
że nie mamy ogólnej, ciągłej konstrukcji wzdłuż trajektorii,
która działałaby dla dowolnego $f\in L^2([0,T])$.
Co ciekawe, \emph{losowość} pozwoli ten brak regularności pokonać.

\begin{przyklad}[Dlaczego punkt próbkowania ma znaczenie]
Weźmy podział $\pi=\{0=t_0<\cdots<t_n=T\}$ i oznaczmy
$\Delta_kW=W_{t_{k+1}}-W_{t_k}$. Z tożsamości
$W_{t_{k+1}}^2-W_{t_k}^2=2W_{t_k}\Delta_kW+(\Delta_kW)^2$
wynika dokładnie
\[
\sum_{k=0}^{n-1}W_{t_k}\Delta_kW
=\frac12W_T^2-\frac12\sum_{k=0}^{n-1}(\Delta_kW)^2.
\]
W tomie II pokazaliśmy, że ostatnia suma zbiega do $T$ w $L^2$
dla podziałów deterministycznych o średnicy dążącej do zera.
Dlatego sumy \emph{lewostronne} mają granicę
$\frac12(W_T^2-T)$ w $L^2$. Gdy w każdej sumie użyjemy
punktu środkowego w \emph{wartościach procesu}, czyli
$\frac12(W_{t_k}+W_{t_{k+1}})$, otrzymamy zamiast tego
$\frac12W_T^2$ dla każdego podziału: odpowiednia suma
jest teleskopowa.

Dwie pozornie podobne reguły próbkowania dają różne odpowiedzi.
Nie jest to usterka rachunku, ale sygnał, że przy całkowaniu względem
procesu Wienera należy \emph{ustalić konwencję}. W całce Itô
stosujemy wartości funkcji podcałkowej dostępne \emph{przed} przyrostem.
\end{przyklad}

\begin{center}
\begin{tikzpicture}[>=Latex,x=1.05cm,y=0.9cm]
\draw[->] (0,0)--(8.3,0) node[right] {$t$};
\draw[->] (0,-1.45)--(0,2.35) node[above] {$w(t)$};
\draw[thick,blue!75!black] (0,0)--(.5,.3)--(1,.15)--(1.5,1.1)--(2,.5)--(2.5,.7)--(3,-.55)--(3.5,-.15)--(4.1,-.8)--(4.6,-.15)--(5.2,.75)--(5.7,.4)--(6.2,1.65)--(6.8,1.1)--(7.4,1.75);
\foreach \x in {0,1.5,3,4.6,6.2,7.4}{\draw[dashed,gray!65] (\x,-1.15)--(\x,2);}
\node[below,scale=.78] at (0,0) {$t_0$};
\node[below,scale=.78] at (1.5,0) {$t_1$};
\node[below,scale=.78] at (3,0) {$t_2$};
\node[below,scale=.78] at (4.6,0) {$t_3$};
\node[below,scale=.78] at (6.2,0) {$t_4$};
\node[below,scale=.78] at (7.4,0) {$t_5$};
\node[align=left,scale=.85] at (4.1,2.05) {Przyrost jest znany dopiero\\po osiągnięciu prawego końca};
\end{tikzpicture}
\par\smallskip\small Schemat przyrostów: reguła lewego końca nie korzysta z przyszłych wartości procesu.
\end{center}

\section{Funkcje schodkowe i pierwszy wzór całkowy}
\begin{definicja}[Deterministyczna funkcja schodkowa]
Funkcją schodkową na $[0,T]$ nazywamy funkcję postaci
$f=\sum_{k=0}^{n-1}a_k\ind{(t_k,t_{k+1}]}$, gdzie
$0=t_0<t_1<\cdots<t_n=T$ i $a_k\in\R$.
Jej wartość w punkcie $0$ nie ma znaczenia dla całek względem
miary Lebesgue'a ani dla konstruowanej całki Wienera.
Przestrzeń takich funkcji oznaczamy przez $\calS_T$.
\end{definicja}

\begin{definicja}[Całka Wienera dla funkcji schodkowej]
Dla $f=\sum_{k=0}^{n-1}a_k\ind{(t_k,t_{k+1}]}\in\calS_T$
definiujemy
\[
I_W(f):=\sum_{k=0}^{n-1}a_k(W_{t_{k+1}}-W_{t_k}).
\tag{1.1}\label{eq:wiener-simple}
\]
Wynikiem jest zmienna losowa z $L^2(\Omega)$.
\end{definicja}

\begin{lemat}[Niezależność reprezentacji]
Dla każdej funkcji $f\in\calS_T$ wartość $I_W(f)$
nie zależy od wyboru podziału ani od zapisu $f$ jako sumy
funkcji wskaźnikowych przedziałów. Odwzorowanie
$I_W:\calS_T\to L^2(\Omega)$ jest liniowe.
\end{lemat}
\begin{proof}
Jeżeli podzielimy $(t_k,t_{k+1}]$ w punkcie $u$, to
$W_{t_{k+1}}-W_{t_k}=(W_u-W_{t_k})+(W_{t_{k+1}}-W_u)$.
Wkład przedziału do sumy nie zmienia się. Dwa przedstawienia
funkcji schodkowej mają wspólne rozdrobnienie podziałów.
Po przejściu do niego współczynniki na każdym przedziale są
równe (poza punktami końcowymi), ponieważ oba przedstawienia
określają tę samą funkcję prawie wszędzie. Stąd równość sum.
Liniowość uzyskujemy przez sprowadzenie dwóch funkcji do
wspólnego podziału i dodanie współczynników.
\end{proof}

\begin{twierdzenie}[Izometria Wienera i prawo Gaussa]\label{th:wiener-isometry}
Dla dowolnych $f,g\in\calS_T$ zachodzą
\[
\E I_W(f)=0,\qquad
\E[I_W(f)I_W(g)]=\int_0^T f(t)g(t)\,\dd t.
\tag{1.2}\label{eq:wienscalar}
\]
Ponadto $I_W(f)$ ma rozkład normalny
$\Normal(0,\int_0^Tf(t)^2\,\dd t)$,
przy czym wariancja zero oznacza rozkład skupiony w zerze.
\end{twierdzenie}
\begin{proof}
Z poprzedniego lematu możemy przyjąć wspólny podział,
tak że $f=\sum a_k\ind{(t_k,t_{k+1}]}$ i
$g=\sum b_k\ind{(t_k,t_{k+1}]}$. Oznaczmy
$\Delta_kW=W_{t_{k+1}}-W_{t_k}$. Ponieważ przyrosty Wienera
mają średnią zero, $\E I_W(f)=0$. Dla $j\ne k$
niezależność i zerowa średnia przyrostów dają
$\E[\Delta_jW\Delta_kW]=0$; dla $j=k$ mamy
$\E(\Delta_kW)^2=t_{k+1}-t_k$. Zatem
\[
\E[I_W(f)I_W(g)]
=\sum_{j,k}a_jb_k\E[\Delta_jW\Delta_kW]
=\sum_k a_kb_k(t_{k+1}-t_k)
=\int_0^Tfg\,\dd t.
\]
Na koniec przyrosty $\Delta_kW$ są niezależnymi zmiennymi
normalnymi, a ich kombinacja liniowa jest normalna.
Obliczona wariancja określa dokładnie jej rozkład.
\end{proof}

\begin{uwaga}[Ograniczenie do funkcji deterministycznych]
W powyższym dowodzie wolno było wyciągnąć współczynniki
$a_k$ i $b_k$ przed wartość oczekiwaną, ponieważ były liczbami.
Dla losowych współczynników nie jest to dopuszczalne.
W rozdziale drugim użyjemy w zamian warunkowania względem
$\calF_{t_k}$ oraz niezależności przyszłego przyrostu od tej
$\sigma$-algebry. Właśnie to prowadzi do przewidywalności.
\end{uwaga}

\section{Domknięcie przestrzeni: definicja dla funkcji kwadratowalnych}
\begin{lemat}[Gęstość funkcji schodkowych]
Przestrzeń $\calS_T$ jest gęsta w $L^2([0,T])$.
\end{lemat}
\begin{proof}
Wskaźniki przedziałów półotwartych generują $\calB([0,T])$.
Ich skończone kombinacje liniowe są gęste w $L^2$:
po pierwsze, wskaźniki zbiorów borelowskich można
aproksymować w $L^2$ wskaźnikami skończonych sum przedziałów
(z regularności miary Lebesgue'a); po drugie,
każdą funkcję z $L^2$ przybliża się funkcjami prostymi,
najpierw przez obcięcie, a potem przez kwantyzację jej wartości.
Po wspólnym rozdrobnieniu przedziałów wszystkie aproksymacje
należą do $\calS_T$.
\end{proof}

\begin{twierdzenie}[Całka Wienera w $L^2$]\label{th:wiener-extension}
Istnieje dokładnie jedno liniowe i ciągłe odwzorowanie
$I_W:L^2([0,T])\to L^2(\Omega)$, które zgadza się z
\eqref{eq:wiener-simple} na $\calS_T$. Jest izometrią:
\[
\norm{I_W(f)}_{L^2(\Omega)}=\norm{f}_{L^2([0,T])}.
\tag{1.3}\label{eq:wiener-iso}
\]
W szczególności $I_W(f)$ jest granicą w $L^2(\Omega)$
zmiennych $I_W(f_n)$ dla dowolnego ciągu $f_n\in\calS_T$
zbiegającego do $f$ w $L^2([0,T])$. Definicja granicy
nie zależy od wyboru ciągu aproksymującego.
\end{twierdzenie}
\begin{proof}
Z izometrii funkcji schodkowych dostajemy dla $m,n$
$\norm{I_W(f_n)-I_W(f_m)}_2=\norm{f_n-f_m}_2$.
Skoro $f_n\to f$ w $L^2([0,T])$, ciąg $(f_n)$ jest Cauchy'ego,
a więc $(I_W(f_n))$ jest Cauchy'ego w zupełnej przestrzeni
$L^2(\Omega)$. Ma zatem granicę $Y$. Jeżeli $g_n\to f$
jest innym ciągiem schodkowym, to
$\norm{I_W(f_n)-I_W(g_n)}_2=\norm{f_n-g_n}_2\to0$,
więc oba ciągi mają tę samą granicę. Kładziemy $I_W(f)=Y$.
Liniowość wynika z liniowości dla funkcji schodkowych i
ciągłości dodawania w $L^2$. Przechodząc do granicy w izometrii
otrzymujemy \eqref{eq:wiener-iso}; dowodzi to także ciągłości.
Jeżeli inne ciągłe liniowe odwzorowanie zgadzałoby się na gęstej
podprzestrzeni $\calS_T$, to przez ciągłość zgadzałoby się
wszędzie.
\end{proof}

\begin{twierdzenie}[Gaussowskość i kowariancja całki]\label{th:wiener-gaussian}
Dla dowolnych $f_1,\ldots,f_m\in L^2([0,T])$ wektor
$(I_W(f_1),\ldots,I_W(f_m))$ jest centrowanym wektorem gaussowskim
z macierzą kowariancji
$\big(\int_0^T f_i(t)f_j(t)\,\dd t\big)_{i,j=1}^m$.
W szczególności
\[
\E e^{iuI_W(f)}=\exp\left(-\frac{u^2}{2}\norm{f}_2^2\right),
\quad u\in\R.
\tag{1.4}
\]
Jeżeli $\int_0^Tfg\,\dd t=0$, to $I_W(f)$ i $I_W(g)$
są niezależne.
\end{twierdzenie}
\begin{proof}
Dla każdej kombinacji liniowej $h=\sum_{j=1}^ma_jf_j$
zmienna $I_W(h)$ jest granicą w $L^2$, a więc według
rozkładu, zmiennych normalnych $I_W(h_n)$ o wariancjach
dążących do $\norm{h}_2^2$. Zbieżność funkcji
charakterystycznych daje $I_W(h)\sim\Normal(0,\norm{h}_2^2)$.
Każda kombinacja liniowa wektora całek jest więc normalna,
co stanowi definicję wspólnej gaussowskości. Wzór na kowariancję
wynika z polaryzacji izometrii:
$\E I_W(f)I_W(g)=\int_0^Tfg\,\dd t$.
Dla wektora gaussowskiego nieskorelowanie dwóch współrzędnych
implikuje niezależność, co daje ostatnie stwierdzenie.
\end{proof}

\begin{wniosek}[Całkowanie po rozłącznych przedziałach]
Jeżeli $f,g\in L^2([0,T])$ mają rozłączne nośniki z dokładnością
do zbioru miary Lebesgue'a zero, to $I_W(f)$ i $I_W(g)$
są niezależne. Dla $0\le a<b\le T$ mamy
$I_W(\ind{(a,b]})=W_b-W_a$.
\end{wniosek}

\begin{przyklad}[Od funkcji stałej do funkcji wykładniczej]
Dla $f\equiv1$ mamy $I_W(f)=W_T\sim\Normal(0,T)$.
Dla $f=\ind{(a,b]}$ otrzymujemy przyrost Wienera.
Dla $f(t)=e^{-\lambda t}$, $\lambda>0$, wynik jest
zmienną normalną o średniej zero i wariancji
\[
\int_0^Te^{-2\lambda t}\,\dd t
=\frac{1-e^{-2\lambda T}}{2\lambda}.
\]
Łatwo w ten sposób znaleźć rozkład całki mimo braku
klasycznego całkowania wzdłuż trajektorii.
\end{przyklad}

\begin{zadanie}[Rozkład całki i niezależność]
Obliczyć prawo losowej zmiennej
$X=\int_0^T(2t-T)\,\dd W_t$ oraz ustalić,
czy jest ona niezależna od $W_T$.

\emph{Rozwiązanie.}
Z twierdzenia \ref{th:wiener-gaussian} zmienna $X$ ma rozkład
normalny $\Normal(0,\sigma^2)$, gdzie
$\sigma^2=\int_0^T(2t-T)^2\,\dd t=T^3/3$.
Ponadto
$\Cov(X,W_T)=\int_0^T(2t-T)\cdot1\,\dd t=0$.
Para $(X,W_T)$ jest wspólnie gaussowska, zatem jej
współrzędne są niezależne. Samo zerowanie się kowariancji
bez wspólnej gaussowskości nie wystarczyłoby do tego wniosku.
\end{zadanie}


\begin{zadanie}[Osobliwość funkcji podcałkowej]
Ustalić, dla jakich $\alpha\in\R$ istnieje całka Wienera
$X_\alpha=\int_0^Tt^{-\alpha}\,\dd W_t$, a dla jakich
nie może powstać jako granica całek obciętych
$\int_\varepsilon^Tt^{-\alpha}\,\dd W_t$
w prawdopodobieństwie, gdy $\varepsilon\downarrow0$.

\emph{Rozwiązanie.} Funkcja $t^{-\alpha}$ należy do
$L^2([0,T])$ wtedy i tylko wtedy, gdy
$\int_0^Tt^{-2\alpha}\,\dd t<\infty$, czyli
$2\alpha<1$. Dla $\alpha<1/2$ całka istnieje,
a jej rozkład jest normalny ze średnią zero i wariancją
$T^{1-2\alpha}/(1-2\alpha)$.

Jeżeli $\alpha\ge1/2$, to dla każdego $\varepsilon>0$
zmienna $X_\varepsilon:=\int_\varepsilon^T t^{-\alpha}\,\dd W_t$
jest normalna z wariancją
$v_\varepsilon=\int_\varepsilon^Tt^{-2\alpha}\,\dd t$,
która dąży do nieskończoności. Ponieważ
$X_\varepsilon/\sqrt{v_\varepsilon}\sim\Normal(0,1)$,
dla każdego ustalonego $R>0$ zachodzi
$\Prob(|X_\varepsilon|\le R)
=\Prob(|Z|\le R/\sqrt{v_\varepsilon})\to0$.
Rodzina rozkładów $(X_\varepsilon)$ nie jest ciasna,
a więc nie może zbiegać w prawdopodobieństwie do
skończonej zmiennej losowej. W ten sposób warunek
$f\in L^2$ otrzymuje także konkretne znaczenie:
zapobiega nieograniczonemu wzrostowi wariancji.
\end{zadanie}

\begin{przyklad}[Ziarno wykładnicze i początek procesu Ornsteina--Uhlenbecka]
Dla $\lambda>0$ i $t\in[0,T]$ rozważmy
$X_t:=\int_0^te^{-\lambda(t-s)}\,\dd W_s$.
Dla każdej ustalonej chwili $t$ funkcja
$s\mapsto e^{-\lambda(t-s)}\ind{(0,t]}(s)$
jest deterministyczna i należy do $L^2([0,T])$.
Stąd $X_t$ jest centrowaną zmienną normalną, a
\[
\Var(X_t)=\int_0^te^{-2\lambda(t-s)}\,\dd s
=\frac{1-e^{-2\lambda t}}{2\lambda}.
\]
Dla $0\le s\le t$ izometria spolaryzowana daje
\[
\Cov(X_s,X_t)
=\int_0^s e^{-\lambda(s-u)}e^{-\lambda(t-u)}\,\dd u
=e^{-\lambda(t-s)}\frac{1-e^{-2\lambda s}}{2\lambda}.
\]
Otrzymaliśmy bez rozwiązywania równania stochastycznego
konkretny proces gaussowski z pamięcią wykładniczą.
W dalszych rozdziałach pokażemy, że jest to rozwiązanie
liniowego równania stochastycznego z dryfem
$-\lambda X_t$.
\end{przyklad}

\section{Proces całkowy dla deterministycznej funkcji}
\begin{definicja}[Proces całek Wienera]
Dla $f\in L^2([0,T])$ i $t\in[0,T]$ kładziemy
$M_t^f:=I_W(f\ind{(0,t]})$.
Zmienna $M_t^f$ jest dobrze określona w $L^2(\Omega)$.
\end{definicja}
\begin{twierdzenie}[Martyngał, kowariancja i przyrosty]
Proces $(M_t^f)_{0\le t\le T}$ jest martyngałem względem
$(\calF_t)$ i ma ciągłą modyfikację. Dla $0\le s\le t\le T$
zachodzą
\[
\E[M_t^f\mid\calF_s]=M_s^f,\qquad
\E[M_s^fM_t^f]=\int_0^sf(u)^2\,\dd u.
\tag{1.5}
\]
Przyrost $M_t^f-M_s^f$ jest niezależny od $\calF_s$
i ma rozkład $\Normal(0,\int_s^tf(u)^2\,\dd u)$.
\end{twierdzenie}
\begin{proof}
Najpierw załóżmy, że $f$ jest schodkowa. Przyrost
$M_t^f-M_s^f$ jest deterministyczną kombinacją przyrostów
procesu Wienera na $(s,t]$. Jest więc niezależny od $\calF_s$
i ma średnią zero. Zatem własność martyngałowa zachodzi.
Dla ogólnego $f$ wybierzmy $f_n\in\calS_T$ z
$\norm{f_n-f}_2\to0$. Z izometrii wynika
$M_t^{f_n}\to M_t^f$ w $L^2$ dla każdego $t$.
Warunkowa wartość oczekiwana jest kontrakcją w $L^1$,
więc przechodząc do granicy otrzymujemy
$\E[M_t^f\mid\calF_s]=M_s^f$. Jednocześnie rozkład
i niezależność przyrostu od $\calF_s$ wynikają z granicy
w $L^2$ oraz faktu, że zmienne niezależne od danej
$\sigma$-algebry zachowują tę własność przy zbieżności
według prawdopodobieństwa (wystarczy sprawdzić równość
oczekiwań $\E[g(Y_n)\ind A]=\E g(Y_n)\Prob(A)$ dla
$A\in\calF_s$, a następnie przejść do granicy dla
granicznie ograniczonych ciągłych $g$). Wzór na kowariancję
wynika z \eqref{eq:wienscalar} po domknięciu.

Ciągłość trajektorii wymaga osobnego argumentu. Dla
schodkowej $f_n$ proces $M^{f_n}$ jest ciągły, a
nierówność maksymalna Dooba i izometria dają
\[
\E\sup_{0\le t\le T}|M_t^{f_n}-M_t^{f_m}|^2
\le4\norm{f_n-f_m}_2^2.
\]
Możemy więc wybrać podciąg, dla którego
$\sum_j\big(\E\sup_t|M_t^{f_{n_{j+1}}}-M_t^{f_{n_j}}|^2\big)^{1/2}<\infty$.
Nierówność Cauchy'ego--Schwarza i twierdzenie o monotonicznej
zbieżności pokazują, że suma supremów różnic jest skończona
prawie na pewno. Wobec tego podciąg zbiega jednostajnie
prawie na pewno do procesu o ciągłych trajektoriach.
Dla każdego ustalonego $t$ granica zgadza się z $M_t^f$
w $L^2$, a więc jest jego modyfikacją.
\end{proof}

\begin{uwaga}[Zmiana czasu]
Funkcja $A_f(t):=\int_0^tf(u)^2\,\dd u$ jest ciągła
i niemalejąca. Rozkłady skończeniewymiarowe procesu $M^f$
są takie same jak procesu $B_{A_f(t)}$, gdzie $B$
jest standardowym ruchem Browna na odpowiednio rozszerzonej
przestrzeni czasu. Jest to równość \emph{praw procesów},
a nie stwierdzenie $M_t^f=B_{A_f(t)}$ dla ustalonego
wcześniej, niezależnego ruchu Browna.
\end{uwaga}

\section{Całkowanie przez części i granica sum}
\begin{twierdzenie}[Deterministyczne całkowanie przez części]
Jeżeli $f:[0,T]\to\R$ jest funkcją absolutnie ciągłą
i $f'\in L^2([0,T])$, to
\[
\int_0^T f(t)\,\dd W_t
=f(T)W_T-\int_0^T f'(t)W_t\,\dd t
\quad\text{prawie na pewno}.
\tag{1.6}\label{eq:ibp-det}
\]
Całka po prawej jest zwykłą całką Lebesgue'a względem czasu.
\end{twierdzenie}
\begin{proof}
Dla funkcji schodkowej wyznaczonej wartościami $f(t_k)$
suma Wienera spełnia dyskretne całkowanie przez części:
\[
\sum_{k=0}^{n-1}f(t_k)\Delta_kW
=f(T)W_T-\sum_{k=0}^{n-1}W_{t_{k+1}}(f(t_{k+1})-f(t_k))
\]
jest dokładną tożsamością teleskopową, ponieważ $W_0=0$.
Wystarczy zatem zbadać granice obu stron. Funkcje $f_n(t)=f(t_k)$ dla
$t\in(t_k,t_{k+1}]$ zbiegają jednostajnie do $f$,
a więc w $L^2([0,T])$. Lewa strona zbiega zatem do
$I_W(f)$ w $L^2(\Omega)$.
Ponieważ $f$ jest absolutnie ciągła,
$\Delta_k f=\int_{t_k}^{t_{k+1}}f'(u)\,\dd u$,
a więc druga suma jest równa
$\int_0^T W_{r_n(u)}f'(u)\,\dd u$,
gdzie $r_n(u)$ jest prawym końcem przedziału zawierającego $u$.
Ciągłość trajektorii $W$ zapewnia jednostajną zbieżność
$W_{r_n(u)}\to W_u$ prawie na pewno. Ponieważ
$\int_0^T|f'(u)|\,\dd u<\infty$, suma zbiega niemal
na pewno do $\int_0^TW_uf'(u)\,\dd u$.
Granice w prawdopodobieństwie są jednoznaczne, co daje tezę.
\end{proof}

\begin{przyklad}[Ważona całka z procesu Wienera]
Dla $f(t)=t$ otrzymujemy
$\int_0^Tt\,\dd W_t=T W_T-\int_0^TW_t\,\dd t$.
Losowa zmienna po lewej ma rozkład
$\Normal(0,\int_0^Tt^2\,\dd t)=\Normal(0,T^3/3)$.
Nie mylimy jej ze zwykłą całką $\int_0^TW_t\,\dd t$.
Co ciekawe, obie zmienne mają ten sam rozkład normalny
$\Normal(0,T^3/3)$, a nawet tę samą kowariancję
$T^2/2$ z $W_T$. Nie są jednak równe:
na mocy \eqref{eq:ibp-det} ich różnica to
$\int_0^T(2t-T)\,\dd W_t$, mająca wariancję
$T^3/3>0$. To przykład pokazujący, że jednakowy
rozkład nie oznacza równości zmiennych losowych.
\end{przyklad}

\section{Miara Wienera a całka po trajektoriach}
Dotychczas całkowaliśmy funkcję czasu względem losowych
przyrostów. Istnieje jeszcze inny rodzaj całkowania:
całkowanie funkcjonałów całej ścieżki względem \emph{miary
Wienera}. Nie wolno tych pojęć utożsamiać.

\begin{definicja}[Miara Wienera i całka po trajektoriach]
Niech $C_0([0,T])$ oznacza przestrzeń funkcji ciągłych
$w:[0,T]\to\R$ o $w(0)=0$, z normą supremum i borelowskim
$\sigma$-ciałem. Niech $\mu_W$ będzie rozkładem procesu Wienera
jako zmiennej losowej o wartościach w $C_0([0,T])$:
$\mu_W(A):=\Prob(W\in A)$ dla $A\in\calB(C_0([0,T]))$.
Dla mierzalnego i całkowalnego funkcjonału
$F:C_0([0,T])\to\R$ jego \emph{całką względem miary Wienera}
nazywamy
\[
\int_{C_0([0,T])}F(w)\,\dd\mu_W(w)=\E[F(W)].
\tag{1.7}
\]
\end{definicja}

\begin{uwaga}[Różne obiekty po dwóch stronach znaku całki]
W $\int_0^Tf(t)\,\dd W_t$ zmienna $t$ przebiega przedział
czasu, a wynikiem jest zmienna losowa. W
$\int_{C_0([0,T])}F(w)\,\dd\mu_W(w)$ przebiegamy zbiór
całych funkcji $w$ i otrzymujemy liczbę. Funkcjonał $F$
może zależeć na przykład od maksimum trajektorii,
jej pola pod wykresem albo wartości w wielu chwilach.
Zapis $\dd\mu_W(w)$ jest ścisłym zapisem całki po trajektoriach
w czasie euklidesowym. Formalna całka Feynmana o oscylacyjnym
czynniku $\exp(iS/\hbar)$ wymaga innej konstrukcji.
\end{uwaga}

\begin{zadanie}[Całka funkcyjna dla pola pod trajektorią]\label{zad:area-char}
Obliczyć dla $u\in\R$ całkę
\[
\int_{C_0([0,T])}
\exp\left(iu\int_0^T w(s)\,\dd s\right)\,\dd\mu_W(w).
\]
\emph{Rozwiązanie.}
Kładziemy $A_T=\int_0^TW_s\,\dd s$.
Z \eqref{eq:ibp-det} zastosowanego do $f(s)=T-s$
wynika $A_T=I_W(T-\cdot)$, ponieważ $(T-s)'=-1$
i $f(T)=0$. Jest to zmienna normalna o średniej zero
i wariancji
$\int_0^T(T-s)^2\,\dd s=T^3/3$.
Zatem jej funkcja charakterystyczna wynosi
\[
\E e^{iuA_T}=\exp\left(-\frac{u^2T^3}{6}\right).
\]
Otrzymaliśmy całkę po nieskończeniewymiarowej przestrzeni
trajektorii, sprowadzając ją do funkcji charakterystycznej
jednej normalnej zmiennej losowej.
\end{zadanie}

\begin{przyklad}[Całka cylindryczna]
Niech $0<t_1<\cdots<t_m\le T$ i
$F(w)=\varphi(w(t_1),\ldots,w(t_m))$ dla ograniczonej
funkcji borelowskiej $\varphi:\R^m\to\R$.
Wówczas $\int F\,\dd\mu_W$ jest zwykłą
$m$-wymiarową całką względem wspólnego rozkładu
$(W_{t_1},\ldots,W_{t_m})$. Jest to rozkład gaussowski
z macierzą kowariancji $K_{ij}=\min(t_i,t_j)$.
Dla przykładów nieskończeniewymiarowych naturalną metodą
jest aproksymowanie funkcjonału takimi funkcjonałami
cylindrycznymi, a następnie uzasadnienie zbieżności całek.
\end{przyklad}

\section{Dlaczego całka Wienera nie jest punktowym funkcjonałem ciągłym?}
Izometria \eqref{eq:wiener-iso} zachodzi pomiędzy
\emph{przestrzeniami klas funkcji}. Nie mówi, że dla prawie
każdego $\omega$ odwzorowanie
$f\mapsto I_W(f)(\omega)$ jest ograniczonym funkcjonałem
liniowym na $L^2([0,T])$. Takie stwierdzenie byłoby fałszywe.

\begin{twierdzenie}[Brak jednoczesnej ścieżkowej ciągłości]
Nie istnieje zbiór $\Omega_0\subset\Omega$ o
$\Prob(\Omega_0)=1$ i rodzina ograniczonych liniowych
funkcjonałów $L_\omega:L^2([0,T])\to\R$, $\omega\in\Omega_0$,
taka że dla każdego $f\in L^2([0,T])$ równość
$L_\omega(f)=I_W(f)(\omega)$ zachodzi niemal na pewno
(w sensie wybranych wersji zmiennych $I_W(f)$).
\end{twierdzenie}
\begin{proof}
Niech $(e_n)_{n\ge1}$ będzie układem ortonormalnym zupełnym
w $L^2([0,T])$; taka baza istnieje, ponieważ jest to
separowalna przestrzeń Hilberta. Z twierdzenia
\ref{th:wiener-gaussian} wynika, że
$Z_n:=I_W(e_n)$ są niezależnymi zmiennymi o rozkładzie
$\Normal(0,1)$. Ponieważ baza jest przeliczalna, możemy
wybrać wspólny zbiór pełnego prawdopodobieństwa, na którym
$L_\omega(e_n)=Z_n(\omega)$ dla wszystkich $n$.
Jeżeli $L_\omega$ jest ograniczonym funkcjonałem liniowym,
na mocy twierdzenia Riesza istnieje $h_\omega\in L^2([0,T])$
taka, że $L_\omega(f)=\langle f,h_\omega\rangle$.
Nierówność Bessela daje wtedy
$\sum_{n\ge1}Z_n(\omega)^2\le\norm{h_\omega}_2^2<\infty$.
Tymczasem mocne prawo wielkich liczb dla niezależnych
zmiennych $Z_n^2$ zapewnia
$n^{-1}\sum_{k=1}^nZ_k^2\to\E Z_1^2=1$ prawie na pewno.
Oznacza to, że suma $\sum Z_n^2$ jest nieskończona
prawie na pewno, sprzeczność.
\end{proof}

\begin{uwaga}[Intuicja: biały szum]
Zapis heurystyczny $\dd W_t=\xi(t)\,\dd t$ sugeruje,
że całkujemy $f$ względem \emph{białego szumu} $\xi$.
Powyższe twierdzenie wyjaśnia, dlaczego $\xi$ nie może
być zwykłą losową funkcją z $L^2([0,T])$. Można jednak
zdefiniować biały szum jako losowy obiekt uogólniony,
który działa na funkcjach testowych; ścisłą teorię
dystrybucji losowych rozwiniemy w dalszej części tomu.
\end{uwaga}

\section{Podsumowanie konstrukcji i zadania}
Konstrukcja całki Wienera opiera się na trzech zasadach:
niezależnych normalnych przyrostach, izometrii w $L^2$
i zupełności tej przestrzeni. Umożliwia ona całkowanie
funkcji deterministycznych $f\in L^2$, lecz nie wystarcza
dla losowej funkcji podcałkowej zależnej od wartości
procesu w chwili $t$. Tę trudność rozwiążemy przez
warunkowanie względem informacji dostępnej do czasu $t$.

\begin{zadanie}[Niezależne całki z niepokrywających się przedziałów]
Dla $0<a<b<c<T$ wyznaczyć rozkład wektora
$(X,Y)$, gdzie $X=\int_a^b t\,\dd W_t$,
$Y=\int_b^c (t-b)\,\dd W_t$. Czy jego współrzędne
są niezależne?

\emph{Rozwiązanie.} Jest to wektor centrowany gaussowski.
Wariancje wynoszą $\Var(X)=(b^3-a^3)/3$ oraz
$\Var(Y)=(c-b)^3/3$, a kowariancja jest równa zero,
gdyż iloczyn funkcji podcałkowych znika niemal wszędzie.
Stąd $X$ i $Y$ są niezależne. Wspólny punkt $b$ nie
wpływa na wynik, ponieważ ma zerową miarę Lebesgue'a.
\end{zadanie}

\chapter{Całka Itô: losowe funkcje podcałkowe}
\label{ch:ito-construction}
\section{Dlaczego funkcja podcałkowa musi znać tylko przeszłość?}
W rozdziale pierwszym współczynnik $a_k$ w sumie
$\sum a_k\Delta_kW$ był ustaloną liczbą. Jeżeli chcemy
modelować strategię inwestycyjną, siłę losowego zakłócenia
albo przyrost procesu zależny od aktualnego położenia,
to naturalny współczynnik ma postać $H_{t_k}(\omega)$.
Aby rachunek nie wykorzystywał przyszłości, musimy żądać,
by $H_{t_k}$ był mierzalny względem $\calF_{t_k}$.

W przeciwnym razie suma mogłaby znać znak następnego
przyrostu i osiągać dodatnią średnią nawet wtedy, gdy
podstawowy proces Wienera jest martyngałem. Na przykład
$H_{t_k}=\operatorname{sgn}(\Delta_kW)$
uzależnia decyzję od przyrostu, który jeszcze nie nastąpił:
wtedy każdy składnik $H_{t_k}\Delta_kW=|\Delta_kW|$
jest nieujemny. Taka strategia nie jest przewidywalna.

\begin{definicja}[Proces prosty przewidywalny]
Proces $H:[0,T]\times\Omega\to\R$ nazywamy prostym
przewidywalnym, jeżeli istnieje deterministyczny podział
$0=t_0<\cdots<t_n=T$ i zmienne losowe $\xi_k$
mierzalne względem $\calF_{t_k}$, takie że
\[
H_t=\sum_{k=0}^{n-1}\xi_k\ind{(t_k,t_{k+1}]}(t),
\qquad 0\le t\le T.
\tag{2.1}\label{eq:ito-step}
\]
W przestrzeni $\calS_{\mathrm{pred}}^2$ rozpatrujemy
te procesy proste, dla których $\E\xi_k^2<\infty$
na każdym przedziale o dodatniej długości.
\end{definicja}

\begin{uwaga}[Wspólna filtracja ma znaczenie]
Jeżeli do filtracji dołączylibyśmy już w chwili zerowej
wartość $W_T$, to $W$ nie byłby procesem Wienera
\emph{względem powiększonej filtracji}: przyrost
$W_T-W_0$ nie byłby niezależny od $\calF_0$.
Dlatego każda definicja całki stochastycznej zawiera
zarówno proces Wienera, jak i filtrację, względem której
jest on procesem Wienera.
\end{uwaga}

\begin{center}
\begin{tikzpicture}[>=Latex,x=1cm,y=0.9cm]
\draw[->] (0,0)--(9.2,0) node[right] {$t$};
\draw[->] (0,-.35)--(0,2.8) node[above] {$H_t$};
\fill[blue!12] (.2,0) rectangle (2,1.35);
\fill[blue!20] (2,0) rectangle (4.15,2.15);
\fill[blue!12] (4.15,0) rectangle (6.4,.8);
\fill[blue!20] (6.4,0) rectangle (8.45,1.8);
\draw[thick,blue!65!black] (.2,1.35)--(2,1.35) (2,2.15)--(4.15,2.15)
(4.15,.8)--(6.4,.8) (6.4,1.8)--(8.45,1.8);
\foreach \x/\name in {0.2/{t_0},2/{t_1},4.15/{t_2},6.4/{t_3},8.45/{t_4}}{
 \draw[dashed,gray] (\x,0)--(\x,2.5);
 \node[below,scale=.83] at (\x,0) {$\name$};}
\node[scale=.86,blue!55!black] at (1.05,1.68) {$\xi_0$};
\node[scale=.86,blue!55!black] at (3.05,2.48) {$\xi_1$};
\node[scale=.86,blue!55!black] at (5.3,1.13) {$\xi_2$};
\node[scale=.86,blue!55!black] at (7.5,2.12) {$\xi_3$};
\end{tikzpicture}
\par\smallskip\small Przewidywalny proces prosty: wysokość przedziału jest znana przy jego lewym końcu.
\end{center}

\section{Całka prostego procesu i warunkowa izometria}
\begin{definicja}[Całka Itô dla procesów prostych]
Dla $H$ postaci \eqref{eq:ito-step} i $t\in[0,T]$ kładziemy
\[
I_t(H):=\int_0^t H_s\,\dd W_s
:=\sum_{k=0}^{n-1}\xi_k
\big(W_{t\wedge t_{k+1}}-W_{t\wedge t_k}\big).
\tag{2.2}\label{eq:ito-simple}
\]
Zmienna $I_t(H)$ jest $\calF_t$-mierzalna.
\end{definicja}

\begin{lemat}[Poprawność definicji]
Wzór \eqref{eq:ito-simple} nie zależy od wyboru
reprezentacji $H$ jako procesu prostego (z dokładnością
do równości $\dd t\otimes\dd\Prob$-prawie wszędzie).
Dla prostych procesów $H,K$ i $a,b\in\R$ zachodzi
$I_t(aH+bK)=aI_t(H)+bI_t(K)$ niemal na pewno.
\end{lemat}
\begin{proof}
Przy rozdrobnieniu deterministycznego podziału współczynnik
$\xi_k$ na każdym nowym przedziale pozostaje mierzalny
względem $\calF$ jego lewego końca, bo filtracja jest rosnąca.
Przyrost procesu Wienera dzieli się na sumę przyrostów,
więc wartość całki nie zmienia się. Dwie reprezentacje
procesu sprowadzamy do wspólnego podziału. Na każdym
przedziale ich współczynniki są równe niemal na pewno,
ponieważ norma różnicy w przestrzeni produktowej jest zerowa:
$\E\int_0^T|H-K|^2\,\dd t=0$.
Po usunięciu skończenie wielu zbiorów zerowych
współczynniki są równe jednocześnie. Stąd niezależność
od reprezentacji i liniowość.
\end{proof}

\begin{twierdzenie}[Martyngałowość i izometria Itô dla procesów prostych]
\label{th:ito-simple-iso}
Niech $H,K\in\calS_{\mathrm{pred}}^2$. Wtedy
$I(H)=(I_t(H))_{0\le t\le T}$ jest ciągłym
martyngałem z $I_0(H)=0$ i zachodzą
\[
\E I_t(H)=0,\qquad
\E[I_t(H)I_t(K)]
=\E\int_0^tH_sK_s\,\dd s,
\tag{2.3}\label{eq:ito-iso-simple}
\]
\[
\E\big[(I_t(H)-I_s(H))^2\mid\calF_s\big]
=\E\left[\left.\int_s^tH_u^2\,\dd u\,\right|\calF_s\right]
\quad(0\le s\le t\le T).
\tag{2.4}\label{eq:ito-conditional-var}
\]
\end{twierdzenie}
\begin{proof}
Ustalmy wspólny podział dla $H,K$ i początkowo przyjmijmy
$t=T$. Zapiszmy $H=\sum\xi_k\ind{(t_k,t_{k+1}]}$,
$K=\sum\eta_k\ind{(t_k,t_{k+1}]}$ i
$\Delta_kW=W_{t_{k+1}}-W_{t_k}$.
Dla każdej $k$ mamy
$\E[\Delta_kW\mid\calF_{t_k}]=0$ i
$\E[(\Delta_kW)^2\mid\calF_{t_k}]=\Delta t_k$,
ponieważ przyrost jest niezależny od $\calF_{t_k}$.
Stąd
$\E[\xi_k\Delta_kW\mid\calF_{t_k}]=0$ i
$\E[\xi_k\eta_k(\Delta_kW)^2]=\E[\xi_k\eta_k]\Delta t_k$.

Dla $j<k$ zmienna $\xi_j\eta_k\Delta_jW$ jest
$\calF_{t_k}$-mierzalna, zatem
\[
\E[\xi_j\eta_k\Delta_jW\Delta_kW]
=\E\big[\xi_j\eta_k\Delta_jW\,
\E(\Delta_kW\mid\calF_{t_k})\big]=0.
\]
Identycznie znikają wyrazy o $k<j$.
Po wymnożeniu iloczynu dwóch sum pozostają więc
wyłącznie wyrazy diagonalne:
\[
\E[I_T(H)I_T(K)]
=\sum_k\E(\xi_k\eta_k)\Delta t_k
=\E\int_0^TH_sK_s\,\dd s.
\]
Gdy $t<T$, wystarczy uciąć procesy na $(0,t]$,
co dowodzi \eqref{eq:ito-iso-simple}. Analogicznie,
aby pokazać własność martyngałową, dzielimy $I_t(H)$
na $I_s(H)$ i sumę przyrostów z $(s,t]$. Każdy taki
składnik ma warunkową średnią zero względem informacji
w swoim lewym końcu; iterowane warunkowanie daje
$\E[I_t(H)-I_s(H)\mid\calF_s]=0$. Ciągłość trajektorii
wynika ze skończonej sumy ciągłych procesów przyrostowych.

Na koniec zastosujmy identyczne rozwinięcie kwadratu
do całki na przedziale $(s,t]$, ale tym razem z warunkowaniem
względem $\calF_s$. Wyrazy mieszane znikają po iterowanym
warunkowaniu, a składniki diagonalne mają warunkową
wartość oczekiwaną równą
$\E[\xi_k^2\Delta t_k\mid\calF_s]$.
Zsumowanie daje \eqref{eq:ito-conditional-var}.
\end{proof}

\begin{uwaga}[Uzasadnienie całkowalności iloczynów]
Jeżeli $\E\xi_k^2$ i $\E\eta_k^2$ są skończone,
to $\E|\xi_k\eta_k|<\infty$ z nierówności
Cauchy'ego--Schwarza. Same składniki całek mają skończone
drugie momenty dzięki warunkowej wariancji przyszłego
przyrostu. W konsekwencji wszystkie zastosowania
warunkowej wartości oczekiwanej w dowodzie są uprawnione.
\end{uwaga}

\section{Procesy przewidywalne i przestrzeń całkowania}
W praktyce proces $H$ nie musi mieć skończonej liczby zmian.
Na przykład $H_t=W_t$ jest ciągły i zależy od całej
przeszłości. Aby przybliżać takie procesy przez
$\calS_{\mathrm{pred}}^2$, potrzebujemy dokładnie określonej
mierzalności względem czasu i zdarzeń losowych.

\begin{definicja}[Przewidywalne $\sigma$-ciało]
Na przestrzeni $[0,T]\times\Omega$ definiujemy
\emph{przewidywalne $\sigma$-ciało} $\calP_T$ jako
$\sigma$-ciało generowane przez zbiory
$\{0\}\times A$, $A\in\calF_0$, oraz
$(s,t]\times A$, gdzie $0\le s<t\le T$ i $A\in\calF_s$.
Proces $H$ nazywamy \emph{przewidywalnym}, gdy odwzorowanie
$(t,\omega)\mapsto H_t(\omega)$ jest
$\calP_T/\calB(\R)$-mierzalne.
Definiujemy
\[
\calH_T^2:=\left\{H\text{ przewidywalny}:
\E\int_0^T H_t^2\,\dd t<\infty\right\}\big/\sim,
\quad H\sim K\iff H=K\quad
(\dd t\otimes\Prob)\text{-p.w.}
\tag{2.5}
\]
Norma ma postać
$\norm{H}_{\calH_T^2}:=(\E\int_0^TH_t^2\,\dd t)^{1/2}$.
\end{definicja}

\begin{uwaga}[Oznaczenia i mierzalność]
Produktem miar w \eqref{eq:ito-iso-simple} jest
$\dd t\otimes\dd\Prob$, w zwykłym znaczeniu miary
Lebesgue'a na czasie i miary probabilistycznej na $\Omega$.
Nie każdy proces adaptowany jest automatycznie przewidywalny.
Natomiast każdy adaptowany proces o ciągłych albo
lewostronnie ciągłych trajektoriach jest przewidywalny:
można przybliżać go punktowo procesami stałymi na
przedziałach, używając wartości w lewych końcach.
\end{uwaga}

\begin{lemat}[Gęstość prostych procesów przewidywalnych]
\label{lem:pred-dense}
Przestrzeń $\calS_{\mathrm{pred}}^2$ jest gęsta w
$\calH_T^2$.
\end{lemat}
\begin{proof}
Przestrzeń $\calH_T^2$ jest zamkniętą podprzestrzenią
$L^2([0,T]\times\Omega,\calB([0,T])\otimes\calF,
\dd t\otimes\Prob)$ z warunkiem mierzalności względem
$\calP_T$, a więc jest zupełna.
Rozważmy rodzinę prostokątów $(s,t]\times A$ z
$A\in\calF_s$. Ich wskaźniki odpowiadają dokładnie
prostym procesom przewidywalnym
$\ind A\ind{(s,t]}$. Rodzina tych prostokątów, wraz
ze zbiorami z chwili zerowej, jest zamknięta na
skończone przecięcia: jeśli $s\le u$, to
$\big((s,t]\times A\big)\cap\big((u,v]\times B\big)$
jest albo pusta, albo równa
$(u,\min\{t,v\}]\times(A\cap B)$, gdzie
$A\cap B\in\calF_u$. Tak więc jest to $\pi$-układ
generujący $\calP_T$.

Niech $\calD$ będzie klasą zbiorów $C\in\calP_T$,
dla których wskaźnik $\ind C$ należy do domknięcia
w $L^2(\dd t\otimes\Prob)$ przestrzeni liniowej
rozpiętej przez wskaźniki wspomnianych prostokątów.
Ponieważ miara produktowa ma skończoną masę $T$,
$\calD$ jest klasą Dynkina: zawiera całą przestrzeń
(z dokładnością do przekroju $\{0\}\times\Omega$ o
zerowej mierze produktowej), jest zamknięta na różnice
zbiorów zagnieżdżonych oraz na przeliczalne rozłączne
sumy, co wynika ze zbieżności wskaźników w $L^2$
i z domknięcia. Twierdzenie o $\pi$- i $\lambda$-układach
daje $\calD=\calP_T$. Wobec tego wskaźniki dowolnych
zbiorów przewidywalnych dają się aproksymować w $L^2$
przez skończone liniowe kombinacje wskaźników prostokątów.
Na koniec każdą funkcję z $L^2(\calP_T)$ przybliżamy
funkcjami prostymi. Otrzymane aproksymacje są
prostymi procesami przewidywalnymi o ograniczonych
współczynnikach, a zatem należą do
$\calS_{\mathrm{pred}}^2$.
\end{proof}

\begin{twierdzenie}[Konstrukcja całki Itô w $L^2$]\label{th:ito-extension}
Dla każdego $H\in\calH_T^2$ istnieje jednoznaczny
(z dokładnością do nierozróżnialności) ciągły,
adaptowany martyngał $I(H)=(I_t(H))_{0\le t\le T}$
z $I_0(H)=0$, taki że dla każdego ciągu
$H^n\in\calS_{\mathrm{pred}}^2$ z
$\norm{H^n-H}_{\calH_T^2}\to0$ zachodzi
\[
\E\sup_{0\le t\le T}|I_t(H^n)-I_t(H)|^2\longrightarrow0.
\tag{2.6}\label{eq:ito-process-limit}
\]
Odwzorowanie $H\mapsto I_T(H)$ jest liniową izometrią
$\calH_T^2\to L^2(\Omega)$ oraz dla dowolnych
$H,K\in\calH_T^2$ i $t\in[0,T]$ zachodzi
\[
\E\left|\int_0^tH_s\,\dd W_s\right|^2
=\E\int_0^tH_s^2\,\dd s,
\qquad
\E\left[\int_0^tH_s\,\dd W_s\int_0^tK_s\,\dd W_s\right]
=\E\int_0^tH_sK_s\,\dd s.
\tag{2.7}\label{eq:ito-full-iso}
\]
\end{twierdzenie}
\begin{proof}
Gęstość z lematu \ref{lem:pred-dense} zapewnia istnienie
$H^n\to H$ w $\calH_T^2$. Dla procesów prostych
$M^n:=I(H^n)$ nierówność maksymalna Dooba i
\eqref{eq:ito-iso-simple} dają
\[
\E\sup_{t\le T}|M_t^n-M_t^m|^2
\le4\E|M_T^n-M_T^m|^2
=4\norm{H^n-H^m}_{\calH_T^2}^2.
\tag{2.8}\label{eq:ito-doob-cauchy}
\]
W szczególności $(M^n)$ jest ciągiem Cauchy'ego
w normie $\norm{M}_{\mathbb S^2}:=(\E\sup_{t\le T}|M_t|^2)^{1/2}$.
Aby nie opierać się niejawnie na zupełności przestrzeni
procesów, wybierzmy podciąg $(M^{n_j})$, dla którego
$\sum_j\norm{M^{n_{j+1}}-M^{n_j}}_{\mathbb S^2}<\infty$.
Z nierówności Jensena i zbieżności monotonicznej
$\E\sum_j\sup_t|M_t^{n_{j+1}}-M_t^{n_j}|<\infty$.
Zatem suma supremów jest skończona prawie na pewno,
a podciąg zbiega jednostajnie na $[0,T]$ poza jednym
zbiorem zerowym. Jego granica ma ciągłe trajektorie.
Niech $M$ oznacza ten proces, z wartością zero na
wyjątkowym zbiorze (który można wybrać w $\calF_0$
przez zupełność filtracji).

Ponieważ procesy $M^{n_j}$ są adaptowane, ich granice
punktowe są $\calF_t$-mierzalne. Z nierówności
\eqref{eq:ito-doob-cauchy} wynika, że cały ciąg
$(M^n)$ zbiega do $M$ w normie $\mathbb S^2$.
Dla $s\le t$ mamy
$\E[M_t^n\mid\calF_s]=M_s^n$; zbieżność w $L^2$
i kontrakcja warunkowej wartości oczekiwanej pozwalają
przejść do granicy i otrzymać $\E[M_t\mid\calF_s]=M_s$.
Tym samym $M$ jest martyngałem.

Gdybyśmy wybrali inny ciąg aproksymujący $K^n\to H$,
to \eqref{eq:ito-doob-cauchy} dla różnicy pokazuje,
że oba procesy graniczne mają zerową normę odległości
w $\mathbb S^2$. Są więc nierozróżnialne.
Izometrię \eqref{eq:ito-full-iso} otrzymujemy,
przechodząc do granicy w odpowiedniej równości dla
procesów prostych, najpierw dla kwadratów, a następnie
stosując polaryzację. Dokładnie to samo daje
\eqref{eq:ito-process-limit}.
\end{proof}

\begin{wniosek}[Oszacowanie maksymalne]
Dla każdego $H\in\calH_T^2$ zachodzi
\[
\E\sup_{0\le t\le T}\left|\int_0^tH_s\,\dd W_s\right|^2
\le4\E\int_0^TH_s^2\,\dd s.
\tag{2.9}\label{eq:ito-max}
\]
\end{wniosek}
\begin{proof}
Jest to nierówność maksymalna Dooba dla martyngału
z twierdzenia \ref{th:ito-extension} wraz z izometrią.
\end{proof}

\begin{uwaga}[Co daje oszacowanie maksymalne?]
Sama izometria w chwili $T$ kontroluje tylko
$\E|I_T(H)|^2$. Oszacowanie \eqref{eq:ito-max}
kontroluje \emph{całą ścieżkę} całki stochastycznej.
Pozwala więc uzyskać zbieżność jednostajną na zwartym
przedziale w prawdopodobieństwie, a po wyborze podciągu
również zbieżność jednostajną prawie na pewno. To właśnie
uzasadnia wybór ciągłej wersji procesu całkowego.
\end{uwaga}


\begin{lemat}[Ciągły proces adaptowany jest przewidywalny]
Jeżeli $H$ jest rzeczywistym, adaptowanym procesem
o ciągłych trajektoriach, to jest przewidywalny.
Jeśli ponadto $\E\int_0^TH_t^2\,\dd t<\infty$,
to należy do $\calH_T^2$.
\end{lemat}
\begin{proof}
Dla każdego $n$ wybierzmy deterministyczny podział
$0=t_0^{(n)}<\cdots<t_{k_n}^{(n)}=T$ o średnicy
dążącej do zera i zdefiniujmy
$H_t^{(n)}=\sum_k H_{t_k^{(n)}}
\ind{(t_k^{(n)},t_{k+1}^{(n)}]}(t)$,
a w punkcie zero przyjmijmy $H_0^{(n)}=H_0$.
Każdy $H^{(n)}$ jest przewidywalny, bo
$H_{t_k^{(n)}}$ jest $\calF_{t_k^{(n)}}$-mierzalny.
Dla każdej ciągłej trajektorii mamy
$H_t^{(n)}\to H_t$ dla każdego $t$ dzięki
ciągłości w czasie. Granica punktowa funkcji
mierzalnych względem $\calP_T$ pozostaje
$\calP_T$-mierzalna. Druga część wynika wprost
z definicji przestrzeni $\calH_T^2$.
\end{proof}

\begin{przyklad}[Aproksymacja $H_t=g(W_t)$ i jej dokładność]
Niech $g:\R\to\R$ będzie funkcją globalnie Lipschitzowską
z pewną stałą $L\ge0$. Proces $H_t=g(W_t)$ jest
adaptowany i ciągły, a także należy do $\calH_T^2$,
ponieważ $|g(x)|\le|g(0)|+L|x|$ i
$\E W_t^2=t$. Na równomiernej siatce $t_k=kT/n$
przybliżmy $H$ przez
$H_t^{(n)}=g(W_{t_k})$ dla $t\in(t_k,t_{k+1}]$.
Wówczas
\[
\E\int_0^T|H_t^{(n)}-H_t|^2\,\dd t
\le L^2\sum_{k=0}^{n-1}\int_{t_k}^{t_{k+1}}
\E|W_t-W_{t_k}|^2\,\dd t
=\frac{L^2T^2}{2n}.
\]
Izometria Itô daje błąd średniokwadratowy całki
w chwili $T$ nie większy niż $L^2T^2/(2n)$,
a nierówność Dooba daje analogiczną kontrolę
całego procesu, z czterokrotnie większą stałą.
Takie oszacowanie jest punktem wyjścia
numerycznej aproksymacji całek stochastycznych.
\end{przyklad}

\section{Podstawowe własności wynikające z konstrukcji}
\begin{twierdzenie}[Lokalność w czasie i zależność od informacji]
\label{th:ito-properties}
Niech $H,K\in\calH_T^2$, $0\le s\le t\le T$.
Wówczas:
\begin{enumerate}[label=\textup{(\roman*)}]
\item $\displaystyle\int_0^tH_u\,\dd W_u
=\int_0^T H_u\ind{(0,t]}(u)\,\dd W_u$;
\item jeśli $H=K$ $\dd t\otimes\Prob$-prawie wszędzie
na $[0,t]\times\Omega$, to $I_u(H)=I_u(K)$ dla wszystkich
$u\le t$ poza jednym zbiorem zerowym;
\item $\displaystyle I_t(H)-I_s(H)=\int_0^TH_u\ind{(s,t]}(u)\,\dd W_u$;
\item jeśli $A\in\calF_s$, to
$\displaystyle\ind A\big(I_t(H)-I_s(H)\big)
=\int_s^t\ind A H_u\,\dd W_u$;
\item zachodzi warunkowa izometria
\[
\E\big[(I_t(H)-I_s(H))^2\mid\calF_s\big]
=\E\left[\left.\int_s^tH_u^2\,\dd u\,\right|\calF_s\right].
\]
\end{enumerate}
\end{twierdzenie}
\begin{proof}
Dla procesów prostych (i) i (iii) wynikają z samej
definicji sum. Dla ogólnych procesów przechodzimy
do granicy w $\calH_T^2$; pomnożenie procesu przez
wskaźnik deterministycznego przedziału nie zwiększa
normy. W (ii) izometria stosowana do $(H-K)\ind{(0,t]}$
daje równość całek w każdej chwili $u\le t$;
ciągłość obu procesów rozszerza ją z chwil wymiernych
na wszystkie chwile poza jednym zbiorem zerowym.

W (iv) proces $\ind A H_u\ind{(s,t]}(u)$ jest
przewidywalny, ponieważ $A\in\calF_s$ i na
$(s,t]$ informacja $A$ jest już dostępna. Dla procesu
prostego wzór wynika z przeniesienia wskaźnika $\ind A$
przed skończoną sumę. Dla ogólnego $H$ aproksymujemy
przez procesy proste i korzystamy z izometrii,
gdyż mnożenie przez $\ind A$ jest kontrakcją
w $\calH_T^2$.

W (v) dla procesu prostego mamy \eqref{eq:ito-conditional-var}.
Dla $H^n\to H$ w $\calH_T^2$ przyrosty całek
zbiegają w $L^2$, więc ich kwadraty zbiegają w $L^1$:
$\E|X_n^2-X^2|\le\norm{X_n-X}_2(\norm{X_n}_2+\norm X_2)$.
Jednocześnie $\int_s^t(H_u^n)^2\,\dd u$ zbiega
w $L^1(\Omega)$ do $\int_s^tH_u^2\,\dd u$,
co wynika z tej samej nierówności Cauchy'ego--Schwarza
zastosowanej w przestrzeni produktowej.
Ponieważ warunkowa wartość oczekiwana jest kontrakcją
w $L^1$, przechodzimy do granicy w tożsamości
warunkowej dla procesów prostych.
\end{proof}

\begin{twierdzenie}[Kwadrat całki i przewidywalny nawias]
\label{th:ito-bracket}
Dla $H\in\calH_T^2$ proces
\[
Q_t:=\left(\int_0^tH_s\,\dd W_s\right)^2
-\int_0^tH_s^2\,\dd s,\qquad 0\le t\le T,
\tag{2.10}\label{eq:ito-square-mg}
\]
jest martyngałem. Jeżeli przez $\langle M\rangle$
oznaczamy przewidywalny nawias ciągłego martyngału
kwadratowo całkowalnego z tomu II, to dla
$M_t=\int_0^tH_s\,\dd W_s$ mamy
$\langle M\rangle_t=\int_0^tH_s^2\,\dd s$.
W szczególności dla ciągłego martyngału
$[M]=\langle M\rangle$, więc
$[M]_t=\int_0^tH_s^2\,\dd s$.
\end{twierdzenie}
\begin{proof}
Własność martyngałową $M$ oraz warunkową izometrię
ustaliliśmy powyżej. Dla $s\le t$ zapisujemy
$M_t=M_s+(M_t-M_s)$ i rozwijamy kwadrat.
Po warunkowaniu względem $\calF_s$ składnik liniowy
$2M_s\E[M_t-M_s\mid\calF_s]$ znika. Pozostaje
\[
\E[M_t^2\mid\calF_s]
=M_s^2+\E\left[\left.\int_s^tH_u^2\,\dd u\,\right|\calF_s\right].
\]
Odjęcie $\E[\int_0^tH_u^2\,\dd u\mid\calF_s]$
daje $\E[Q_t\mid\calF_s]=Q_s$. Wszystkie wyrazy
są całkowalne na mocy izometrii.
Proces $A_t=\int_0^tH_s^2\,\dd s$ jest adaptowany,
ciągły i niemalejący, a więc przewidywalny.
Z definicji i jednoznaczności przewidywalnego nawiasu
ciągłego martyngału kwadratowo całkowalnego wynika
$\langle M\rangle=A$. W tomie II wykazaliśmy,
że dla ciągłych lokalnych martyngałów wariacja
kwadratowa i nawias przewidywalny są tym samym
ciągłym procesem; stąd ostatnia równość.
\end{proof}

\begin{uwaga}[Nie jest to jeszcze pełny wzór Itô]
Tożsamość \eqref{eq:ito-square-mg} identyfikuje
martyngałowy kompensator kwadratu. W kolejnym rozdziale
wyprowadzimy wzór dla ogólnego $F(t,X_t)$, a także
rachunek różniczkowy z poprawką wariacji kwadratowej.
Nie powinniśmy traktować symboli $(\dd W_t)^2=\dd t$
jako zwykłej algebry nieskończenie małych wielkości;
są one skrótem dla ścisłych twierdzeń o granicach sum.
\end{uwaga}


\begin{przyklad}[Martyngał bez niezależnych przyrostów]
Proces $M_t:=\int_0^tW_u\,\dd W_u$
jest martyngałem, ale jego przyrosty nie są niezależne
od przeszłości. Z warunkowej izometrii dla $s<t$
wynika
\[
\E[(M_t-M_s)^2\mid\calF_s]
=\E\left[\left.\int_s^tW_u^2\,\dd u\,\right|\calF_s\right]
=\int_s^t\big(W_s^2+(u-s)\big)\,\dd u
=(t-s)W_s^2+\frac{(t-s)^2}{2}.
\]
W drugim przejściu użyliśmy rozkładu
$W_u=W_s+(W_u-W_s)$ oraz niezależności
przyrostu od $\calF_s$. Prawa strona jest
rzeczywiście losowa dla $s>0$. Gdyby
$M_t-M_s$ był niezależny od $\calF_s$, jego
warunkowy drugi moment byłby stałą, co przeczy
otrzymanemu wzorowi. Ta obserwacja odróżnia
martyngałowość od własności niezależnych przyrostów.
\end{przyklad}

\section{Zatrzymywanie i całka stochastyczna}
Przy zatrzymywaniu strategii na losowym horyzoncie nie
wystarcza zwykłe pomnożenie przez wskaźnik zdarzenia
$\{\tau\le t\}$. Potrzebujemy procesu, który w chwili $u$
wskazuje, czy strategia jeszcze działa. Właśnie taki proces
jest przewidywalny.

\begin{twierdzenie}[Zatrzymanie całki Itô]\label{th:ito-stopping}
Niech $H\in\calH_T^2$ i niech $\tau$ będzie czasem
zatrzymania o wartościach w $[0,T]$. Wówczas proces
$H_u\ind{(0,\tau]}(u)$ jest przewidywalny i należy
do $\calH_T^2$, a dla wszystkich $t\in[0,T]$
zachodzi, poza jednym zbiorem zerowym,
\[
\int_0^{t\wedge\tau}H_u\,\dd W_u
=\int_0^t H_u\ind{\{u\le\tau\}}\,\dd W_u.
\tag{2.11}\label{eq:stopped-ito}
\]
\end{twierdzenie}
\begin{proof}
Proces $u\mapsto\ind{\{u\le\tau\}}$ jest adaptowany
i lewostronnie ciągły dla każdej trajektorii na $(0,T]$,
więc jest przewidywalny. Iloczyn z przewidywalnym $H$
też jest przewidywalny, a jego norma w $\calH_T^2$
nie przekracza normy $H$.

Najpierw niech $H$ będzie procesem prostym. Ustalmy
ciąg czasów $\tau_n:=2^{-n}\lceil2^n\tau\rceil\wedge T$,
gdzie w razie potrzeby ostatni przedział siatki skracamy
do $T$. Każdy $\tau_n$ jest czasem zatrzymania o skończenie
wielu wartościach, a $\tau_n\downarrow\tau$.
Po wspólnym rozdrobnieniu siatki procesu $H$ i siatki
czasu $\tau_n$ można zapisać proces
$H\ind{(0,\tau_n]}$ jako prosty przewidywalny:
współczynnik na odcinku $(u,v]$ zawiera wskaźnik
$\{\tau_n\ge v\}=\{\tau_n>u\}$, który jest
$\calF_u$-mierzalny. Z definicji sum wynika dokładnie
$I_t(H\ind{(0,\tau_n]})=I_{t\wedge\tau_n}(H)$.
Ponieważ $H$ jest kwadratowo całkowalny,
\[
\E\int_0^T H_u^2
|\ind{\{u\le\tau_n\}}-\ind{\{u\le\tau\}}|^2\,\dd u
\longrightarrow0
\]
na mocy zbieżności zdominowanej względem
$\dd t\otimes\dd\Prob$; pojedyncza chwila $u=\tau$
nie wpływa na całkę po czasie. Izometria Itô pozwala
przejść do granicy po prawej, a ciągłość trajektorii
$I(H)$ po lewej. Dostajemy tezę dla prostego $H$.

W przypadku ogólnym wybierzmy proste $H^m\to H$
w $\calH_T^2$. Po zatrzymaniu zachodzi ta sama
zbieżność w $\calH_T^2$, gdyż mnożenie przez
$\ind{(0,\tau]}$ jest kontrakcją.
Nierówność maksymalna \eqref{eq:ito-max} daje zbieżność
całych procesów $I(H^m)\to I(H)$ oraz
$I(H^m\ind{(0,\tau]})\to I(H\ind{(0,\tau]})$
w normie $\mathbb S^2$. Tożsamość dla $H^m$
przechodzi zatem do granicy. Równość w chwilach wymiernych
i ciągłość procesów zapewniają wspólny zbiór pełnego
prawdopodobieństwa, na którym zachodzi dla wszystkich $t$.
\end{proof}

\begin{zadanie}[Zatrzymany ruch Browna]
Niech $a>0$, a proces Wienera będzie określony na
$[0,\infty)$. Niech $\tau_a=\inf\{t\ge0:W_t=a\}$.
Dla ustalonego $T>0$
wyznaczyć $\int_0^T\ind{\{t\le\tau_a\}}\,\dd W_t$
i jego wartość oczekiwaną.

\emph{Rozwiązanie.} Z twierdzenia \ref{th:ito-stopping}
i równości $\int_0^t1\,\dd W_s=W_t$ otrzymujemy
$\int_0^T\ind{\{t\le\tau_a\}}\,\dd W_t=W_{T\wedge\tau_a}$.
Ponieważ czas $T\wedge\tau_a$ jest ograniczony,
próbkowanie opcjonalne lub bezpośrednia martyngałowość
całki daje $\E W_{T\wedge\tau_a}=0$.
Nie wolno bez dodatkowych warunków przechodzić tu do
$T\to\infty$ pod wartością oczekiwaną:
$W_{\tau_a}=a$ prawie na pewno, ale $\E W_{\tau_a}=a$,
a nie zero. Przy granicy traci się odpowiednią
jednostajną całkowalność.
\end{zadanie}

\section{Lokalna całkowalność: wyjście poza przestrzeń kwadratowalną}
Warunek $\E\int_0^T H_t^2\,\dd t<\infty$ jest
wygodny do konstrukcji, ale w wielu zastosowaniach
okazuje się zbyt mocny. Proces może mieć dobrze
zdefiniowaną całkę mimo nieskończonej wartości oczekiwanej
energii, jeśli energia każdej pojedynczej trajektorii
jest skończona.

\begin{definicja}[Całkowalność lokalna względem procesu Wienera]
Mówimy, że przewidywalny proces $H$ należy do
$\calH_{\mathrm{loc},T}^2$, gdy
$\int_0^T H_t^2\,\dd t<\infty$ prawie na pewno.
Nie zakładamy skończoności wartości oczekiwanej tej całki.
\end{definicja}

\begin{twierdzenie}[Konstrukcja całki Itô przez lokalizację]\label{th:ito-local}
Dla każdego $H\in\calH_{\mathrm{loc},T}^2$
istnieje jednoznaczny, z dokładnością do nierozróżnialności,
ciągły lokalny martyngał $I(H)$, taki że dla
$\tau_n:=\inf\{t\in[0,T]:\int_0^t H_s^2\,\dd s\ge n\}\wedge T$
zachodzi
\[
I_{t\wedge\tau_n}(H)
=\int_0^tH_s\ind{\{s\le\tau_n\}}\,\dd W_s,
\qquad 0\le t\le T.
\tag{2.12}
\]
Procesy po prawej są zdefiniowane przez
poprzednią konstrukcję w $L^2$, a $\tau_n\uparrow T$
prawie na pewno. Całka zachowuje liniowość i własność
zatrzymywania.
\end{twierdzenie}
\begin{proof}
Proces $A_t=\int_0^tH_s^2\,\dd s$ jest adaptowany,
ciągły i niemalejący, więc $\tau_n$ jest czasem
zatrzymania. Ciągłość $A$ zapewnia
$A_{\tau_n}\le n$, chyba że $\tau_n=T$ i $A_T<n$;
w obu przypadkach nierówność pozostaje prawdziwa.
Zatem
\[
\E\int_0^T H_s^2\ind{\{s\le\tau_n\}}\,\dd s
=\E A_{\tau_n}\le n.
\]
W konsekwencji każdy proces
$H^{(n)}:=H\ind{(0,\tau_n]}$ należy do $\calH_T^2$.
Dla $m\ge n$ mamy
$H^{(m)}\ind{(0,\tau_n]}=H^{(n)}$ niemal wszędzie.
Twierdzenie \ref{th:ito-stopping} daje więc
$I_{t\wedge\tau_n}(H^{(m)})=I_t(H^{(n)})$ dla wszystkich
$t$ poza wspólnym zbiorem zerowym.
Można zatem skleić ciągłe procesy $I(H^{(n)})$
w jedną ciągłą trajektorię: definiujemy ją w chwili $t$
przez $I_t(H^{(n)})$, jeśli $t\le\tau_n$, co jest
niezależne od dostatecznie dużego $n$.
Ponieważ $A_T<\infty$ prawie na pewno, dla każdego
$\omega$ poza zbiorem zerowym przy dostatecznie dużym
$n$ mamy $\tau_n(\omega)=T$; sklejanie obejmuje więc
cały przedział. Otrzymany proces jest lokalnym martyngałem,
bo każdy zatrzymany proces $I_{\cdot\wedge\tau_n}(H)$
jest martyngałem kwadratowo całkowalnym.
Jednoznaczność wynika z jednoznaczności wszystkich
zatrzymanych procesów. Liniowość otrzymujemy,
lokalizując równocześnie dwa integrandy przez
wspólną sekwencję czasów opartą na sumie ich energii.
\end{proof}

\begin{uwaga}[Kres własności całki]
Dla lokalnie kwadratowo całkowalnego procesu $H$ całka
$I_t(H)$ \emph{nie musi} należeć do $L^2(\Omega)$ ani
być prawdziwym martyngałem. Jej podstawową własnością
jest lokalna martyngałowość. Nie wolno więc bezwarunkowo
stosować wzoru $\E I_t(H)=0$ ani izometrii w postaci
z nieskończoną prawą stroną. Poprawny jest zawsze
wzór lokalny, po zatrzymaniu na $\tau_n$.
\end{uwaga}

\begin{przyklad}[Proces lokalnie całkowalny, lecz nie w $\calH_T^2$]
Niech $Z$ będzie zmienną losową niezależną od $W$,
$Z<\infty$ prawie na pewno, ale $\E Z^2=\infty$.
Rozszerzmy filtrację przez ujawnienie $Z$ w chwili $0$
i zdefiniujmy $H_t=Z$. Proces jest przewidywalny i
$\int_0^TH_t^2\,\dd t=TZ^2<\infty$ prawie na pewno,
lecz jego wartość oczekiwana jest nieskończona.
Całka lokalna ma postać $I_t(H)=ZW_t$.
Proces zatrzymujemy po przekroczeniu poziomu energii;
w tej prostej sytuacji można też bezpośrednio lokalizować
przez zdarzenia $\{|Z|\le n\}\in\calF_0$.
\end{przyklad}

\section{Obliczanie całek: trzy szczegółowe przykłady}
\begin{zadanie}[Podstawowa całka $\int W\,\dd W$]\label{zad:W-dW}
Obliczyć $\int_0^T W_t\,\dd W_t$ bez używania wzoru Itô.

\emph{Rozwiązanie.}
Proces $H_t=W_t$ jest adaptowany i ciągły,
a więc przewidywalny. Ponadto
$\E\int_0^TW_t^2\,\dd t=\int_0^Tt\,\dd t=T^2/2<\infty$.
Na równomiernej siatce $t_k=kT/n$ rozważmy proces
$H_t^{(n)}=W_{t_k}$ dla $t\in(t_k,t_{k+1}]$.
Z niezależności przyrostów
\[
\E\int_0^T|H_t^{(n)}-W_t|^2\,\dd t
=\sum_k\int_{t_k}^{t_{k+1}}(t-t_k)\,\dd t
=\frac{T^2}{2n}\longrightarrow0.
\]
Izometria Itô daje więc zbieżność sum
$\sum_kW_{t_k}\Delta_kW$ do szukanej całki w $L^2$.
Z tożsamości algebraicznej
\[
2\sum_kW_{t_k}\Delta_kW
=W_T^2-\sum_k(\Delta_kW)^2
\]
oraz zbieżności wariacji kwadratowej z tomu II
$\sum_k(\Delta_kW)^2\to T$ w $L^2$ otrzymujemy
\[
\int_0^TW_t\,\dd W_t=\frac12(W_T^2-T)
\quad\text{prawie na pewno}.
\tag{2.13}\label{eq:W-dW}
\]
Warto zauważyć, że średnia prawej strony wynosi zero,
a jej wariancja jest równa
$\frac14\Var(W_T^2)=T^2/2$, zgodnie z izometrią.
\end{zadanie}

\begin{zadanie}[Całka $\int W^2\,\dd W$ bez wzoru Itô]
Obliczyć $\int_0^T W_t^2\,\dd W_t$ bez stosowania
rachunku Itô do funkcji $x\mapsto x^3$.

\emph{Rozwiązanie.} Proces $W_t^2$ jest przewidywalny,
a $\E\int_0^T W_t^4\,\dd t
=3\int_0^T t^2\,\dd t=T^3<\infty$.
Dla siatki równomiernej $t_k=k\Delta$, $\Delta=T/n$,
przyrostów $\Delta_kW=W_{t_{k+1}}-W_{t_k}$ mamy
dokładną tożsamość algebraiczną
\[
W_T^3
=3\sum_kW_{t_k}^2\Delta_kW
+3\sum_kW_{t_k}(\Delta_kW)^2
+\sum_k(\Delta_kW)^3.
\tag{2.14}\label{eq:ito-cube-expand}
\]
Najpierw udowodnimy zbieżność pierwszej sumy do całki Itô.
Na przedziale $[t_k,t_{k+1}]$, pisząc
$W_t=W_{t_k}+\delta_t$ z przyrostem niezależnym od
$\calF_{t_k}$, otrzymujemy
$\E(W_t^2-W_{t_k}^2)^2
=4t_k(t-t_k)+3(t-t_k)^2$.
Wobec tego
\[
\E\int_0^T|W_t^2-W_{t_k(t)}^2|^2\,\dd t
\le 4T^2\Delta+3T\Delta^2\longrightarrow0,
\]
gdzie $t_k(t)$ oznacza lewy koniec przedziału
zawierającego $t$. Izometria daje pożądaną zbieżność
pierwszej sumy w $L^2$.

Drugi składnik rozdzielamy na
\[
\sum_kW_{t_k}(\Delta_kW)^2
=\sum_kW_{t_k}\Delta+R_n,\qquad
R_n:=\sum_kW_{t_k}[(\Delta_kW)^2-\Delta].
\]
Poszczególne składniki $R_n$ są różnicami
martyngałowymi, więc są ortogonalne w $L^2$. Ponieważ
$\E[(\Delta_kW)^2-\Delta]^2=2\Delta^2$,
\[
\E R_n^2
=2\Delta^2\sum_k\E W_{t_k}^2
\le2T^2\Delta\longrightarrow0.
\]
Suma Riemanna $\sum_kW_{t_k}\Delta$ dąży do
$\int_0^TW_t\,\dd t$ w $L^2$ (na przykład
z nierówności Cauchy'ego--Schwarza oraz ciągłości
średniokwadratowej $W$). Wreszcie składniki
$(\Delta_kW)^3$ są niezależne, mają średnią zero
i szóste momenty $15\Delta^3$. Stąd
$\E|\sum_k(\Delta_kW)^3|^2=15n\Delta^3=15T\Delta^2\to0$.
Przechodząc w \eqref{eq:ito-cube-expand} do granicy w $L^2$
otrzymujemy
\[
\int_0^TW_t^2\,\dd W_t
=\frac13W_T^3-\int_0^TW_t\,\dd t.
\tag{2.15}\label{eq:W2-dW}
\]
Ta tożsamość jest zapowiedzią wzoru Itô dla
funkcji $x^3$ i pokazuje, skąd bierze się dodatkowy
wyraz względem zwykłego rachunku różniczkowego.
\end{zadanie}

\begin{zadanie}[Strategia zależna od pierwszej obserwacji]
Załóżmy, że inwestor obserwuje $W_{T/2}$, a następnie
w drugiej połowie przedziału utrzymuje pozycję równą
$1$, jeśli $W_{T/2}>0$, i zero w przeciwnym przypadku.
Wyznaczyć prawo końcowego zysku
$X=\int_0^T\ind{\{W_{T/2}>0\}}\ind{(T/2,T]}(t)\,\dd W_t$.

\emph{Rozwiązanie.} Funkcja podcałkowa jest prostym
procesem przewidywalnym: decyzja jest
$\calF_{T/2}$-mierzalna. Z definicji
$X=\ind{\{W_{T/2}>0\}}(W_T-W_{T/2})$.
Przyrost $Y=W_T-W_{T/2}\sim\Normal(0,T/2)$ jest
niezależny od $W_{T/2}$. Z symetrii rozkładu normalnego
$\Prob(W_{T/2}>0)=1/2$. Zatem rozkład $X$ jest mieszaniną
$\frac12\delta_0+\frac12\Normal(0,T/2)$.
Nie jest rozkładem normalnym, ponieważ ma atom
w zerze o masie $1/2$. Jednocześnie
$\E X=0$ i $\Var(X)=(1/2)(T/2)=T/4$,
co potwierdza izometria Itô.
\end{zadanie}

\section{Przewidywalność nie jest kosmetycznym założeniem}
\begin{przyklad}[Integrand antycypujący]
Przypuśćmy błędnie, że moglibyśmy w całej konstrukcji
brać losowy współczynnik $\xi=W_T$ niezależnie od
czasu. Formalna suma dawałaby
$\sum_kW_T\Delta_kW=W_T^2$.
Jeśli nierozważnie użyć izometrii Itô, otrzymalibyśmy
$\E W_T^4\overset{?}=\E\int_0^TW_T^2\,\dd t=T^2$.
Lewa strona wynosi jednak $3T^2$.
Sprzeczność bierze się stąd, że $W_T$ nie jest
$\calF_{t_k}$-mierzalny dla $t_k<T$.
Nie oznacza to, że nigdy nie można całkować
antycypująco: istnieje na przykład całka Skorochoda,
lecz ma inną konstrukcję i inne własności.
\end{przyklad}

\begin{przyklad}[Dlaczego całka Itô nie musi mieć rozkładu normalnego]
Niech $Z\in\{0,1\}$ będzie zmienną Bernoulliego
z parametrem $1/2$, niezależną od procesu Wienera,
i niech $Z$ będzie znana w chwili $0$.
Wtedy $H_t=Z$ jest przewidywalny, zaś
$\int_0^TH_t\,\dd W_t=ZW_T$.
Prawo tej całki jest mieszaniną
$\frac12\delta_0+\frac12\Normal(0,T)$.
Choć każdy składnik sumy w definicji całki ma
normalny przyrost warunkowo względem przeszłości,
całka z \emph{losowym} integrandem nie musi być
normalna. Teza o gaussowskości całki z rozdziału
pierwszego obowiązywała dla funkcji deterministycznych.
\end{przyklad}


% Uzupełnienie audytowe rozdzialu 2: osobne uzasadnienia dwoch przejsc granicznych.
\section{Dwa kryteria przybliżania całek i ich ograniczenia}
\begin{lemat}[Lewostronna aproksymacja procesu ciągłego]
\label{lem:step-approx-ito}
Niech $H$ będzie ciągłym, adaptowanym procesem na
$[0,T]$, a $\sigma$ przewidywalnym procesem spełniającym
$\int_0^T\sigma_s^2\,\dd s<\infty$ prawie na pewno.
Niech $\pi_n$ będą deterministycznymi podziałami
$[0,T]$ o średnicy dążącej do zera, a
$H_s^{(n)}=H_{t_k}$ dla $s\in(t_k,t_{k+1}]$.
Wówczas, w sensie całek zlokalizowanych,
\[
\sup_{t\le T}\left|
\int_0^t(H_s^{(n)}-H_s)\sigma_s\,\dd W_s\right|
\longrightarrow0\quad\text{według prawdopodobieństwa}.
\]
W szczególności sumy lewostronne
$\sum_kH_{t_k}(W_{t_{k+1}\wedge t}-W_{t_k\wedge t})$
zbiegają do $\int_0^tH_s\,\dd W_s$ jednostajnie
według prawdopodobieństwa.
\end{lemat}
\begin{proof}
Ponieważ każda trajektoria $H$ jest jednostajnie
ciągła na $[0,T]$, mamy
$\norm{H^{(n)}-H}_\infty\to0$ prawie na pewno.
Wprowadźmy czas zatrzymania
$\tau_m=\inf\{t:\int_0^t\sigma_s^2\,\dd s\ge m\}
\wedge T$ oraz $\rho_m=\inf\{t:|H_t|\ge m\}\wedge T$.
Rozważmy $\theta_m=\tau_m\wedge\rho_m$.
Na $[0,\theta_m]$ mamy
$\int_0^{\theta_m}\sigma_s^2\,\dd s\le m$ oraz
$|H_s|\le m$; dla procesów $H^{(n)}$ po przekroczeniu
$\theta_m$ nie całkujemy, a ich wartości na
$[0,\theta_m]$ są ograniczone przez $m$.
Zatem na mocy izometrii i zbieżności zdominowanej
na przestrzeni $\Omega$,
\[
\E\int_0^{\theta_m}|H_s^{(n)}-H_s|^2
\sigma_s^2\,\dd s
\le \E\left[
\Big(\sup_{s\le\theta_m}|H_s^{(n)}-H_s|\Big)^2
\int_0^{\theta_m}\sigma_s^2\,\dd s\right]\longrightarrow0,
\]
ponieważ supremum na zatrzymanym przedziale jest
ograniczone przez $2m$ i zbiega do zera prawie na pewno.
Nierówność Dooba zapewnia
zbieżność zatrzymanych całek w $L^2$ normy supremum.
Ponieważ $\Prob(\theta_m<T)\to0$, usunięcie
lokalizacji daje tezę. Dla $\sigma\equiv1$
pierwsza równość wprost identyfikuje granicę
sum lewostronnych.
\end{proof}

\begin{zadanie}[Dlaczego zbieżność w $L^2$ nie daje zbieżności
ścieżkowej dla całego ciągu?]
Niech $(Y_n)$ będzie ciągiem niezależnych zdarzeń
z $\Prob(Y_n=1)=1/n$ i $\Prob(Y_n=0)=1-1/n$
($n\ge2$). Wyznaczyć rodzaje zbieżności $Y_n$ do zera.

\emph{Rozwiązanie.} Mamy $\E|Y_n|^2=1/n\to0$,
a zatem $Y_n\to0$ w $L^2$ i według prawdopodobieństwa.
Ponieważ $\sum_n\Prob(Y_n=1)=\infty$ i zdarzenia są
niezależne, drugi lemat Borela--Cantellego daje
$\Prob(Y_n=1\ \text{nieskończenie często})=1$.
Zatem nie ma zbieżności do zera prawie na pewno.
Wyjaśnia to, dlaczego w konstrukcji całki Itô potrzebny
był wybór odpowiedniego podciągu do uzyskania ciągłej
wersji procesu: z samej zbieżności w $L^2$ całego
ciągu nie wolno wyciągać wniosku o zbieżności
ścieżkowej całego ciągu.
\end{zadanie}


\section{Co udowodniliśmy, a co pozostaje do zrobienia?}
Całka Itô jest obecnie określona dla wszystkich
procesów przewidywalnych o skończonej oczekiwanej energii,
a przez lokalizację również dla procesów o energii
skończonej niemal na pewno. Udowodniliśmy izometrię,
martyngałowość, ciągłość, zatrzymywanie i identyfikację
wariacji kwadratowej. Na tej podstawie można już
formułować stochastyczne równania całkowe
$X_t=X_0+\int_0^tb_s\,\dd s+\int_0^t\sigma_s\,\dd W_s$.
Nie udowodniliśmy jeszcze ogólnego wzoru Itô,
nie skonstruowaliśmy całki Stratonovicha i nie badaliśmy
istnienia oraz jednoznaczności rozwiązań SDE.
Te zagadnienia będą przedmiotem kolejnych rozdziałów.

\begin{uwaga}[Konwencje na dalszą część tomu]
Przez $\int_0^tH_s\,\dd W_s$ będziemy rozumieli
całkę Itô, o ile wyraźnie nie zaznaczymy innej
konwencji. Skrót $\dd X_t=b_t\,\dd t+\sigma_t\,\dd W_t$
zawsze będziemy interpretować przez odpowiednie
równanie całkowe. Całki po trajektoriach z miarą
Wienera i całki względem pojedynczego procesu
pozostają odrębnymi konstrukcjami, nawet gdy
w szczególnych przykładach łączy je wzór Feynmana--Kaca.
\end{uwaga}
% ============================================================================
% ROZDZIAL 3 — WZOR ITO
% ============================================================================
\chapter{Wzór Itô: rachunek różniczkowy z wariacją kwadratową}
\label{ch:ito-formula}

\section{Motywacja: dlaczego zwykła reguła łańcuchowa zawodzi?}
Niech $W$ będzie procesem Wienera względem filtracji ustalonej
we wprowadzeniu. Gdybyśmy mogli bez zmian stosować klasyczną
regułę łańcuchową, uzyskalibyśmy
$W_t^2=2\int_0^tW_s\,\dd W_s$. Tymczasem w rozdziale
\ref{ch:ito-construction} wyprowadziliśmy ścisłą tożsamość
$W_t^2=2\int_0^tW_s\,\dd W_s+t$.
Dodatkowy składnik nie jest błędem przybliżenia: wynika z tego,
że suma kwadratów przyrostów $W$ nie znika w granicy, lecz
zbiega do długości przedziału.

W rachunku deterministycznym, przy różniczkowaniu $F(x)$,
wyraz drugiego rzędu jest rzędu $(\Delta x)^2$ i po
zsumowaniu zanika dla trajektorii o skończonej wariacji.
Dla ruchu Browna typowy przyrost na odcinku długości
$\Delta t$ jest rzędu $\sqrt{\Delta t}$, a więc jego
kwadrat jest rzędu $\Delta t$. Suma takich wyrazów nie
zanika. Celem rozdziału jest zastąpienie tej heurystyki
ścisłym twierdzeniem.

\section{Proces Itô i jego wariacja kwadratowa}
\paragraph{Konwencja zbieżności.}
Dla procesów $Y^{(n)}$ i $Y$ na $[0,T]$ piszemy
$Y^{(n)}\to Y$ \emph{ucp} (jednostajnie na przedziałach
zwartych według prawdopodobieństwa), jeżeli dla każdego
$\varepsilon>0$ zachodzi
$\Prob(\sup_{0\le t\le T}|Y_t^{(n)}-Y_t|>\varepsilon)
\to0$; gdy czas przebiega $[0,\infty)$, wymaga się
tej własności dla każdego skończonego $T$.

\begin{definicja}[Proces Itô]
Niech $T>0$, $W$ będzie rzeczywistym procesem Wienera,
$X_0$ będzie rzeczywistą zmienną $\calF_0$-mierzalną,
$b$ będzie procesem progresywnie mierzalnym, a $\sigma$
procesem przewidywalnym. Zakładamy, że prawie na pewno
$\int_0^T|b_s|\,\dd s<\infty$ oraz
$\int_0^T\sigma_s^2\,\dd s<\infty$.
Proces
\[
X_t=X_0+\int_0^tb_s\,\dd s+\int_0^t\sigma_s\,\dd W_s,
\qquad 0\le t\le T,
\tag{3.1}\label{eq:ito-process}
\]
nazywamy \emph{procesem Itô}. Druga całka jest całką
skonstruowaną w rozdziale \ref{ch:ito-construction}
(przez lokalizację, jeśli to konieczne).
\end{definicja}

\begin{uwaga}[Dlaczego nie zakładamy skończonych wartości oczekiwanych?]
Do dowodu lokalnego wzoru Itô wystarcza skończoność całek
z poprzedniej definicji dla niemal każdej trajektorii.
Zatrzymując proces, możemy na każdym ograniczonym przedziale
kontrolować zarówno $|X|$, jak i energię $\int\sigma^2$,
zmienność $\int|b|$ oraz pochodne funkcji, do której
stosujemy wzór. Dopiero po wykonaniu rachunku zdejmujemy
zatrzymanie. Nie wolno zatem bez uzasadnienia zakładać, że
całki występujące we wzorze mają skończone wartości oczekiwane.
\end{uwaga}

\begin{lemat}[Wariacja kwadratowa procesu Itô]
\label{lem:ito-qv}
Niech $X$ będzie procesem \eqref{eq:ito-process}. Dla
każdego ciągu deterministycznych podziałów $\pi_n$ przedziału
$[0,T]$ o średnicy dążącej do zera sumy
$\sum_{[u,v]\in\pi_n,\,v\le t}(X_v-X_u)^2$
zbiegają do $\int_0^t\sigma_s^2\,\dd s$ jednostajnie
w prawdopodobieństwie na $[0,T]$ (w skrócie: ucp).
W szczególności $[X]_t=\int_0^t\sigma_s^2\,\dd s$.
\end{lemat}
\begin{proof}
Piszemy $X=X_0+A+M$, gdzie
$A_t=\int_0^tb_s\,\dd s$ jest ciągłym procesem
skończonej wariacji, a $M_t=\int_0^t\sigma_s\,\dd W_s$
ciągłym lokalnym martyngałem. Na dowolnym podziale mamy
\[
\sum(\Delta X)^2-\sum(\Delta M)^2
=2\sum\Delta A\,\Delta M+\sum(\Delta A)^2.
\]
Dla ustalonej trajektorii ciągłość $A$ zapewnia
$\max|\Delta A|\to0$, a stąd
\[
\sum|\Delta A|^2
\le (\max|\Delta A|)\operatorname{TV}_{[0,T]}(A)
\longrightarrow0.
\]
Z nierówności Cauchy'ego--Schwarza otrzymujemy
\[
\left|\sum\Delta A\,\Delta M\right|
\le\Big(\sum(\Delta A)^2\Big)^{1/2}
\Big(\sum(\Delta M)^2\Big)^{1/2}.
\]
Drugi czynnik jest ograniczony w prawdopodobieństwie,
ponieważ sumy kwadratów przyrostów ciągłego lokalnego
martyngału zbiegają ucp do $[M]_T$; wynik ten udowodniono
w tomie II. Oszacowania stosują się również do sum
uciętych w dowolnej chwili $t$, z supremum po $t$.
Otrzymujemy więc zbieżność ucp sum dla $X$ i $M$ do
tej samej granicy. Z rozdziału \ref{ch:ito-construction}
wiemy, że $[M]_t=\int_0^t\sigma_s^2\,\dd s$.
\end{proof}

\begin{lemat}[Ważone sumy kwadratów przyrostów]
\label{lem:weighted-qv}
Niech $X$ będzie ciągłym procesem o wariacji kwadratowej
$[X]$ ciągłej i niemalejącej, przy czym sumy kwadratów
przyrostów na deterministycznych podziałach o średnicy
dążącej do zera zbiegają do $[X]$ ucp. Niech $G$ będzie
ciągłym procesem o trajektoriach ograniczonych na $[0,T]$.
Wówczas dla takich podziałów
\[
\sum_{[u,v]\in\pi_n,\,v\le t}G_u(X_v-X_u)^2
\ \xrightarrow[\,n\to\infty\,]{\mathrm{ucp}}\
\int_0^tG_s\,\dd[X]_s.
\tag{3.2}\label{eq:weighted-qv}
\]
Po prawej stronie jest zwykła całka Lebesgue'a--Stieltjesa
względem niemalejącej, ciągłej trajektorii $[X]$.
\end{lemat}
\begin{proof}
Ustalmy najpierw deterministyczny podział pomocniczy
$0=r_0<\cdots<r_m=T$ i zastąpmy $G$ procesem
$G^{(m)}_s=G_{r_j}$ dla $s\in(r_j,r_{j+1}]$.
Dla takiej funkcji wynik wynika z ucp-zbieżności
sum kwadratów na kolejnych podprzedziałach i z
ciągłości $[X]$. Jeżeli siatka $\pi_n$ nie zawiera
wszystkich $r_j$, dołączamy je: wpływ skończenie wielu
rozdzielonych przyrostów jest zaniedbywalny, gdyż
$\sup|\Delta X|\to0$, a całkowite sumy kwadratów są
ograniczone w prawdopodobieństwie.

Pozostaje oszacować błąd zastąpienia $G$ przez $G^{(m)}$.
Jego moduł w supremum po $t$ nie przekracza
$\norm{G-G^{(m)}}_\infty\sum_{[u,v]\in\pi_n}(X_v-X_u)^2$,
a dla całki granicznej jest ograniczony przez
$\norm{G-G^{(m)}}_\infty[X]_T$. Z ciągłości trajektorii
$G$ otrzymujemy $\norm{G-G^{(m)}}_\infty\to0$ prawie
na pewno przy zagęszczaniu pomocniczej siatki. Ponadto
$[X]_T$ i sumy kwadratów są ograniczone w prawdopodobieństwie.
Najpierw wybieramy $m$ dostatecznie duże, a następnie $n$;
daje to \eqref{eq:weighted-qv}.
\end{proof}

\section{Jednowymiarowy wzór Itô}
\begin{twierdzenie}[Wzór Itô dla funkcji klasy $C^{1,2}$]
\label{th:ito1}
Niech $X$ będzie procesem \eqref{eq:ito-process}, a
$F:[0,T]\times\R\to\R$ funkcją klasy $C^{1,2}$,
tzn. $F$, $\partial_tF$, $\partial_xF$ i
$\partial_{xx}F$ są ciągłe na $[0,T]\times\R$
(jednostronnie na końcach przedziału czasu). Wtedy
z prawdopodobieństwem jeden, jednocześnie dla wszystkich
$t\in[0,T]$, zachodzi
\[
\begin{split}
F(t,X_t)=F(0,X_0)
&+\int_0^t\partial_tF(s,X_s)\,\dd s
+\int_0^t\partial_xF(s,X_s)b_s\,\dd s\\
&+\int_0^t\partial_xF(s,X_s)\sigma_s\,\dd W_s
+\frac12\int_0^t\partial_{xx}F(s,X_s)\sigma_s^2\,\dd s.
\end{split}
\tag{3.3}\label{eq:ito1}
\]
Stochastyczna całka po prawej stronie jest rozumiana lokalnie.
\end{twierdzenie}
\begin{proof}
\emph{Krok 1: lokalizacja.} Dla liczb całkowitych $m\ge1$
zatrzymajmy proces przy wyjściu $X$ z $[-m,m]$ lub
przekroczeniu przez $\int_0^t(|b_s|+\sigma_s^2)\,\dd s$
poziomu $m$. Ciągłość trajektorii i skończoność całek
zapewniają, że czasy zatrzymania rosną do $T$ prawie na pewno.
Na zbiorze zatrzymanych trajektorii $X$ pozostaje w
przedziale zwartym, a $\partial_tF$, $\partial_xF$,
$\partial_{xx}F$ są tam ograniczone i jednostajnie ciągłe.
Wystarczy wykazać tożsamość do każdego z czasów zatrzymania
i następnie usunąć lokalizację. Możemy więc w dalszej
części dowodu pracować z ograniczonym zakresem $X$
i ze skończonymi, kontrolowanymi energiami.

\emph{Krok 2: rozwinięcie na podziale.} Weźmy podział
$\pi_n=\{0=t_0<\dots<t_{N_n}=t\}$ o średnicy
dążącej do zera. Oznaczmy $\Delta_kX=X_{t_{k+1}}-X_{t_k}$
i $\Delta t_k=t_{k+1}-t_k$. Rozbijamy przyrost funkcji na zmianę wartości przestrzennej
w chwili $t_k$ oraz zmianę czasu przy 
ustalonej \emph{nowej} wartości przestrzennej:
\[
F(t_{k+1},X_{t_{k+1}})-F(t_k,X_{t_k})
=\bigl[F(t_k,X_{t_{k+1}})-F(t_k,X_{t_k})\bigr]
 +\bigl[F(t_{k+1},X_{t_{k+1}})-F(t_k,X_{t_{k+1}})\bigr].
\]
Jednostajna ciągłość pochodnych na odpowiednim zwartym
prostokącie i wzór Taylora dają
\[
\begin{split}
F(t_{k+1},X_{t_{k+1}})-F(t_k,X_{t_k})
={}&\partial_tF(t_k,X_{t_k})\Delta t_k
+\partial_xF(t_k,X_{t_k})\Delta_kX\\
&+\frac12\partial_{xx}F(t_k,X_{t_k})(\Delta_kX)^2+r_{k,n},
\end{split}
\]
gdzie suma reszt $\sum_k r_{k,n}$ dąży do zera
w prawdopodobieństwie. Uzasadnijmy to bez ukrytego
założenia o istnieniu pochodnej mieszanej $F_{tx}$.
Niech $\omega_t(\delta)$ i $\omega_{xx}(\delta)$
będą modułami ciągłości, odpowiednio, $F_t$ i $F_{xx}$
na ustalonym zwartym prostokącie zawierającym
zatrzymaną trajektorię. Z twierdzenia o wartości
średniej w zmiennej $t$ i wzoru Taylora w zmiennej $x$
wynika
\[
\sum_k|r_{k,n}|\le
T\,\omega_t\!\left(|\pi_n|+\max_k|\Delta_kX|\right)
+\tfrac12\omega_{xx}\!\left(\max_k|\Delta_kX|\right)
\sum_k(\Delta_kX)^2.
\]
Rzeczywiście, w składniku czasowym odejmujemy
$F_t(t_k,X_{t_k})\Delta t_k$ od
$\int_{t_k}^{t_{k+1}}F_t(s,X_{t_{k+1}})\,\dd s$;
różnica jest kontrolowana przez pierwszy moduł ciągłości.
W składniku przestrzennym stosujemy jednowymiarowy
wzór Taylora z resztą kontrolowaną przez drugi moduł.
Z ciągłości $X$ mamy $\max_k|\Delta_kX|\to0$
prawie na pewno, a z lematu \ref{lem:ito-qv} sumy
kwadratów przyrostów są ograniczone w prawdopodobieństwie.
Oba składniki oszacowania dążą więc do zera w
prawdopodobieństwie. To rozumowanie wymaga dokładnie
klasy $C^{1,2}$, a nie silniejszej regularności.

\emph{Krok 3: identyfikacja granic.} Sumy
$\sum_k\partial_tF(t_k,X_{t_k})\Delta t_k$
zbiegają do $\int_0^t\partial_tF(s,X_s)\,\dd s$
ścieżkowo, ponieważ integrand jest ciągły.
Dla części liniowej piszemy
$\Delta_kX=\int_{t_k}^{t_{k+1}}b_s\,\dd s
+\int_{t_k}^{t_{k+1}}\sigma_s\,\dd W_s$.
Funkcje schodkowe
$H_s^{(n)}=\partial_xF(t_k,X_{t_k})$ dla
$s\in(t_k,t_{k+1}]$ są przewidywalne. Ich różnica
od ciągłego, adaptowanego procesu
$H_s=\partial_xF(s,X_s)$ dąży jednostajnie do zera
na każdej zatrzymanej trajektorii. Zatem
$\int H^{(n)}b\,\dd s\to\int Hb\,\dd s$
ścieżkowo, a po odpowiedniej lokalizacji
$\E\int|H^{(n)}-H|^2\sigma^2\,\dd s\to0$ przez
zbieżność zdominowaną. Izometria Itô wraz z
nierównością maksymalną Dooba daje zbieżność
odpowiednich całek stochastycznych.

Dla wyrazów kwadratowych stosujemy lemat
\ref{lem:weighted-qv} z
$G_s=\partial_{xx}F(s,X_s)$, otrzymując
\[
\sum_kG_{t_k}(\Delta_kX)^2
\longrightarrow\int_0^tG_s\sigma_s^2\,\dd s
\]
w prawdopodobieństwie. Sumując wszystkie granice,
otrzymujemy \eqref{eq:ito1} dla ustalonego $t$.
Obie strony mają ciągłe trajektorie. Równość w chwilach
wymiernych (oraz w $T$) daje ją jednocześnie dla wszystkich
$t$, poza jednym zbiorem zerowym. Na końcu usuwamy
lokalizację, pozwalając poziomom zatrzymania rosnąć do nieskończoności.
\end{proof}

\begin{uwaga}[Znaczenie wzoru różniczkowego]
Często zapisujemy \eqref{eq:ito1} symbolicznie jako
$\dd F(t,X_t)=F_t\,\dd t+F_x\,\dd X_t
+\frac12F_{xx}\,\dd[X]_t$. Nie jest to równość
między zwykłymi różniczkami: jej ścisłym znaczeniem
jest tożsamość całkowa \eqref{eq:ito1}.
\end{uwaga}

\section{Przykłady: kwadrat, potęgi, wykładnik i martyngały}
\begin{zadanie}[Potęgi procesu Wienera]
Dla $n\in\N$, $n\ge2$, obliczyć rozkład na składnik
martyngałowy i składnik skończonej wariacji procesu $W_t^n$.

\emph{Rozwiązanie.} Dla $F(x)=x^n$ mamy
$F_x(x)=nx^{n-1}$ oraz $F_{xx}(x)=n(n-1)x^{n-2}$.
Wzór Itô daje
\[
W_t^n=n\int_0^tW_s^{n-1}\,\dd W_s
+\frac{n(n-1)}2\int_0^tW_s^{n-2}\,\dd s.
\]
Przy $n=2$ dostajemy $W_t^2-t=2\int_0^tW_s\,\dd W_s$;
przy $n=3$ otrzymujemy $W_t^3-3\int_0^tW_s\,\dd s
=3\int_0^tW_s^2\,\dd W_s$. Nie wolno utożsamiać
składnika całkowego z martyngałem jednostajnie
całkowalnym na całym $[0,\infty)$ bez sprawdzenia
jego momentów. Na dowolnym skończonym horyzoncie
potrzebna całkowalność wynika z momentów Gaussa.
\end{zadanie}

\begin{twierdzenie}[Wykładnik stochastyczny dla stałego parametru]
Niech $\theta\in\R$ i $W$ będzie procesem Wienera na
$[0,T]$. Wówczas $Z_t=\exp(\theta W_t-\tfrac12\theta^2t)$
jest dodatnim martyngałem i
\[
Z_t=1+\theta\int_0^t Z_s\,\dd W_s.
\tag{3.4}\label{eq:exp-mg}
\]
\end{twierdzenie}
\begin{proof}
Stosując twierdzenie \ref{th:ito1} do
$F(t,x)=\exp(\theta x-\frac12\theta^2t)$, zauważamy,
że $F_t+\frac12F_{xx}=0$, więc otrzymujemy
\eqref{eq:exp-mg} (początkowo lokalnie).
Dla $0\le s<t$ niezależny od $\calF_s$ przyrost
$W_t-W_s$ ma rozkład normalny $\Normal(0,t-s)$.
Stąd
$\E[Z_t\mid\calF_s]=Z_s\E\exp(\theta(W_t-W_s)
-\frac12\theta^2(t-s))=Z_s$.
Wszystkie wartości oczekiwane są skończone, a
$\E Z_t=1$.
\end{proof}

\begin{zadanie}[Wykładnik czasu i równanie ciepła]
Dla $u(t,x)=x^2+(T-t)$, $0\le t\le T$,
wykazać, że $u(t,W_t)$ jest martyngałem i
wyjaśnić związek z równaniem ciepła.

\emph{Rozwiązanie.} Mamy $u_t=-1$, $u_{xx}=2$,
więc $u_t+\frac12u_{xx}=0$. Wzór Itô daje
$\dd u(t,W_t)=2W_t\,\dd W_t$ i tym samym
$u(t,W_t)=T+2\int_0^tW_s\,\dd W_s$.
Jest to martyngał kwadratowo całkowalny na $[0,T]$.
W ogólności, jeśli $u\in C^{1,2}$ spełnia
$u_t+\frac12u_{xx}=0$ oraz wystarczające warunki
całkowalności, to $u(t,W_t)$ jest martyngałem.
Pokazuje to probabilistyczną rolę równania ciepła;
pełny wzór Feynmana--Kaca pojawi się później.
\end{zadanie}

\section{Iloczyny, całkowanie przez części i kowariacja}
\begin{twierdzenie}[Wzór na iloczyn procesów Itô]
\label{th:ito-product}
Niech $X,Y$ będą procesami Itô, napędzanymi tym samym
$W$, o współczynnikach odpowiednio $(b,\sigma)$ oraz
$(c,\rho)$, spełniających warunki z definicji procesu Itô.
Wówczas z prawdopodobieństwem jeden, dla wszystkich $t$,
\[
X_tY_t=X_0Y_0+\int_0^tX_sc_s\,\dd s
+\int_0^tY_sb_s\,\dd s
+\int_0^t(X_s\rho_s+Y_s\sigma_s)\,\dd W_s
+\int_0^t\sigma_s\rho_s\,\dd s.
\tag{3.5}\label{eq:ito-product}
\]
W szczególności $[X,Y]_t=\int_0^t\sigma_s\rho_s\,\dd s$.
\end{twierdzenie}
\begin{proof}
Tożsamość $4xy=(x+y)^2-(x-y)^2$ pozwala wyprowadzić
wzór na iloczyn z jednowymiarowego wzoru Itô dla
$X+Y$ i $X-Y$: stosujemy \eqref{eq:ito1} do kwadratu,
odejmujemy odpowiednie równości i dzielimy przez cztery.
Daje to składniki liniowe oraz kowariację
$\int\sigma\rho\,\dd s$. Alternatywnie, z definicji
kowariacji i polaryzacji $[X,Y]=\frac14([X+Y]-[X-Y])$,
a następnie $[X+Y]-[X-Y]=4\int\sigma\rho\,\dd s$.
\end{proof}

\begin{zadanie}[Całkowanie przez części a momenty]
Niech $W$ będzie procesem Wienera. Obliczyć
$\int_0^t sW_s\,\dd W_s$.

\emph{Rozwiązanie.} Stosujemy wzór na iloczyn do
$X_s=s$ i $Y_s=W_s^2$, dla którego
$\dd Y_s=2W_s\,\dd W_s+\dd s$.
Ponieważ $[s,W^2]_t=0$, otrzymujemy
$\dd(sW_s^2)=W_s^2\,\dd s+s(2W_s\,\dd W_s+\dd s)$.
Zatem
\[
\int_0^t sW_s\,\dd W_s
=\frac12\left(tW_t^2-\int_0^t(W_s^2+s)\,\dd s\right).
\]
Sprawdzenie: wartość oczekiwana prawej strony wynosi
$\frac12(t^2-\int_0^t2s\,\dd s)=0$, zgodnie z
martyngałowością całki.
\end{zadanie}

\section{Wzór wielowymiarowy i kierunek dyfuzji}
\begin{lemat}[Całki względem różnych współrzędnych Wienera]
\label{lem:multi-orthogonality}
Niech $W=(W^1,\ldots,W^m)$ będzie standardowym
$m$-wymiarowym procesem Wienera względem filtracji
$(\calF_t)$, to znaczy wektor jego przyrostów na
$(s,t]$ jest niezależny od $\calF_s$ i ma macierz
kowariancji $(t-s)I_m$. Jeśli $H,K$ są rzeczywistymi
procesami przewidywalnymi, dla których
$\E\int_0^T(H_s^2+K_s^2)\,\dd s<\infty$, to
\[
\E\left[
\left(\int_0^tH_s\,\dd W_s^i\right)
\left(\int_0^tK_s\,\dd W_s^j\right)\right]
=\delta_{ij}\E\int_0^tH_sK_s\,\dd s.
\]
Ponadto kowariacja ciągłych martyngałów
$M^i=\int H\,\dd W^i$ i
$N^j=\int K\,\dd W^j$ wynosi
$[M^i,N^j]_t=\delta_{ij}\int_0^tH_sK_s\,\dd s$.
W wersji lokalnie całkowalnej ostatni wzór nadal
obowiązuje dzięki lokalizacji.
\end{lemat}
\begin{proof}
Dla prostych procesów przewidywalnych przechodzimy do
wspólnego deterministycznego podziału. Iloczyn dwóch
sum zawiera wyrazy z różnych przedziałów i z tego
samego przedziału. Pierwsze mają wartość oczekiwaną
zero dzięki iterowanemu warunkowaniu; dla drugich
\[
\E[\Delta_kW^i\Delta_kW^j\mid\calF_{t_k}]
=\delta_{ij}(t_{k+1}-t_k).
\]
Jest to dokładnie podany wzór kowariancji wartości
końcowych. Przy $i\ne j$ ten sam rachunek po
warunkowaniu pokazuje, że iloczyn $M^iN^j$ jest
martyngałem. Przy $i=j$ proces
$M^iN^i-\int HK\,\dd s$ jest martyngałem,
co wynika z analogicznego rozwinięcia przyrostów.
Jednoznaczność takiej korekty wynika z faktu, że
ciągły lokalny martyngał o skończonej wariacji jest
stały: różnica dwóch korekt byłaby jednocześnie
lokalnym martyngałem i procesem skończonej wariacji.
To identyfikuje kowariację.
Dla ogólnych integrandów stosujemy gęstość procesów
prostych, izometrię Itô i nierówność Cauchy'ego--Schwarza.
W wersji lokalnej najpierw obcinamy energię obu
integrandów, a potem usuwamy zatrzymanie.
\end{proof}

\begin{twierdzenie}[Wielowymiarowy wzór Itô]
\label{th:ito-multi}
Niech $W=(W^1,\ldots,W^m)$ będzie standardowym
$m$-wymiarowym procesem Wienera względem wspólnej
filtracji. Niech $X=(X^1,\ldots,X^d)$ będzie procesem
postaci
\[
X_t^i=X_0^i+\int_0^t b_s^i\,\dd s
+\sum_{k=1}^{m}\int_0^t\sigma_s^{ik}\,\dd W_s^k,
\quad i=1,\ldots,d,
\]
gdzie $X_0$ jest $\calF_0$-mierzalny, $b$ jest
progresywnie mierzalny, $\sigma$ przewidywalna i
$\int_0^T(\norm{b_s}+\norm{\sigma_s}_{\mathrm F}^2)\,\dd s
<\infty$ prawie na pewno. Jeśli
$F\in C^{1,2}([0,T]\times\R^d)$, to z prawdopodobieństwem
jeden, dla wszystkich $t\in[0,T]$,
\[
\begin{split}
F(t,X_t)=F(0,X_0)&+\int_0^t\left(
F_t+\sum_{i=1}^d b_s^iF_{x_i}
+\frac12\sum_{i,j=1}^d(\sigma_s\sigma_s^{\mathsf T})_{ij}
F_{x_ix_j}\right)(s,X_s)\,\dd s\\
&+\sum_{k=1}^m\int_0^t
\left(\sum_{i=1}^dF_{x_i}(s,X_s)\sigma_s^{ik}\right)
\,\dd W_s^k.
\end{split}
\tag{3.6}\label{eq:ito-multi}
\]
\end{twierdzenie}
\begin{proof}
Stosujemy wielowymiarowy wzór Taylora drugiego rzędu
na deterministycznym podziale do funkcji $F(t,X_t)$.
Kontrola reszt odbywa się tak jak w dowodzie twierdzenia
\ref{th:ito1}, po lokalizacji na zwartych obszarach.
Dla wszystkich $i,j$ mamy z polaryzacji wariacji
kwadratowej
\[
[X^i,X^j]_t
=\sum_{k=1}^m\int_0^t\sigma_s^{ik}\sigma_s^{jk}\,\dd s.
\]
Istotnie, niezależność współrzędnych $W$ daje
$[W^k,W^\ell]_t=\delta_{k\ell}t$, a izometria Itô,
lokalizacja i aproksymacja współczynników procesami
prostymi pozwalają wyprowadzić tę samą relację dla
całek stochastycznych. Zastosowanie ważonej wersji
lematu \ref{lem:weighted-qv} do sum
$\sum F_{x_ix_j}(t_k,X_{t_k})\Delta_kX^i\Delta_kX^j$
(dla $i\ne j$ przez polaryzację) daje dokładnie
składnik drugiego rzędu z macierzą $\sigma\sigma^{\mathsf T}$.
Część pierwszego rzędu zbiega do całek względem czasu
i procesów Wienera z izometrii i lokalizacji, jak w
dowodzie jednowymiarowym. Otrzymujemy tezę.
\end{proof}

\begin{przyklad}[Ruch Browna i funkcje harmoniczne]
Niech $B$ będzie standardowym $d$-wymiarowym ruchem
Browna, $d\ge2$, i $h\in C^2(\R^d)$ spełnia
$\Delta h=0$. Stosując \eqref{eq:ito-multi} dla
$X=B$, $b=0$, $\sigma=I$, otrzymujemy
$h(B_t)=h(B_0)+\int_0^t\nabla h(B_s)\cdot\dd B_s$.
Jest to lokalny martyngał; po zatrzymaniu przy wyjściu
z kuli staje się kwadratowo całkowalnym martyngałem,
gdyż $\nabla h$ jest ograniczony na odpowiednim
zwartym obszarze. Ten przykład tłumaczy związek
ruchu Browna z funkcjami harmonicznymi.
\end{przyklad}

\begin{zadanie}[Kwadrat odległości w $\R^d$]
Dla $d$-wymiarowego ruchu Browna $B$ wyprowadzić
równanie dla $R_t^2=\norm{B_t}^2$ i obliczyć jego
wartość oczekiwaną, jeśli $B_0=x$ jest deterministyczny.

\emph{Rozwiązanie.} Z \eqref{eq:ito-multi} dla
$F(x)=\sum_i x_i^2$ dostajemy
\[
\norm{B_t}^2=\norm{x}^2+2\int_0^t B_s\cdot\dd B_s+d\,t.
\]
Po wzięciu wartości oczekiwanych mamy
$\E\norm{B_t}^2=\norm{x}^2+dt$. Ostatnia równość
może też zostać otrzymana przez zsumowanie
$\E(B_t^i)^2=x_i^2+t$ dla poszczególnych współrzędnych.
\end{zadanie}

\section{Czas lokalny procesu Wienera: pierwszy kontakt}
Wzór Itô wymaga klasy $C^{1,2}$, tymczasem
funkcja $x\mapsto|x-a|$ nie jest różniczkowalna w $a$.
W przybliżeniu jej druga pochodna koncentruje się
w punkcie $a$. Wartość graniczna daje nowy obiekt:
\emph{czas lokalny}, który mierzy intensywność
przebywania procesu w otoczeniu danego poziomu.

\begin{twierdzenie}[Wzór Tanaki i czas lokalny w ustalonym poziomie]
\label{th:tanaka}
Niech $W$ będzie standardowym procesem Wienera, $a\in\R$.
Przyjmujemy $\operatorname{sgn}(x)=1$ dla $x>0$,
$-1$ dla $x<0$ oraz $0$ dla $x=0$. Istnieje ciągły,
niemalejący, nieujemny proces $(L_t^a)_{0\le t\le T}$,
zaczynający się w zerze, taki że prawie na pewno,
jednocześnie dla wszystkich $t\in[0,T]$,
\[
|W_t-a|=|W_0-a|+\int_0^t\operatorname{sgn}(W_s-a)
\,\dd W_s+L_t^a.
\tag{3.7}\label{eq:tanaka}
\]
Ponadto dla każdego $t$ zachodzi zbieżność w
prawdopodobieństwie
\[
\frac1{2\varepsilon}\int_0^t
\ind{\{|W_s-a|<\varepsilon\}}\,\dd s
\ \longrightarrow\ L_t^a
\qquad(\varepsilon\downarrow0).
\tag{3.8}\label{eq:local-time-occ}
\]
Proces $L^a$ nazywamy \emph{czasem lokalnym} ruchu Browna
na poziomie $a$, w przyjętej normalizacji.
\end{twierdzenie}
\begin{proof}
Zaczynamy od funkcji wypukłej
\[
f_\varepsilon(x)=
\begin{cases}
|x-a|,&|x-a|\ge\varepsilon,\\
\displaystyle\frac{(x-a)^2}{2\varepsilon}+
\frac\varepsilon2,&|x-a|<\varepsilon.
\end{cases}
\]
Jest ona klasy $C^1$, jej pochodna jest Lipschitzowska,
a druga pochodna istnieje poza dwoma punktami
$a\pm\varepsilon$ i wynosi
$f_\varepsilon''(x)=\varepsilon^{-1}
\ind{\{|x-a|<\varepsilon\}}$.
Aby legalnie stosować wzór Itô, przybliżamy
$f_\varepsilon$ splotami z gładkim jądrem
mollifikującym. Pochodne pierwszego rzędu zbiegają
jednostajnie, a drugie — prawie wszędzie i są
jednostajnie ograniczone (przy ustalonym
$\varepsilon$). Ponieważ rozkład $W_s$ ma gęstość
przy każdym $s>0$, a $\int_0^T\Prob(W_s=a\pm\varepsilon)
\,\dd s=0$, zbieżność zdominowana daje
zbieżność składników drugiego rzędu w $L^1(\Omega)$.
Izometria Itô daje zbieżność całek z pochodnych
pierwszego rzędu w $L^2(\Omega)$.
Przechodząc do granicy w gładkim wzorze Itô, uzyskujemy
\[
f_\varepsilon(W_t)=f_\varepsilon(W_0)
+\int_0^t f_\varepsilon'(W_s)\,\dd W_s
+\frac1{2\varepsilon}\int_0^t
\ind{\{|W_s-a|<\varepsilon\}}\,\dd s.
\tag{3.9}\label{eq:tanaka-eps}
\]
Teraz $|f_\varepsilon(x)-|x-a||\le\varepsilon/2$ dla
każdego $x$, zaś $f_\varepsilon'(x)\to
\operatorname{sgn}(x-a)$ poza punktem $a$, przy czym
$|f_\varepsilon'|\le1$. Ponieważ
$\E\int_0^T\ind{\{W_s=a\}}\,\dd s=0$,
z twierdzenia Fubiniego i zbieżności zdominowanej
mamy
$\E\int_0^T|f_\varepsilon'(W_s)-
\operatorname{sgn}(W_s-a)|^2\,\dd s\to0$.
Izometria i nierówność Dooba zapewniają zbieżność
całek stochastycznych jednostajnie w $L^2$.
Z \eqref{eq:tanaka-eps} wynika więc zbieżność ucp
nieujemnych i niemalejących procesów z lewej strony
\eqref{eq:local-time-occ} do ciągłego procesu
\[
L_t^a:=|W_t-a|-|W_0-a|
-\int_0^t\operatorname{sgn}(W_s-a)\,\dd W_s.
\]
Wybieramy podciąg zbieżny jednostajnie prawie na pewno;
jego granica jest niemalejąca i nieujemna. Dostajemy
\eqref{eq:tanaka} oraz \eqref{eq:local-time-occ}.
\end{proof}

\begin{zadanie}[Oczekiwany czas lokalny w zerze]
Dla procesu Wienera $W_0=0$ obliczyć $\E L_t^0$.

\emph{Rozwiązanie.} Biorąc wartości oczekiwane w
\eqref{eq:tanaka}, zauważamy, że całka
$\int_0^t\operatorname{sgn}(W_s)\,\dd W_s$ jest
kwadratowo całkowalnym martyngałem o średniej zero,
gdyż integrand jest ograniczony. Zatem
$\E L_t^0=\E|W_t|=\sqrt{2t/\pi}$.
Wartość jest dodatnia, mimo że $\Prob(W_s=0)=0$
dla każdej ustalonej chwili $s>0$: czas lokalny
mierzy nagromadzenie wizyt \emph{w dowolnie małych
otoczeniach} zera, a nie miarę Lebesgue'a zbioru
chwil, w których $W_s=0$.
\end{zadanie}

\begin{uwaga}[Czego jeszcze nie dowiedliśmy?]
W tym rozdziale skonstruowaliśmy $L_t^a$ dla każdego
\emph{ustalonego} $a$. Wspólnie ciągłą wersję w
parametrach $(t,a)$, pełny wzór okupacyjny i
analogiczne wyniki dla ogólnych półmartyngałów
przedstawimy po wprowadzeniu odpowiednich narzędzi.
Nie wolno uważać \eqref{eq:local-time-occ} za
jednostajną w $a$ zbieżność bez dodatkowego dowodu.
\end{uwaga}


% ============================================================================
% ROZDZIAL 4 — STRATONOVICH
% ============================================================================
\chapter{Całka Stratonovicha i jej związek z całką Itô}
\label{ch:stratonovich}

\section{Dwie reguły sumowania, dwa rachunki}
W klasycznym rachunku różniczkowym funkcja złożona
spełnia regułę łańcuchową bez dodatkowych wyrazów.
W rachunku Itô pojawia się poprawka zależna od wariacji
kwadratowej. Można jednak zdefiniować inną całkę,
która w przypadku ciągłych procesów Itô zachowuje
zwykłą regułę łańcuchową. Jest nią \emph{całka
Stratonovicha}.

Trzeba tu odróżnić dwa pytania. Pierwsze dotyczy
\emph{istnienia granicy} sum, w których współczynnik
wybieramy na środku przedziału. Drugie dotyczy tego,
\emph{co dokładnie rozumiemy przez środek}: w poniższej
konstrukcji użyjemy średniej arytmetycznej wartości
procesu podcałkowego na obu końcach przedziału.
Takie sumy są symetryczne, lecz w przeciwieństwie do
sum Itô korzystają z przyszłej wartości integrandu.
Nie jest to strategia przewidywalna w rozumieniu
rozdziału \ref{ch:ito-construction}.

\begin{przyklad}[Najprostsza różnica]
Dla procesu Wienera $W_0=0$ rozważmy
\[
S_\pi^{\mathrm L}=\sum_kW_{t_k}\Delta_kW,
\qquad
S_\pi^{\mathrm S}=
\frac12\sum_k(W_{t_k}+W_{t_{k+1}})\Delta_kW.
\]
Drugie wyrażenie jest dokładnie teleskopowe:
$S_\pi^{\mathrm S}=\frac12W_T^2$ dla każdego podziału.
Natomiast $S_\pi^{\mathrm L}\to\frac12(W_T^2-T)$
w $L^2$. Zatem odpowiednia poprawka wynosi $T/2$.
\emph{To nie wybór dowolnej estetycznej konwencji}:
obydwa rachunki odpowiadają różnym konstrukcjom
matematycznym i różnym sposobom aproksymowania hałasu.
\end{przyklad}

\begin{center}
\begin{tikzpicture}[>=Latex,x=1.0cm,y=.85cm]
\draw[->] (0,0)--(8.5,0) node[right] {$t$};
\draw[->] (0,-.35)--(0,2.5) node[above] {$W_t$};
\draw[thick,blue!70!black]
 (0,0)--(.6,.3)--(1.2,1.15)--(1.8,.9)--(2.4,1.6)
 --(3,1.0)--(3.6,.65)--(4.2,1.35)--(4.8,1.1)
 --(5.4,1.9)--(6,1.4)--(6.6,2.1)--(7.2,1.5)--(7.8,1.85);
\foreach \x in {0,1.8,3.6,5.4,7.8}
 {\draw[dashed,gray!70] (\x,0)--(\x,2.25);}
\fill[blue!65!black] (1.8,.9) circle (2pt);
\fill[blue!65!black] (3.6,.65) circle (2pt);
\draw[<->,orange!75!black,thick] (1.8,2.3)--(3.6,2.3)
 node[midway,above,scale=.85] {jeden przedział};
\node[scale=.81,anchor=north] at (1.8,-.12) {$t_k$};
\node[scale=.81,anchor=north] at (3.6,-.12) {$t_{k+1}$};
\node[scale=.78,align=left,anchor=west] at (5.0,2.42)
 {Itô: lewy koniec\\Stratonovich: średnia końców};
\end{tikzpicture}
\par\smallskip\small Porównanie reguł próbkowania na tym samym odcinku trajektorii.
\end{center}

\section{Kowariacja i zbieżność sum symetrycznych}
\begin{definicja}[Kowariacja ciągłych procesów Itô]
Niech $H$ i $X$ będą ciągłymi procesami Itô,
$H=H_0+\int_0^.\alpha_s\,\dd s+
\int_0^.\beta_s\,\dd W_s$ oraz
$X=X_0+\int_0^.b_s\,\dd s+
\int_0^.\sigma_s\,\dd W_s$,
z lokalnymi warunkami całkowalności jak w
rozdziale \ref{ch:ito-formula}. Definiujemy
\[
[H,X]_t:=\int_0^t\beta_s\sigma_s\,\dd s.
\tag{4.1}\label{eq:HX-bracket}
\]
Jest to dobrze określony ciągły proces skończonej
wariacji, ponieważ
$\int_0^T|\beta_s\sigma_s|\,\dd s
\le(\int_0^T\beta_s^2\,\dd s)^{1/2}
(\int_0^T\sigma_s^2\,\dd s)^{1/2}<\infty$
prawie na pewno.
\end{definicja}

\begin{twierdzenie}[Zbieżność symetrycznych sum]
\label{th:strat-sums}
Niech $X,H$ będą ciągłymi procesami Itô, jak w
poprzedniej definicji. Niech $\pi_n$ będą
\emph{deterministycznymi} podziałami $[0,T]$
o średnicy dążącej do zera. Dla $t\in[0,T]$
niech
\[
S_t^{(n)}:=\frac12\sum_k
(H_{t_k}+H_{t_{k+1}})
\bigl(X_{t_{k+1}\wedge t}-X_{t_k\wedge t}\bigr)
\ind{\{t_k<t\}},
\]
przy czym w składniku przeciętym przez $t$
zastępujemy $H_{t_{k+1}}$ wartością $H_t$.
Wtedy
\[
S_t^{(n)}\ \xrightarrow[n\to\infty]{\mathrm{ucp}}
\int_0^tH_s\,\dd X_s+\frac12[H,X]_t,
\tag{4.2}\label{eq:strat-sum}
\]
gdzie $\int H\,\dd X:=\int Hb\,\dd s+
\int H\sigma\,\dd W$ oznacza całkę Itô.
\end{twierdzenie}
\begin{proof}
Dla $t=T$ rozkładamy sumę na część lewostronną i
część poprawkową:
\[
S_T^{(n)}=\sum_kH_{t_k}\Delta_kX
+\frac12\sum_k\Delta_kH\,\Delta_kX.
\tag{4.3}\label{eq:strat-decomp}
\]
Dla składnika pierwszego definiujemy
$H_s^{(n)}=H_{t_k}$ na $(t_k,t_{k+1}]$.
Jest to proces przewidywalny. Z ciągłości $H$ mamy
$\sup_s|H_s^{(n)}-H_s|\to0$ ścieżkowo. Część
całki względem skończonej wariacji spełnia
\[
\left|\int_0^T(H_s^{(n)}-H_s)b_s\,\dd s\right|
\le\norm{H^{(n)}-H}_\infty\int_0^T|b_s|\,\dd s
\longrightarrow0
\]
prawie na pewno. Dla części stochastycznej
lokalizujemy energię $\int_0^T\sigma_s^2\,\dd s$
i supremum $|H|$. Wtedy zbieżność zdominowana
i izometria Itô zapewniają
$\E|\int_0^T(H_s^{(n)}-H_s)\sigma_s\,\dd W_s|^2\to0$
po zatrzymaniu. Usuwając lokalizację, otrzymujemy
zbieżność w prawdopodobieństwie do $\int_0^TH_s\,\dd X_s$.

W drugim składniku \eqref{eq:strat-decomp} korzystamy
z kowariacji. Z polaryzacji i lematu
\ref{lem:ito-qv} zastosowanego do $H+X$ i $H-X$
wynika ucp-zbieżność
$\sum_{t_{k+1}\le t}\Delta_kH\Delta_kX\to[H,X]_t$.
Połączenie obydwu granic daje \eqref{eq:strat-sum}
w chwili końcowej. Dla uciętych sum wszystkie powyższe
oszacowania zachodzą także w supremum po $t$, ponieważ
całki stochastyczne kontroluje nierówność maksymalna
Dooba, a sumy kwadratów przyrostów zbiegają ucp.
Różnica spowodowana ostatnim, niepełnym odcinkiem
jest zaniedbywalna wskutek ciągłości $H,X$ i
ograniczoności odpowiednich sum w prawdopodobieństwie.
Otrzymujemy zbieżność ucp na całym $[0,T]$.
\end{proof}

\begin{definicja}[Całka Stratonovicha]
Dla ciągłych procesów Itô $H,X$ z poprzedniego
twierdzenia definiujemy
\[
\int_0^tH_s\circ\dd X_s
:=\int_0^tH_s\,\dd X_s+\frac12[H,X]_t.
\tag{4.4}\label{eq:strat-def}
\]
Symbol $\circ\dd X_s$ oznacza całkowanie
Stratonovicha, natomiast $\dd X_s$ bez kółka
oznacza całkowanie Itô. Równość
\eqref{eq:strat-sum} pokazuje, że definicja zgadza
się z granicą symetrycznych sum Riemanna.
\end{definicja}

\begin{uwaga}[Granice stosowalności]
Wzór \eqref{eq:strat-def} zdefiniowaliśmy dla
\emph{ciągłych procesów Itô}, a nie dla każdego
procesu przewidywalnego. Dla ogólnych
półmartyngałów ciągłych można go rozszerzyć przez
kowariację; dla procesów ze skokami trzeba ponadto
rozróżnić kilka nierównoważnych konwencji (np.
całkę kanoniczną Marcusa). Nie będziemy ich utożsamiać.
\end{uwaga}

\section{Reguła łańcuchowa bez członu drugiego rzędu}
\begin{twierdzenie}[Reguła łańcuchowa Stratonovicha]
\label{th:strat-chain}
Niech $X$ będzie jednowymiarowym ciągłym procesem
Itô. Dla każdej funkcji $F\in C^3(\R)$ zachodzi
prawie na pewno, dla wszystkich $t\in[0,T]$,
\[
F(X_t)-F(X_0)=\int_0^t F'(X_s)\circ\dd X_s.
\tag{4.5}\label{eq:strat-chain}
\]
Jeśli $F\in C^{1,3}([0,T]\times\R)$, a pochodna mieszana
$F_{tx}$ istnieje i jest ciągła, to
\[
F(t,X_t)-F(0,X_0)
=\int_0^tF_t(s,X_s)\,\dd s
+\int_0^tF_x(s,X_s)\circ\dd X_s.
\tag{4.6}\label{eq:strat-chain-time}
\]
\end{twierdzenie}
\begin{proof}
Ze wzoru Itô mamy
$F(X_t)-F(X_0)=\int_0^tF'(X_s)\,\dd X_s
+\frac12\int_0^tF''(X_s)\,\dd[X]_s$.
Ponadto, stosując wzór Itô do $F'(X)$
(co wyjaśnia założenie $F\in C^3$),
dostajemy, że współczynnik przy $\dd W_s$ w
rozkładzie $F'(X_s)$ jest równy
$F''(X_s)\sigma_s$. Stąd
$[F'(X),X]_t=\int_0^tF''(X_s)\sigma_s^2\,\dd s
=\int_0^tF''(X_s)\,\dd[X]_s$.
Po podstawieniu definicji \eqref{eq:strat-def}
człony drugiego rzędu znoszą się i otrzymujemy
\eqref{eq:strat-chain}. W przypadku funkcji zależnej
od czasu stosujemy tę samą argumentację do
$H_s=F_x(s,X_s)$: skończona wariacja składnika
$\int F_{xt}\,\dd s$ nie wnosi kowariacji.
\end{proof}

\begin{zadanie}[Trzy elementarne całki]
Obliczyć, bez pozostawiania symboli całek po prawej
stronie, wyrażenia
$\int_0^TW_s\circ\dd W_s$,
$\int_0^TW_s^2\circ\dd W_s$
oraz $\int_0^Te^{W_s}\circ\dd W_s$.

\emph{Rozwiązanie.} Z reguły łańcuchowej dla
$F(x)=x^2/2$, $F(x)=x^3/3$ i $F(x)=e^x$
(odpowiednio) otrzymujemy
\[
\int_0^TW_s\circ\dd W_s=\frac12W_T^2,
\quad \int_0^TW_s^2\circ\dd W_s=\frac13W_T^3,
\quad \int_0^Te^{W_s}\circ\dd W_s=e^{W_T}-1.
\]
Dla pierwszej całki porównanie z Itô daje
$\int_0^T W_s\circ\dd W_s
=\int_0^TW_s\,\dd W_s+T/2$.
Dla drugiej różnica wynosi $\int_0^TW_s\,\dd s$;
nie jest zatem stałą deterministyczną.
\end{zadanie}

\begin{twierdzenie}[Iloczyn i całkowanie przez części]
Dla dwóch ciągłych procesów Itô $X,Y$ zachodzi
\[
X_tY_t-X_0Y_0
=\int_0^tX_s\circ\dd Y_s+
\int_0^tY_s\circ\dd X_s.
\tag{4.7}\label{eq:strat-product}
\]
\end{twierdzenie}
\begin{proof}
Ze wzoru na iloczyn Itô (twierdzenie
\ref{th:ito-product}) otrzymujemy
$\dd(XY)=X\,\dd Y+Y\,\dd X+\dd[X,Y]$.
Z kolei \eqref{eq:strat-def} daje
$\int X\circ\dd Y=\int X\,\dd Y+\frac12[X,Y]$
oraz analogiczną równość po zamianie $X,Y$.
Dodanie obu wyrażeń daje dokładnie
\eqref{eq:strat-product}.
\end{proof}

\begin{lemat}[Jednoznaczność rozkładu na martyngał i skończoną wariację]
\label{lem:uniqueness-mart-fv}
Niech $M,N$ będą ciągłymi lokalnymi martyngałami zaczynającymi się
w zerze, a $A,B$ ciągłymi, adaptowanymi procesami o skończonej
wariacji na każdym przedziale zwartym, także zaczynającymi się
w zerze. Jeżeli $M_t+A_t=N_t+B_t$ dla wszystkich $t$ prawie
na pewno (z dokładnością do nierozróżnialności), to $M=N$ i
$A=B$ z tą samą dokładnością.
\end{lemat}
\begin{proof}
Proces $D=M-N=B-A$ jest ciągłym lokalnym martyngałem
i jednocześnie ma skończoną wariację. Dla dowolnej
trajektorii ciągłego procesu o skończonej wariacji na $[0,T]$
\[
\sum_{[u,v]\in\pi}(D_v-D_u)^2
\le\max_{[u,v]\in\pi}|D_v-D_u|\,\operatorname{TV}_{[0,T]}(D)
\longrightarrow0
\]
gdy średnica deterministycznego podziału dąży do zera.
Z definicji wariacji kwadratowej ciągłego lokalnego
martyngału wynika więc $[D]_t=0$ dla każdego $t$.
Zatrzymajmy $D$ tak, aby $D^\tau$ i jego kwadrat były
całkowalne. Proces $(D^\tau)^2-[D^\tau]$
jest martyngałem, zatem
$\E(D_{t\wedge\tau})^2=0$.
Otrzymujemy $D_{t\wedge\tau}=0$ prawie na pewno.
Usuwamy lokalizację i korzystamy z ciągłości trajektorii,
aby uzyskać $D\equiv0$, co dowodzi obu równości.
\end{proof}

\section{Konwersja równań Itô i Stratonovicha}
Sama funkcja $\sigma$ nie określa procesu, jeśli nie
wiemy, w jakiej konwencji jest całka. Rozważmy
formalnie równania
\[
\dd X_t=b_I(X_t)\,\dd t+\sigma(X_t)\,\dd W_t,
\qquad
\dd X_t=b_S(X_t)\,\dd t+\sigma(X_t)\circ\dd W_t.
\]
Współczynniki dryfu $b_I$ i $b_S$ nie są na ogół
identyczne.

\begin{twierdzenie}[Konwersja w jednym wymiarze]
\label{th:ito-strat-conversion}
Niech $b_S,\sigma\in C^2(\R)$ mają lokalnie ograniczone
pochodne, a $X$ będzie ciągłym procesem Itô o
wartościach rzeczywistych, dla którego
\[
X_t=X_0+\int_0^tb_S(X_s)\,\dd s
+\int_0^t\sigma(X_s)\circ\dd W_s
\tag{4.8}\label{eq:StratSDE}
\]
przy czym wszystkie całki istnieją lokalnie.
Wówczas równoważnie $X$ spełnia równanie Itô
\[
X_t=X_0+\int_0^t
\left[b_S(X_s)+\frac12\sigma(X_s)\sigma'(X_s)\right]\dd s
+\int_0^t\sigma(X_s)\,\dd W_s.
\tag{4.9}\label{eq:ItoSDE}
\]
Odwrotnie, każde rozwiązanie Itô \eqref{eq:ItoSDE}
(mające wskazaną regularność) spełnia \eqref{eq:StratSDE}.
\end{twierdzenie}
\begin{proof}
Najpierw wyeliminujmy pozorną kolistość rozumowania.
Z założenia $X$ jest procesem Itô, zatem istnieją
procesy $\widetilde b$ i $\widetilde\sigma$, dla których
$\dd X_t=\widetilde b_t\,\dd t+
\widetilde\sigma_t\,\dd W_t$. Po rozwinięciu całki
Stratonovicha z \eqref{eq:StratSDE} otrzymujemy
\[
X_t=X_0+\int_0^t b_S(X_s)\,\dd s
+\int_0^t\sigma(X_s)\,\dd W_s
+\tfrac12[\sigma(X),W]_t.
\]
Ostatni składnik ma skończoną wariację. Na mocy lematu \ref{lem:uniqueness-mart-fv}
jednoznaczność rozkładu na ciągły lokalny martyngał
i proces o skończonej wariacji daje
$\int_0^.\widetilde\sigma_s\,\dd W_s
=\int_0^.\sigma(X_s)\,\dd W_s$.
Z izometrii po lokalizacji wnioskujemy
$\widetilde\sigma_s=\sigma(X_s)$ dla
$(\lambda_1\otimes\Prob)$-prawie wszystkich $(s,\omega)$.
Teraz wzór Itô dla $\sigma(X)$ pokazuje, że jego
współczynnik dyfuzji wynosi
$\sigma'(X_s)\sigma(X_s)$; zatem
$[\sigma(X),W]_t=\int_0^t\sigma'(X_s)\sigma(X_s)\,\dd s$.
Definicja całki Stratonovicha daje więc
\[
\int_0^t\sigma(X_s)\circ\dd W_s
=\int_0^t\sigma(X_s)\,\dd W_s
+\frac12\int_0^t\sigma'(X_s)\sigma(X_s)\,\dd s.
\]
Po podstawieniu do \eqref{eq:StratSDE} otrzymujemy
\eqref{eq:ItoSDE}. W drugą stronę odejmujemy tę samą
poprawkę. Lokalizacja pozwala kontrolować współczynniki
bez zakładania ich globalnej ograniczoności.
\end{proof}

\begin{zadanie}[Geometryczny ruch Browna w obu konwencjach]
Niech $X_0=x_0>0$, $\mu\in\R$, $\gamma\in\R$.
Pokazać, że proces
$X_t=x_0\exp((\mu-\gamma^2/2)t+\gamma W_t)$
rozwiązuje równanie Itô
$\dd X_t=\mu X_t\,\dd t+\gamma X_t\,\dd W_t$
i podać równoważne równanie Stratonovicha.

\emph{Rozwiązanie.} Stosujemy wzór Itô do
$F(t,w)=x_0\exp((\mu-\gamma^2/2)t+\gamma w)$.
Ponieważ $F_t=(\mu-\gamma^2/2)F$,
$F_w=\gamma F$ i $F_{ww}=\gamma^2F$, wynika stąd
$\dd X_t=\mu X_t\,\dd t+\gamma X_t\,\dd W_t$.
Równoważny zapis Stratonovicha brzmi
\[
\dd X_t=(\mu-\tfrac12\gamma^2)X_t\,\dd t
+\gamma X_t\circ\dd W_t.
\]
Przy zerowym dryfie \emph{Stratonovicha}, czyli
$\dd X_t=\gamma X_t\circ\dd W_t$, rozwiązaniem jest
$x_0e^{\gamma W_t}$, nie
$x_0e^{\gamma W_t-\gamma^2t/2}$.
\end{zadanie}

\begin{przyklad}[Szum addytywny]
Jeśli $\sigma(x)\equiv\gamma$ jest stała, to
$\sigma\sigma'=0$. Równania
$\dd X_t=-\kappa X_t\,\dd t+\gamma\,\dd W_t$
oraz
$\dd X_t=-\kappa X_t\,\dd t+\gamma\circ\dd W_t$
są równoważne. Dotyczy to m.in. procesu
Ornsteina--Uhlenbecka z szumem addytywnym.
Różnica między konwencjami pojawia się dopiero przy
odpowiedniej zależności współczynnika dyfuzji od stanu.
\end{przyklad}

\section{Wielowymiarowa poprawka i znaczenie geometryczne}
\begin{twierdzenie}[Konwersja wielowymiarowa]
Niech $W=(W^1,\ldots,W^m)$ będzie $m$-wymiarowym
ruchem Browna, $X$ ciągłym procesem Itô w $\R^d$,
a $b_S:\R^d\to\R^d$ i
$\sigma:\R^d\to\R^{d\times m}$ będą klasy $C^2$
z lokalnie ograniczonymi pochodnymi.
Załóżmy, że $X$ spełnia w sensie całkowym
\[
\dd X_t^i=b_S^i(X_t)\,\dd t
+\sum_{k=1}^m\sigma_{ik}(X_t)\circ\dd W_t^k,
\quad i=1,\ldots,d.
\]
Wówczas $X$ spełnia równanie Itô o tym samym
współczynniku dyfuzji oraz dryfie
\[
b_I^i(x)=b_S^i(x)+\frac12\sum_{j=1}^d
\sum_{k=1}^m\partial_{x_j}\sigma_{ik}(x)\,
\sigma_{jk}(x).
\tag{4.10}\label{eq:strat-multi-correction}
\]
\end{twierdzenie}
\begin{proof}
Podobnie jak w dowodzie jednowymiarowym, zaczynamy
od dowolnego rozkładu Itô procesu $X$ o macierzy
dyfuzji $\widetilde\sigma$. Zapis całkowy równania
Stratonovicha różni się od całki Itô z macierzą
$\sigma(X)$ wyłącznie procesem o skończonej wariacji.
Lemat \ref{lem:uniqueness-mart-fv}, a następnie izometria
Itô, dają $\widetilde\sigma_s=\sigma(X_s)$ prawie
wszędzie względem $\dd s\otimes\dd\Prob$.
Dopiero teraz, dla ustalonego $i,k$, stosujemy
wielowymiarowy wzór Itô do $H_t=\sigma_{ik}(X_t)$.
Współczynnik przy $\dd W_t^\ell$ w jego części
martyngałowej jest równy
$\sum_j\partial_{x_j}\sigma_{ik}(X_t)\sigma_{j\ell}(X_t)$.
Zatem, z definicji kowariacji,
$[\sigma_{ik}(X),W^k]_t=
\int_0^t\sum_j\partial_{x_j}\sigma_{ik}(X_s)
\sigma_{jk}(X_s)\,\dd s$.
Sumujemy to po $k$ i dodajemy połowę kowariacji do
całek Itô. To jest dokładnie różnica dryfów podana
we wzorze \eqref{eq:strat-multi-correction}.
\end{proof}

\begin{przyklad}[Pole wektorowe na płaszczyźnie]
Niech $V(x,y)=(-y,x)$ i rozważmy równanie
Stratonovicha $\dd X_t=V(X_t)\circ\dd W_t$
dla $X_t=(X_t^1,X_t^2)$. Macierz pochodnej $DV$
ma postać $\left(\begin{smallmatrix}0&-1\\1&0\end{smallmatrix}\right)$,
zatem $DV(x,y)V(x,y)=(-x,-y)$.
Równoważne równanie Itô to
\[
\dd X_t=-\frac12X_t\,\dd t+V(X_t)\,\dd W_t.
\]
Z reguły łańcuchowej Stratonovicha dla $\norm{x}^2$
wynika $\dd\norm{X_t}^2=2X_t\cdot V(X_t)\circ\dd W_t=0$.
Stąd $\norm{X_t}=\norm{X_0}$ dla wszystkich $t$:
trajektoria porusza się po okręgu. W konwencji Itô
ten sam wynik wymaga poprawki od wariacji kwadratowej:
$-\norm{X_t}^2\,\dd t+\norm{V(X_t)}^2\,\dd t=0$.
To elementarny przykład geometrycznego znaczenia
rachunku Stratonovicha.
\end{przyklad}

\section{Aproksymacje gładkimi trajektoriami: elementarny przypadek Wong--Zakaia}
\begin{zadanie}[Dlaczego wygładzenie szumu prowadzi do Stratonovicha?]
Niech $W^{(n)}$ będzie liniową interpolacją
trajektorii Wienera na deterministycznym podziale
$\pi_n$ o średnicy dążącej do zera. Dla $x_0>0$
i $\gamma\in\R$ rozwiązać klasyczne równanie
Riemanna--Stieltjesa
\[
X_t^{(n)}=x_0+\gamma\int_0^t
X_s^{(n)}\,\dd W_s^{(n)}
\]
i znaleźć granicę jednostajną na $[0,T]$.

\emph{Rozwiązanie.} Każda trajektoria $W^{(n)}$
jest odcinkami liniowa, a więc ma skończoną wariację
na $[0,T]$. Klasyczna reguła łańcuchowa daje
$X_t^{(n)}=x_0\exp(\gamma W_t^{(n)})$.
Ponieważ prawie każda trajektoria $W$ jest
jednostajnie ciągła na $[0,T]$, jej interpolacje
liniowe zbiegają jednostajnie do $W$. Ciągłość
funkcji wykładniczej na zwartym zbiorze zawierającym
wartości wszystkich interpolacji zapewnia
\[
\sup_{t\le T}|X_t^{(n)}-x_0e^{\gamma W_t}|
\longrightarrow0\qquad\text{prawie na pewno}.
\]
Proces graniczny rozwiązuje
$\dd X_t=\gamma X_t\circ\dd W_t$,
a nie $\dd X_t=\gamma X_t\,\dd W_t$.
Udowodniliśmy zatem elementarną wersję zjawiska
Wong--Zakaia. Ogólne twierdzenie dla wielowymiarowych,
nieliniowych układów wymaga dodatkowej teorii i
nie wynika samo z tego przykładu.
\end{zadanie}

\section{Co otrzymaliśmy i jaką teorię rozwiniemy dalej?}
Całka Stratonovicha nie jest prostym przemianowaniem
całki Itô. Łączy się z nią inna reguła sumowania i
inne przesunięcie współczynnika dryfu. Jest wygodna
przy zmianie współrzędnych, w problemach geometrycznych
i przy aproksymowaniu ruchu Browna gładkimi ścieżkami.
Natomiast całka Itô naturalnie współpracuje z
warunkową wartością oczekiwaną, martyngałami i
rachunkiem probabilistycznym. Obydwie konwencje są
ściśle powiązane wzorem korekcyjnym.

Dalsze rozdziały poświęcimy istnieniu, jednoznaczności
i własnościom rozwiązań stochastycznych równań
różniczkowych, wzorowi Feynmana--Kaca oraz całkom
po przestrzeni trajektorii względem miary Wienera.


% ============================================================================
% ROZDZIAL 5 — SDE: konstrukcja, istnienie i jednoznacznosc
% ============================================================================
\chapter{Stochastyczne równania różniczkowe: konstrukcja rozwiązań}
\label{ch:sde-existence}
\paragraph{Odsyłacze międzytomowe.} Numeryczne zastosowania SDE łączą się z aproksymacją z rozdziału~\ref{ch:sde-properties} oraz symulacyjnymi narzędziami tomu IV, rozdziału~\ref{IV-ch:monte-carlo}.


\section{Dlaczego równanie stochastyczne nie jest zwykłym równaniem różniczkowym?}
Rozważmy równanie $\dd X_t=b(t,X_t)\,\dd t+\sigma(t,X_t)\,\dd W_t$.
Gdyby $W$ miał gładkie trajektorie, można by próbować zastąpić
ostatni człon przez $\sigma(t,X_t)\dot W_t\,\dd t$.
Wiemy jednak z tomu II, że trajektorie procesu Wienera
są niemal na pewno nieróżniczkowalne. Ponadto całka Itô
została zdefiniowana przez domknięcie w przestrzeni procesów
przewidywalnych, a nie jako klasyczna całka wzdłuż trajektorii.
Zapis różniczkowy jest więc skrótem równania całkowego.

W teorii RÓR zadane pole wektorowe i punkt początkowy wyznaczają
trajektorię. Tu do tych danych dochodzi \emph{realizacja szumu}.
Pytamy osobno, czy rozwiązanie istnieje, czy jest jednoznaczne
przy ustalonym szumie, czy jego rozkład jest jednoznaczny i czy
może eksplodować w skończonym czasie. Rozróżnienia te nie są
formalnością: brak odpowiedniej regularności współczynników
może zniszczyć jednoznaczność nawet w równaniu deterministycznym.

\begin{center}
\begin{tikzpicture}[>=Latex,node distance=6mm]
\node[draw,rounded corners,align=center,fill=blue!5,minimum width=31mm] (data)
  {$X_0$, $b$, $\sigma$\\oraz proces $W$};
\node[draw,rounded corners,align=center,fill=orange!9,right=13mm of data,minimum width=37mm] (eq)
  {$X_t=X_0+\int_0^t b(s,X_s)\,\dd s$\\$\displaystyle+\int_0^t\sigma(s,X_s)\,\dd W_s$};
\node[draw,rounded corners,align=center,fill=green!8,right=13mm of eq,minimum width=32mm] (sol)
  {ciągły proces\\adaptowany $X$};
\draw[->,thick] (data)--(eq);
\draw[->,thick] (eq)--(sol);
\end{tikzpicture}
\par\smallskip\small Równanie stochastyczne jest warunkiem na \emph{cały proces}, nie wzorem na pochodną jego trajektorii.
\end{center}

\section{Definicje i dokładne znaczenie całek}
\begin{definicja}[Rozwiązanie silne na ustalonej przestrzeni]
\label{def:sde-strong}
Niech $T>0$ i niech $(\Omega,\calF,\Prob)$ będzie
przestrzenią probabilistyczną wyposażoną w filtrację
$(\calF_t)_{0\le t\le T}$ spełniającą warunki zwykłe.
Niech $W=(W^1,\ldots,W^m)$ będzie procesem Wienera
$m$-wymiarowym względem tej filtracji. Niech $X_0$ będzie
$\calF_0$-mierzalną zmienną o wartościach w $\R^d$.
Niech $b:[0,T]\times\R^d\to\R^d$ i
$\sigma:[0,T]\times\R^d\to\R^{d\times m}$
będą funkcjami borelowskimi. \emph{Silnym rozwiązaniem} równania
\[
\dd X_t=b(t,X_t)\,\dd t+\sigma(t,X_t)\,\dd W_t,
\qquad X_0\text{ zadane},
\tag{5.1}\label{eq:sde}
\]
nazywamy ciągły, adaptowany proces $X$ o wartościach w $\R^d$,
dla którego prawie na pewno
$\int_0^T (\norm{b(s,X_s)}+\norm{\sigma(s,X_s)}_{\mathrm F}^2)\,\dd s<\infty$
oraz, z prawdopodobieństwem jeden jednocześnie dla wszystkich $t$,
\[
X_t=X_0+\int_0^tb(s,X_s)\,\dd s
       +\int_0^t\sigma(s,X_s)\,\dd W_s.
\]
Druga całka jest sumą $m$ całek Itô po kolumnach macierzy $\sigma$.
\end{definicja}

\begin{uwaga}[Mierzalność i rola filtracji]
Ponieważ $X$ jest ciągły i adaptowany, jest przewidywalny.
Borelowskość $b$ i $\sigma$ zapewnia, że
$(t,\omega)\mapsto b(t,X_t(\omega))$ i
$(t,\omega)\mapsto\sigma(t,X_t(\omega))$ są przewidywalne,
a więc całki w definicji mają sens. Ważne jest, że nie wystarcza
samo istnienie procesu Wienera jako rodziny zmiennych losowych:
jego przyszłe przyrosty muszą być niezależne od aktualnej filtracji.
\end{uwaga}

\begin{definicja}[Jednoznaczność ścieżkowa i jednoznaczność co do prawa]
Dla ustalonych danych $b,\sigma$ mówimy o
\emph{jednoznaczności ścieżkowej}, gdy dowolne dwa rozwiązania
silne $X,Y$ zbudowane na tej samej filtrowanej przestrzeni,
z tym samym procesem Wienera i tym samym warunkiem początkowym
($X_0=Y_0$ prawie na pewno), są nierozróżnialne:
$\Prob(X_t=Y_t\ \text{dla wszystkich }t\in[0,T])=1$.
\emph{Jednoznaczność co do prawa} oznacza, że dowolne dwa
rozwiązania słabe o tym samym rozkładzie początkowym mają
ten sam rozkład na $C([0,T];\R^d)$.
\end{definicja}

\begin{definicja}[Rozwiązanie słabe]
Niech $\mu$ będzie miarą probabilistyczną na $\R^d$.
\emph{Rozwiązaniem słabym} \eqref{eq:sde} z rozkładem
początkowym $\mu$ nazywamy pewną filtrowaną przestrzeń
probabilistyczną, pewien $m$-wymiarowy proces Wienera $W$
względem jej filtracji oraz ciągły proces adaptowany $X$
spełniający równanie całkowe i $\Law(X_0)=\mu$.
W tej definicji przestrzeń probabilistyczna i szum nie są
a priori zadane. Słowo „słaby” nie oznacza, że równanie
obowiązuje tylko po scałkowaniu z funkcją testową.
\end{definicja}

\section{Podstawowe oszacowanie dla całek stochastycznych}
\begin{lemat}[Oszacowanie supremum]
\label{lem:sde-sup}
Niech $T>0$, $A$ będzie procesem progresywnie mierzalnym
o wartościach w $\R^d$, a $B$ przewidywalnym o wartościach
w $\R^{d\times m}$. Załóżmy, że
$\E\int_0^T(\norm{A_s}^2+\norm{B_s}_{\mathrm F}^2)\,\dd s<\infty$.
Wówczas dla $0\le t\le T$
\[
\E\sup_{r\le t}\norm{\int_0^r A_s\,\dd s+
\int_0^r B_s\,\dd W_s}^2
\le 2t\int_0^t\E\norm{A_s}^2\,\dd s
+8\int_0^t\E\norm{B_s}_{\mathrm F}^2\,\dd s.
\tag{5.2}\label{eq:sde-basic-est}
\]
\end{lemat}
\begin{proof}
Z nierówności $(u+v)^2\le2u^2+2v^2$ dla norm wektorowych
wystarczy osobno oszacować obie całki. Dla pierwszej,
przez Cauchy'ego--Schwarza, dla każdego $r\le t$ mamy
$\norm{\int_0^r A_s\,\dd s}^2
\le r\int_0^r\norm{A_s}^2\,\dd s
\le t\int_0^t\norm{A_s}^2\,\dd s$.
Dla drugiej, stosując nierówność Dooba w $L^2$
do kwadratu normy martyngału (lub do jego współrzędnych),
uzyskujemy
$\E\sup_{r\le t}\norm{\int_0^rB_s\,\dd W_s}^2
\le4\E\norm{\int_0^tB_s\,\dd W_s}^2$.
Izometria Itô daje $\E\norm{\int_0^tB_s\,\dd W_s}^2
=\E\int_0^t\norm{B_s}_{\mathrm F}^2\,\dd s$.
Po połączeniu oszacowań otrzymujemy tezę.
\end{proof}

\begin{uwaga}
Oszacowanie kontroluje \emph{całą trajektorię} przy pomocy
norm współczynników, a nie tylko wartość w ustalonej chwili.
Właśnie dlatego metoda Picarda zbiega w przestrzeni
ciągłych procesów, co jest silniejsze niż zbieżność
punktowa dla każdego $t$.
\end{uwaga}

\section{Istnienie i jednoznaczność przy globalnym warunku Lipschitza}
\begin{twierdzenie}[Twierdzenie o istnieniu i jednoznaczności silnego rozwiązania]
\label{th:sde-global}
Niech $T>0$, niech $W$ i filtracja będą jak w definicji
\ref{def:sde-strong}, a $X_0\in L^2(\Omega;\R^d)$ będzie
$\calF_0$-mierzalny. Załóżmy, że borelowskie współczynniki
$b:[0,T]\times\R^d\to\R^d$ i
$\sigma:[0,T]\times\R^d\to\R^{d\times m}$ spełniają:
\begin{enumerate}[label=(\roman*)]
\item dla pewnego $L\ge0$ i wszystkich $t\in[0,T]$ oraz
$x,y\in\R^d$ mamy
$\norm{b(t,x)-b(t,y)}+\norm{\sigma(t,x)-\sigma(t,y)}_{\mathrm F}
\le L\norm{x-y}$;
\item dla pewnego $K<\infty$ i wszystkich $t\in[0,T]$
mamy $\norm{b(t,0)}+\norm{\sigma(t,0)}_{\mathrm F}\le K$.
\end{enumerate}
Wówczas równanie \eqref{eq:sde} ma dokładnie jedno silne
rozwiązanie w sensie nierozróżnialności. Ponadto dla pewnej
stałej $C=C(T,L,K)$ zachodzi
\[
\E\sup_{0\le t\le T}\norm{X_t}^2
\le C(1+\E\norm{X_0}^2).
\tag{5.3}\label{eq:sde-moment}
\]
\end{twierdzenie}
\begin{proof}
\emph{Krok 1: liniowy wzrost.} Z założeń i nierówności
trójkąta otrzymujemy, dla pewnego $K_1=K_1(L,K)$,
$\norm{b(t,x)}^2+\norm{\sigma(t,x)}_{\mathrm F}^2
\le K_1(1+\norm{x}^2)$, jednostajnie względem $t$.

\emph{Krok 2: iteracje Picarda.} Ustalmy $X^{(0)}_t=X_0$
i dla $n\ge0$ zdefiniujmy
\[
X_t^{(n+1)}=X_0+\int_0^tb(s,X_s^{(n)})\,\dd s
+\int_0^t\sigma(s,X_s^{(n)})\,\dd W_s.
\tag{5.4}\label{eq:sde-picard}
\]
Przez indukcję wszystkie iteracje są ciągłymi procesami
adaptowanymi i należą do przestrzeni
$L^2(\Omega;C([0,T];\R^d))$; w kroku indukcyjnym
stosujemy liniowy wzrost i lemat \ref{lem:sde-sup}.
Nie postulujemy zatem istnienia całek przed wykazaniem
ich całkowalności.

\emph{Krok 3: szybkość zbieżności różnic.}
Niech $D_n(t)=\E\sup_{r\le t}
\norm{X_r^{(n+1)}-X_r^{(n)}}^2$.
Z lematu \ref{lem:sde-sup} i warunku Lipschitza,
po zastosowaniu nierówności $r\le T$, otrzymujemy
$D_n(t)\le C_0\int_0^t D_{n-1}(s)\,\dd s$
dla $n\ge1$ i $C_0=L^2(2T+8)$.
Pierwszy przyrost jest całkowalny:
$D_0(T)<\infty$ na mocy oszacowania zastosowanego
do $b(s,X_0),\sigma(s,X_0)$. Indukcja daje
\[
D_n(T)\le D_0(T)\frac{(C_0T)^n}{n!}.
\tag{5.5}\label{eq:picard-factorial}
\]
Zatem
$\sum_{n=0}^\infty\|X^{(n+1)}-X^{(n)}\|_{L^2(\Omega;C)}
\le\sqrt{D_0(T)}\sum_{n=0}^\infty(C_0T)^{n/2}/\sqrt{n!}<\infty$.
Przestrzeń $L^2(\Omega;C([0,T];\R^d))$ jest zupełna,
więc istnieje jej element $X$ będący granicą iteracji.
Możemy wybrać podciąg zbieżny jednostajnie niemal
na pewno; wynika to np. z nierówności Markowa i
lematu Borela--Cantellego. W tej wersji $X$ ma
ciągłe trajektorie i jest adaptowany.

\emph{Krok 4: przejście do granicy w równaniu.}
Z warunku Lipschitza i lematu \ref{lem:sde-sup}
wynika, że różnica między procesami całkowymi po
podstawieniu $X^{(n)}$ i $X$ zbiega do zera
w normie $L^2(\Omega;C)$, ponieważ jest ograniczona
przez $C_0\int_0^T\E\norm{X_s^{(n)}-X_s}^2\,\dd s$.
Granica spełnia więc równanie całkowe. Liniowy
wzrost oraz $X\in L^2(\Omega;C)$ gwarantują
wszystkie warunki całkowalności z definicji.

\emph{Krok 5: jednoznaczność.} Jeśli $X,Y$ są
rozwiązaniami o tym samym warunku początkowym i
przyrostach Wienera, oszacowanie z kroku 3 daje
$u(t):=\E\sup_{r\le t}\norm{X_r-Y_r}^2
\le C_0\int_0^t u(s)\,\dd s$.
Lemat Gronwalla implikuje $u(t)=0$ dla każdego $t$,
co oznacza nierozróżnialność.
Jeśli $X,Y$ są jedynie rozwiązaniami w sensie definicji
\ref{def:sde-strong}, bez założenia kwadratowej całkowalności
supremów, nie można stosować powyższego oszacowania wprost.
Dla $R>0$ wprowadzamy wspólny czas zatrzymania
$\tau_R=\inf\{t\in[0,T]:\norm{X_t}\vee\norm{Y_t}\ge R\}\wedge T$.
Ponieważ procesy są ciągłe, a $X_0=Y_0$ prawie na pewno,
na $[0,\tau_R]$ integrandy różnic są ograniczone przez
stałą zależną od $R$. Stosujemy lemat \ref{lem:sde-sup}
do różnicy całek zatrzymanych, uzyskując
$u_R(t):=\E\sup_{r\le t}\norm{X_{r\wedge\tau_R}-Y_{r\wedge\tau_R}}^2
\le C_0\int_0^t u_R(s)\,\dd s$.
Lemat Gronwalla daje $u_R(t)=0$. Dla prawie każdej trajektorii
ciągłe procesy $X,Y$ są ograniczone na $[0,T]$, więc
$\tau_R=T$ dla wszystkich dostatecznie dużych $R$.
Przechodząc do granicy po rosnącym ciągu promieni,
otrzymujemy nierozróżnialność na całym $[0,T]$.

\emph{Krok 6: ograniczenie momentów.}
Z równania i nierówności
$\norm{u+v+w}^2\le3(\norm u^2+\norm v^2+\norm w^2)$,
po zastosowaniu lematu \ref{lem:sde-sup} do obu części,
otrzymujemy
\[
\E\sup_{r\le t}\norm{X_r}^2
\le3\E\norm{X_0}^2+C_1\int_0^t(1+
\E\sup_{u\le s}\norm{X_u}^2)\,\dd s,
\]
gdzie $C_1$ zależy tylko od $T,L,K$.
Z lematu Gronwalla wynika \eqref{eq:sde-moment}.
\end{proof}

\begin{uwaga}[Dlaczego w dowodzie pojawia się silnia?]
Dla odwzorowania Picarda na całym przedziale $[0,T]$
nie musimy otrzymać kontrakcji w jednym kroku.
Każda kolejna iteracja wprowadza jednak dodatkowe
całkowanie po czasie. Objętość $n$-wymiarowego
simpleksu $0<s_n<\dots<s_1<T$ wynosi $T^n/n!$;
to właśnie źródło silni w \eqref{eq:picard-factorial}.
\end{uwaga}

\begin{zadanie}[Dlaczego regularność współczynników ma znaczenie?]
Rozważyć równanie $\dd X_t=2\sqrt{|X_t|}\,\dd t$
z warunkiem $X_0=0$, bez składnika szumowego.

\emph{Rozwiązanie.} Dla każdego $c\ge0$ proces
deterministyczny $X_t=(t-c)_+^2$ spełnia równanie:
dla $t\le c$ obie strony są zerowe, a dla $t>c$
mamy $\dot X_t=2(t-c)=2\sqrt{X_t}$. Istnieje
więc nieskończenie wiele rozwiązań. Współczynnik
$b(x)=2\sqrt{|x|}$ nie jest lokalnie Lipschitzowski
w zerze. Nawet bez szumu warunku tego nie można
w ogólności po prostu pominąć.
\end{zadanie}

\begin{lemat}[Mierzalna konstrukcja iteracji Picarda]
\label{lem:picard-measurable}
Niech współczynniki $b,\sigma$ spełniają założenia
twierdzenia \ref{th:sde-global}, a $\mu$ będzie rozkładem
na $\R^d$ o skończonym drugim momencie.
Dla ustalonego $T>0$ iteracje Picarda można skonstruować
tak, że prawo procesu $X^{(n)}$ jest jednoznacznie
wyznaczone przez prawo pary $(X_0,W)$ na
$\R^d\times C_0([0,T];\R^m)$, niezależnie od
filtrowanej przestrzeni probabilistycznej.
\end{lemat}
\begin{proof}
Dla $n=0$ proces jest stały $X_0$, więc wniosek jest oczywisty.
Załóżmy, że proces $X^{(n)}$ otrzymaliśmy jako mierzalny
funkcjonał pary $(X_0,W)$ (z dokładnością do zbioru miary
zero wyznaczonego przez jej prawo). Proces
$H_t=\sigma(t,X_t^{(n)})$ jest przewidywalny i należy
do $L^2(\dd t\otimes\Prob)$; analogicznie
$b(t,X_t^{(n)})$ jest całkowalny po czasie.
Dla $h=T/N$ i $k\ge1$ określamy
$H_k^{(h)}=h^{-1}\int_{(k-1)h}^{kh}H_s\,\dd s$,
a następnie proces schodkowy równy $H_k^{(h)}$
na $[kh,(k+1)h)$, z wartością zero na $[0,h)$.
Każda wartość $H_k^{(h)}$ jest $\calF_{kh}$-mierzalna;
więc przybliżenie jest przewidywalne i jego całka
Itô jest skończoną sumą
$\sum_k H_k^{(h)}(W_{t\wedge t_{k+1}}-W_{t\wedge t_k})$.
Jest to funkcjonał mierzalny pary $(X_0,W)$.

Dla dowolnego procesu z klasy $L^2(\dd t\otimes\Prob)$
przesunięte średnie na poprzednim przedziale zbiegają
do oryginalnego procesu w tej normie. Najpierw sprawdzamy
to dla funkcji schodkowych w czasie, a następnie dla
wszystkich funkcji $L^2$ przez ich gęstość i fakt, że
uśrednianie jest kontrakcją w $L^2$ (Jensen oraz Fubini).
Izometria Itô i nierówność maksymalna Dooba dają stąd
zbieżność całek w $L^2(\Omega;C([0,T]))$.
Dla całki względem czasu mierzalność funkcjonału
jest bezpośrednia z aproksymacji całek Lebesgue'a.
Wybór podciągu zbieżnego niemal na pewno pozwala
zidentyfikować całe $X^{(n+1)}$ jako mierzalny
funkcjonał tej samej pary. Wszystkie konstrukcje
odbywają się względem ustalonego wspólnego prawa
$(X_0,W)$, więc prawo iteracji nie zależy od
konkretnej realizacji przestrzeni probabilistycznej.
Indukcja kończy dowód.
\end{proof}

\begin{wniosek}[Jednoznaczność co do prawa]
\label{cor:sde-law}
Przy założeniach twierdzenia \ref{th:sde-global},
niech $\mu$ będzie rozkładem probabilistycznym
na $\R^d$ o skończonym drugim momencie.
Wówczas wszystkie rozwiązania słabe równania
\eqref{eq:sde} z $\Law(X_0)=\mu$ mają ten sam
rozkład na $C([0,T];\R^d)$.
\end{wniosek}
\begin{proof}
Dla każdego rozwiązania słabego $X_0$ jest
$\calF_0$-mierzalny, a $W$ jest procesem Wienera
względem tej filtracji. Cała trajektoria $W$
jest więc niezależna od $X_0$, a wspólne prawo pary
$(X_0,W)$ jest zawsze równe
$\mu\otimes\mu_W$, gdzie $\mu_W$ jest miarą Wienera.
Lemat \ref{lem:picard-measurable} daje mierzalną
konstrukcję kolejnych iteracji Picarda na podstawie
tej pary, a zatem ich prawo nie zależy od wyboru
rozwiązania słabego. W szczególności nie zakładamy
nieudowodnionej zależności całki Itô od całej trajektorii
szumu: uzyskujemy ją jako granicę jawnych sum
z przewidywalnymi współczynnikami. Na mocy
\eqref{eq:picard-factorial} iteracje zbiegają w
prawdopodobieństwie w przestrzeni trajektorii
do rozwiązania. Ponieważ rozkład granicy w
przestrzeni metrycznej jest wyznaczony przez rozkłady
przybliżeń, otrzymujemy tezę.
\end{proof}

\begin{przyklad}[Ten sam rozkład nie oznacza tych samych trajektorii]
Dla równania $\dd X_t=\dd W_t$, $X_0=0$,
jedynym silnym rozwiązaniem przy zadanym szumie jest
$X_t=W_t$. Na innych przestrzeniach możemy wziąć
różne ruchy Browna; odpowiednie rozwiązania nie są
funkcjami tego samego wyniku $\omega$, lecz mają
identyczne prawa na przestrzeni trajektorii.
To ilustruje różnicę między jednoznacznością ścieżkową
a jednoznacznością co do prawa.
\end{przyklad}

\section{Stabilność względem warunku początkowego}
\begin{twierdzenie}[Stabilność w normie średniokwadratowej]
\label{th:sde-stability}
Niech $b,\sigma$ spełniają założenia twierdzenia
\ref{th:sde-global}, a $X,Y$ będą rozwiązaniami
na tej samej przestrzeni, napędzanymi tym samym $W$,
z warunkami początkowymi $X_0,Y_0\in L^2$.
Wówczas dla pewnego $C=C(T,L)$
\[
\E\sup_{t\le T}\norm{X_t-Y_t}^2
\le C\E\norm{X_0-Y_0}^2.
\tag{5.6}\label{eq:sde-stability}
\]
\end{twierdzenie}
\begin{proof}
Odejmujemy równania całkowe. Ponieważ
$X_t-Y_t=X_0-Y_0+\int_0^t(b(s,X_s)-b(s,Y_s))\,\dd s
+\int_0^t(\sigma(s,X_s)-\sigma(s,Y_s))\,\dd W_s$,
z lematu \ref{lem:sde-sup} i nierówności
$\norm{u+v}^2\le2\norm u^2+2\norm v^2$
otrzymujemy
$u(t)\le2\E\norm{X_0-Y_0}^2+C_2\int_0^t u(s)\,\dd s$,
gdzie $u(t)=\E\sup_{r\le t}\norm{X_r-Y_r}^2$.
Lemat Gronwalla daje \eqref{eq:sde-stability}.
\end{proof}

\begin{przyklad}[Co mierzy oszacowanie stabilności?]
Jeśli dla tego samego procesu Wienera rozwiązujemy
dwa równania o identycznych współczynnikach, ale
z warunkami początkowymi oddalonymi o $\varepsilon$,
to średni kwadrat \emph{największej w czasie} różnicy
między trajektoriami jest nie większy od $C\varepsilon^2$.
Wynik nie mówi, że \emph{dwie niezależnie zaszumione}
trajektorie są bliskie. Sprzężenie polegające na
użyciu tego samego $W$ jest częścią założenia.
\end{przyklad}

\section{Warunki lokalne i zagadnienie eksplozji}
\begin{definicja}[Czas eksplozji]
Dla procesu ciągłego $X$ określonego do losowej chwili
$\zeta$ i wartościowanego w $\R^d$ definiujemy
$\tau_R=\inf\{t<\zeta:\norm{X_t}\ge R\}$
(przyjmujemy $\inf\varnothing=\infty$).
Jeżeli $\tau_R\uparrow\zeta$, to $\zeta$ nazywamy
\emph{czasem życia} procesu. Proces jest
\emph{nieeksplozywny} na $[0,T]$, jeśli $\Prob(\zeta>T)=1$.
\end{definicja}

\begin{twierdzenie}[Lokalny Lipschitz i liniowy wzrost]
\label{th:sde-local}
Niech $T>0$, niech $X_0\in L^2(\Omega;\R^d)$,
a $b,\sigma$ będą borelowskie i dla każdego $R>0$
istnieje $L_R<\infty$, takie że dla wszystkich
$t\in[0,T]$, $\norm x,\norm y\le R$ zachodzi
$\norm{b(t,x)-b(t,y)}+\norm{\sigma(t,x)-\sigma(t,y)}_{\mathrm F}
\le L_R\norm{x-y}$.
Załóżmy także istnienie $K$ z własnością
$\norm{b(t,x)}^2+\norm{\sigma(t,x)}_{\mathrm F}^2
\le K(1+\norm{x}^2)$ dla wszystkich $t,x$.
Wówczas \eqref{eq:sde} ma jedno silne rozwiązanie
na $[0,T]$, jest ono nieeksplozywne i spełnia
oszacowanie momentów analogiczne do \eqref{eq:sde-moment}.
\end{twierdzenie}
\begin{proof}
\emph{Obcięcie.} Niech $\rho_R(x)$ będzie rzutem
na domkniętą kulę promienia $R$; jest on $1$-Lipschitzowski
i równy identyczności w tej kuli. Definiujemy
$b_R(t,x)=b(t,\rho_R(x))$,
$\sigma_R(t,x)=\sigma(t,\rho_R(x))$.
Nowe współczynniki są globalnie Lipschitzowskie
ze stałą $L_R$ i spełniają ten sam warunek liniowego
wzrostu z tą samą stałą $K$, ponieważ
$\norm{\rho_R(x)}\le\norm x$.
Z twierdzenia \ref{th:sde-global} otrzymujemy
globalne rozwiązania $X^{(R)}$.

\emph{Zgodność rozwiązań.} Dla $S>R$ rozwiązania
$X^{(R)}$ i $X^{(S)}$ mają identyczne współczynniki,
gdy oba przebywają w kuli promienia $R$.
Stosujemy oszacowanie różnicy z dowodu jednoznaczności
do procesów zatrzymanych przy pierwszym opuszczeniu tej kuli
przez którykolwiek z nich. Ponieważ różnica początkowa
wynosi zero, Gronwall daje zgodność do tej chwili;
a więc również zgodność czasów pierwszego wyjścia.
Można zatem skleić lokalne rozwiązania na przedziale
przed granicą tych czasów.

\emph{Wykluczenie eksplozji.} Oszacowanie momentów z
kroku 6 poprzedniego dowodu zależy od stałej liniowego
wzrostu $K$, nie od stałej Lipschitza $L_R$.
Zatem $\sup_R\E\sup_{t\le T}\norm{X_t^{(R)}}^2
\le C(T,K)(1+\E\norm{X_0}^2)$.
Dla czasu wyjścia $\tau_R$ z kuli o promieniu $R$
wynika z nierówności Markowa
\[
\Prob(\tau_R\le T)\le
\frac{C(T,K)(1+\E\norm{X_0}^2)}{R^2}.
\]
Prawa strona dąży do zera. Stąd z prawdopodobieństwem
jeden lokalne rozwiązanie można przedłużyć na cały
$[0,T]$ i otrzymujemy jednoznaczność przez zgodność
obciętych równań.
\end{proof}

\begin{zadanie}[Eksplozja bez warunku liniowego wzrostu]
Rozwiązać równanie $\dd X_t=X_t^2\,\dd t$,
$X_0=x_0>0$, i wyznaczyć czas jego eksplozji.

\emph{Rozwiązanie.} Równanie jest deterministyczne.
Rozdzielenie zmiennych prowadzi do
$X_t=x_0/(1-x_0t)$ dla $t<1/x_0$.
W chwili $t\uparrow1/x_0$ rozwiązanie dąży do
nieskończoności, zatem $\zeta=1/x_0$. Funkcja
$b(x)=x^2$ jest lokalnie Lipschitzowska, ale nie
spełnia globalnego warunku liniowego wzrostu.
To pokazuje, dlaczego w poprzednim twierdzeniu
występują dwa odrębne założenia.
\end{zadanie}

\section{Rozwiązane modele: wykładniczy i liniowy}
\begin{zadanie}[Geometryczny ruch Browna: rozwiązanie i rozkład]
\label{zad:gbm}
Niech $a,c\in\R$ i $x_0>0$. Rozwiązać
\[
\dd X_t=aX_t\,\dd t+cX_t\,\dd W_t,
\qquad X_0=x_0,
\]
i wyznaczyć $\E X_t$, $\Var(X_t)$ oraz
$\Prob(X_t>K)$ dla ustalonego $K>0$ i $t>0$.

\emph{Rozwiązanie.} Ponieważ współczynniki są
Lipschitzowskie, rozwiązanie istnieje i jest jednoznaczne.
Szukamy postaci $X_t=x_0e^{Y_t}$, gdzie
$Y_t=(a-c^2/2)t+cW_t$. Wzór Itô dla funkcji
wykładniczej daje
$\dd X_t=X_t(a-c^2/2+c^2/2)\,\dd t+cX_t\,\dd W_t$.
A zatem
\[
X_t=x_0\exp\bigl((a-\tfrac12c^2)t+cW_t\bigr).
\tag{5.7}\label{eq:gbm-solution}
\]
Zmienne $W_t$ mają rozkład $\Normal(0,t)$, a więc
$X_t$ ma rozkład lognormalny (dla $c\ne0$).
Dzięki $\E e^{uW_t}=e^{u^2t/2}$ otrzymujemy
$\E X_t=x_0e^{at}$,
$\E X_t^2=x_0^2e^{(2a+c^2)t}$ oraz
$\Var(X_t)=x_0^2e^{2at}(e^{c^2t}-1)$.
Dla $c\ne0$, po standaryzacji zmiennej normalnej,
\[
\Prob(X_t>K)=1-\Phi\!\left(
\frac{\log(K/x_0)-(a-c^2/2)t}{|c|\sqrt t}\right),
\]
gdzie $\Phi$ jest dystrybuantą $\Normal(0,1)$.
Gdy $c=0$, proces jest deterministyczny i powyższe
prawdopodobieństwo wynosi $\ind{\{x_0e^{at}>K\}}$.
W szczególności $X_t>0$ prawie na pewno dla wszystkich $t$,
co nie jest oczywiste z samej postaci różniczkowej.
\end{zadanie}

\begin{zadanie}[Proces Ornsteina--Uhlenbecka]
\label{zad:ou}
Niech $\theta>0$, $m\in\R$, $\eta>0$ i
$X_0=x_0\in\R$. Rozwiązać
$\dd X_t=\theta(m-X_t)\,\dd t+\eta\,\dd W_t$
i obliczyć rozkład $X_t$ oraz $\Cov(X_s,X_t)$ dla $s\le t$.

\emph{Rozwiązanie.} Wprowadzamy czynnik całkujący
$e^{\theta t}$. Wzór na iloczyn Itô i zerowa
kowariacja funkcji czasu z $X$ dają
$\dd(e^{\theta t}(X_t-m))=\eta e^{\theta t}\,\dd W_t$.
Stąd
\[
X_t=m+(x_0-m)e^{-\theta t}
+\eta\int_0^t e^{-\theta(t-u)}\,\dd W_u.
\tag{5.8}\label{eq:ou-solution}
\]
Ostatnia całka jest gaussowska, ma średnią zero
i wariancję $\eta^2\int_0^t e^{-2\theta(t-u)}\,\dd u
=\eta^2(1-e^{-2\theta t})/(2\theta)$.
Dla $s\le t$ rozbijamy całkę w chwili $s$.
Przyrost od $s$ do $t$ jest niezależny od $\calF_s$,
a część wcześniejsza jest równa
$e^{-\theta(t-s)}(X_s-\E X_s)$. Dlatego
\[
\Cov(X_s,X_t)=\frac{\eta^2}{2\theta}
 e^{-\theta(t-s)}(1-e^{-2\theta s}).
\]
Z rozkładu widać, że gdy $t\to\infty$,
$X_t$ zbiega według rozkładu do
$\Normal(m,\eta^2/(2\theta))$.
Nie oznacza to, że $X_t$ zbiega prawie na pewno
do stałej $m$: w stanie równowagi szum nadal działa.
\end{zadanie}

\begin{przyklad}[Liniowe równanie o zmiennym dryfie]
Dla deterministycznych, ograniczonych funkcji $a,g$
i $h\in L^2([0,T])$ rozwiązanie
$\dd X_t=(a(t)X_t+g(t))\,\dd t+h(t)\,\dd W_t$
ma postać
\[
X_t=e^{A(t)}\left(X_0+\int_0^te^{-A(s)}g(s)\,\dd s
+\int_0^te^{-A(s)}h(s)\,\dd W_s\right),
\qquad A(t)=\int_0^t a(s)\,\dd s.
\]
Wystarczy zastosować wzór na iloczyn do
$e^{-A(t)}X_t$. Jeżeli $X_0$ jest deterministyczny,
proces jest gaussowski: jest sumą funkcji deterministycznej
i całki Wienera z deterministycznym integrandem.
\end{przyklad}

\section{Podsumowanie i granica teorii}
Wykazaliśmy istnienie i jednoznaczność przy globalnym
warunku Lipschitza, a następnie przy lokalnym warunku
Lipschitza i liniowym wzroście. Otrzymaliśmy także
stabilność i pierwsze modele jawne. Nie twierdzimy,
że współczynniki muszą być globalnie Lipschitzowskie,
aby SDE miało rozwiązanie: jest to użyteczne kryterium
wystarczające. W następnym rozdziale wykorzystamy
rozwiązania do budowania procesów Markowa, ich generatorów,
porównywania trajektorii i badania aproksymacji numerycznych.
% ============================================================================
% ROZDZIAL 6 — SDE: wlasnosci, generatory, porownanie, aproksymacja
% ============================================================================
\chapter{Własności rozwiązań SDE, procesy Markowa i aproksymacje}
\label{ch:sde-properties}
\paragraph{Odsyłacze międzytomowe.} Estymację parametrów procesów zależnych i dobór modeli omawia tom IV, rozdziały~\ref{IV-ch:time-series} i~\ref{IV-ch:model-selection}.


\section{Od równania do ewolucji prawdopodobieństwa}
W poprzednim rozdziale skonstruowaliśmy proces $X$ na tej
samej przestrzeni co zadany proces Wienera. Chcemy teraz
zrozumieć jego \emph{rozkład} i sposób, w jaki zależy od
warunku początkowego. W tym celu potrzebujemy dwóch
perspektyw. Pierwsza śledzi trajektorie z użyciem wspólnego
szumu; druga opisuje ewolucję wartości oczekiwanych
funkcji testowych. Ich punktem wspólnym jest wzór Itô.

Wszystkie twierdzenia tej sekcji dotyczą rozwiązania
\[
X_t^x=x+\int_0^t b(X_s^x)\,\dd s
+\int_0^t\sigma(X_s^x)\,\dd W_s,\qquad t\ge0,
\tag{6.1}\label{eq:timehom-sde}
\]
gdzie $b:\R^d\to\R^d$, $\sigma:\R^d\to\R^{d\times m}$
są globalnie Lipschitzowskie, a $W$ jest $m$-wymiarowym
procesem Wienera względem filtracji spełniającej warunki
zwykłe. Na każdym skończonym przedziale proces jest
skonstruowany przez twierdzenie \ref{th:sde-global};
jego rozszerzenia na rosnące przedziały są zgodne
dzięki jednoznaczności ścieżkowej.

\section{Własność Markowa wynikająca z jednoznaczności silnej}
\begin{definicja}[Półgrupa przejścia]
Dla $t\ge0$, $x\in\R^d$ i ograniczonej funkcji
borelowskiej $f:\R^d\to\R$ oznaczamy
\[
P_tf(x)=\E[f(X_t^x)],\qquad
P_t(x,A)=\Prob(X_t^x\in A),\quad A\in\calB(\R^d).
\tag{6.2}\label{eq:sde-semigroup}
\]
Wartości te nie zależą od konkretnej przestrzeni
probabilistycznej, na której skonstruowano rozwiązanie:
rozwiązania otrzymane metodą Picarda z identycznych
danych mają te same rozkłady.
\end{definicja}

\begin{lemat}[Zależność rozkładu od warunku początkowego]
\label{lem:sde-kernel}
Dla każdego $t\ge0$ funkcja $x\mapsto P_tf(x)$
jest ciągła, jeżeli $f:\R^d\to\R$ jest ograniczona
i ciągła. Ponadto dla każdej ograniczonej funkcji
borelowskiej $f$ odwzorowanie $x\mapsto P_tf(x)$
jest borelowskie.
\end{lemat}
\begin{proof}
Niech $x_n\to x$ i sprzęgnijmy rozwiązania
$X^{x_n},X^x$ przez ten sam proces $W$.
Z twierdzenia \ref{th:sde-stability} otrzymujemy
$\E\sup_{s\le t}\norm{X_s^{x_n}-X_s^x}^2\to0$.
Stąd $X_t^{x_n}\to X_t^x$ według prawdopodobieństwa.
Jeśli $f$ jest ograniczona i ciągła, to
$f(X_t^{x_n})\to f(X_t^x)$ według prawdopodobieństwa,
a z ograniczoności wynika zbieżność wartości oczekiwanych.
Dla ogólnej funkcji borelowskiej rozpatrujemy klasę
ograniczonych funkcji $f$, dla których $P_tf$ jest
borelowska. Zawiera ona ograniczone funkcje ciągłe;
pozostaje też zamknięta na ograniczone granice punktowe,
ponieważ zbieżność zdominowana działa dla każdej
ustalonej wartości $x$. Ograniczone funkcje ciągłe
wyznaczają borelowskie $\sigma$-ciało na $\R^d$.
Twierdzenie o klasie monotonicznej, zastosowane do
funkcji, daje zatem wszystkie ograniczone funkcje
borelowskie. Jest to zarazem uzasadnienie, że
$P_t(x,A)$ jest jądrem probabilistycznym, a nie
tylko formalnym symbolem.
\end{proof}

\begin{twierdzenie}[Własność Markowa i równania Chapmana--Kołmogorowa]
\label{th:sde-markov}
Niech $X$ będzie rozwiązaniem \eqref{eq:timehom-sde}
z $\calF_0$-mierzalnym $X_0\in L^2$, a $W$ procesem
Wienera względem tej filtracji. Wówczas dla dowolnych
$s,t\ge0$ i ograniczonej funkcji borelowskiej $f$
\[
\E[f(X_{s+t})\mid\calF_s]=P_tf(X_s)
\quad\text{prawie na pewno}.
\tag{6.3}\label{eq:sde-markov}
\]
W szczególności $P_0f=f$ i $P_sP_tf=P_{s+t}f$.
Własność ta dotyczy \emph{całej przeszłości} zawartej
w $\calF_s$, a nie tylko pojedynczej zmiennej $X_s$.
\end{twierdzenie}
\begin{proof}
Ustalmy $s$ i rozważmy szum przesunięty
$\widehat W_u=W_{s+u}-W_s$. Jest on procesem Wienera
niezależnym od $\calF_s$. Na przedziale $[s,s+t]$
proces $Z_u=X_{s+u}$ spełnia równanie
\[
Z_u=X_s+\int_0^u b(Z_v)\,\dd v
+\int_0^u\sigma(Z_v)\,\dd\widehat W_v.
\]
Porównajmy go z iteracjami Picarda uruchomionymi
w chwili $s$: $Z^{(0)}_u=X_s$ i
$Z^{(n+1)}_u=X_s+\int_0^u b(Z_v^{(n)})\,\dd v
+\int_0^u\sigma(Z_v^{(n)})\,\dd\widehat W_v$.
Dla deterministycznego $x$ analogiczne iteracje
zaczynają się w $x$ i korzystają wyłącznie z
$\widehat W$. Dla dowolnej ograniczonej funkcji
borelowskiej $f$ wartość warunkowa
$\E[f(Z_t^{(n)})\mid\calF_s]$ zależy wyłącznie od
$X_s$: w kroku indukcyjnym korzystamy z niezależności
całego procesu $\widehat W$ od $\calF_s$ oraz z
mierzalności iteracji względem $X_s$ i $\widehat W$.
Dokładniej, dla skończenie wielu prostych przybliżeń
całek Itô jest to zwykła własność warunkowania
funkcji dwóch niezależnych zmiennych losowych;
przejście do granicy w $L^2$ uzasadnia tę samą
własność dla iteracji Picarda.

Z twierdzenia \ref{th:sde-global} iteracje zbiegają do
$Z$ w $L^2(\Omega;C([0,t];\R^d))$.
Dla ograniczonego ciągłego $f$ wynika stąd
zbieżność $f(Z_t^{(n)})\to f(Z_t)$ w $L^1$.
Odpowiednie prawa dla deterministycznego $x$
zbiegać będą do $P_tf(x)$; zatem
$\E[f(Z_t)\mid\calF_s]=P_tf(X_s)$.
Przejście do ograniczonych funkcji borelowskich
uzyskujemy przez klasę monotoniczną, używając
lematu \ref{lem:sde-kernel} dla mierzalności
prawej strony. Na końcu korzystamy z tower property:
$P_{s+t}f(x)=\E[P_tf(X_s^x)]=P_s(P_tf)(x)$.
\end{proof}

\begin{twierdzenie}[Mocna własność Markowa]
\label{th:sde-strongmarkov}
Przy założeniach twierdzenia \ref{th:sde-markov}
niech $\tau$ będzie prawie na pewno skończonym
czasem zatrzymania. Dla każdego $t\ge0$ i dowolnej
ograniczonej funkcji borelowskiej $f$ zachodzi
\[
\E[f(X_{\tau+t})\mid\calF_\tau]=P_tf(X_\tau)
\quad\text{prawie na pewno}.
\tag{6.4}\label{eq:sde-strongmarkov}
\]
\end{twierdzenie}
\begin{proof}
Najpierw niech $\tau$ przyjmuje skończenie wiele
wartości $s_1,\ldots,s_N$. Mnożąc
\eqref{eq:sde-markov} przez wskaźnik
$\ind{\{\tau=s_k\}}\in\calF_{s_k}$ i sumując,
otrzymujemy warunkową równość dla $\tau$.

Następnie niech $\tau\le R$ prawie na pewno.
Przybliżamy ją z góry czasami
$\tau_n=2^{-n}\lceil2^n\tau\rceil$,
które są czasami zatrzymania dla filtracji
prawostronnie ciągłej, przyjmują skończenie wiele
wartości i spełniają $\tau_n\downarrow\tau$.
Dla $A\in\calF_\tau$ mamy $A\in\calF_{\tau_n}$,
zatem
\[
\E[\ind A f(X_{\tau_n+t})]
=\E[\ind A P_tf(X_{\tau_n})].
\]
Jeżeli $f$ jest ograniczona i ciągła, to lemat
\ref{lem:sde-kernel}, ciągłość trajektorii $X$
i zbieżność zdominowana pozwalają przejść do granicy,
dając identyczną równość z $\tau$.
Dla ograniczonej funkcji borelowskiej $f$ stosujemy
ponownie twierdzenie o klasie monotonicznej.
Ograniczenie $\tau\le R$ usuwamy przez zastosowanie
wyniku do $\tau\wedge R$ i przejście $R\to\infty$;
na zdarzeniu $\{\tau\le R\}$ odpowiednie tożsamości
pokrywają się, zaś $\Prob(\tau>R)\to0$.
\end{proof}

\begin{przyklad}[Co proces Markowa naprawdę „pamięta”?]
Dla procesu Ornsteina--Uhlenbecka z zadania
\ref{zad:ou}, przy ustalonych $s,t\ge0$,
\[
\E[X_{s+t}\mid\calF_s]
=m+(X_s-m)e^{-\theta t},
\quad
\Var(X_{s+t}\mid\calF_s)
=\frac{\eta^2}{2\theta}(1-e^{-2\theta t}).
\]
To wynika z rozbicia całki stochastycznej w chwili
$s$ i niezależności przyszłych przyrostów. Znajomość
całej trajektorii do chwili $s$ nie zmienia tych
wzorów po ustaleniu $X_s$. Nie jest to jednak
niezależność kolejnych wartości procesu.
\end{przyklad}

\section{Generator dyfuzji i wzór Dynkina}
\begin{definicja}[Operator generatora]
Dla funkcji $f\in C_c^2(\R^d)$ określamy
\[
Lf(x)=\sum_{i=1}^d b_i(x)\partial_i f(x)
+\frac12\sum_{i,j=1}^d a_{ij}(x)\partial_{ij}f(x),
\qquad a(x)=\sigma(x)\sigma(x)^{\mathsf T}.
\tag{6.5}\label{eq:sde-generator}
\]
Operator $L$ nazywamy \emph{generatorem różniczkowym}
dyfuzji. Nie zakładamy w tej definicji, że funkcja
$f$ leży już w dziedzinie generatora silnej półgrupy
w przestrzeni funkcji ograniczonych.
\end{definicja}

\begin{twierdzenie}[Wzór Dynkina w skończonym horyzoncie]
\label{th:sde-dynkin}
Niech $X^x$ będzie rozwiązaniem \eqref{eq:timehom-sde}
z deterministycznym $x\in\R^d$ i niech
$f\in C_c^2(\R^d)$. Wówczas proces
\[
M_t^f=f(X_t^x)-f(x)-\int_0^tLf(X_s^x)\,\dd s
\tag{6.6}\label{eq:dynkin-mart}
\]
jest kwadratowo całkowalnym martyngałem na każdym
skończonym przedziale. Dla dowolnego ograniczonego
czasu zatrzymania $\tau$ otrzymujemy
\[
\E f(X_\tau^x)=f(x)+\E\int_0^\tau Lf(X_s^x)\,\dd s.
\tag{6.7}\label{eq:dynkin-optional}
\]
\end{twierdzenie}
\begin{proof}
Wzór Itô \ref{th:ito-multi} zastosowany do $f(X_t)$
daje
$\dd f(X_t)=Lf(X_t)\,\dd t+
\nabla f(X_t)\sigma(X_t)\,\dd W_t$.
Ponieważ $f$ ma zwarty nośnik, $\nabla f$ i
$D^2 f$ są ograniczone i zerują się poza pewnym
zbiorem zwartym. Współczynniki $b,\sigma$ są
ciągłe i lokalnie ograniczone, więc $Lf$ oraz
$\nabla f\,\sigma$ są \emph{globalnie ograniczone}.
W szczególności
$\E\int_0^T\norm{\nabla f(X_s)\sigma(X_s)}^2\,\dd s
\le T\sup_x\norm{\nabla f(x)\sigma(x)}^2<\infty$.
Stąd całka stochastyczna jest kwadratowo całkowalnym
martyngałem, co daje \eqref{eq:dynkin-mart}.
Dla ograniczonego $\tau$ stosujemy twierdzenie
o próbkowaniu opcjonalnym, a do skończenie zmiennej
w czasie całki względem $s$ — Fubiniego.
\end{proof}

\begin{wniosek}[Generator jako pochodna średniej]
\label{cor:sde-gen-limit}
Przy założeniach twierdzenia \ref{th:sde-dynkin},
dla każdej funkcji $f\in C_c^2(\R^d)$ i każdego
$x\in\R^d$ zachodzi
\[
\lim_{t\downarrow0}\frac{P_tf(x)-f(x)}{t}=Lf(x).
\tag{6.8}\label{eq:genlimit}
\]
\end{wniosek}
\begin{proof}
Z \eqref{eq:dynkin-mart} po wzięciu oczekiwanych mamy
$(P_tf(x)-f(x))/t=(1/t)\int_0^t\E Lf(X_s^x)\,\dd s$.
Funkcja $Lf$ jest ograniczona i ciągła, a
$X_s^x\to x$ prawie na pewno, gdy $s\downarrow0$.
Z twierdzenia o zbieżności zdominowanej wynika
$\E Lf(X_s^x)\to Lf(x)$.
Średnia Cesàro takiej funkcji również dąży do
$Lf(x)$, co dowodzi tezy.
\end{proof}

\begin{uwaga}[Co mierzy generator?]
Dla bardzo małego $h>0$ mamy w przybliżeniu
$\E[f(X_{t+h})\mid X_t=x]-f(x)\approx hLf(x)$.
Pierwsze pochodne w \eqref{eq:sde-generator}
opisują dryf, a pochodne drugiego rzędu są
konsekwencją wariacji kwadratowej szumu.
To związek z równaniami parabolicznymi, który
będzie rozwijany w rozdziałach o półgrupach i
wzorze Feynmana--Kaca.
\end{uwaga}

\begin{zadanie}[Wyjście dyfuzji z przedziału: generator i warunek brzegowy]
\label{zad:exit-drift}
Niech $a>0$, $\mu\in\R$, $\sigma_0>0$ i niech
$X_t=x+\mu t+\sigma_0W_t$, gdzie $x\in(0,a)$.
Dla $\tau=\inf\{t\ge0:X_t\notin(0,a)\}$ obliczyć
$\Prob_x(X_\tau=a)$.

\emph{Rozwiązanie.} W pierwszym kroku dowodzimy, że
$\tau<\infty$ prawie na pewno. Ustalmy $\delta>0$.
Niezależnie od położenia procesu w $(0,a)$ na początku
przedziału długości $\delta$, zdarzenie
$\{\mu\delta+\sigma_0(W_{s+\delta}-W_s)>a\}$
wymusza opuszczenie $(0,a)$ do chwili $s+\delta$.
Ma ono prawdopodobieństwo $q>0$, niezależne od
$\calF_s$. Iterowane warunkowanie daje
$\Prob_x(\tau>n\delta)\le(1-q)^n$.
W szczególności $\E_x\tau<\infty$ i $X_\tau\in\{0,a\}$.

Generator ma postać $L=\mu\partial_x+
(\sigma_0^2/2)\partial_{xx}$. Poszukujemy funkcji
$u\in C^2([0,a])$ spełniającej
$Lu=0$ na $(0,a)$, $u(0)=0$, $u(a)=1$.
Jeśli $\mu\ne0$, równanie
$u''+(2\mu/\sigma_0^2)u'=0$ po jednokrotnym
scałkowaniu daje
$u'(z)=C\exp(-2\mu z/\sigma_0^2)$;
po uwzględnieniu wartości brzegowych otrzymujemy
\[
u(x)=\frac{1-\exp(-2\mu x/\sigma_0^2)}
{1-\exp(-2\mu a/\sigma_0^2)}.
\]
Dla $\mu=0$ równanie $u''=0$ daje $u(x)=x/a$.

Aby wykazać, że to rzeczywiście szukane prawdopodobieństwo,
stosujemy wzór Itô do $u(X_{t\wedge\tau})$.
Ponieważ $Lu=0$ przed wyjściem, jest to lokalny
martyngał. Jest on jednak ograniczony przez $1$,
więc jest prawdziwym martyngałem i
$\E_xu(X_{t\wedge\tau})=u(x)$.
Dzięki prawie pewnej skończoności $\tau$ i ograniczoności
$u$ możemy zastosować zbieżność zdominowaną, uzyskując
$u(x)=\E_xu(X_\tau)=\Prob_x(X_\tau=a)$.
Nie wyciągamy zatem wyniku z samego formalnego rozwiązania
równania różniczkowego: wskazaliśmy, dlaczego jego
rozwiązanie przedstawia prawdopodobieństwo zdarzenia.
\end{zadanie}

\section{Równanie w przód: rozkład prawdopodobieństwa jako rozwiązanie słabe}
\begin{twierdzenie}[Słabe równanie Fokkera--Plancka]
\label{th:sde-fp}
Niech $X$ będzie rozwiązaniem \eqref{eq:timehom-sde}
z $X_0\in L^2(\Omega;\R^d)$, a
$\mu_t=\Law(X_t)$ dla $t\ge0$. Wówczas dla każdej
$f\in C_c^2(\R^d)$
\[
\int_{\R^d}f(x)\,\dd\mu_t(x)
=\int_{\R^d}f(x)\,\dd\mu_0(x)
+\int_0^t\int_{\R^d}Lf(x)\,\dd\mu_s(x)\,\dd s.
\tag{6.9}\label{eq:fp-weak}
\]
Jeśli ponadto dla $t>0$ rozkład $\mu_t$ ma gęstość
$\rho(t,\cdot)$, funkcja $\rho$ jest klasy $C^1$
w zmiennej $t$, a funkcje $b_i\rho$ i $a_{ij}\rho$
mają ciągłe pochodne przestrzenne odpowiednio rzędu
pierwszego i drugiego, przy lokalnej całkowalności
umożliwiającej różniczkowanie całek funkcji testowych,
to $\rho$ spełnia klasycznie równanie
\[
\partial_t\rho
=-\sum_i\partial_i(b_i\rho)
+\tfrac12\sum_{i,j}\partial_{ij}(a_{ij}\rho),
\qquad a=\sigma\sigma^{\mathsf T}.
\tag{6.10}\label{eq:fp-formal}
\]
Bez tych dodatkowych założeń \eqref{eq:fp-weak}
jest właściwym, matematycznie ścisłym sformułowaniem.
\end{twierdzenie}
\begin{proof}
Ze wzoru Itô i oszacowania w dowodzie twierdzenia
\ref{th:sde-dynkin} wynika
$\E f(X_t)-\E f(X_0)=\E\int_0^t Lf(X_s)\,\dd s$.
Ponieważ $Lf$ jest ograniczone, twierdzenie Fubiniego
daje \eqref{eq:fp-weak} po zapisaniu wartości
oczekiwanych jako całek względem rozkładów $\mu_s$.
W przypadku istnienia regularnej gęstości podstawiamy
$\dd\mu_s(x)=\rho(s,x)\,\dd\lambda_d(x)$.
Dla każdej funkcji testowej o zwartym nośniku możemy
różniczkować lewą stronę względem czasu (na dowolnym
zwartym podprzedziale dodatnich chwil). Po scałkowaniu
przez części w przestrzeni dostajemy
$\int f\,\partial_t\rho\,\dd\lambda_d
=\int f[-\partial_i(b_i\rho)+\frac12\partial_{ij}(a_{ij}\rho)]
\,\dd\lambda_d$ (z sumowaniem po indeksach).
Wszystkie te funkcje są ciągłe na rozpatrywanym obszarze.
Lemat podstawowy rachunku wariacyjnego, stosowany do
$C_c^\infty(\R^d)$, utożsamia ich wartości punktowo.
To uzasadnia postać klasyczną \eqref{eq:fp-formal},
a nie tylko jej zapis formalny.
\end{proof}

\begin{zadanie}[Ruch Browna i równanie ciepła]
Niech $X_t=x+W_t$ w $\R^d$. Wykazać, że jego gęstość
spełnia równanie ciepła z generatorem $\tfrac12\Delta$.

\emph{Rozwiązanie.} Tutaj $b=0$, $\sigma=I_d$,
a więc $Lf=\tfrac12\Delta f$. Dla $t>0$ gęstość
\[
p_t(x,y)=\frac1{(2\pi t)^{d/2}}
\exp\!\left(-\frac{\norm{y-x}^2}{2t}\right)
\]
spełnia $\partial_t p_t=\frac12\Delta_y p_t$:
bezpośrednie różniczkowanie daje
$\partial_t p_t=p_t(-d/(2t)+\norm{y-x}^2/(2t^2))$
oraz
$\Delta_y p_t=p_t(-d/t+\norm{y-x}^2/t^2)$.
Tożsamość jest zatem sprawdzona wprost.
Co ważne, w chwili $t=0$ prawo procesu jest miarą
Diraca $\delta_x$, nie zwykłą funkcją-gęstością.
Stąd naturalne jest słabe sformułowanie
\eqref{eq:fp-weak}.
\end{zadanie}

\begin{zadanie}[Rozkład stacjonarny procesu Ornsteina--Uhlenbecka]
Niech $X$ spełnia równanie z zadania \ref{zad:ou}
ze stałymi $\theta,\eta>0$ i $m\in\R$.
Znaleźć rozkład niezmienniczy i sprawdzić jego
niezmienniczość.

\emph{Rozwiązanie.} Z \eqref{eq:ou-solution}
wynika, że dla deterministycznego $X_0=x$ rozkład
$X_t$ jest normalny o średniej
$m+(x-m)e^{-\theta t}$ i wariancji
$\eta^2(1-e^{-2\theta t})/(2\theta)$.
Kandydatem na rozkład niezmienniczy jest
$\pi=\Normal(m,\eta^2/(2\theta))$.
Jeżeli $X_0\sim\pi$ i $X_0$ jest niezależny od
przyszłych przyrostów $W$ (co wynika z założeń
filtracji), to \eqref{eq:ou-solution} jest sumą
niezależnych zmiennych normalnych, a jej średnia
wynosi $m$ i wariancja
\[
\frac{\eta^2}{2\theta}e^{-2\theta t}
+\frac{\eta^2}{2\theta}(1-e^{-2\theta t})
=\frac{\eta^2}{2\theta}.
\]
Stąd $\Law(X_t)=\pi$ dla każdego $t\ge0$.
To sprawdza niezmienniczość. Zbieżność rozkładów
przy deterministycznym punkcie startowym, otrzymana
w poprzednim rozdziale, wyjaśnia fizyczną intuicję
powrotu do równowagi.
\end{zadanie}

\section{Porównywanie rozwiązań w jednym wymiarze}
\begin{twierdzenie}[Zasada porównania]
\label{th:sde-comparison}
Niech $b_1,b_2,\sigma:[0,T]\times\R\to\R$
będą borelowskie, globalnie Lipschitzowskie w
zmiennej przestrzennej jednostajnie względem czasu
i ograniczone w zerze. Zakładamy, że
$b_1(t,x)\le b_2(t,x)$ dla wszystkich $t,x$.
Niech $X,Y$ będą silnymi rozwiązaniami napędzanymi
tym samym procesem Wienera, odpowiednio,
\[
\dd X_t=b_1(t,X_t)\,\dd t+\sigma(t,X_t)\,\dd W_t,
\quad
\dd Y_t=b_2(t,Y_t)\,\dd t+\sigma(t,Y_t)\,\dd W_t,
\]
z $X_0,Y_0\in L^2(\Omega)$ i $X_0\le Y_0$ prawie
na pewno. Wówczas
$\Prob(X_t\le Y_t\ \text{dla wszystkich }t\in[0,T])=1$.
\end{twierdzenie}
\begin{proof}
Niech $D_t=X_t-Y_t$. Wybierzmy gładką, niemalejącą
funkcję $g:\R\to[0,1]$, dla której $g(z)=0$
przy $z\le0$ oraz $g(z)=1$ przy $z\ge1$.
Istnieje stała $C_g<\infty$ taka, że
$0\le g'(z)\le C_g\ind{\{0<z<1\}}$.
Dla $\varepsilon>0$ określmy
$\psi_\varepsilon(z)=\int_0^{z^+}g(u/\varepsilon)\,\dd u$.
Jest to funkcja klasy $C^2$, wypukła, spełniająca
\[
0\le\psi_\varepsilon'(z)\le\ind{\{z>0\}},\quad
0\le z^+-\psi_\varepsilon(z)\le\varepsilon,\quad
0\le\psi_\varepsilon''(z)\le
\frac{C_g}{\varepsilon}\ind{\{0<z<\varepsilon\}}.
\]
Możemy więc zastosować zwykły wzór Itô bez
odwoływania się do niesformułowanej wersji dla
funkcji z nieciągłą drugą pochodną. Po wzięciu
wartości oczekiwanych
całka stochastyczna znika, ponieważ
$\psi_\varepsilon'$ jest ograniczona, a
$\sigma(t,X_t)-\sigma(t,Y_t)$ jest kwadratowo
całkowalna. Otrzymujemy
\[
\begin{split}
\E\psi_\varepsilon(D_t)
=&\E\psi_\varepsilon(D_0)
+\E\int_0^t\psi_\varepsilon'(D_s)
[b_1(s,X_s)-b_2(s,Y_s)]\,\dd s\\
&+\frac12\E\int_0^t\psi_\varepsilon''(D_s)
[\sigma(s,X_s)-\sigma(s,Y_s)]^2\,\dd s.
\end{split}
\]
Pierwszy składnik jest zerowy, gdyż $D_0\le0$.
Na zbiorze $D_s>0$ mamy
$b_1(s,X_s)-b_2(s,Y_s)
\le b_1(s,X_s)-b_1(s,Y_s)\le L D_s^+$.
Dla wyrazu drugiego rzędu warunek Lipschitza daje
\[
\psi_\varepsilon''(D_s)
[\sigma(s,X_s)-\sigma(s,Y_s)]^2
\le \frac{C_gL^2}{\varepsilon}D_s^2
\ind{\{0<D_s<\varepsilon\}}\le C_gL^2\varepsilon.
\]
Wobec tego
$\E\psi_\varepsilon(D_t)
\le L\int_0^t\E D_s^+\,\dd s+
\tfrac12 C_gL^2\varepsilon t$.
Ponieważ $0\le D_t^+-\psi_\varepsilon(D_t)\le\varepsilon$,
po przejściu $\varepsilon\downarrow0$ mamy
$\E D_t^+\le L\int_0^t\E D_s^+\,\dd s$.
Lemat Gronwalla daje $\E D_t^+=0$ dla każdego $t$.
Na przeliczalnym zbiorze gęstym i dzięki ciągłości
trajektorii otrzymujemy żądaną nierówność
jednocześnie dla wszystkich chwil.
\end{proof}

\begin{uwaga}
Kluczowym założeniem jest \emph{wspólny} współczynnik
dyfuzji $\sigma$, a nie tylko porządek dryfów.
Jeśli różne szumy działają na oba procesy, nawet
początkowy porządek nie musi się utrzymać.
\end{uwaga}

\section{Metoda Eulera--Maruyamy: pierwszy schemat numeryczny}
\begin{definicja}[Schemat Eulera--Maruyamy]
\label{def:em}
Niech $T>0$, $h=T/N$ i $t_k=kh$, $k=0,\ldots,N$.
Dla funkcji $b,\sigma$ z rozdziału
\ref{ch:sde-existence} przyjmujemy
$\bar X_0=X_0$ oraz
\[
\bar X_{k+1}=\bar X_k+b(t_k,\bar X_k)h
+\sigma(t_k,\bar X_k)(W_{t_{k+1}}-W_{t_k}).
\tag{6.11}\label{eq:em}
\]
Ciągły interpolant \emph{stochastyczny} definiujemy
przez $X_t^h=X_0+\int_0^t b(\eta_h(s),X^h_{\eta_h(s)})\,\dd s
+\int_0^t\sigma(\eta_h(s),X^h_{\eta_h(s)})\,\dd W_s$,
gdzie $\eta_h(s)=t_k$ dla $s\in[t_k,t_{k+1})$.
W punktach siatki $X_{t_k}^h=\bar X_k$.
\end{definicja}

\begin{uwaga}[Dwa rodzaje błędu]
Błąd \emph{silny} kontroluje różnicę $X^h-X$
na tej samej realizacji szumu. Błąd \emph{słaby}
dotyczy różnicy wartości oczekiwanych
$\E f(X_T^h)-\E f(X_T)$ dla funkcji testowych.
Poniżej dowodzimy wyłącznie oszacowania silnego;
błędy słabe wymagają odrębnych założeń i analizy.
\end{uwaga}

\begin{twierdzenie}[Silna zbieżność metody Eulera--Maruyamy]
\label{th:em-rate}
Niech $X$ będzie rozwiązaniem SDE z twierdzenia
\ref{th:sde-global}, a $X^h$ interpolantem
\ref{def:em}, napędzanym tym samym procesem Wienera
i z tym samym $X_0\in L^2$. Zakładamy dodatkowo,
że istnieje $K$ takie, iż dla wszystkich $s,t\in[0,T]$
i $x\in\R^d$
\[
\norm{b(t,x)-b(s,x)}+
\norm{\sigma(t,x)-\sigma(s,x)}_{\mathrm F}
\le K(1+\norm{x})|t-s|^{1/2}.
\tag{6.12}\label{eq:time-holder}
\]
Wówczas dla pewnej stałej $C=C(T,L,K,\E\norm{X_0}^2)$,
niezależnej od $N$,
\[
\E\sup_{t\le T}\norm{X_t-X_t^h}^2\le Ch.
\tag{6.13}\label{eq:em-strong}
\]
Oznacza to silny rząd $1/2$ w normie
$\|\sup_{t\le T}\norm{X_t-X_t^h}\|_{L^2}$.
\end{twierdzenie}
\begin{proof}
\emph{Krok 1: jednostajne ograniczenie momentów schematu.}
Z definicji interpolantu, liniowego wzrostu
współczynników i lematu \ref{lem:sde-sup} otrzymujemy
\[
\E\sup_{r\le t}\norm{X_r^h}^2
\le C_0(1+\E\norm{X_0}^2)
+C_1\int_0^t\E\norm{X^h_{\eta_h(s)}}^2\,\dd s.
\]
Ponieważ $\eta_h(s)\le s$, ostatnia wartość
oczekiwana jest nie większa niż
$\E\sup_{u\le s}\norm{X_u^h}^2$.
Z lematu Gronwalla
\[
\sup_{N\ge1}\E\sup_{t\le T}\norm{X_t^h}^2<\infty.
\tag{6.14}\label{eq:em-moment}
\]

\emph{Krok 2: przyrost wewnątrz pojedynczego kroku.}
Dla $s\in[t_k,t_{k+1})$, warunkując względem
$\calF_{t_k}$ i korzystając z izometrii,
otrzymujemy
\[
\begin{split}
\E\norm{X_s^h-X_{t_k}^h}^2
&\le2(s-t_k)^2\E\norm{b(t_k,X_{t_k}^h)}^2\\
&\quad+2(s-t_k)\E\norm{\sigma(t_k,X_{t_k}^h)}_{\mathrm F}^2
\le C_2h,
\end{split}
\]
gdzie ostatni krok wynika z liniowego wzrostu i
\eqref{eq:em-moment}. W szczególności
$\int_0^T\E\norm{X_s^h-X_{\eta_h(s)}^h}^2\,\dd s
\le C_2Th$.

\emph{Krok 3: porównanie równań całkowych.}
Oznaczmy $e_t=X_t-X_t^h$. Odejmując równania,
a następnie stosując lemat \ref{lem:sde-sup},
rozkładamy różnice współczynników na trzy składniki:
\[
\begin{split}
\norm{b(s,X_s)-b(\eta_h(s),X^h_{\eta_h(s)})}
&\le L\norm{X_s-X_s^h}
+L\norm{X_s^h-X^h_{\eta_h(s)}}\\
&\quad+K(1+\norm{X^h_{\eta_h(s)}})\sqrt h,
\end{split}
\]
i analogiczną nierówność dla macierzy $\sigma$.
Po podniesieniu do kwadratu, scałkowaniu i
wykorzystaniu kroków 1--2 otrzymujemy
\[
\E\sup_{r\le t}\norm{e_r}^2
\le C_3\int_0^t\E\sup_{u\le s}\norm{e_u}^2\,\dd s
+C_4h.
\]
Lemat Gronwalla daje \eqref{eq:em-strong}.
Oszacowanie jest jednostajne na całym przedziale,
a nie tylko w punktach siatki.
\end{proof}

\begin{zadanie}[Schemat Eulera dla procesu Ornsteina--Uhlenbecka]
Dla równania $\dd X_t=\theta(m-X_t)\,\dd t+\eta\,\dd W_t$
z deterministycznym $X_0=x_0$ wyznaczyć rekurencję
Eulera, jej średnią i wariancję po $k$ krokach.

\emph{Rozwiązanie.} Ze wzoru \eqref{eq:em}
\[
\bar X_{k+1}=(1-\theta h)\bar X_k+\theta hm
+\eta\Delta_kW.
\]
Wprowadźmy $a=1-\theta h$. Po odjęciu $m$
i iteracji otrzymujemy
$\bar X_k-m=a^k(x_0-m)+\eta\sum_{j=0}^{k-1}
a^{k-1-j}\Delta_jW$. Przyrosty są niezależne,
mają wariancję $h$, więc
\[
\E\bar X_k=m+a^k(x_0-m),\qquad
\Var(\bar X_k)=\eta^2h\sum_{j=0}^{k-1}a^{2j}.
\]
Jeżeli $|a|<1$, to wariancja graniczna schematu
przy stałym $h$ wynosi $\eta^2h/(1-a^2)$,
co na ogół nie jest równe dokładnej wariancji
stacjonarnej $\eta^2/(2\theta)$.
To pokazuje, dlaczego dyskretyzacja może zmieniać
własności długoczasowe, nawet gdy zbiega na każdym
ustalonym przedziale wraz z $h\downarrow0$.
\end{zadanie}

\begin{center}
\begin{tikzpicture}[>=Latex,x=1cm,y=.7cm]
\draw[->] (0,0)--(10,0) node[right] {$t$};
\draw[->] (0,-.3)--(0,3) node[above] {$X_t$};
\draw[thick,blue!70!black] (0,.6)--(.6,.95)--(1.2,.72)--(1.8,1.48)
--(2.4,1.3)--(3,1.65)--(3.6,1.12)--(4.2,1.85)
--(4.8,1.43)--(5.4,1.7)--(6,2.35)--(6.6,1.9)
--(7.2,2.55)--(7.8,2.06)--(8.4,2.53)--(9,2.18);
\draw[very thick,orange!85!black] (0,.6)--(1.8,1.25)--(3.6,1.05)
--(5.4,1.65)--(7.2,2.32)--(9,2.2);
\foreach \x in {0,1.8,3.6,5.4,7.2,9} \draw[dashed,gray!55] (\x,0)--(\x,2.8);
\node[blue!70!black,anchor=west,scale=.82] at (5.8,2.95)
{trajektoria procesu (schemat)};
\node[orange!75!black,anchor=west,scale=.82] at (5.9,.6)
{przybliżenie na siatce};
\end{tikzpicture}
\par\smallskip\small Ilustracja idei dyskretyzacji; nie przedstawia obliczonej próbki konkretnego SDE.
\end{center}

\section{Kryterium nieeksplozji poza przypadkiem liniowego wzrostu}
\begin{twierdzenie}[Przykład kryterium Lapunowa]
\label{th:lyapunov-example}
Niech $W$ będzie jednowymiarowym procesem Wienera,
a $X$ lokalnym rozwiązaniem
\[
\dd X_t=-X_t^3\,\dd t+\dd W_t,\qquad X_0=x\in\R.
\]
Współczynniki są lokalnie Lipschitzowskie, lecz dryf
nie ma liniowego wzrostu. Mimo tego rozwiązanie jest
nieeksplozywne na każdym skończonym przedziale.
\end{twierdzenie}
\begin{proof}
Z obcięcia współczynników, jak w dowodzie twierdzenia
\ref{th:sde-local}, otrzymujemy zgodne lokalne
rozwiązania do czasów $\tau_R=\inf\{t:\abs{X_t}\ge R\}$.
Stosujemy wzór Itô do $V(x)=x^2$ i procesu
zatrzymanego w chwili $\tau_R$:
\[
\abs{X_{t\wedge\tau_R}}^2
=x^2+\int_0^{t\wedge\tau_R}(1-2X_s^4)\,\dd s
+2\int_0^{t\wedge\tau_R}X_s\,\dd W_s.
\]
Na przedziale przed $\tau_R$ integrand ostatniej
całki jest ograniczony przez $R$, a więc całka
ma średnią zero. Dostajemy
$\E\abs{X_{t\wedge\tau_R}}^2\le x^2+t$.
Na zdarzeniu $\{\tau_R\le t\}$ ciągłość trajektorii
daje $\abs{X_{t\wedge\tau_R}}=R$, zatem
$R^2\Prob(\tau_R\le t)\le x^2+t$.
Po przejściu $R\to\infty$ wnioskujemy, że czas
życia lokalnego rozwiązania jest nieskończony
prawie na pewno.
\end{proof}

\begin{uwaga}
Powyższy argument jest przykładem metody funkcji
Lapunowa. Ujemny składnik $-2X^4$ przeciwdziała
ucieczce rozwiązania do nieskończoności. W ogólnej
wersji poszukuje się nieujemnej funkcji $V$, która
rośnie przy $\norm{x}\to\infty$ i spełnia odpowiednią
nierówność dla generatora $LV$.
\end{uwaga}

\section{Dwa przykłady pokazujące granice twierdzeń}
\begin{przyklad}[Silna i słaba jednoznaczność dla równania deterministycznego]
Równanie $\dd X_t=2\sqrt{|X_t|}\,\dd t$, $X_0=0$,
ma zarówno rozwiązanie $X_t\equiv0$, jak i
$X_t=(t-c)_+^2$ dla każdego $c\ge0$.
Możemy nadać dowolne rozkłady losowemu opóźnieniu $c$,
otrzymując różne prawa na przestrzeni trajektorii.
Nie ma zatem ani jednoznaczności ścieżkowej, ani
jednoznaczności co do prawa. Oba pojęcia są różne,
choć w tym przykładzie zawodzą jednocześnie.
\end{przyklad}
\begin{przyklad}[Dlaczego sama zbieżność rozkładów nie daje zbieżności silnej?]
Niech $W$ i $\widehat W$ będą niezależnymi ruchami Browna
na tej samej przestrzeni. Dla dowolnego $t>0$ mają identyczny
rozkład $\Normal(0,t)$, ale
$\E|W_t-\widehat W_t|^2=2t$.
Nie można więc wnioskować o błędzie silnym metody
numerycznej ze zbieżności jej rozkładów jednowymiarowych.
W twierdzeniu \ref{th:em-rate} proces i przybliżenie
korzystają z \emph{tych samych} przyrostów Wienera.
\end{przyklad}

\section{Podsumowanie i dalsza droga}
W rozdziałach 5--6 zbudowaliśmy rozwiązania SDE i
wyprowadziliśmy ich podstawowe własności bez
zakładania z góry teorii rozwiązań. Globalny warunek
Lipschitza zapewnia istnienie silne i jednoznaczność
ścieżkową, a lokalny warunek wraz z odpowiednią
kontrolą wzrostu — brak eksplozji. Własność Markowa
i generator wiążą proces z równaniami różniczkowymi
cząstkowymi, natomiast zasada porównania i metoda
Eulera--Maruyamy opisują zależność od danych i
możliwość obliczeń numerycznych.

W dalszej części tomu przejdziemy do słabych
rozwiązań SDE, problemów martyngałowych, wzoru
Feynmana--Kaca, zmiany miary Girsanowa oraz
bardziej ogólnego rachunku dla półmartyngałów.

% ============================================================================
% ROZDZIAL 7 — Problem martyngalowy i Feynman--Kac
% ============================================================================
\chapter{Problem martyngałowy i twierdzenie Feynmana--Kaca}
\label{ch:martingale-fk}

\section{Od generatora do prawa procesu: motywacja}

W rozdziale \ref{ch:sde-properties} rozpoczęliśmy od procesu $X$, a następnie
wyznaczyliśmy jego generator $L$. Można odwrócić to rozumowanie:
\emph{czy operator $L$ i rozkład początkowy wyznaczają prawo procesu?}
Nie należy utożsamiać tego pytania z istnieniem silnego rozwiązania SDE.
Silne rozwiązanie wymaga zadania procesu Wienera i jego filtracji,
natomiast prawo procesu można próbować scharakteryzować bez wskazywania
konkretnego źródła szumu.

Niech $b:\R^d\to\R^d$ i $\sigma:\R^d\to\R^{d\times m}$ będą
ciągłe, a $a=\sigma\sigma^{\mathsf T}$. Rozważamy operator
\[
Lf(x)=\sum_i b_i(x)\partial_i f(x)
 +\frac12\sum_{i,j}a_{ij}(x)\partial_{ij}f(x),
 \qquad f\in C_c^2(\R^d).
\]
W równaniu SDE składnik $\int\nabla f(X_s)\sigma(X_s)\,\dd W_s$
ma średnią warunkową zero. Istotą konstrukcji jest więc fakt, że
$f(X_t)-f(X_0)-\int_0^tLf(X_s)\,\dd s$ jest martyngałem.

\begin{definicja}[Problem martyngałowy dla operatora $L$]
\label{def:mp}
Niech $\mu$ będzie miarą probabilistyczną na $\R^d$ i niech
$\Omega_C=C([0,\infty);\R^d)$ będzie przestrzenią ciągłych trajektorii
z jej $\sigma$-ciałem borelowskim (dla topologii zbieżności jednostajnej
na przedziałach zwartych). Niech $Y_t(\omega)=\omega(t)$ i
$\calF_t^Y=\sigma(Y_s:0\le s\le t)$.
Mówimy, że prawo probabilistyczne $\mathbf Q$ na $\Omega_C$
\emph{rozwiązuje problem martyngałowy $(L,\mu)$}, jeśli
$\Law_{\mathbf Q}(Y_0)=\mu$ i dla każdej $f\in C_c^2(\R^d)$
proces
\[
M_t^f=f(Y_t)-f(Y_0)-\int_0^tLf(Y_s)\,\dd s
\]
jest martyngałem względem $(\calF_t^Y)$ na każdym skończonym
przedziale czasu. Zakładamy przy tym całkowalność występujących
wyrażeń. Gdy $b,\sigma$ są lokalnie ograniczone, $Lf$ jest ograniczone
ze względu na zwarty nośnik $f$, więc całkowalność jest automatyczna.
\end{definicja}

\begin{uwaga}[Co jest, a czego nie ma w definicji?]
W definicji nie występuje proces Wienera. Warunek
martyngałowy opisuje \emph{warunkowe średnie przyrostów funkcji
stanu}. Własność ta jest silniejsza od równości
$\E f(Y_t)-\E f(Y_0)=\int_0^t\E Lf(Y_s)\,\dd s$:
ta ostatnia nie zawiera informacji o przeszłości $\calF_s^Y$.
Różne prawa mogą spełniać te same słabe równania dla rozkładów
jednowymiarowych; problem martyngałowy kontroluje całe prawo trajektorii.
\end{uwaga}

\begin{twierdzenie}[SDE rozwiązuje problem martyngałowy]
\label{th:sde-to-mp}
Niech $b,\sigma$ będą globalnie Lipschitzowskie i niech $X$ będzie
jedynym silnym rozwiązaniem
$\dd X_t=b(X_t)\,\dd t+\sigma(X_t)\,\dd W_t$
z $\Law(X_0)=\mu$ oraz $\E\norm{X_0}^2<\infty$.
Wówczas prawo procesu $X$ na $C([0,\infty);\R^d)$ rozwiązuje
problem martyngałowy $(L,\mu)$.
\end{twierdzenie}
\begin{proof}
Dla ustalonego $f\in C_c^2$ wzór Itô daje
$M_t^f=\int_0^t\nabla f(X_s)\sigma(X_s)\,\dd W_s$.
Ponieważ $\nabla f$ znika poza zwartym zbiorem, a $\sigma$ jest
tam ograniczona, integrand jest globalnie ograniczony.
Dlatego całka jest prawdziwym, kwadratowo całkowalnym
martyngałem względem filtracji napędzającej. Musimy jeszcze
uzasadnić przejście do mniejszej filtracji naturalnej.
Dla $s\le t$ i ograniczonej zmiennej $G$, mierzalnej względem
$\sigma(X_u:u\le s)$, mamy $\E[(M_t^f-M_s^f)G]=0$,
gdyż $G$ jest $\calF_s$-mierzalna. Jest to definicyjny
warunek martyngałowy względem filtracji naturalnej.
Transportując tożsamość przez odwzorowanie
$\omega\mapsto(X_t(\omega))_{t\ge0}$ otrzymujemy
wymagany warunek dla prawa kanonicznego.
\end{proof}

\section{Przykład fundamentalny: problem martyngałowy ruchu Browna}

Powyższe twierdzenie jest kierunkiem łatwiejszym: z równania
otrzymujemy problem martyngałowy. Pokażemy teraz na ważnym
przykładzie, że problem potrafi także \emph{wyznaczyć proces}.

\begin{lemat}[Rozszerzenie funkcji testowej przez obcięcie]
\label{lem:mp-cutoff}
Załóżmy, że $Y$ jest ciągłym procesem o wartościach w $\R^d$,
którego prawo rozwiązuje problem martyngałowy dla
$L=\frac12\Delta$ i rozkładu początkowego $\delta_x$.
Jeżeli $f\in C^2(\R^d)$ ma ograniczone wartości,
pochodne pierwszego i drugiego rzędu, to
$f(Y_t)-f(Y_0)-\frac12\int_0^t\Delta f(Y_s)\,\dd s$
jest martyngałem na każdym przedziale $[0,T]$.
\end{lemat}
\begin{proof}
Wybierzmy $\chi_R\in C_c^\infty(\R^d)$ równą $1$ na
kuli promienia $R$ i $0$ poza kulą promienia $2R$.
Funkcja $f_R=\chi_R f$ jest klasy $C_c^2$.
Niech $\tau_R=\inf\{t\ge0:\norm{Y_t}\ge R\}$.
Dla $t\le\tau_R$ zarówno $f_R(Y_t)=f(Y_t)$,
jak i $\Delta f_R(Y_t)=\Delta f(Y_t)$.
Stąd proces zatrzymany w $\tau_R$ jest martyngałem,
przez opcjonalne zatrzymywanie w ograniczonym czasie
$T\wedge\tau_R$. Dla każdej ciągłej trajektorii
$\sup_{0\le t\le T}\norm{Y_t}<\infty$, więc
$\tau_R>T$ od pewnego $R$ (zależnego od trajektorii).
Ponadto na $[0,T]$ zmienne
$|f(Y_{t\wedge\tau_R})-f(Y_0)
-\frac12\int_0^{t\wedge\tau_R}\Delta f(Y_s)\,\dd s|$
są ograniczone przez
$2\norm f_\infty+\frac T2\norm{\Delta f}_\infty$.
Dla każdego $A\in\calF_s^Y$ przechodzimy więc
do granicy w tożsamości
$\E[\ind A(M_{t\wedge\tau_R}^f-M_{s\wedge\tau_R}^f)]=0$
za pomocą zbieżności zdominowanej. Otrzymujemy tezę.
\end{proof}

\begin{twierdzenie}[Jednoznaczność problemu martyngałowego dla ruchu Browna]
\label{th:mp-bm}
Niech $x\in\R^d$ i $L=\frac12\Delta$ na $C_c^2(\R^d)$.
Jedynym prawem na $C([0,\infty);\R^d)$ rozwiązującym
problem martyngałowy $(L,\delta_x)$ jest prawo procesu
$x+W$, gdzie $W$ jest standardowym $d$-wymiarowym
ruchem Browna.
\end{twierdzenie}
\begin{proof}
Istnienie wynika z twierdzenia \ref{th:sde-to-mp},
przy $b=0$, $\sigma=I_d$. Dla jednoznaczności niech
$Y$ będzie procesem kanonicznym pod dowolnym innym
rozwiązaniem $\mathbf Q$. Ustalmy $u\in\R^d$ i
$f_u(y)=e^{i\langle u,y\rangle}$. W rzeczywistym
problemie martyngałowym stosujemy osobno części rzeczywistą
i urojoną. Pochodne $f_u$ są ograniczone, zatem lemat
\ref{lem:mp-cutoff} daje, że
\[
Z_t=f_u(Y_t)+\frac{\norm u^2}{2}\int_0^t f_u(Y_s)\,\dd s
\]
jest martyngałem zespolonym. Zróżniczkowanie iloczynu
$e^{\norm u^2t/2}f_u(Y_t)$ z deterministyczną funkcją
czasu (całkowanie przez części) pokazuje, że
$H_t=e^{\norm u^2t/2}f_u(Y_t)$ jest martyngałem.
Na $[0,T]$ jest on ograniczony przez $e^{\norm u^2T/2}$,
nie ma więc problemu z lokalizacją.
Dla $s\le t$ otrzymujemy
\[
\E_{\mathbf Q}[e^{i\langle u,Y_t-Y_s\rangle}\mid\calF_s^Y]
=e^{-\norm u^2(t-s)/2}.
\]
Lewa strona jest warunkową funkcją charakterystyczną
przyrostu. Prawo przyrostu jest więc normalne
$\Normal(0,(t-s)I_d)$ i nie zależy od $\calF_s^Y$.
Jednoznaczność funkcji charakterystycznych i
argument z klasą monotoniczną dowodzą niezależności
przyrostu od całej przeszłości. Ponieważ $Y_0=x$ i
trajektorie są ciągłe z definicji przestrzeni kanonicznej,
$Y-x$ jest standardowym ruchem Browna.
Prawo takiego procesu jest wyznaczone przez rozkłady
skończeniewymiarowe, co kończy dowód.
\end{proof}

\begin{przyklad}[Dlaczego testujemy wiele funkcji?]
Sam warunek $\E(Y_t-Y_s\mid\calF_s)=0$ nie wystarcza
do rozpoznania procesu Wienera: spełnia go również
wiele procesów skokowych i ciągłych martyngałów o innej
wariacji kwadratowej. Funkcje
$e^{i\langle u,x\rangle}$ kodują wszystkie
funkcje charakterystyczne przyrostów i pozwalają
rozpoznać całe ich rozkłady.
\end{przyklad}

\section{Półgrupy, potencjały i motywacja formuły Feynmana--Kaca}

Niech $X^x$ rozwiązuje jednorodne SDE o globalnie
Lipschitzowskich współczynnikach, jak w
rozdziale \ref{ch:sde-properties}. Dla ograniczonej
funkcji borelowskiej $f$ i $t\ge0$ znamy już półgrupę
$P_tf(x)=\E_xf(X_t)$. Jeżeli cząstka może zostać
\emph{zabita} z intensywnością zależną od położenia,
to każdą trajektorię trzeba dodatkowo ważyć przez
$e^{-\int_0^tV(X_s)\,\dd s}$. Prowadzi to do
nowej rodziny operatorów.

\begin{definicja}[Półgrupa Feynmana--Kaca]
\label{def:fk-semigroup}
Niech $V:\R^d\to\R$ będzie ograniczoną funkcją
borelowską. Dla ograniczonej borelowskiej $f$ określamy
\[
(T_t^Vf)(x)=\E_x\left[
\exp\left(-\int_0^tV(X_s)\,\dd s\right)f(X_t)\right].
\]
Całka w wykładniku jest zwykłą całką względem czasu.
Cały nawias oznacza funkcjonał losowej trajektorii.
\end{definicja}

\begin{twierdzenie}[Własność półgrupowa i oszacowanie]
\label{th:fk-semigroup}
Przy powyższych założeniach $T_0^V=I$,
$T_{s+t}^V=T_s^VT_t^V$ na ograniczonych
funkcjach borelowskich oraz
$\norm{T_t^Vf}_\infty\le e^{t\norm V_\infty}\norm f_\infty$.
Jeśli $V\ge0$, to $T_t^V$ jest dodatnią kontrakcją:
$0\le f\le1$ implikuje $0\le T_t^Vf\le1$.
\end{twierdzenie}
\begin{proof}
Oszacowanie wynika z nierówności
$|e^{-\int_0^t V(X_s)\,\dd s}|\le e^{t\norm V_\infty}$.
Dla własności półgrupowej rozbijamy
$\int_0^{s+t}V(X_r)\,\dd r$ na odcinki
$[0,s]$ i $[s,s+t]$. Warunkowo względem $\calF_s$
dalszy fragment procesu, po ustaleniu $X_s$, ma
prawo procesu startującego w $X_s$.
Stosujemy więc własność Markowa nie tylko do
jednego punktu końcowego, lecz do ograniczonego
mierzalnego funkcjonału całego fragmentu przyszłej trajektorii.
Wyjaśnijmy przejście, które jest istotne przy borelowskim $V$.
Dla każdej skończonej rodziny chwil $r_1,\ldots,r_k\in[0,t]$
własność Markowa i iterowanie wartości oczekiwanej warunkowej
dają stosowną tożsamość dla funkcji
$F(w)=\varphi(w(r_1),\ldots,w(r_k))$.
Zbiory takich funkcji cylindrycznych generują borelowskie
$\sigma$-ciało na $C([0,t];\R^d)$; twierdzenie o klasie
monotonicznej rozszerza tożsamość na wszystkie ograniczone
borelowskie funkcjonały $F$. Funkcjonał
$F(w)=\exp(-\int_0^tV(w(r))\,\dd r)f(w(t))$
jest borelowski: odwzorowanie $(r,w)\mapsto V(w(r))$
jest wspólnie mierzalne, a całkowanie ograniczonej funkcji
mierzalnej względem $r$ zachowuje mierzalność względem $w$
(z twierdzenia Fubiniego). Nie trzeba więc zakładać
zbieżności sum Riemanna dla potencjału jedynie borelowskiego.
Dostajemy
\[
\E_x\left[e^{-\int_0^{s+t}V(X_r)\,\dd r}f(X_{s+t})
\mid\calF_s\right]
=e^{-\int_0^sV(X_r)\,\dd r}(T_t^Vf)(X_s).
\]
Po wzięciu wartości oczekiwanych dostajemy
$T_{s+t}^Vf=T_s^V(T_t^Vf)$. Dodatniość dla $V\ge0$
jest bezpośrednia.
\end{proof}

\section{Wzór Feynmana--Kaca dla klasycznych rozwiązań PDE}

Poniższe twierdzenie ma celowo precyzyjny zakres.
Najpierw \emph{weryfikujemy} reprezentację dla
istniejącego klasycznego rozwiązania równania wstecznego.
Samo istnienie rozwiązania $C^{1,2}$ dla ogólnych
współczynników wymaga odrębnych twierdzeń teorii PDE.

\begin{twierdzenie}[Feynman--Kac: postać wsteczna z wymuszeniem]
\label{th:feynman-kac}
Niech $T>0$, $b:\R^d\to\R^d$,
$\sigma:\R^d\to\R^{d\times m}$ będą globalnie
Lipschitzowskie, $V:\R^d\to\R$ oraz
$q:[0,T]\times\R^d\to\R$ będą ograniczonymi funkcjami
ciągłymi, a $g:\R^d\to\R$ będzie ograniczona
i ciągła. Niech $u\in C^{1,2}([0,T)\times\R^d)
\cap C([0,T]\times\R^d)$ będzie ograniczona,
a jej pochodne pierwszego i drugiego rzędu względem
$x$ oraz pierwsza pochodna względem $t$ będą
ograniczone na $[0,T)\times\R^d$.
Załóżmy, że $u$ spełnia
\[
\partial_tu(t,x)+Lu(t,x)-V(x)u(t,x)+q(t,x)=0
\quad (t<T),\qquad u(T,x)=g(x).
\]
Wówczas dla każdego $t\in[0,T]$ i $x\in\R^d$
\[
\begin{split}
u(t,x)=\E\Big[&e^{-\int_t^T V(X_r^{t,x})\,\dd r}
                         g(X_T^{t,x})\\
 &+\int_t^T e^{-\int_t^s V(X_r^{t,x})\,\dd r}
                         q(s,X_s^{t,x})\,\dd s\Big],
\end{split}
\tag{7.1}\label{eq:fk}
\]
gdzie $(X_s^{t,x})_{s\ge t}$ jest rozwiązaniem SDE
startującym w $x$ w chwili $t$.
W klasie funkcji o podanej regularności i ograniczoności
rozwiązanie równania wstecznego jest jednoznaczne.
\end{twierdzenie}
\begin{proof}
Niech $X_s=X_s^{t,x}$ oraz
$A_s=\exp(-\int_t^sV(X_r)\,\dd r)$ dla $s\in[t,T]$.
Jest to ciągły proces o skończonej wariacji, spełniający
$\dd A_s=-V(X_s)A_s\,\dd s$, $A_t=1$.
Dla $s<T$ stosujemy wzór Itô do $u(s,X_s)$
i regułę iloczynu dla $A_su(s,X_s)$.
Kowariacja $A$ z ciągłym lokalnym martyngałem wynosi zero,
więc
\[
\dd[A_su(s,X_s)]
=A_s(\partial_su+Lu-Vu)(s,X_s)\,\dd s
 +A_s\nabla u(s,X_s)\sigma(X_s)\,\dd W_s.
\]
Po podstawieniu PDE dryf wynosi $-A_sq(s,X_s)\,\dd s$.
Dlatego proces
$A_su(s,X_s)+\int_t^sA_rq(r,X_r)\,\dd r$
jest lokalnym martyngałem. W rzeczywistości jest
prawdziwym martyngałem: wszystkie jego składniki poza
całką stochastyczną są ograniczone przez
$e^{T\norm V_\infty}(\norm u_\infty+T\norm q_\infty)$,
a sama całka ma z równania tę samą całkowalną,
a nawet ograniczoną na $[t,T)$ różnicę tych składników.
Po zatrzymaniu i zastosowaniu zbieżności zdominowanej
usuwamy lokalizację. Biorąc wartości oczekiwane
w chwili $T-\varepsilon$ i przechodząc $\varepsilon\downarrow0$
przez ciągłość $u$, $X$ i ograniczoność $A,q$, dostajemy
\eqref{eq:fk}.
Jeśli $u_1,u_2$ spełniają te same warunki, to oba
mają tę samą reprezentację prawostronną, a więc są równe.
\end{proof}

\begin{uwaga}[Wersja z potencjałem zależnym od czasu]
Ten sam dowód działa dla ograniczonej ciągłej funkcji
$V:[0,T]\times\R^d\to\R$, po zamianie
$V(X_s)$ na $V(s,X_s)$ w całce oraz w PDE.
Nie trzeba przy tym zakładać, że $V$ jest nieujemna;
ograniczoność zapewnia całkowalność czynnika wykładniczego.
\end{uwaga}

\section{Przykłady rachunkowe i rozwiązane zadania}

\begin{zadanie}[Równanie ciepła i jądro gaussowskie]
Niech $g\in C_b^2(\R)$, przy czym $g''$ jest
jednostajnie ciągła. Rozwiązać problem
$\partial_tu+\frac12\partial_{xx}u=0$ dla $t<T$
i $u(T,x)=g(x)$.

\emph{Rozwiązanie.} Weźmy $X_s^{t,x}=x+W_s-W_t$.
Kandydatem jest
\[
u(t,x)=\E g(x+W_{T-t})
=\int_\R g(y)\frac{e^{-(y-x)^2/(2(T-t))}}
{\sqrt{2\pi(T-t)}}\,\dd\lambda_1(y)
\quad(t<T).
\]
Dla $h>0$ można różniczkować pod całką względem $x$,
gdyż $g',g''$ są ograniczone: $u_x=\E g'(x+W_h)$,
$u_{xx}=\E g''(x+W_h)$.
Z własności półgrupy ciepła oraz jednostajnej ciągłości
$g''$ wynika
$\partial_h\E g(x+W_h)=\frac12\E g''(x+W_h)$.
Można to też otrzymać z wzoru Itô i zbieżności
zdominowanej po podzieleniu przez mały przyrost czasu.
Ponieważ $h=T-t$, mamy $u_t=-\frac12u_{xx}$.
Jednostajna ciągłość $g$ zapewnia warunek końcowy.
Zatem $u$ jest klasycznym rozwiązaniem z twierdzenia
\ref{th:feynman-kac}; otrzymaliśmy zarówno istnienie,
jak i jednoznaczność w rozważanej klasie.
\end{zadanie}

\begin{zadanie}[Stały potencjał: zabijanie trajektorii]
Niech $\kappa\ge0$, $V(x)=\kappa$, $q=0$ i $g\equiv1$.
Wykazać, że rozwiązaniem równania
$\partial_tu+\frac12\Delta u-\kappa u=0$,
$u(T,x)=1$, jest $u(t,x)=e^{-\kappa(T-t)}$.

\emph{Rozwiązanie.} Dla każdej trajektorii
$\int_t^TV(X_s)\,\dd s=\kappa(T-t)$.
Wzór \eqref{eq:fk} daje więc
$u(t,x)=\E e^{-\kappa(T-t)}=e^{-\kappa(T-t)}$.
Weryfikacja PDE jest bezpośrednia:
$u_t=\kappa u$, $\Delta u=0$, stąd
$u_t-\kappa u=0$. W interpretacji probabilistycznej
$u$ jest prawdopodobieństwem przeżycia niezależnego
czasu wykładniczego o intensywności $\kappa$ od $t$ do $T$.
\end{zadanie}

\begin{zadanie}[Równanie Poissona i średni czas wyjścia]
Niech $B_t$ będzie standardowym ruchem Browna w $\R^d$,
a $\tau_R=\inf\{t\ge0:\norm{x+B_t}\ge R\}$ dla $\norm x<R$.
Wyznaczyć $\E_x\tau_R$ i wyjaśnić, czemu równanie
$Lu=-1$ z warunkiem $u=0$ na sferze daje tę wielkość.

\emph{Rozwiązanie.} Funkcja
$u(x)=(R^2-\norm x^2)/d$ spełnia
$\frac12\Delta u=-1$ w kuli oraz $u=0$ na jej brzegu.
Z wzoru Itô dla procesu zatrzymanego wynika
$u(x+B_{t\wedge\tau_R})=u(x)-(t\wedge\tau_R)+M_t$,
gdzie $M_t$ jest całką stochastyczną z ograniczonym
integrandem na zatrzymanym procesie (bo $\nabla u$ jest
ograniczony w domkniętej kuli). Stąd
$\E(t\wedge\tau_R)\le u(x)$,
ponieważ $u\ge0$ na kuli. Zatem $\E\tau_R<\infty$.
Po przejściu $t\to\infty$ przez zbieżność zdominowaną
w pierwszym składniku oraz monotoniczną w drugim
otrzymujemy $u(x)=\E_x\tau_R$.
To ten sam mechanizm co we wzorze Feynmana--Kaca:
obserwowaną wielkość przedstawiamy przez funkcjonał
procesu i identyfikujemy odpowiedni problem brzegowy.
\end{zadanie}

\begin{center}
\begin{tikzpicture}[>=Latex,node distance=7mm]
\node[draw,rounded corners,fill=blue!6,align=center,minimum width=31mm] (sde)
{równanie SDE\\proces $X$};
\node[draw,rounded corners,fill=orange!8,align=center,right=13mm of sde,minimum width=34mm] (gen)
{generator $L$\\problem martyngałowy};
\node[draw,rounded corners,fill=green!7,align=center,right=13mm of gen,minimum width=31mm] (pde)
{równanie PDE\\wzór Feynmana--Kaca};
\draw[->,thick] (sde)--(gen);
\draw[->,thick] (gen)--(pde);
\draw[->,bend left=35,thin,dashed] (pde.north) to node[above,scale=.78]
{reprezentacja rozwiązania} (sde.north);
\end{tikzpicture}
\par\smallskip\small Związki między procesem, jego generatorem i równaniem wstecznym.
\end{center}

\begin{uwaga}[Prawidłowy zakres wyniku o rozwiązaniach PDE]
Twierdzenie \ref{th:feynman-kac} ma charakter weryfikacyjny:
każde rozwiązanie w wymienionej klasie ma podaną reprezentację.
Sama mierzalność lub ograniczoność potencjału nie zapewnia
istnienia klasycznego rozwiązania $C^{1,2}$ bez dodatkowej
teorii regularności PDE. Podobnie równanie słabe dla rozkładów
jednowymiarowych nie zastępuje dobrze postawionego problemu
martyngałowego dla prawa całego procesu.
\end{uwaga}

\section{Granice uzyskanej teorii}
Wykazaliśmy, że silne rozwiązanie SDE spełnia problem martyngałowy,
a dla operatora $\frac12\Delta$ udowodniliśmy także kierunek odwrotny
na poziomie prawa procesu. W pełnej teorii związek między dobrze
postawionym problemem martyngałowym a słabymi rozwiązaniami SDE
wymaga dalszych narzędzi, m.in. konstrukcji całek względem ogólnych
ciągłych lokalnych martyngałów. Rozwiniemy je po wprowadzeniu
półmartyngałów. Wzór Feynmana--Kaca został udowodniony jako
twierdzenie weryfikacyjne dla klasycznych rozwiązań PDE;
nie przypisujemy mu bez dodatkowych twierdzeń regularności,
której nie udowodniliśmy.


% ============================================================================
% ROZDZIAL 8 — Zmiana miary i Girsanov
% ============================================================================
\chapter{Zmiana miary probabilistycznej i twierdzenie Girsanowa}
\label{ch:girsanov}

\section{Jedna trajektoria, dwa prawa: dlaczego zmieniamy miarę?}

W równaniu $\dd X_t=b(X_t)\,\dd t+\sigma(X_t)\,\dd W_t$
część deterministyczna i losowa odgrywają różne role.
Nasuwa się pytanie, czy można zmienić \emph{prawdopodobieństwa
przypisywane trajektoriom} tak, aby część dryfu została
włączona do nowego procesu Wienera. Odpowiedź daje
\emph{twierdzenie Girsanowa}. Nie zmieniamy tu wartości $X_t(\omega)$
ani samego zbioru trajektorii; zmieniamy miarę, według której
liczymy oczekiwania.

Niech $(\Omega,\calF,(\calF_t)_{0\le t\le T},\Prob)$
będzie filtrowaną przestrzenią probabilistyczną spełniającą
warunki zwykłe. Jeżeli $Z_T>0$ prawie na pewno i
$\E_{\Prob}Z_T=1$, to możemy zdefiniować nową miarę
$\mathbf Q$ na $\calF_T$ przez
$\mathbf Q(A)=\E_{\Prob}[Z_T\ind A]$.
Tak więc dla całkowalnej względem $\mathbf Q$ zmiennej $Y$
\[
\E_{\mathbf Q}Y=\E_{\Prob}(Z_TY).
\]
Gęstość $Z_T$ \emph{przeważa} trajektorie: zdarzenia,
na których $Z_T$ jest duże, otrzymują większą masę.

\begin{lemat}[Warunkowy wzór Bayesa]
\label{lem:bayes-change}
Załóżmy, że $Z_T>0$ $\Prob$-prawie na pewno,
$\E_{\Prob}Z_T=1$, a $\mathbf Q=Z_T\cdot\Prob$
na $\calF_T$. Niech $Z_t=\E_{\Prob}[Z_T\mid\calF_t]$.
Wówczas $Z_t>0$ $\Prob$-prawie na pewno i dla
$0\le t\le T$ oraz $Y\in L^1(\mathbf Q)$
\[
\E_{\mathbf Q}[Y\mid\calF_t]
=\frac{\E_{\Prob}[Z_TY\mid\calF_t]}{Z_t}
\quad\mathbf Q\text{-prawie na pewno}.
\tag{8.1}\label{eq:bayes-measure}
\]
\end{lemat}
\begin{proof}
Ponieważ $Z_T>0$ prawie na pewno, gdyby
$A=\{Z_t=0\}\in\calF_t$ miało dodatnie prawdopodobieństwo,
to $\E[Z_T\ind A]=\E[Z_t\ind A]=0$,
co przeczyłoby dodatniości $Z_T$. Zatem mianownik
jest dodatni prawie na pewno. Oznaczmy prawą stronę
przez $U$. Dla każdego $A\in\calF_t$ mamy
\[
\E_{\mathbf Q}[\ind A U]
=\E_{\Prob}[Z_T\ind A U]
=\E_{\Prob}[Z_t\ind A U]
=\E_{\Prob}[\ind A\E_{\Prob}(Z_TY\mid\calF_t)]
=\E_{\mathbf Q}[\ind A Y].
\]
W drugim przejściu użyliśmy $\calF_t$-mierzalności $U$;
dla nieograniczonych zmiennych najpierw rozpatrujemy
część dodatnią i ujemną albo obcięcia i granicę monotoniczną.
Ponadto $\E_{\mathbf Q}|U|\le\E_{\mathbf Q}|Y|$.
Jest to dokładnie definicja wartości oczekiwanej warunkowej.
\end{proof}

\section{Wykładnik stochastyczny: skąd pochodzi gęstość?}

Niech $W=(W^1,\ldots,W^m)$ będzie $m$-wymiarowym
procesem Wienera, a $\theta=(\theta_t)_{0\le t\le T}$
procesem przewidywalnym o wartościach w $\R^m$.
Na razie załóżmy, że istnieje deterministyczna stała
$K<\infty$ taka, iż $\norm{\theta_t}\le K$
dla $\dd t\otimes\Prob$-prawie wszystkich $(t,\omega)$.
Definiujemy
\[
M_t=-\int_0^t\langle\theta_s,\dd W_s\rangle,
\qquad
Z_t=\exp\left(M_t-\frac12\int_0^t\norm{\theta_s}^2\,\dd s\right).
\tag{8.2}\label{eq:girsanov-density}
\]
Zastosowanie wzoru Itô do $e^{M_t-[M]_t/2}$
pokazuje, że
\[
\dd Z_t=-Z_t\langle\theta_t,\dd W_t\rangle,
\qquad Z_0=1.
\tag{8.3}\label{eq:density-sde}
\]
Poprawka $-\frac12[M]_t$ jest niezbędna: bez niej
wykładnik zwykłej całki martyngałowej ma na ogół
niezerowy dryf.

\begin{lemat}[Wykładnik ograniczonego integrandu jest martyngałem]
\label{lem:bounded-exponential}
Jeśli $\theta$ jest przewidywalny i
$\norm{\theta}\le K$ $\dd t\otimes\Prob$-prawie wszędzie,
to $Z$ z \eqref{eq:girsanov-density} jest ściśle dodatnim,
ciągłym, jednostajnie całkowalnym martyngałem na $[0,T]$.
Dla każdego $p>1$ zachodzi
\[
\sup_{0\le t\le T}\E Z_t^p
\le\exp\left(\frac{p(p-1)}2 K^2T\right),
\tag{8.4}\label{eq:z-moments}
\]
a w szczególności $\E Z_T=1$.
\end{lemat}
\begin{proof}
Najpierw ustalmy $p=2$. Wybierzmy czasy
$\tau_n=\inf\{t\le T:Z_t\ge n\}\wedge T$.
Z wzoru Itô dla $Z^2$ otrzymujemy
\[
Z_{t\wedge\tau_n}^2
=1-2\int_0^{t\wedge\tau_n}Z_s^2\langle\theta_s,\dd W_s\rangle
+\int_0^{t\wedge\tau_n}Z_s^2\norm{\theta_s}^2\,\dd s.
\]
Po zatrzymaniu całka stochastyczna ma zerową średnią,
a $\E Z_{t\wedge\tau_n}^2\le
1+K^2\int_0^t\E Z_{s\wedge\tau_n}^2\,\dd s$.
Gronwall daje oszacowanie jednostajne w $n$.
Ponieważ ciągłe trajektorie $Z$ są ograniczone na
$[0,T]$, mamy $\tau_n=T$ od pewnego losowego $n$;
Fatou i oszacowanie dają $\E Z_t^2\le e^{K^2T}$.
Rodzina zatrzymanych martyngałów $Z^{\tau_n}$
jest zatem jednostajnie całkowalna w każdej ustalonej
chwili (ma jednostajnie ograniczone drugie momenty).
Przechodząc do granicy w warunkowych tożsamościach
martyngałowych otrzymujemy $\E(Z_t\mid\calF_s)=Z_s$.
W szczególności $\E Z_t=1$.

Dla dowolnego $p>1$ zapisujemy
\[
Z_t^p=\underbrace{\exp\left(
-p\int_0^t\langle\theta_s,\dd W_s\rangle
-\frac{p^2}2\int_0^t\norm{\theta_s}^2\,\dd s\right)}_{Z_t^{(p)}}
\exp\left(\frac{p(p-1)}2\int_0^t\norm{\theta_s}^2\,\dd s\right).
\]
Proces $Z^{(p)}$ jest martyngałem na mocy właśnie
udowodnionego przypadku stosowanego do ograniczonego
integrandu $p\theta$; stąd
$\E Z_t^p\le e^{p(p-1)K^2T/2}\E Z_t^{(p)}$.
Wreszcie proces $Z_t=\E[Z_T\mid\calF_t]$ jest
jednostajnie całkowalny jako rodzina wartości oczekiwanych
warunkowych jednej zmiennej całkowalnej.
\end{proof}

\begin{uwaga}[Granica twierdzenia]
Założenie ograniczoności $\theta$ jest wystarczające,
lecz niekonieczne. W dalszej teorii używa się m.in.
warunku Novikova
$\E\exp(\frac12\int_0^T\norm{\theta_s}^2\,\dd s)<\infty$.
W tej iteracji pełny dowód przeprowadzamy przy
ograniczonym integrandzie; nie utożsamiamy lokalnej
martyngałowości wykładnika z zachowaniem jego średniej.
\end{uwaga}

\section{Twierdzenie Girsanowa z pełnym dowodem}

\begin{twierdzenie}[Girsanow przy ograniczonym dryfie]
\label{th:girsanov-bounded}
Niech $T>0$, niech $W$ będzie $m$-wymiarowym
procesem Wienera względem filtracji $(\calF_t)$
spełniającej warunki zwykłe i niech $\theta$ będzie
przewidywalnym procesem o wartościach w $\R^m$,
$\norm{\theta_t}\le K$ dla $\dd t\otimes\Prob$-prawie
wszystkich $(t,\omega)$. Niech $Z$ określa
\eqref{eq:girsanov-density} i niech $\mathbf Q$ na
$\calF_T$ będzie zadana przez $\dd\mathbf Q=Z_T\,\dd\Prob$.
Wówczas proces
\[
\widetilde W_t=W_t+\int_0^t\theta_s\,\dd s,
\qquad 0\le t\le T,
\tag{8.5}\label{eq:shifted-wiener}
\]
jest standardowym $m$-wymiarowym procesem Wienera
względem tej samej filtracji, ale pod miarą $\mathbf Q$.
Ponadto $\Prob$ i $\mathbf Q$ są równoważne na
$\calF_T$ (mają te same zbiory miary zero).
\end{twierdzenie}
\begin{proof}
\emph{Krok 1: zmiana miary jest prawidłowa.}
Z lematu \ref{lem:bounded-exponential} otrzymujemy
$Z_T>0$ i $\E Z_T=1$. Jest więc to gęstość miary
probabilistycznej. Ścisła dodatniość zapewnia
równoważność miar.

\emph{Krok 2: warunkowa funkcja charakterystyczna.}
Ustalmy $u\in\R^m$ i oznaczmy
\[
H_t=\exp\left(i\langle u,\widetilde W_t\rangle
                    +\frac12\norm u^2t\right).
\]
Na mocy wzoru Itô dla rzeczywistej i urojonej
części (albo przez bezpośredni rachunek pochodnych)
$\dd H_t=H_t(i\langle u,\dd W_t\rangle
+i\langle u,\theta_t\rangle\,\dd t)$.
Ze wzoru na iloczyn i \eqref{eq:density-sde} mamy
\[
\dd(Z_tH_t)=Z_tH_t\langle iu-\theta_t,\dd W_t\rangle.
\tag{8.6}\label{eq:girsanov-product}
\]
Wynika to z \emph{dokładnego skasowania} dwóch
członów dryfu: $+iZ_tH_t\langle u,\theta_t\rangle\,\dd t$
i $-iZ_tH_t\langle u,\theta_t\rangle\,\dd t$
pochodzącego z kowariacji.

\emph{Krok 3: lokalny martyngał jest prawdziwy.}
Na $[0,T]$ moduł $H_t$ jest ograniczony przez
$e^{\norm u^2T/2}$. Proces $Z$ ma ograniczone
drugie momenty z \eqref{eq:z-moments}; przy
ograniczonym $\theta$ kwadrat integrandu w
\eqref{eq:girsanov-product} ma skończoną wartość
oczekiwaną po scałkowaniu w czasie. Zatem $ZH$
jest (zespolonym) kwadratowo całkowalnym martyngałem.
Dla $s\le t$ dostajemy
\[
\E_{\Prob}[Z_tH_t\mid\calF_s]=Z_sH_s.
\]
Ponieważ $Z_t=\E_{\Prob}[Z_T\mid\calF_t]$,
wzór Bayesa \eqref{eq:bayes-measure} zastosowany do
$H_t$ daje $\E_{\mathbf Q}[H_t\mid\calF_s]=H_s$.
Po podzieleniu przez dodatni $\calF_s$-mierzalny
czynnik $e^{\norm u^2s/2}e^{i\langle u,\widetilde W_s\rangle}$
otrzymujemy
\[
\E_{\mathbf Q}[e^{i\langle u,\widetilde W_t-\widetilde W_s\rangle}
\mid\calF_s]
=e^{-\norm u^2(t-s)/2}.
\tag{8.7}\label{eq:girsanov-char}
\]
\emph{Krok 4: identyfikacja ruchu Browna.}
Prawa strona jest funkcją charakterystyczną rozkładu
$\Normal(0,(t-s)I_m)$ i nie zależy od $\calF_s$.
Jednoznaczność funkcji charakterystycznych i
argument z klasą monotoniczną dają niezależność
przyrostu od $\calF_s$. Proces $\widetilde W$ ma
ciągłe trajektorie i $\widetilde W_0=0$, zatem
spełnia definicję procesu Wienera pod $\mathbf Q$.
\end{proof}

\begin{uwaga}[Jak odczytywać znak?]
W naszej konwencji gęstość zawiera
$-\int\theta\,\dd W$, a nowy ruch Browna ma postać
$\widetilde W=W+\int\theta\,\dd t$.
Spotyka się także przeciwną konwencję, w której
w wykładniku stoi znak dodatni. Wtedy znak przed
dryfem w nowym procesie Wienera również się zmienia.
Nie ma sprzeczności, ale znaki muszą być uzgodnione.
\end{uwaga}

\section{Jak zmienia się równanie stochastyczne?}

\begin{wniosek}[Transformacja dryfu w SDE]
\label{cor:girsanov-sde}
Załóżmy warunki twierdzenia \ref{th:girsanov-bounded}
i niech proces $X$ spełnia pod $\Prob$, dla
$t\in[0,T]$, równanie
\[
X_t=X_0+\int_0^tb_s\,\dd s+\int_0^t\sigma_s\,\dd W_s,
\]
gdzie $b$ jest progresywnie mierzalny,
$\int_0^T\norm{b_s}\,\dd s<\infty$ prawie na pewno,
a $\sigma$ jest przewidywalny i
$\int_0^T\norm{\sigma_s}_{\mathrm F}^2\,\dd s<\infty$
prawie na pewno. Wtedy pod $\mathbf Q$ ten sam proces $X$
spełnia
\[
X_t=X_0+\int_0^t(b_s-\sigma_s\theta_s)\,\dd s
        +\int_0^t\sigma_s\,\dd\widetilde W_s.
\tag{8.8}\label{eq:girsanov-sde}
\]
\end{wniosek}
\begin{proof}
Tożsamość dla prostych procesów przewidywalnych wynika
z $W_t=\widetilde W_t-\int_0^t\theta_s\,\dd s$
po zastosowaniu liniowości sum całkowych.
Dla ogólnego $\sigma$ wprowadźmy czasy
$\tau_n=\inf\{t:\int_0^t\norm{\sigma_s}_{\mathrm F}^2\,\dd s\ge n\}\wedge T$.
Na przedziale zatrzymanym proces $\sigma$ należy do
$L^2(\dd t\otimes\Prob)$; ponadto zmiana miary może
pozbawić nas globalnego $L^2(\mathbf Q)$, ale przez
zatrzymanie dodatkowo w chwilach, gdy $Z$ lub $1/Z$
opuszczają przedział $[1/n,n]$, otrzymujemy równoważne
normy $L^2$ na zatrzymanym przedziale.
Przybliżając $\sigma$ prostymi procesami w obu normach,
otrzymujemy tożsamość całek najpierw lokalnie, a następnie
na całym $[0,T]$, dzięki wyczerpaniu przedziału przez
czasy lokalizujące. Ponieważ $\theta$ jest ograniczony,
całka $\int_0^t\sigma_s\theta_s\,\dd s$ jest określona
na mocy nierówności Cauchy'ego--Schwarza.
\end{proof}


\begin{uwaga}[Dokładne znaczenie lokalizacji przy zmianie miary]
W dowodzie wniosku \ref{cor:girsanov-sde} trzeba zachować
równoważność norm pod dwiema miarami, ale tylko na
zatrzymanych przedziałach. Dla ustalonego $n$ wybieramy
$\rho_n=\inf\{t\le T:Z_t\notin[1/n,n]\}\wedge T$.
Proces $Z$ jest ciągły i ściśle dodatni; na każdej trajektorii
$\inf_{t\le T}Z_t>0$ i $\sup_{t\le T}Z_t<\infty$,
a zatem $\rho_n=T$ od pewnego losowego $n$.
Jeśli $H$ jest przewidywalny i zeruje się po $\rho_n$,
to na $\calF_{\rho_n}$ zachodzi
$\dd\mathbf Q=Z_{\rho_n}\,\dd\Prob$, skąd
\[
  \frac1n\E_{\Prob}\int_0^{\rho_n}|H_s|^2\,\dd s
 \le \E_{\mathbf Q}\int_0^{\rho_n}|H_s|^2\,\dd s
 \le n\E_{\Prob}\int_0^{\rho_n}|H_s|^2\,\dd s.
\]
Istotne jest zastosowanie gęstości \emph{zatrzymanej}
$Z_{\rho_n}=\E[Z_T\mid\calF_{\rho_n}]$,
a nie gęstości końcowej $Z_T$ bez kontroli.
Przy dodatkowym zatrzymaniu kwadratowego czasu całkowania
powyższe porównanie umożliwia aproksymację tymi samymi
procesami prostymi pod obiema miarami i uzasadnia
podstawienie $\dd W_t=\dd\widetilde W_t-\theta_t\,\dd t$.
\end{uwaga}

\begin{uwaga}[Co pozostaje niezmienione?]
Równoważna zmiana miary nie zmienia własności
ścieżkowych zachodzących prawie na pewno, np. ciągłości
lub wariacji kwadratowej konkretnego procesu.
Zmienia natomiast rozkłady trajektorii i ich wartości
oczekiwane. Wzór \eqref{eq:girsanov-sde} pokazuje,
jak przesunięcie szumu zmienia dryf.
\end{uwaga}

\begin{przyklad}[Wspólny dryf i różne miary: dlaczego horyzont ma znaczenie?]
Dla stałego $c\ne0$ na każdym $[0,T]$ miary prawa
$W$ i $W+ct$ są równoważne, a ich gęstość jest dana
przez wykładnik $e^{cW_T-c^2T/2}$. Nie istnieje jednak
jedna gęstość na całym nieskończonym horyzoncie,
która czyniłaby te prawa równoważnymi na
$C([0,\infty);\R)$: z mocnego prawa wielkich liczb
dla przyrostów w chwilach całkowitych wynika
$W_n/n\to0$ prawie na pewno, podczas gdy
$(W_n+cn)/n\to c$. Zdarzenie
$\{\lim_{n\to\infty}w(n)/n=0\}$ ma więc prawdopodobieństwo
$1$ pod pierwszym prawem, a $0$ pod drugim.
Równoważność na każdym skończonym przedziale nie oznacza
równoważności na całej półprostej.
\end{przyklad}

\section{Zadania rozwiązane i kontrprzykłady intuicyjne}

\begin{zadanie}[Ruch Browna z dryfem jako ruch Browna pod nową miarą]
Niech $W$ będzie procesem Wienera pod $\Prob$, $c\in\R$
i $X_t=W_t+ct$ na $[0,T]$. Skonstruować miarę,
pod którą $X$ jest ruchem Browna, i porównać
prawdopodobieństwo zdarzenia $\{X_T>a\}$.

\emph{Rozwiązanie.} Wybieramy $\theta_t=c$ i
\[
Z_T=\exp(-cW_T-\tfrac12c^2T),\qquad
\dd\mathbf Q=Z_T\,\dd\Prob.
\]
Z twierdzenia Girsanowa $\widetilde W_t=W_t+ct=X_t$
jest procesem Wienera pod $\mathbf Q$. Stąd
\[
\Prob(X_T>a)=1-\Phi\left(\frac{a-cT}{\sqrt T}\right),
\qquad
\mathbf Q(X_T>a)=1-\Phi\left(\frac a{\sqrt T}\right).
\]
Są to różne liczby mimo tego, że $X_t(\omega)$
jest dosłownie tą samą funkcją na obu przestrzeniach
probabilistycznych. Na tym polega zmiana miary, a nie
„usunięcie dryfu z trajektorii”.
\end{zadanie}

\begin{zadanie}[Usuwanie ograniczonego dryfu przez zmianę miary]
Niech $b:\R\to\R$ będzie ograniczoną funkcją Lipschitzowską,
$\sigma_0>0$ i niech $X$ spełnia
$\dd X_t=b(X_t)\,\dd t+\sigma_0\,\dd W_t$, $X_0=x$.
Znaleźć miarę $\mathbf Q$, pod którą $X_t-x$
jest ruchem Browna przeskalowanym przez $\sigma_0$.

\emph{Rozwiązanie.} Chcemy, aby dryf w
\eqref{eq:girsanov-sde} zniknął, więc ustawiamy
$\theta_t=b(X_t)/\sigma_0$.
Proces $\theta$ jest przewidywalny i ograniczony,
ponieważ $X$ jest ciągły i adaptowany. Definiujemy
$Z_T$ przez \eqref{eq:girsanov-density}. Wówczas
$\widetilde W=W+\int\theta\,\dd t$ jest procesem
Wienera pod $\mathbf Q$, a
$\dd X_t=(b(X_t)-\sigma_0\theta_t)\,\dd t
+\sigma_0\,\dd\widetilde W_t=\sigma_0\,\dd\widetilde W_t$.
Stąd $X_t=x+\sigma_0\widetilde W_t$ jako procesy
na całym $[0,T]$. Jeśli dryf $b$ nie byłby ograniczony,
nie moglibyśmy bez dalszego sprawdzenia zastosować
udowodnionej wersji twierdzenia.
\end{zadanie}

\begin{zadanie}[Wartość oczekiwana pod inną miarą]
Niech $W$ będzie jednowymiarowym procesem Wienera,
a $\mathbf Q$ będzie określona przez
$Z_T=e^{-cW_T-c^2T/2}$, gdzie $c\in\R$.
Obliczyć $\E_{\mathbf Q}W_T$ i sprawdzić wynik
dwiema metodami.

\emph{Rozwiązanie.} Ponieważ $\widetilde W_t=W_t+ct$
jest ruchem Browna pod $\mathbf Q$,
$\E_{\mathbf Q}W_T=-cT$.
Z drugiej strony
$\E_{\mathbf Q}W_T=\E_{\Prob}(W_Te^{-cW_T-c^2T/2})$.
Dla $W_T\sim\Normal(0,T)$ różniczkujemy funkcję
tworzącą momenty $\E e^{uW_T}=e^{u^2T/2}$:
$\E(W_Te^{uW_T})=uT e^{u^2T/2}$.
Po podstawieniu $u=-c$ otrzymujemy
$(-cT)e^{c^2T/2}e^{-c^2T/2}=-cT$.
Dwie metody dają ten sam wynik, a druga pokazuje
jak w praktyce wykorzystuje się gęstość Radona--Nikodyma.
\end{zadanie}

\begin{przyklad}[Dlaczego zmiana miary musi zależeć od całej filtracji?]
Jeżeli $\theta_t$ zależy od $W_s$ dla $s>t$, to nie
jest procesem przewidywalnym i całka Itô
$\int\theta\,\dd W$ w ogólności nie ma sensu w
naszej konstrukcji. Próba „odwrócenia znaku dryfu”
na podstawie przyszłego wyniku nie jest poprawną
strategią Girsanowa. Konieczna jest adaptacja
do bieżącej informacji.
\end{przyklad}

\begin{center}
\begin{tikzpicture}[>=Latex,node distance=8mm]
\node[draw,rounded corners,fill=blue!6,minimum width=35mm,align=center] (p)
{$\Prob$\\$W$ jest procesem Wienera};
\node[draw,rounded corners,fill=orange!8,minimum width=36mm,align=center,right=18mm of p] (q)
{$\mathbf Q=Z_T\cdot\Prob$\\$\widetilde W=W+\int\theta\,\dd t$ jest procesem Wienera};
\draw[->,very thick] (p)--node[above,align=center,scale=.85]
{ważenie trajektorii\\przez $Z_T$}(q);
\end{tikzpicture}
\par\smallskip\small Zmiana miary przekształca prawo szumu, nie jego realizacje.
\end{center}

\section{Co dalej?}
Udowodniona wersja twierdzenia Girsanowa ma przejrzysty
zakres: ograniczony przewidywalny dryf, skończony horyzont,
ciągły szum Wienera. Wersje z integrandami nieograniczonymi,
warunkami Novikova i Kazamakiego oraz z miarami losowymi
skoków wymagają odrębnych twierdzeń.
W kolejnych rozdziałach wrócimy do zmiany miary przy
rozważaniu przestrzeni trajektorii i twierdzenia
Camerona--Martina, a później przy procesach Lévy'ego
oraz równaniach ze skokami.

% ===========================================================================
% ROZDZIAL 9 — Miara Wienera i Cameron--Martin
% ===========================================================================
\chapter{Miara Wienera i twierdzenie Camerona--Martina}
\label{ch:cameron-martin}

\begin{uwaga}[O wersjach funkcjonału gaussowskiego]
Wzór Camerona--Martina zawiera $I(h')$, czyli element
$L^2(C_0([0,T]),\mu_T)$, a nie kanonicznie wyznaczoną
liczbę dla każdej trajektorii $w$. Wybór dowolnego
borelowskiego reprezentanta tej klasy daje tę samą
pochodną Radona--Nikodyma prawie wszędzie. W ogólności
nie istnieje ciągły funkcjonał liniowy na całej
przestrzeni $C_0([0,T])$, który realizowałby $I(h')$
dla dowolnego $h'\in L^2$.
\end{uwaga}

\section{Dlaczego miara na przestrzeni trajektorii jest innym obiektem niż całka Itô?}

Całka $\int_0^T H_t\,\dd W_t$ jest zmienną losową, a w przypadku
całkowania na zmiennym przedziale --- procesem. Natomiast dla funkcjonału
$F:C_0([0,T])\to\R$ zapis $\int F(w)\,\dd\mu_T(w)$ jest zwykłą całką
Lebesgue'a względem \emph{miary na przestrzeni funkcji}. Są to dwa
różne rodzaje całkowania. Pierwszy wykorzystuje przewidywalność i przyrosty
procesu, drugi jest uśrednieniem po wszystkich możliwych trajektoriach.

Na przykład $F(w)=\exp(-\int_0^T w(s)^2\,\dd s)$ jest funkcjonałem
trajektorii: dla każdej ciągłej $w$ całka po czasie istnieje klasycznie.
Wówczas $\int F\,\dd\mu_T$ jest liczbą z przedziału $(0,1]$.
Nie ma tu całki stochastycznej. W innym przykładzie
$F(w)=w(T)^2$ dostaniemy $\int F\,\dd\mu_T=T$.

\begin{definicja}[Kanoniczna przestrzeń Wienera]
\label{def:canonical-wiener}
Niech $T>0$ i
$C_0([0,T])=\{w\in C([0,T];\R):w(0)=0\}$ z normą
$\|w\|_\infty=\sup_{0\le t\le T}|w(t)|$ i
borelowskim $\sigma$-ciałem $\mathcal B_T$.
\emph{Miara Wienera} $\mu_T$ jest prawem standardowego procesu Wienera
na $C_0([0,T])$. Na przestrzeni kanonicznej proces współrzędnych
$B_t(w)=w(t)$ jest procesem Wienera pod miarą $\mu_T$.
\end{definicja}

W tomie II skonstruowaliśmy proces Wienera o ciągłych trajektoriach;
zatem odwzorowanie $\omega\mapsto(W_t(\omega))_{t\in[0,T]}$ jest
mierzalne jako odwzorowanie w $C_0([0,T])$ (wystarcza mierzalność
współrzędnych w chwilach wymiernych). Obraz miary probabilistycznej
przez to odwzorowanie jest $\mu_T$. Konstrukcja nie wymaga formalnej
„miary $\prod_t\dd w(t)$”.

\begin{twierdzenie}[Rozkłady skończeniewymiarowe i jednoznaczność miary]
\label{th:wiener-cyl}
Niech $0<t_1<\cdots<t_n\le T$, $x_0=0$, $t_0=0$, oraz
$p_t(z)=(2\pi t)^{-1/2}\exp(-z^2/(2t))$ dla $t>0$.
Dla każdego ograniczonego borelowskiego $g:\R^n\to\R$ zachodzi
\[
\int_{C_0([0,T])}g(w(t_1),\ldots,w(t_n))\,\dd\mu_T(w)
=\int_{\R^n}g(x_1,\ldots,x_n)
\prod_{j=1}^n p_{t_j-t_{j-1}}(x_j-x_{j-1})
\,\dd\lambda_n(x).
\]
Miara $\mu_T$ jest jedyną miarą probabilistyczną na
$C_0([0,T])$, dla której zachodzą te tożsamości.
\end{twierdzenie}
\begin{proof}
Dla procesu Wienera wektor przyrostów
$(W_{t_1},W_{t_2}-W_{t_1},\ldots,W_{t_n}-W_{t_{n-1}})$
ma niezależne współrzędne normalne z wariancjami
$t_1,t_2-t_1,\ldots,t_n-t_{n-1}$. Liniowa transformacja
$(y_1,\ldots,y_n)\mapsto(y_1,y_1+y_2,\ldots,y_1+\cdots+y_n)$
ma wyznacznik równy $1$. Wzór na gęstość wynika więc ze
zmiany zmiennych w $\R^n$.

Dla jednoznaczności wystarczy zauważyć, że funkcje
$e_q(w)=w(q)$, gdzie $q\in[0,T]\cap\Q$ (oraz $q=T$),
generują $\mathcal B_T$. Istotnie, dla $w,v\in C_0([0,T])$
z ciągłości wynika
$\|w-v\|_\infty=\sup_{q\in D}|w(q)-v(q)|$, gdzie
$D=([0,T]\cap\Q)\cup\{T\}$. Otwarte kule w normie
jednostajnej są zatem mierzalne względem
$\sigma(e_q:q\in D)$. Ponieważ $C_0([0,T])$ jest
separowalną przestrzenią metryczną, każde otwarte podzbiory
są przeliczalnymi sumami takich kul. Z drugiej strony
wszystkie $e_q$ są ciągłe. Stąd równość $\sigma$-ciał.
Zbiory cylindryczne w wymiernych chwilach tworzą
$\pi$-układ generujący $\mathcal B_T$; dwie miary
zgodne na nim są równe z twierdzenia o $\pi$- i $\lambda$-układach.
\end{proof}

\begin{przyklad}[Uśrednianie funkcjonału trajektorii]
Dla $F(w)=\int_0^T w(t)^2\,\dd t$ twierdzenie Tonellego daje
\[
\int F(w)\,\dd\mu_T(w)
=\int_0^T\E B_t^2\,\dd t
=\int_0^Tt\,\dd t=\frac{T^2}{2}.
\]
Wykorzystaliśmy tutaj \emph{dwie} miary: Lebesgue'a w czasie i
Wienera na przestrzeni trajektorii. Ich kolejność można zamienić,
ponieważ funkcja podcałkowa jest nieujemna.
\end{przyklad}

\section{Deterministyczna całka Wienera jako gaussowski funkcjonał trajektorii}

\begin{twierdzenie}[Izometria i funkcje charakterystyczne]
\label{th:pw-map}
Dla każdego $f\in L^2([0,T])$ istnieje zmienna losowa
$I(f)\in L^2(C_0([0,T]),\mu_T)$, jednoznaczna w $L^2$,
taka że dla funkcji schodkowej
$f=\sum_{j=0}^{n-1}c_j\ind{(t_j,t_{j+1}]}$ zachodzi
$I(f)(w)=\sum_j c_j(w(t_{j+1})-w(t_j))$.
Odwzorowanie $I:L^2([0,T])\to L^2(\mu_T)$ jest liniowe i
\[
\E_{\mu_T}[I(f)I(g)]=\int_0^T f(t)g(t)\,\dd t,
\qquad
\E_{\mu_T}e^{iuI(f)}=e^{-u^2\|f\|_2^2/2}.
\]
Dla dowolnych $f_1,\ldots,f_k\in L^2([0,T])$ wektor
$(I(f_1),\ldots,I(f_k))$ jest gaussowski.
\end{twierdzenie}
\begin{proof}
Dla schodkowych $f,g$ pierwsza tożsamość wynika z niezależności
przyrostów i z faktu, że kowariancja dwóch przyrostów ruchu Browna
jest równa długości przecięcia odpowiadających im przedziałów.
W szczególności $\|I(f)-I(g)\|_{L^2(\mu_T)}=\|f-g\|_2$.
Jeśli $f_n$ są schodkowe i zbiegają w $L^2$ do $f$,
to $I(f_n)$ jest ciągiem Cauchy'ego w $L^2(\mu_T)$;
jego granica nie zależy od wybranej aproksymacji.
Wzór na kowariancję przechodzi do granicy przez nierówność
Cauchy'ego--Schwarza. Każda kombinacja
$\sum_j u_jI(f_j)=I(\sum_j u_jf_j)$ jest granicą w $L^2$
zmiennych normalnych o wariancjach zbiegających do
$\|\sum_j u_jf_j\|_2^2$. Zbieżność funkcji charakterystycznych
i twierdzenie o ciągłości dowodzą normalności granicy.
To samo rozumowanie dla wszystkich liniowych kombinacji
daje wspólną gaussowskość wektora.
\end{proof}

\begin{uwaga}[Gdzie ukryta jest subtelność mierzalności?]
Dla każdej pojedynczej funkcji $f$ otrzymujemy klasę
równoważności $I(f)$, a nie jeden kanoniczny wzór punktowy,
który działałby jednocześnie dla wszystkich $f$ i wszystkich
trajektorii $w$. Nie możemy więc pisać bez wyjaśnienia
$I(f)(w)=\int f(t)w'(t)\,\dd t$: typowa trajektoria
Wienera nie posiada pochodnej. W dowodach poniżej każda
występująca całka stochastyczna jest zdefiniowana jako
granica w $L^2$, a wzory dotyczą równości prawie na pewno.
\end{uwaga}

\section{Przestrzeń Camerona--Martina: które przesunięcia są dopuszczalne?}

\begin{definicja}[Przestrzeń Camerona--Martina]
\label{def:CMspace}
Niech
\[
\calH_T=\left\{h\in C_0([0,T]):
\begin{array}{l}
 h\text{ jest absolutnie ciągła oraz}\\[-2pt]
 h'\in L^2([0,T])
\end{array}\right\}.
\]
Wprowadzamy iloczyn skalarny
$\langle h,k\rangle_{\calH_T}=\int_0^T h'(s)k'(s)\,\dd s$
oraz normę $\|h\|_{\calH_T}=\|h'\|_2$.
\end{definicja}

\begin{twierdzenie}[Struktura Hilberta i jądro reprodukujące]
\label{th:cm-hilbert}
Przestrzeń $\calH_T$ jest przestrzenią Hilberta,
a odwzorowanie $D:h\mapsto h'$ jest izometrią
$\calH_T$ na $L^2([0,T])$. Ponadto
\[
|h(t)-h(s)|\le\sqrt{|t-s|}\,\|h\|_{\calH_T}
\quad(s,t\in[0,T]),
\]
a funkcjonał wartości w chwili $t$ jest reprezentowany
przez $k_t(s)=\min(s,t)$:
$h(t)=\langle h,k_t\rangle_{\calH_T}$.
\end{twierdzenie}
\begin{proof}
Z definicji absolutnej ciągłości i podstawowego twierdzenia
rachunku całkowego dla całki Lebesgue'a mamy
$h(t)=\int_0^t h'(r)\,\dd r$. Odwrotnie każdej
$f\in L^2([0,T])$ odpowiada funkcja
$h_f(t)=\int_0^t f(r)\,\dd r$, która jest absolutnie
ciągła, zaczyna się od zera i ma pochodną $f$ prawie
wszędzie. Otrzymujemy izometryczną bijekcję z zupełną
przestrzenią $L^2$, a więc zupełność $\calH_T$.
Nierówność wynika z Cauchy'ego--Schwarza:
$|\int_s^t h'(r)\,\dd r|\le\sqrt{|t-s|}\,\|h'\|_2$.
Wreszcie $k_t'(r)=\ind{(0,t)}(r)$ niemal wszędzie, zatem
$\langle h,k_t\rangle_{\calH_T}=\int_0^t h'(r)\,\dd r=h(t)$.
\end{proof}

Pojawienie się $\min(s,t)$ nie jest przypadkowe. To jednocześnie
kowariancja procesu Wienera $\E B_sB_t$ i jądro reprodukujące
przestrzeni funkcji, o które wolno przesuwać miarę Wienera
bez zmiany jej zbiorów miary zero. Ten fakt stanie się treścią
następnego twierdzenia, a nie założeniem.

\begin{lemat}[Aproksymacja w przestrzeni Camerona--Martina]
\label{lem:cm-approx}
Dla $h\in\calH_T$ i podziału dyadycznego
$t_j^{(n)}=jT2^{-n}$ niech $h_n$ będzie interpolacją liniową
wartości $h(t_j^{(n)})$. Wtedy $h_n\in\calH_T$,
$h_n'$ jest schodkowa oraz
$\|h_n'-h'\|_2\to0$ i $\|h_n-h\|_\infty\to0$.
\end{lemat}
\begin{proof}
Na przedziale $(t_j^{(n)},t_{j+1}^{(n)})$ pochodna $h_n$
jest średnią funkcji $h'$ na tym przedziale:
\[
 h_n'|_{(t_j^{(n)},t_{j+1}^{(n)})}
 =\frac{1}{t_{j+1}^{(n)}-t_j^{(n)}}
 \int_{t_j^{(n)}}^{t_{j+1}^{(n)}}h'(s)\,\dd s.
\]
Operator zastępowania funkcji przez jej średnie przedziałowe
jest kontrakcją w $L^2$ z nierówności Jensena.
Dla funkcji ciągłych średnie zbiegają jednostajnie,
a funkcje ciągłe są gęste w $L^2$. Stąd dla każdego
$f\in L^2$ przybliżenia przedziałowe zbiegają w $L^2$ do $f$.
Biorąc $f=h'$ dostajemy pierwszą zbieżność. Druga wynika
z $\sup_t|h_n(t)-h(t)|\le\sqrt T\|h_n'-h'\|_2$.
\end{proof}

\section{Twierdzenie Camerona--Martina: gęstość przesuniętej miary}

Dla $h\in C_0([0,T])$ oznaczmy przez $\tau_h(w)=w+h$
przesunięcie trajektorii, a przez
$\mu_T^h=\mu_T\circ\tau_h^{-1}$ prawo $B+h$.
Nie należy mylić $\mu_T^h$ z przesunięciem samego argumentu
w całce Itô. Teraz przesuwamy \emph{całą miarę}.

\begin{lemat}[Skończeniewymiarowa translacja gaussowska]
\label{lem:gauss-shift}
Niech $Y\sim\Normal(0,D)$ w $\R^n$, gdzie $D$ jest
symetryczną dodatnio określoną macierzą, i niech
$a\in\R^n$. Dla każdej ograniczonej borelowskiej
$F:\R^n\to\R$ mamy
\[
\E F(Y+a)=\E\left[F(Y)\exp\left(
 a^{\mathsf T}D^{-1}Y-\frac12a^{\mathsf T}D^{-1}a
\right)\right].
\]
\end{lemat}
\begin{proof}
Gęstość $Y+a$ w punkcie $y$ jest proporcjonalna do
$\exp[-\frac12(y-a)^{\mathsf T}D^{-1}(y-a)]$.
Iloraz tej gęstości przez gęstość $Y$ jest równy
$\exp(a^{\mathsf T}D^{-1}y-\frac12a^{\mathsf T}D^{-1}a)$.
Pomnożenie przez $F$ i scałkowanie kończy dowód.
\end{proof}

\begin{twierdzenie}[Cameron--Martin na przedziale skończonym]
\label{th:cameron-martin}
Niech $T>0$ i $h\in\calH_T$.
Miary $\mu_T^h$ i $\mu_T$ są równoważne, a
\[
\frac{\dd\mu_T^h}{\dd\mu_T}(w)
=\exp\left(I(h')(w)-\frac12\|h'\|_2^2\right)
\quad\text{dla }\mu_T\text{-prawie każdego }w.
\tag{9.1}\label{eq:cm-density}
\]
Innymi słowy, dla każdego ograniczonego borelowskiego
funkcjonału $F:C_0([0,T])\to\R$ zachodzi
\[
\int F(w+h)\,\dd\mu_T(w)
=\int F(w)\exp\left(I(h')(w)-\frac12\|h'\|_2^2\right)
\,\dd\mu_T(w).
\tag{9.2}\label{eq:cm-identity}
\]
\end{twierdzenie}
\begin{proof}
\emph{Krok 1: przesunięcia łamanymi.} Załóżmy najpierw, że
$h'$ jest funkcją schodkową stałą na przedziałach podziału
$0=t_0<\cdots<t_m=T$. Dla funkcjonału cylindrycznego
$F(w)=\varphi(w(s_1),\ldots,w(s_k))$ uzupełnijmy podział
o wszystkie $s_j$. Wektor niezależnych przyrostów $B$ na
nowym podziale ma diagonalną macierz kowariancji, której
przekątna składa się z długości przedziałów.
Przesunięcie przez $h$ dodaje do każdego przyrostu
$\Delta h_j=\int_{t_j}^{t_{j+1}}h'(s)\,\dd s$.
Lemat \ref{lem:gauss-shift} zastosowany do tych przyrostów
daje gęstość
\[
\exp\left(\sum_j\frac{\Delta h_j}{\Delta t_j}\Delta B_j
 -\frac12\sum_j\frac{(\Delta h_j)^2}{\Delta t_j}\right)
=\exp\left(I(h')-\frac12\|h'\|_2^2\right).
\]
Tożsamość \eqref{eq:cm-identity} jest więc prawdziwa dla
wszystkich ograniczonych borelowskich funkcji cylindrycznych.

\emph{Krok 2: aproksymacja dowolnego $h\in\calH_T$.}
Niech $h_n$ będą interpolantami z lematu \ref{lem:cm-approx},
$f_n=h_n'$, a $f=h'$. Z izometrii Wienera
$I(f_n)\to I(f)$ w $L^2(\mu_T)$ i w prawdopodobieństwie,
a $\|f_n\|_2^2\to\|f\|_2^2$. Wykładniki
$Z_n=\exp(I(f_n)-\|f_n\|_2^2/2)$ zbiegają więc
w prawdopodobieństwie do
$Z=\exp(I(f)-\|f\|_2^2/2)$.
Ponadto dla zmiennej normalnej $I(f_n)$ obliczamy
$\E Z_n^2=e^{\|f_n\|_2^2}$; ta rodzina drugich
momentów jest ograniczona. Zatem $(Z_n)$ jest jednostajnie
całkowalna i $Z_n\to Z$ w $L^1(\mu_T)$.

Jeśli $F$ jest ograniczoną ciągłą funkcją cylindryczną,
to $F(B+h_n)\to F(B+h)$ punktowo, gdyż $h_n\to h$
jednostajnie. Zbieżność zdominowana daje zbieżność
lewych stron tożsamości dla $h_n$, a zbieżność $Z_n$
w $L^1$ daje zbieżność prawych stron.
Wobec tego \eqref{eq:cm-identity} zachodzi dla $h$ i
wszystkich ograniczonych ciągłych funkcji cylindrycznych.
Funkcje te rozdzielają miary na $C_0([0,T])$:
wyznaczają rozkłady skończeniewymiarowe, a te wyznaczają
miarę na $\mathcal B_T$ na mocy twierdzenia \ref{th:wiener-cyl}.
Zatem dwie miary po obu stronach \eqref{eq:cm-identity}
są równe, co rozszerza wzór na wszystkie ograniczone
borelowskie $F$.

\emph{Krok 3: równoważność.} Wykładnik $Z$ jest ściśle
dodatni i skończony $\mu_T$-prawie na pewno.
Ponieważ $\dd\mu_T^h=Z\,\dd\mu_T$, otrzymujemy
$\mu_T^h\ll\mu_T$; dodatniość $Z$ implikuje też
$\mu_T\ll\mu_T^h$.
\end{proof}

\begin{uwaga}[Znak w wykładniku i związek z Girsanowem]
W gęstości prawa $B+h$ względem prawa $B$ stoi
$+I(h')$. Jest to zgodne z rozdziałem \ref{ch:girsanov},
jeśli uwzględnimy kierunek, w którym zmieniamy miarę.
Prawo przesuniętej trajektorii pod starą miarą jest tym
samym, co odpowiednio przeważone prawo nieprzesuniętej
trajektorii pod nową miarą. Mylenie tych dwóch kierunków
prowadzi do błędów znaku.
\end{uwaga}

\section{Przesunięcia spoza przestrzeni Camerona--Martina}

\begin{twierdzenie}[Kryterium singularności przesunięcia]
\label{th:cm-singular}
Niech $h\in C_0([0,T])$. Jeśli $h\notin\calH_T$,
to miary $\mu_T^h$ i $\mu_T$ są wzajemnie singularne.
Dokładniej, istnieje zbiór borelowski $A\subset C_0([0,T])$
taki, że $\mu_T(A)=0$ i $\mu_T^h(A)=1$.
\end{twierdzenie}
\begin{proof}
Niech $t_j^{(n)}=jT2^{-n}$ i zdefiniujmy funkcję
schodkową $f_n$ jako nachylenia interpolantu $h$ na
przedziałach dyadycznych, bez zakładania istnienia $h'$.
Oznaczmy energię dyskretną
\[
E_n=\|f_n\|_2^2=
\sum_{j=0}^{2^n-1}\frac{[h(t_{j+1}^{(n)})-h(t_j^{(n)})]^2}
{t_{j+1}^{(n)}-t_j^{(n)}}.
\]
\emph{Krok 1: $E_n\to\infty$.} Każda średnia $f_n$ na
przedziale siatki jest średnią $f_{n+1}$ na dwóch
podprzedziałach. Z nierówności Jensena wynika
$E_n\le E_{n+1}$. Gdyby $(E_n)$ było ograniczone,
to $(f_n)$ byłoby ograniczone w przestrzeni Hilberta $L^2$.
Istniałby słabo zbieżny podciąg $f_{n_k}\rightharpoonup f$.
Dla każdej dyadycznej chwili $t$ oraz dostatecznie dużego $k$
mamy $\int_0^t f_{n_k}(r)\,\dd r=h(t)$.
Przechodząc do granicy słabej przy funkcji testowej
$\ind{[0,t]}\in L^2$ dostajemy $\int_0^t f=h(t)$.
Obie strony są ciągłe w $t$, a chwile dyadyczne są gęste,
więc równość zachodzi dla wszystkich $t$.
Oznaczałoby to $h\in\calH_T$, sprzecznie z założeniem.
Stąd $E_n\uparrow\infty$.

\emph{Krok 2: test rozdzielający prawa.} Niech
\[
G_n(w)=\sum_{j=0}^{2^n-1}\frac{\Delta_j h}{\Delta_j t}
\,[w(t_{j+1}^{(n)})-w(t_j^{(n)})].
\]
Jest to ciągła funkcja na $C_0([0,T])$.
Pod $\mu_T$ zmienna $G_n$ ma rozkład $\Normal(0,E_n)$,
ponieważ jest kombinacją niezależnych przyrostów Wienera.
Natomiast $G_n(w+h)=G_n(w)+E_n$, więc pod $\mu_T^h$
ma rozkład $\Normal(E_n,E_n)$.

\emph{Krok 3: Borel--Cantelli.} Wybierzmy rosnący podciąg
$n_k$ taki, że $E_{n_k}\ge k^3$.
Z nierówności Czebyszewa
\[
\mu_T\{G_{n_k}>E_{n_k}/2\}\le\frac4{E_{n_k}},\qquad
\mu_T^h\{G_{n_k}\le E_{n_k}/2\}\le\frac4{E_{n_k}}.
\]
Prawa strona jest sumowalna po $k$. Pierwszy lemat
Borela--Cantellego oznacza, że pod $\mu_T$ nierówność
$G_{n_k}>E_{n_k}/2$ zachodzi co najwyżej skończenie wiele razy,
a pod $\mu_T^h$ zachodzi od pewnego miejsca zawsze.
Zbiór
$A=\bigcup_{K\ge1}\bigcap_{k\ge K}\{G_{n_k}>E_{n_k}/2\}$
ma zatem $\mu_T(A)=0$ i $\mu_T^h(A)=1$.
\end{proof}

\begin{wniosek}[Pełna dychotomia przesunięć]
Dla każdego $h\in C_0([0,T])$ zachodzi dokładnie jedna
z dwóch sytuacji: $h\in\calH_T$ i wtedy $\mu_T^h\sim\mu_T$,
albo $h\notin\calH_T$ i wtedy $\mu_T^h\perp\mu_T$.
\end{wniosek}

\section{Przykłady rachunkowe i zadania rozwiązane}

\begin{zadanie}[Przesunięcie liniowe i zdarzenie końcowe]
Niech $h(t)=ct$, $c\in\R$, i $a\in\R$.
Wyznaczyć gęstość prawa $B+h$ oraz
$\mu_T^h\{w:w(T)>a\}$.

\emph{Rozwiązanie.} Ponieważ $h'=c$,
$\|h\|_{\calH_T}^2=c^2T$ i $I(h')=cB_T$.
Stąd $\dd\mu_T^h/\dd\mu_T=e^{cB_T-c^2T/2}$.
Dla zdarzenia końcowego możemy liczyć bezpośrednio:
$B_T+cT\sim\Normal(cT,T)$, a więc
\[
\mu_T^h\{w(T)>a\}=1-\Phi\left(\frac{a-cT}{\sqrt T}\right).
\]
Dla kontroli można też scałkować gęstość przeważoną:
$\int_a^\infty e^{cy-c^2T/2}p_T(y)\,\dd y$.
Po uzupełnieniu kwadratu wynik jest ten sam.
\end{zadanie}

\begin{zadanie}[Które funkcje potęgowe należą do przestrzeni Camerona--Martina?]
Niech $h(t)=t^\alpha$ na $[0,T]$, gdzie $\alpha>0$.
Wyznaczyć zakres $\alpha$, dla którego przesunięcie
miary Wienera przez $h$ zachowuje zbiory miary zero.

\emph{Rozwiązanie.} Dla $t>0$ mamy
$h'(t)=\alpha t^{\alpha-1}$. Funkcja $h$ jest
absolutnie ciągła dla $\alpha>0$; pozostaje sprawdzić
\[
\|h'\|_2^2=\alpha^2\int_0^T t^{2\alpha-2}\,\dd t.
\]
Całka jest skończona dokładnie wtedy, gdy
$2\alpha-2>-1$, czyli $\alpha>1/2$.
Dla $\alpha>1/2$ miary są równoważne, a dla
$0<\alpha\le1/2$ wzajemnie singularne.
Ciągłość $h$ nie wystarcza: kryterium dotyczy jej
\emph{energii kwadratowej}, nie samego modułu ciągłości.
\end{zadanie}

\begin{zadanie}[Maksimum ruchu Browna z dryfem]
Niech $a>0$, $c\in\R$ oraz
$M_T^c=\sup_{0\le s\le T}(B_s+cs)$.
Obliczyć $\Prob(M_T^c\ge a)$.

\emph{Rozwiązanie.} Twierdzenie \ref{th:cameron-martin}
z $h(t)=ct$ daje
\[
\Prob(M_T^c\ge a)=\E\bigl[
\ind{\{\sup_{s\le T}B_s\ge a\}}e^{cB_T-c^2T/2}\bigr].
\]
Na zdarzeniu $\{B_T\ge a\}$ całkujemy względem
$B_T\sim\Normal(0,T)$; na zdarzeniu
$\{\sup_{s\le T}B_s\ge a,\ B_T<a\}$
zasada odbicia daje gęstość końcową $p_T(2a-y)$
na $y<a$. Stąd
\[
\begin{split}
\Prob(M_T^c\ge a)
&=\int_a^\infty e^{cy-c^2T/2}p_T(y)\,\dd y\\
&\quad+\int_{-\infty}^{a}e^{cy-c^2T/2}p_T(2a-y)\,\dd y\\
&=\Phi\left(\frac{cT-a}{\sqrt T}\right)
+e^{2ca}\Phi\left(\frac{-a-cT}{\sqrt T}\right).
\end{split}
\]
Dla $c=0$ otrzymujemy klasyczny wynik
$2[1-\Phi(a/\sqrt T)]$.
\end{zadanie}

\begin{przyklad}[Entropia względna dwóch miar trajektorii]
Dla $h\in\calH_T$ iloraz gęstości opisuje także koszt
zmiany prawa trajektorii. Definiujemy entropię względną
$D(\mu_T^h\|\mu_T)=\int\log(\dd\mu_T^h/\dd\mu_T)
\,\dd\mu_T^h$. Z \eqref{eq:cm-density} i faktu, że
$I(h')$ pod przesuniętą miarą ma średnią $\|h'\|_2^2$,
dostajemy $D(\mu_T^h\|\mu_T)=\frac12\|h\|_{\calH_T}^2$.
Można ten ostatni fakt wykazać najpierw dla schodkowych
$h'$ przez sumy przyrostów, a następnie przejść do granicy
w $L^2$. Tym samym norma Camerona--Martina ma również
interpretację informacyjną.
\end{przyklad}

\begin{uwaga}[Most Browna a przesunięcie deterministyczne]
Prawo mostu Browna zakończonego w zerze w chwili $T$
jest singularne względem $\mu_T$, ponieważ pod pierwszym
prawem $w(T)=0$ prawie na pewno, a pod drugim
$\mu_T\{w(T)=0\}=0$. Nie ma sprzeczności z twierdzeniem
Camerona--Martina: most nie powstaje przez dodanie jednej
ustalonej funkcji $h$, lecz przez zmianę struktury kowariancji.
Na każdym krótszym przedziale $[0,t]$, $t<T$,
jego prawo jest jednak absolutnie ciągłe względem
prawa ruchu Browna, z gęstością
$p_{T-t}(-B_t)/p_T(0)$; wynika to z warunkowania
w końcowej chwili i wzoru na gęstości przejścia.
\end{uwaga}

\section{Rysunek: przestrzeń trajektorii i kierunki dopuszczalnych przesunięć}
\begin{center}
\begin{tikzpicture}[>=Latex,scale=0.94]
\draw[->] (0,0)--(9.0,0) node[right] {$t$};
\draw[->] (0,-1.6)--(0,2.4) node[above] {$w(t)$};
\draw[blue!70!black,very thick] plot[smooth] coordinates
{(0,0) (0.6,0.3) (1.0,-0.2) (1.4,0.45) (2.1,-0.10)
(2.8,0.7) (3.2,0.25) (4.0,0.9) (4.7,0.1) (5.3,0.65)
(6.1,0.4) (6.7,1.1) (7.3,0.35) (8.2,1.05)};
\draw[orange!85!black,very thick,dashed] plot[smooth] coordinates
{(0,0) (0.6,0.48) (1.0,0.10) (1.4,0.87) (2.1,0.53)
(2.8,1.54) (3.2,1.21) (4.0,2.1) (4.7,1.51) (5.3,2.18)
(6.1,2.14) (6.7,2.0) (7.3,1.93) (8.2,2.25)};
\node[blue!65!black,anchor=west] at (5.3,-1.1) {$w$};
\node[orange!70!black,anchor=west] at (5.3,-1.48) {$w+h$};
\end{tikzpicture}
\par\smallskip\small Schemat idei przesuwania trajektorii. Nie jest to losowa
symulacja; rysunek ilustruje różnicę między dodaniem funkcji $h$
a zmianą miary na przestrzeni wszystkich trajektorii.
\end{center}

\section{Podsumowanie i dalsza perspektywa}
Ustaliliśmy nie tylko \emph{jak} przeważać miarę Wienera,
ale także \emph{kiedy jest to możliwe bez zmiany zbiorów miary zero}.
Kryterium dokładnie wyznacza przestrzeń Camerona--Martina.
Nie jest to tożsamość z całką Itô: całka $I(h')$ pojawia się
we wzorze na pochodną Radona--Nikodyma \emph{dwóch miar}.
W ostatniej części tomu III wykorzystamy tę konstrukcję
w analizie całek po trajektoriach i półgrup Schrödingera.

% ===========================================================================
% ROZDZIAL 10 — Procesy Levy'ego
% ===========================================================================
\chapter{Procesy Lévy'ego i ich wykładniki charakterystyczne}
\label{ch:levy}
\paragraph{Odsyłacze międzytomowe.} Wnioskowanie o zliczeniach i intensywnościach pojawia się w tomie IV, rozdziale~\ref{IV-ch:glm}; proces Poissona przedstawia tom II, rozdział~\ref{II-ch:poisson-process}.


\section{Dlaczego ruch Browna i proces Poissona są członkami jednej rodziny?}

Ruch Browna ma ciągłe trajektorie, a proces Poissona rośnie
skokowo. Łączą je jednak dwie własności: przyrosty na rozłącznych
przedziałach są niezależne, a rozkład przyrostu zależy tylko od
długości przedziału. Naturalne jest więc pytanie, jak wyglądają
\emph{wszystkie} procesy o tych własnościach.
Odpowiedź prowadzi do procesów Lévy'ego: zawierają one
ruchy Browna z dryfem, procesy złożone Poissona, procesy gamma,
procesy stabilne i liczne modele z nieskończenie wieloma skokami.

\begin{definicja}[Proces Lévy'ego]
\label{def:levy}
Niech $d\ge1$. Proces $X=(X_t)_{t\ge0}$ o wartościach w
$\R^d$, określony na przestrzeni probabilistycznej
$(\Omega,\calF,\Prob)$, jest \emph{procesem Lévy'ego}, jeśli:
\begin{enumerate}
\item $X_0=0$ prawie na pewno;
\item dla $0\le t_0<t_1<\cdots<t_n$ zmienne
$X_{t_1}-X_{t_0},\ldots,X_{t_n}-X_{t_{n-1}}$ są niezależne;
\item dla $s,t\ge0$ zachodzi
$X_{s+t}-X_s\overset{\mathrm d}=X_t$;
\item $X_t\to X_s$ według prawdopodobieństwa, gdy $t\to s$;
\item trajektorie są prawostronnie ciągłe i mają granice lewostronne
(własność \emph{càdlàg}) niemal na pewno.
\end{enumerate}
\end{definicja}

Ostatnie dwa warunki mają odrębne role. Ciągłość stochastyczna
opisuje zachowanie rozkładów przy małych zmianach czasu;
własność càdlàg mówi o wyborze \emph{jednej wspólnej wersji}
trajektorii. Każdy proces z pierwszymi czterema własnościami
ma modyfikację càdlàg; dowód ogólnego twierdzenia o regularizacji
nie jest tutaj potrzebny, gdyż od początku pracujemy z taką wersją.

\begin{lemat}[Brak skoku w ustalonej deterministycznej chwili]
\label{lem:levy-no-fixed-jumps}
Niech $X$ będzie procesem Lévy'ego. Dla każdego ustalonego
$t>0$ zachodzi $\Prob(X_t\ne X_{t-})=0$.
\end{lemat}
\begin{proof}
Dla $n$ tak dużych, że $1/n<t$, różnica
$X_t-X_{t-1/n}$ ma prawo $X_{1/n}$. Z ciągłości
stochastycznej w zerze wynika zbieżność tej różnicy do
zera według prawdopodobieństwa. Jednocześnie z istnienia
granicy lewostronnej wynika zbieżność prawie na pewno do
$X_t-X_{t-}$. Granica według prawdopodobieństwa jest
jednoznaczna, więc $X_t-X_{t-}=0$ prawie na pewno.
\end{proof}

Nie jest to sprzeczne ze skokowością procesu Poissona:
może on mieć skoki w \emph{losowych} chwilach, lecz
prawdopodobieństwo skoku w dowolnej chwili ustalonej z góry
jest równe zero.

\section{Półgrupa splotowa i wykładnik charakterystyczny}

Niech $\mu_t=\Law(X_t)$ i
$\varphi_t(u)=\E e^{i\langle u,X_t\rangle}$,
$u\in\R^d$. Niezależność i stacjonarność przyrostów dają
$\mu_{s+t}=\mu_s*\mu_t$ oraz
$\varphi_{s+t}(u)=\varphi_s(u)\varphi_t(u)$.
Ta ostatnia równość jest odpowiednikiem równania wykładniczego.

\begin{twierdzenie}[Wykładnik charakterystyczny procesu Lévy'ego]
\label{th:levy-exponent}
Dla każdego procesu Lévy'ego $X$ istnieje dokładnie jedna
ciągła funkcja $\Psi:\R^d\to\mathbb C$ z $\Psi(0)=0$
taka, że dla wszystkich $t\ge0$ i $u\in\R^d$ zachodzi
\[
\E e^{i\langle u,X_t\rangle}=e^{t\Psi(u)}.
\tag{10.1}\label{eq:levy-exponent}
\]
Ponadto $\Re\Psi(u)\le0$ i $\Psi(-u)=\overline{\Psi(u)}$.
\end{twierdzenie}
\begin{proof}
Ustalmy $u$. Z ciągłości stochastycznej wynika
$\varphi_t(u)\to1$ przy $t\downarrow0$.
Funkcja $\varphi_t(u)$ nie może być zerem dla żadnego $t>0$:
gdyby $\varphi_t(u)=0$, to z
$\varphi_t(u)=[\varphi_{t/n}(u)]^n$ wynikałoby
$\varphi_{t/n}(u)=0$ dla wszystkich $n$, sprzecznie z
granicą w zerze. Ciągła niezerowa funkcja
$t\mapsto\varphi_t(u)$ na przedziale $[0,\infty)$
ma ciągły logarytm $g_u(t)$ z $g_u(0)=0$.
Ponieważ $e^{g_u(s+t)-g_u(s)-g_u(t)}=1$,
wyrażenie w wykładniku przyjmuje wartości w $2\pi i\Z$;
przez ciągłość i wartość zero dla $s=t=0$ jest wszędzie zerem.
Funkcja $g_u$ jest więc ciągła i addytywna, stąd
$g_u(t)=t g_u(1)$. Kładziemy $\Psi(u)=g_u(1)$.

Dla różnych $u$ wybór logarytmu można uzgodnić ciągle
w otoczeniu zera w czasie: jeżeli $u$ pozostaje w ustalonym
małym otoczeniu $u_0$, to $\varphi_t(u)$ jest bliska
$\varphi_t(u_0)$ jednostajnie dla $t$ ze zwartego przedziału,
dzięki wspólnej ciągłości $(t,u)\mapsto\varphi_t(u)$.
Ciągły dobór logarytmu na takim odcinku daje ciągłość
$u\mapsto g_u(1)=\Psi(u)$; możliwa zmiana gałęzi jest
wielokrotnością $2\pi i$, ale zostaje ustalona przez
warunek $g_u(0)=0$. Równość $\varphi_t(0)=1$ daje
$\Psi(0)=0$. Wreszcie $|\varphi_t(u)|\le1$ i
$\varphi_t(-u)=\overline{\varphi_t(u)}$ implikują
wymienione własności $\Psi$.
Jednoznaczność wynika z ciągłości w $t$ i równości
$e^{t\Psi(u)}=e^{t\widetilde\Psi(u)}$ dla każdego $t$.
\end{proof}

\begin{uwaga}[Nieskończona podzielność]
Dla każdego $n\ge1$ zmienna $X_t$ ma rozkład sumy $n$
niezależnych zmiennych o prawie $\mu_{t/n}$.
Zatem $\mu_t$ jest \emph{nieskończenie podzielna}.
To właśnie klasyfikacja nieskończenie podzielnych praw
stoi za ogólnym wzorem Lévy'ego--Chinczyna.
\end{uwaga}

\begin{przyklad}[Trzy wykładniki]
Dla procesu $X_t=bt+\sigma W_t$ w jednym wymiarze
$\Psi(u)=ibu-\sigma^2u^2/2$.
Dla procesu Poissona o intensywności $\lambda$
$\Psi(u)=\lambda(e^{iu}-1)$.
Dla różnicy niezależnych procesów Poissona
$N_t^+-N_t^-$ o intensywnościach $\lambda_+$ i $\lambda_-$
$\Psi(u)=\lambda_+(e^{iu}-1)+\lambda_-(e^{-iu}-1)$.
Mimo zupełnie innej geometrii trajektorii wszystkie trzy
przykłady mają tę samą wykładniczą zależność od czasu.
\end{przyklad}

\section{Złożony proces Poissona: pełny rachunek}

\begin{definicja}[Złożony proces Poissona]
Niech $N$ będzie procesem Poissona o intensywności
$\lambda>0$, a $(Y_j)_{j\ge1}$ niezależnymi zmiennymi
losowymi o wspólnym prawie $\rho$ na $\R^d\setminus\{0\}$,
niezależnymi od $N$. Proces
$J_t=\sum_{j=1}^{N_t}Y_j$ nazywamy \emph{złożonym
procesem Poissona}. Przyjmujemy, że suma pusta jest zerem.
\end{definicja}

\begin{twierdzenie}[Charakterystyka złożonego procesu Poissona]
\label{th:cpp}
Proces $J$ jest procesem Lévy'ego oraz
\[
\E e^{i\langle u,J_t\rangle}
=\exp\left[t\lambda\left(\int e^{i\langle u,y\rangle}
\,\dd\rho(y)-1\right)\right].
\]
Jeśli $\E\|Y_1\|^2<\infty$, to
$\E J_t=\lambda t\E Y_1$ oraz
$\operatorname{Cov}(J_t)=\lambda t\E(Y_1Y_1^{\mathsf T})$.
\end{twierdzenie}
\begin{proof}
Trajektorie są skokowe i mają na każdym zwartym przedziale
skończenie wiele skoków prawie na pewno.
Niezależność przyrostów wynika ze wspólnej niezależności
liczby i rozmiarów skoków na rozłącznych przedziałach;
stacjonarność z jednorodności zegara Poissona.
Ponieważ $\Prob(N_h>0)=1-e^{-\lambda h}\to0$,
proces jest stochastycznie ciągły.

Warunkując względem $N_t$ i korzystając z niezależności
zmiennych $Y_j$, dla
$\phi(u)=\E e^{i\langle u,Y_1\rangle}$ dostajemy
\[
\E e^{i\langle u,J_t\rangle}
=\sum_{n=0}^\infty e^{-\lambda t}
\frac{(\lambda t)^n}{n!}\phi(u)^n
=e^{\lambda t(\phi(u)-1)}.
\]
Dla momentów najpierw warunkujemy:
$\E(J_t\mid N_t)=N_t\E Y_1$ oraz
$\operatorname{Cov}(J_t\mid N_t)=N_t\operatorname{Cov}(Y_1)$.
Z prawa całkowitej kowariancji wynika
$\lambda t\operatorname{Cov}(Y_1)+
\lambda t(\E Y_1)(\E Y_1)^{\mathsf T}$,
co jest podaną macierzą.
\end{proof}

\section{Miara Lévy'ego i małe skoki}

\begin{definicja}[Miara Lévy'ego]
Miara $\nu$ na $\R^d\setminus\{0\}$ jest
\emph{miarą Lévy'ego}, jeśli
\[
\int_{\R^d\setminus\{0\}}(1\wedge\|y\|^2)\,\dd\nu(y)<\infty.
\tag{10.2}\label{eq:levy-measure}
\]
W szczególności $\nu(\{\|y\|>\varepsilon\})<\infty$
dla każdego $\varepsilon>0$, ale może zachodzić
$\nu(\{0<\|y\|\le1\})=\infty$.
\end{definicja}

Warunek z kwadratem jest subtelny: nie żąda całkowalności
rozmiarów małych skoków. Umożliwia jednak sumowanie
\emph{skompensowanych} małych skoków, ponieważ ich wariancje
są sumowalne. Aby to zobaczyć, skonstruujemy proces z zadanej
miary $\nu$ bez korzystania z nieopracowanej jeszcze teorii
ogólnych losowych miar Poissona.

\begin{lemat}[Suma niezależnych skoków w rozłącznych klasach]
\label{lem:annular}
Niech $\nu$ spełnia \eqref{eq:levy-measure}. Dla dowolnego
zbioru borelowskiego $A$ o $\nu(A)<\infty$, rozłącznego z
pewnym otoczeniem zera, można skonstruować złożony proces
Poissona $J^A$ o wykładniku
$\Psi_A(u)=\int_A(e^{i\langle u,y\rangle}-1)\,\dd\nu(y)$.
Dla rozłącznych zbiorów $A_1,\ldots,A_k$ procesy
$J^{A_1},\ldots,J^{A_k}$ można wybrać niezależnie.
Jeśli $\int_A\|y\|^2\,\dd\nu(y)<\infty$, to
$M_t^A=J_t^A-t\int_Ay\,\dd\nu(y)$ jest
kwadratowo całkowalnym martyngałem i
\[
\E\|M_t^A\|^2=t\int_A\|y\|^2\,\dd\nu(y).
\]
\end{lemat}
\begin{proof}
Przy $\nu(A)>0$ bierzemy intensywność
$\lambda=\nu(A)$ i rozkład skoku
$\rho(B)=\nu(A\cap B)/\nu(A)$, po czym stosujemy
konstrukcję złożonego procesu Poissona. Przy $\nu(A)=0$
wybieramy proces zerowy. Niezależność na skończonej rodzinie
rozłącznych klas uzyskujemy na przestrzeni produktowej.
Własność martyngału dla $M^A$ wynika z niezależności
przyszłych przyrostów i ich średniej zero.
Oszacowanie drugiego momentu otrzymujemy ze wzoru na
kowariancję w twierdzeniu \ref{th:cpp} przez wzięcie śladu.
\end{proof}

\section{Wzór Lévy'ego--Chinczyna i konstrukcja zadanej trójki}

\begin{twierdzenie}[Postać Lévy'ego--Chinczyna]
\label{th:LK}
Proces Lévy'ego $X$ o wartościach w $\R^d$ ma wykładnik
charakterystyczny postaci
\[
\begin{split}
\Psi(u)={}&i\langle b,u\rangle-\frac12u^{\mathsf T}Qu\\
&+\int_{\R^d\setminus\{0\}}
\left(e^{i\langle u,y\rangle}-1
-i\langle u,y\rangle\ind{\{\|y\|\le1\}}\right)\,\dd\nu(y),
\end{split}
\tag{10.3}\label{eq:LK}
\]
gdzie $b\in\R^d$, $Q$ jest symetryczną
nieujemnie określoną macierzą, a $\nu$ jest miarą
Lévy'ego. Przy ustalonej funkcji obcięcia
$y\mapsto y\ind{\{\|y\|\le1\}}$ trójka $(b,Q,\nu)$
jest jednoznaczna. Odwrotnie każda taka trójka wyznacza
prawo procesu Lévy'ego.
\end{twierdzenie}

\begin{uwaga}[Status dowodu klasyfikacyjnego]
Dwa aspekty twierdzenia trzeba rozdzielić. \emph{Kierunek
konstrukcyjny} --- od trójki $(b,Q,\nu)$ do procesu ---
udowodnimy poniżej bez stosowania ogólnej teorii miar losowych.
\emph{Kierunek konieczny} i jednoznaczność trójki
są klasycznym twierdzeniem o reprezentacji funkcji
charakterystycznych praw nieskończenie podzielnych.
Pełny dowód tego wyniku analitycznego wymaga teorii
miar Lévy'ego otrzymywanych z granic małoczasowych
oraz wyodrębnienia części gaussowskiej. Przyjmujemy tutaj
klasyfikację praw nieskończenie podzielnych jako
osobne twierdzenie pomocnicze i nie przedstawiamy poniższej
konstrukcji jako dowodu \emph{obu} kierunków.
Dowód przez miary skoków i dekompozycję Lévy'ego--Itô
będzie naturalnym tematem następnego rozdziału.
\end{uwaga}

\begin{twierdzenie}[Konstrukcja procesu dla dowolnej trójki Lévy'ego]
\label{th:triplet-construction}
Niech $b\in\R^d$, $Q\ge0$ oraz $\nu$ spełnia
\eqref{eq:levy-measure}. Istnieje proces Lévy'ego
$X$, którego wykładnik jest równy prawej stronie
\eqref{eq:LK}. Można go otrzymać jako granicę
jednostajną według prawdopodobieństwa na każdym
przedziale zwartym sum niezależnych procesów skokowych,
ruchu Browna i dryfu.
\end{twierdzenie}
\begin{proof}
\emph{Krok 1: skoki duże.} Niech $A_{\rm big}=\{\|y\|>1\}$.
Miara $\nu(A_{\rm big})$ jest skończona, więc lemat
\ref{lem:annular} daje proces $J^{A_{\rm big}}$. Nie potrzebujemy
żadnego założenia momentowego dla dużych skoków:
w skończonym czasie występuje ich skończenie wiele.

\emph{Krok 2: rozbicie małych skoków.} Dla $k\ge0$ niech
$A_k=\{2^{-(k+1)}<\|y\|\le2^{-k}\}$.
Możemy skonstruować niezależne kompensowane procesy
$M^{A_k}$, ponieważ miary $\nu(A_k)$ są skończone.
Definiujemy sumy $S_t^{(n)}=\sum_{k=0}^n M_t^{A_k}$.
Dla $n>m$, niezależność i średnia zero dają
\[
\E\|S_T^{(n)}-S_T^{(m)}\|^2
=T\sum_{k=m+1}^n\int_{A_k}\|y\|^2\,\dd\nu(y).
\]
Norma procesu $S^{(n)}-S^{(m)}$ jest podmartyngałem;
nierówność maksymalna Dooba w $L^2$ zapewnia
\[
\E\sup_{t\le T}\|S_t^{(n)}-S_t^{(m)}\|^2
\le4T\sum_{k=m+1}^n\int_{A_k}\|y\|^2\,\dd\nu(y).
\tag{10.4}\label{eq:levy-l2-sup}
\]
Z warunku \eqref{eq:levy-measure} prawa strona dąży
do zera przy $m,n\to\infty$. Sumy $S^{(n)}$
są więc Cauchy'ego w $L^2$ normy supremum na
$[0,T]$. Wybierając podciąg z sumowalnymi kwadratami
różnic i stosując Borela--Cantellego otrzymujemy
granicę jednostajną prawie na pewno tego podciągu.
Ponieważ jednostajna granica funkcji càdlàg jest càdlàg,
graniczny proces $S$ ma trajektorie càdlàg.
Wykonując konstrukcję dla $T=1,2,\ldots$ i biorąc podciąg
diagonalny, uzyskujemy jeden proces na $[0,\infty)$.

\emph{Uwaga o rozłączności składników.} Składnik $S$ obejmuje
wszystkie skompensowane klasy $A_k$, także $A_0$.
Nie wolno dodawać procesu $J^{A_0}$ po raz drugi: skutkowałoby
to podwójnym policzeniem skoków klasy $A_0$ oraz błędnym
wykładnikiem charakterystycznym.

\emph{Krok 3: część gaussowska.} Niech $W$ będzie
standardowym $d$-wymiarowym procesem Wienera
niezależnym od wszystkich $J^{A_k}$. Niech $Q^{1/2}$
będzie nieujemnym pierwiastkiem macierzy $Q$ i połóżmy
$X_t=bt+Q^{1/2}W_t+J_t^{A_{\rm big}}+S_t$.
Dla każdej sumy skończonej przyrosty na rozłącznych
przedziałach są niezależne i stacjonarne. Zbieżność w
prawdopodobieństwie ich wektorów zachowuje rozkłady
produktowe i stacjonarność. Dostajemy te same własności
dla $X$. Ciągłość stochastyczna wynika z
$\E\|S_{t+h}-S_t\|^2
=h\int_{\{0<\|y\|\le1\}}\|y\|^2\,\dd\nu(y)$,
ciągłości ruchu Browna oraz
$\Prob(J_{t+h}^{A_{\rm big}}-J_t^{A_{\rm big}}\ne0)\le h\nu(A_{\rm big})$.

\emph{Krok 4: obliczenie wykładnika.} Dla $S^{(n)}$
charakterystyka kompensowanego składnika jest równa
\[
\exp\left[t\sum_{k=0}^n\int_{A_k}
(e^{i\langle u,y\rangle}-1-i\langle u,y\rangle)
\,\dd\nu(y)\right].
\]
Zbieżność całek przy $n\to\infty$ wynika z oszacowania
$|e^{iv}-1-iv|\le v^2/2$ dla rzeczywistego $v$.
Duże skoki mają wykładnik
$t\int_{A_{\rm big}}(e^{i\langle u,y\rangle}-1)\,\dd\nu(y)$,
a część gaussowska $-tu^{\mathsf T}Qu/2$.
Mnożąc funkcje charakterystyczne niezależnych składników
i przechodząc do granicy, otrzymujemy dokładnie
\eqref{eq:LK}. Tym samym konstrukcja jest zakończona.
\end{proof}

\begin{uwaga}[Dlaczego kompensujemy małe, lecz nie duże skoki?]
Dla małych skoków może zachodzić
$\int_{\{\|y\|\le1\}}\|y\|\,\dd\nu(y)=\infty$;
zwykła suma wektorowa skoków nie ma wtedy automatycznego
sensu. Po odjęciu średniego przyrostu każdej skończonej
klasy uzyskujemy rodzinę o sumowalnych \emph{wariancjach}.
Duże skoki występują zaś w skończonej liczbie na każdym
przedziale zwartym i można je sumować bez kompensacji.
\end{uwaga}

\section{Generator procesu Lévy'ego: operator nielokalny}

\begin{twierdzenie}[Generator na funkcjach szybko malejących]
\label{th:levy-generator}
Niech $X$ będzie procesem Lévy'ego o trójce $(b,Q,\nu)$,
a $P_tf(x)=\E f(x+X_t)$. Dla każdej rzeczywistej funkcji
$f$ z przestrzeni Schwartza $\calS(\R^d)$ zachodzi
\[
\lim_{t\downarrow0}\frac{P_tf(x)-f(x)}{t}=Af(x)
\qquad (x\in\R^d),
\]
gdzie
\[
\begin{split}
Af(x)={}&\langle b,\nabla f(x)\rangle
+\frac12\operatorname{tr}(QD^2f(x))\\
&+\int_{\R^d\setminus\{0\}}
\bigl(f(x+y)-f(x)
-\ind{\{\|y\|\le1\}}\langle y,\nabla f(x)\rangle\bigr)
\,\dd\nu(y).
\end{split}
\tag{10.5}\label{eq:levy-generator}
\]
Całka po prawej stronie jest bezwzględnie zbieżna.
\end{twierdzenie}
\begin{proof}
Dla $\|y\|\le1$ rozwinięcie Taylora daje
$|f(x+y)-f(x)-\langle y,\nabla f(x)\rangle|
\le\frac12\|D^2f\|_\infty\|y\|^2$.
Dla $\|y\|>1$ różnica $f(x+y)-f(x)$ jest ograniczona
przez $2\|f\|_\infty$. Zatem całka istnieje.

Przyjmijmy konwencję
$\widehat f(u)=\int_{\R^d}e^{-i\langle u,x\rangle}
 f(x)\,\dd\lambda_d(x)$.
Odwrócenie Fouriera i Fubini dają
\[
P_tf(x)=\frac{1}{(2\pi)^d}\int_{\R^d}
 e^{i\langle u,x\rangle}e^{t\Psi(u)}\widehat f(u)
\,\dd\lambda_d(u).
\]
Z \eqref{eq:LK}, $|e^{iv}-1|\le2$ i
$|e^{iv}-1-iv|\le v^2/2$ wynika
$|\Psi(u)|\le C(1+\|u\|^2)$ dla stałej zależnej
od trójki. Ponieważ $\Re\Psi(u)\le0$,
\[
\left|\frac{e^{t\Psi(u)}-1}{t}\right|
=\left|\Psi(u)\int_0^1e^{st\Psi(u)}\,\dd s\right|
\le|\Psi(u)|.
\]
Funkcja $|\Psi(u)\widehat f(u)|$ jest całkowalna,
więc twierdzenie o zbieżności zdominowanej pozwala
przejść do granicy pod całką Fouriera i otrzymać
$\frac1{(2\pi)^d}\int e^{i\langle u,x\rangle}
\Psi(u)\widehat f(u)\,\dd u$.
Wstawiamy wzór \eqref{eq:LK}. Mnożenie transformaty
przez $iu_j$ odpowiada pochodnej $\partial_j f$, a
przez $-u_i u_j$ pochodnej $\partial_{ij}f$.
Dla składnika skokowego Fubini jest dozwolone dzięki
oszacowaniom z pierwszego akapitu i szybkiemu maleniu
$\widehat f$. Otrzymujemy \eqref{eq:levy-generator}.
\end{proof}

\begin{przyklad}[Dlaczego generator jest nielokalny?]
Dla procesu złożonego Poissona o skokach zawsze równych
$a\in\R$ i intensywności $\lambda$ generator jest równy
$Af(x)=\lambda(f(x+a)-f(x))$.
Wartość $Af(x)$ zależy od $f$ w \emph{innym punkcie} $x+a$;
nie da się jej wyrazić wyłącznie przez skończenie wiele
pochodnych funkcji w punkcie $x$. Właśnie dlatego
równania dla procesów skokowych prowadzą do równań
różniczkowo-całkowych, a nie wyłącznie do PDE.
\end{przyklad}

\section{Mocna własność Markowa procesu Lévy'ego}

\begin{twierdzenie}[Mocna własność Markowa]
\label{th:levy-strong-markov}
Niech $X$ będzie procesem Lévy'ego i niech
$(\calF_t)$ będzie uzupełnioną, prawostronnie ciągłą
filtracją naturalną procesu $X$.
Jeżeli $\tau$ jest czasem zatrzymania ze
$\Prob(\tau<\infty)=1$, to proces
$(X_{\tau+t}-X_\tau)_{t\ge0}$ jest niezależny od
$\calF_\tau$ i ma to samo prawo co $X$.
\end{twierdzenie}
\begin{proof}
Najpierw dla czasu deterministycznego $s$ własność wynika
z niezależności przyrostów i twierdzenia o klasie
monotonicznej: przyszłe przyrosty są niezależne od
$\sigma(X_r:r\le s)$. Aby uwzględnić uzupełnienie i
prawostronną ciągłość filtracji, przybliżamy $s$ z góry
chwilami $s_n\downarrow s$; prawostronna ciągłość
trajektorii i zbieżność zdominowana przenoszą równość
warunkową na $\calF_s=\bigcap_{n}\calF_{s_n}$.

Załóżmy teraz, że $\tau\le K$ i niech
$\tau_n=2^{-n}\lceil2^n\tau\rceil$.
Wtedy $\tau_n$ przyjmuje skończenie wiele wartości,
$\tau_n\downarrow\tau$ oraz $\tau\le\tau_n$.
Dla każdego ograniczonego $\calF_\tau$-mierzalnego $Z$
i każdej ograniczonej ciągłej funkcji
$g:(\R^d)^m\to\R$ niezależność przyrostów w
chwilach deterministycznych daje
\[
\E\Big[Z\,g(X_{\tau_n+t_1}-X_{\tau_n},\ldots,
 X_{\tau_n+t_m}-X_{\tau_n})\Big]
=\E Z\,\E g(X_{t_1},\ldots,X_{t_m})
\]
dla $0\le t_1<\cdots<t_m$, gdyż na
$\{\tau_n=r\}$ zmienna $Z\ind{\{\tau_n=r\}}$
jest $\calF_r$-mierzalna. Przez prawostronną ciągłość
trajektorii $X$ przyrosty po $\tau_n$ zbiegają
prawie na pewno do przyrostów po $\tau$.
Zbieżność zdominowana daje tę samą równość dla $\tau$.
Rozszerzamy ją z ciągłych $g$ na ograniczone borelowskie
funkcje przez klasę monotoniczną. Otrzymujemy niezależność
wszystkich wektorów przyszłych przyrostów od $\calF_\tau$.
Rozkłady skończeniewymiarowe wyznaczają prawo procesu
càdlàg na przestrzeni Skorochoda, więc teza zachodzi
dla całego procesu. Dla nieograniczonego, ale skończonego
prawie na pewno $\tau$ stosujemy argument na zdarzeniach
$\{\tau\le K\}$ i przechodzimy $K\to\infty$.
\end{proof}

\section{Modele: proces gamma i procesy stabilne}

\begin{przyklad}[Proces gamma: nieskończenie wiele skoków, ale skończona wariacja]
Niech $\alpha,\beta>0$ i zdefiniujmy miarę na
$(0,\infty)$ przez
$\dd\nu(x)=\alpha e^{-\beta x}x^{-1}\,\dd x$.
Ponieważ $\int_0^1x\,\dd\nu(x)<\infty$ i
$\nu((1,\infty))<\infty$, można zsumować dodatnie
skoki bez kompensacji i otrzymać niemalejący proces Lévy'ego
$G_t$ z wykładnikiem
$\Psi(u)=\int_0^\infty(e^{iux}-1)\,\dd\nu(x)$.
Różniczkując po $u$ pod całką otrzymujemy
$\Psi'(u)=i\alpha\int_0^\infty e^{-(\beta-iu)x}\,\dd x
=i\alpha/(\beta-iu)$, a $\Psi(0)=0$. Stąd
\[
\Psi(u)=-\alpha\log(1-iu/\beta),\qquad
\E e^{iuG_t}=(1-iu/\beta)^{-\alpha t}.
\]
W konwencji \eqref{eq:LK} jego parametr dryfu wynosi
$b=\int_{(0,1]}x\,\dd\nu(x)$, dzięki czemu
skompensowane małe skoki wraz z dryfem tworzą zwykłą
niemalejącą sumę skoków. Wobec tego $G_t$ ma rozkład gamma z parametrem kształtu
$\alpha t$ i intensywnością $\beta$.
Ponieważ $\nu((0,1))=\infty$, liczba skoków procesu
na dowolnym przedziale $[0,t]$ dla $t>0$ jest prawie
na pewno nieskończona: liczby skoków większych niż
$\varepsilon$ są Poissonowskie o parametrach
$t\nu((\varepsilon,\infty))\to\infty$.
Jednak suma ich rozmiarów jest skończona prawie na pewno,
gdyż oczekiwana suma skoków nie większych niż $1$
ma skończoną wartość $t\int_0^1x\,\dd\nu(x)$.
\end{przyklad}

\begin{przyklad}[Symetryczne procesy stabilne]
Dla $0<\alpha<2$ weźmy w jednym wymiarze
$\dd\nu(y)=c|y|^{-1-\alpha}\,\dd y$, $c>0$,
oraz $b=0$, $Q=0$. Warunek \eqref{eq:levy-measure}
jest spełniony: przy zerze całkujemy $|y|^{1-\alpha}$,
a w nieskończoności $|y|^{-1-\alpha}$.
Symetria znosi urojone części całki w \eqref{eq:LK}, zatem
\[
\Psi(u)=\int_\R(\cos(uy)-1)c|y|^{-1-\alpha}\,\dd y
=-C_\alpha |u|^\alpha,
\quad
C_\alpha=2c\int_0^\infty(1-\cos z)z^{-1-\alpha}\,\dd z>0.
\]
Po podstawieniu $z=|u|y$ otrzymujemy skalowanie
$X_t\overset{\mathrm d}=t^{1/\alpha}X_1$.
Dla $\alpha\ne2$ proces ten nie jest gaussowski, a jego
prawa mają zasadniczo inne zachowanie ogonów niż normalne.
\end{przyklad}

\section{Momenty, wariancje i zadanie z trójką Lévy'ego}

\begin{twierdzenie}[Średnia i kowariancja, jeśli istnieją]
\label{th:levy-moments}
Niech $X$ ma trójkę $(b,Q,\nu)$ w konwencji
\eqref{eq:LK}. Jeśli
$\int_{\{\|y\|>1\}}\|y\|\,\dd\nu(y)<\infty$,
to $\E X_t=t(b+\int_{\{\|y\|>1\}}y\,\dd\nu(y))$.
Jeśli dodatkowo
$\int_{\{\|y\|>1\}}\|y\|^2\,\dd\nu(y)<\infty$,
to
\[
\operatorname{Cov}(X_t)
=t\left(Q+\int_{\R^d\setminus\{0\}}yy^{\mathsf T}\,\dd\nu(y)\right).
\]
\end{twierdzenie}
\begin{proof}
W konstrukcji \ref{th:triplet-construction} małe skoki są
scentrowane. Ich częściowe sumy mają średnią zero i
zbieżność w $L^2$, a więc także w $L^1$.
Wektor dryfu ma średnią $bt$, część gaussowska średnią
zero, natomiast suma dużych skoków ma średnią
$t\int_{\|y\|>1}y\,\dd\nu(y)$ pod pierwszym
założeniem. To daje pierwszy wzór.
Przy drugim założeniu wszystkie trzy losowe składniki
są kwadratowo całkowalne i niezależne. Kowariancja
części gaussowskiej wynosi $tQ$, małych skoków
$t\int_{\|y\|\le1}yy^{\mathsf T}\,\dd\nu(y)$,
a dużych $t\int_{\|y\|>1}yy^{\mathsf T}\,\dd\nu(y)$.
Sumując te macierze otrzymujemy tezę.
Prawo i momenty procesu nie zależą od konkretnej
realizacji użytej w konstrukcji, gdyż charakterystyka
\eqref{eq:LK} wyznacza jednoznacznie jego rozkłady.
\end{proof}

\begin{zadanie}[Skoki dodatnie i ujemne]
Niech $X_t=N_t^+-N_t^-$, gdzie $N^+,N^-$ są niezależnymi
procesami Poissona o intensywnościach $\lambda_+,\lambda_->0$.
Wyznaczyć trójkę Lévy'ego, średnią, wariancję i generator.

\emph{Rozwiązanie.} Miara skoków ma postać
$\nu=\lambda_+\delta_1+\lambda_-\delta_{-1}$.
Przy obcięciu $\ind{\{|y|\le1\}}$ skoki $\pm1$
są traktowane jako małe, zatem
$b=\int_{\{|y|\le1\}}y\,\dd\nu(y)=\lambda_+-\lambda_-$,
$Q=0$. Podstawienie do \eqref{eq:LK} daje
$\Psi(u)=\lambda_+(e^{iu}-1)+\lambda_-(e^{-iu}-1)$.
Z twierdzenia \ref{th:levy-moments}:
$\E X_t=t(\lambda_+-\lambda_-)$ oraz
$\Var(X_t)=t(\lambda_++\lambda_-)$. Wreszcie
$Af(x)=\lambda_+(f(x+1)-f(x))+
\lambda_-(f(x-1)-f(x))$.
To zadanie pokazuje, że parametr $b$ zależy od wybranej
funkcji obcięcia: nie zawsze jest dosłownie „dryfem
trajektorii” w intuicyjnym znaczeniu.
\end{zadanie}

\begin{zadanie}[Jak dobrać kompensator?]
Niech $J_t=\sum_{j=1}^{N_t}Y_j$, gdzie
$N$ jest procesem Poissona o intensywności $\lambda$,
a zmienne $Y_j$ mają średnią $m$ i drugi moment $v_2<\infty$.
Znaleźć martyngał związany z $J$ i jego drugi moment.

\emph{Rozwiązanie.} Przyszłe przyrosty $J$ są
niezależne od przeszłości i mają średnią
$\lambda(t-s)m$, stąd
$M_t=J_t-\lambda tm$ jest martyngałem.
Z twierdzenia \ref{th:cpp} jego wariancja wynosi
$\E M_t^2=\lambda tv_2$.
W szczególności poprawnym kompensatorem jest
$\lambda tm$, a nie $\lambda t\Var(Y_1)$:
drugi moment pojedynczego skoku obejmuje także losowość
jego liczby.
\end{zadanie}

\begin{center}
\begin{tikzpicture}[>=Latex,scale=.93]
\draw[->] (0,0)--(9,0) node[right] {$t$};
\draw[->] (0,-1)--(0,3.1) node[above] {$X_t$};
\draw[blue!65!black,very thick] (0,0)--(1.4,.23)--(1.4,1.35)--
(3.0,1.58)--(3.0,.75)--(4.9,1.05)--(4.9,2.23)--
(6.8,2.52)--(6.8,1.82)--(8.5,2.12);
\draw[orange!75!black,dashed,thick] (0,0)--(8.5,1.34);
\node[anchor=west,scale=.87,blue!65!black] at (5.1,2.86)
{część skokowa z dryfem};
\node[anchor=west,scale=.8,orange!80!black] at (5.8,.42)
{trend średni};
\end{tikzpicture}
\par\smallskip\small Schemat trajektorii procesu złożonego Poissona
z dryfem. Skoki mogą mieć oba znaki; pomarańczowa linia
nie jest realizacją procesu, lecz ilustracją jego trendu.
\end{center}

\section{Granica rozdziału i przygotowanie do całkowania ze skokami}
Konstrukcja z twierdzenia \ref{th:triplet-construction}
pokazuje, jak zbudować proces dla trójki Lévy'ego,
wykorzystując niezależne klasy skoków i kontrolę w $L^2$.
Nie zdefiniowaliśmy jednak jeszcze ogólnych \emph{losowych
miar skoków} ani całek względem ich kompensacji.
Będą one przedmiotem rozdziału 11, gdzie otrzymamy
jedną notację obejmującą zarówno złożone procesy Poissona,
jak i skoki o nieskończonej aktywności.

Równocześnie do całek po trajektoriach wrócimy po omówieniu
pełnej teorii półmartyngałów. Dzięki temu nie będziemy
utożsamiać całki względem miary Wienera z całką Itô
ani traktować oscylacyjnej całki Feynmana jak zwykłej
miary probabilistycznej.

% ===========================================================================
% ROZDZIAL 11 -- Miary skokow, kompensacja, calki Poissona
% ===========================================================================
\chapter{Losowe miary skoków i całkowanie względem miar Poissona}
\label{ch:poisson-random-measure}
\paragraph{Odsyłacze międzytomowe.} Proces Poissona z tomem II (rozdział~\ref{II-ch:poisson-process}) łączy się ze statystyczną estymacją liczby zdarzeń z tomu IV (rozdział~\ref{IV-ch:conjugate-families}).


\section{Dlaczego nie wystarcza suma skoków?}

Złożony proces Poissona zapisaliśmy jako $J_t=\sum_{j=1}^{N_t}Y_j$.
W tym wzorze ukryte są \emph{dwie} informacje: czas wystąpienia
skoku i jego rozmiar. Gdy skoków jest nieskończenie wiele, sumowanie
wszystkich ich rozmiarów może nie mieć sensu, nawet na przedziale
skończonym. Warunek definiujący miarę Lévy'ego gwarantuje tylko
całkowalność \emph{kwadratów} małych skoków. Dlatego rozdzielamy
losowe zliczanie skoków od ich całkowania, a następnie odejmujemy
od części sum odpowiednie przewidywalne średnie.

Niech $E$ będzie przestrzenią polską z borelowskim $\sigma$-ciałem
$\calB(E)$, a $\nu$ miarą $\sigma$-skończoną na $E$.
W przykładach głównych $E=\R^d\setminus\{0\}$ i $\nu$ będzie
miarą Lévy'ego. Pracujemy na ustalonym horyzoncie $T>0$;
$(\calF_t)_{0\le t\le T}$ oznacza filtrację spełniającą warunki
zwykłe. Używamy zapisu $\dd t\,\dd\nu(z)$ dla miary
produktowej $\lambda_1\otimes\nu$.

\begin{definicja}[Losowa miara Poissona]
\label{def:PRM}
Losową miarą Poissona o intensywności $\dd t\,\dd\nu(z)$ na
$(0,T]\times E$ nazywamy odwzorowanie $N$, które dla prawie każdego
$\omega$ jest miarą zliczającą i spełnia:
\begin{enumerate}
\item dla każdego mierzalnego $B\subset(0,T]\times E$ o
$(\lambda_1\otimes\nu)(B)<\infty$ zmienna $N(B)$ ma rozkład
Poissona o parametrze $(\lambda_1\otimes\nu)(B)$;
\item dla każdej skończonej rodziny parami rozłącznych takich
zbiorów $B_1,\ldots,B_m$ zmienne $N(B_1),\ldots,N(B_m)$ są niezależne;
\item $N((0,t]\times A)$ jest $\calF_t$-mierzalna, a
$N((s,t]\times A)$ jest niezależna od $\calF_s$ dla $s<t$
i $\nu(A)<\infty$.
\end{enumerate}
Przez $\widetilde N(\dd t,\dd z)$ oznaczamy formalnie
$N(\dd t,\dd z)-\dd t\,\dd\nu(z)$. Sens całkowy tej różnicy
określimy niżej.
\end{definicja}

\begin{uwaga}[O zbiorach nieskończonej intensywności]
Jeżeli $\nu(A)=\infty$, to $N((0,t]\times A)=\infty$
prawie na pewno dla $t>0$, o ile taka miara Poissona istnieje.
Nie oznacza to, że $\widetilde N((0,t]\times A)$ jest określoną
różnicą dwóch nieskończoności: \emph{nie definiujemy} jej w ten
sposób. Kompensowaną miarę rozumiemy przez całki po zbiorach o
skończonej intensywności i przedłużenie w $L^2$.
\end{uwaga}

\begin{twierdzenie}[Istnienie i konstrukcja przez znaczniki]
\label{th:prm-existence}
Dla każdej $\sigma$-skończonej miary $\nu$ na $E$ istnieje
losowa miara Poissona na $(0,\infty)\times E$ o intensywności
$\dd t\,\dd\nu(z)$, z naturalną uzupełnioną filtracją
zapewniającą warunek niezależności przyrostów z
Definicji~\ref{def:PRM}.
\end{twierdzenie}
\begin{proof}
\emph{Krok 1: skończona miara.} Niech $0<\nu(E)=\lambda<\infty$.
Bierzemy proces Poissona $K$ o intensywności $\lambda$, jego
chwile skoków $S_j$ oraz niezależne znaczniki $Z_j\in E$ o
wspólnym prawie $\rho=\nu/\lambda$, niezależne od $K$.
Definiujemy $N=\sum_{j\ge1}\delta_{(S_j,Z_j)}$.
Dla rozłącznych $A_1,\ldots,A_m\subset E$ i przedziału
$(s,t]$ warunkowa funkcja tworząca wektora
$(N((s,t]\times A_k))_k$ wynosi
$(\sum_{k=1}^m\rho(A_k)u_k+1-\sum_{k=1}^m\rho(A_k))^{K_t-K_s}$.
Po uśrednieniu po rozkładzie Poissona dostajemy
$\exp((t-s)\sum_k\nu(A_k)(u_k-1))$.
Jest to iloczyn funkcji tworzących niezależnych zmiennych
Poissona o parametrach $(t-s)\nu(A_k)$.
Dalsze rozbicie przedziału czasu pokazuje niezależność na
rozłącznych prostokątach czasowo-przestrzennych; z twierdzenia
o klasie monotonicznej rozciągamy to na rozłączne zbiory
mierzalne o skończonej intensywności. Obliczenia pokazują
również niezależność od informacji do chwili $s$.

\emph{Krok 2: miara $\sigma$-skończona.} Wybierzmy rozkład
$E=\bigsqcup_{k\ge1}E_k$ z $\nu(E_k)<\infty$.
Na iloczynie przestrzeni probabilistycznych konstruujemy niezależne
miary $N_k$ o intensywnościach $\dd t\,\nu|_{E_k}$ i kładziemy
$N(B)=\sum_kN_k(B\cap((0,\infty)\times E_k))$.
Dla $(\lambda_1\otimes\nu)(B)<\infty$ suma parametrów Poissona
jest skończona, więc $N(B)<\infty$ prawie na pewno.
Niezależność dla rozłącznych zbiorów przechodzi z sum skończonych
przez zbieżność funkcji tworzących, a adaptację zapewnia naturalna
filtracja wygenerowana przez wszystkie zliczenia do danej chwili.
Uzupełnienie i prawostronna regularizacja tej filtracji zachowują
niezależność przyrostów: dla chwil $s_n\downarrow s$ niezależność
od $\calF^0_{s_n}$ przechodzi na $\cap_n\calF^0_{s_n}$, ponieważ
zliczenia na $(s,s_n]$ dążą do zera według prawdopodobieństwa
na zbiorach o skończonej intensywności.
\end{proof}

\section{Średnia liczby skoków i formuła kompensacji}

\begin{twierdzenie}[Formuła Campbella i kompensacji]
\label{th:campbell}
Niech $N$ będzie losową miarą Poissona z Definicji
\ref{def:PRM}. Dla każdej nieujemnej, mierzalnej funkcji
$g:(0,T]\times E\to[0,\infty]$ mamy
\[
\E\int_{(0,T]\times E}g(t,z)\,N(\dd t,\dd z)
=\int_0^T\int_Eg(t,z)\,\dd\nu(z)\,\dd t.
\]
Jeżeli $H(\omega,t,z)\ge0$ jest przewidywalna względem
$(\calF_t)$ w zmiennej czasowej i borelowska w $z$, to
\[
\E\int H\,\dd N
=\E\int_0^T\int_E H_t(z)\,\dd\nu(z)\,\dd t.
\tag{11.1}\label{eq:compensation}
\]
Równość dopuszcza wartość $+\infty$ po obu stronach.
\end{twierdzenie}
\begin{proof}
Najpierw $g=\ind{(s,t]\times A}$, gdzie $\nu(A)<\infty$:
obie strony wynoszą $(t-s)\nu(A)$. Liniowość daje
wzór dla nieujemnych funkcji prostych, a twierdzenie o
zbieżności monotonicznej dla dowolnych $g$.

Dla losowego $H$ zaczynamy od
$H_r(z)=\xi\ind{(s,t]}(r)\ind A(z)$, gdzie $\xi$ jest
nieujemna i $\calF_s$-mierzalna, a $\nu(A)<\infty$.
Niezależność $N((s,t]\times A)$ od $\calF_s$ implikuje
$\E[\xi N((s,t]\times A)]=\E\xi\,(t-s)\nu(A)$.
Liniowość daje tezę dla sum takich procesów. Ich dodatnie
kombinacje oraz granice monotoniczne generują wszystkie
nieujemne funkcje przewidywalne na
$\Omega\times(0,T]\times E$, ponieważ prostokąty
$F_s\times(s,t]\times A$, $F_s\in\calF_s$,
generują przewidywalne $\sigma$-ciało (po domknięciu
względem prostych funkcji przestrzennych). Zatem
klasa monotoniczna i Tonelli kończą dowód.
\end{proof}

\begin{przyklad}[Ile jest dużych skoków?]
Dla miary Lévy'ego $\nu$ i $\varepsilon>0$ liczba skoków
rozmiaru większego od $\varepsilon$ w czasie $t$ jest
rozkład Poissona o parametrze $t\nu(\{\|z\|>\varepsilon\})$.
Prawdopodobieństwo braku takiego skoku wynosi
$\exp[-t\nu(\{\|z\|>\varepsilon\})]$.
Jeżeli $\nu(\{0<\|z\|\le1\})=\infty$, to dla każdego
$t>0$ występuje nieskończenie wiele małych skoków
prawie na pewno, choć suma ich kwadratów ma skończoną
wartość oczekiwaną, gdyż
$\E\sum_{s\le t}\|\Delta X_s\|^2\ind{\{\|\Delta X_s\|\le1\}}
=t\int_{\|z\|\le1}\|z\|^2\,\dd\nu(z)$ w konstrukcji
procesu Lévy'ego.
\end{przyklad}

\section[Całka skompensowana: konstrukcja w L2]{Całka skompensowana: konstrukcja w przestrzeni $L^2$}

\begin{definicja}[Elementarny proces przewidywalny]
\label{def:prm-step}
Proces $H(\omega,t,z)$ na $[0,T]\times E$ jest
\emph{elementarny}, jeżeli ma postać skończonej sumy
\[
H_t(z)=\sum_{j=0}^{m-1}\sum_{k=1}^{r}
\xi_{jk}\ind{(t_j,t_{j+1}]}(t)\ind{A_k}(z),
\]
gdzie $0=t_0<\cdots<t_m=T$, zbiory $A_k\in\calB(E)$
są parami rozłączne i mają skończoną miarę $\nu$,
a $\xi_{jk}$ są ograniczonymi, $\calF_{t_j}$-mierzalnymi
zmiennymi losowymi. Dla takiego procesu definiujemy
$\int_0^t\int_EH_s(z)\,\widetilde N(\dd s,\dd z)$
przez zastąpienie każdego zliczenia $N((t_j,t_{j+1}]\times A_k)$
jego wartością po odjęciu $(t_{j+1}-t_j)\nu(A_k)$;
na odcinku końcowym stosujemy $t\wedge t_{j+1}$.
\end{definicja}

\begin{twierdzenie}[Izometria skompensowanej całki Poissona]
\label{th:prm-isometry}
Niech $H,K$ będą elementarnymi procesami przewidywalnymi.
Jeżeli $I_t(H)=\int_0^t\int_EH_s(z)\,\widetilde N(\dd s,\dd z)$,
to $I(H)$ jest kwadratowo całkowalnym martyngałem o
prawostronnie ciągłych trajektoriach i
\[
\E I_T(H)I_T(K)
=\E\int_0^T\int_E H_t(z)K_t(z)\,\dd\nu(z)\,\dd t.
\tag{11.2}\label{eq:prm-isometry}
\]
W szczególności $\E I_T(H)^2=\E\int H^2\,\dd t\,\dd\nu$.
\end{twierdzenie}
\begin{proof}
Wspólnie rozdrobnijmy podziały czasu i rozbijmy zbiory
przestrzenne na rozłączne atomy. Oznaczmy
$\Delta\widetilde N_{jk}=N((t_j,t_{j+1}]\times A_k)
-(t_{j+1}-t_j)\nu(A_k)$.
Zmienne te mają średnią zero, wariancję
$(t_{j+1}-t_j)\nu(A_k)$; zmienne odpowiadające
rozłącznym atomom są niezależne. Dla $j<l$ wyraz
$\xi_{jk}\eta_{lq}\Delta\widetilde N_{jk}
\Delta\widetilde N_{lq}$ ma wartość oczekiwaną zero:
wszystkie czynniki poza ostatnim są $\calF_{t_l}$-mierzalne,
a ostatni ma warunkową średnią zero. Dla tego samego $j$
i różnych $k,q$ znika warunkowa kowariancja ze względu
na $\calF_{t_j}$. Zostają wyłącznie wyrazy diagonalne,
z których otrzymujemy
$\E[\xi_{jk}\eta_{jk}(t_{j+1}-t_j)\nu(A_k)]$.
Po zsumowaniu daje to wzór \eqref{eq:prm-isometry}.
Z tej samej warunkowej średniej zero wynika martyngałowość
względem filtracji. Na każdym przedziale elementarny
proces całkowy jest skończoną sumą kompensowanych
procesów zliczających, więc ma trajektorie càdlàg.
\end{proof}

\begin{twierdzenie}[Rozszerzenie na $L^2$ procesów przewidywalnych]
\label{th:prm-extend}
Niech $H$ będzie procesem przewidywalnym i
$\E\int_0^T\int_E |H_t(z)|^2\,\dd\nu(z)\,\dd t<\infty$.
Istnieje jednoznaczny (z dokładnością do nierozróżnialności)
càdlàg martyngał $I(H)$, będący granicą ucp całek
procesów elementarnych $H^{(n)}$ takich, że
$\E\int|H^{(n)}-H|^2\,\dd t\,\dd\nu\to0$.
Zachodzi izometria \eqref{eq:prm-isometry} dla wszystkich
procesów tej klasy i dla każdego $t\in[0,T]$.
\end{twierdzenie}
\begin{proof}
Gęstość procesów elementarnych otrzymujemy następująco.
Wystarczy najpierw rozważyć funkcje ograniczone, wspierane
na zbiorach skończonej miary $\Prob\otimes\dd t\otimes\nu$;
obcięcia wartości i lokalizacja w przestrzeni $E$ pozwalają
przejść do tego przypadku. Prostokąty przewidywalne
$F_s\times(s,t]\times A$ generują odpowiednie $\sigma$-ciało,
a więc ich liniowe kombinacje są gęste w $L^2$ z
argumentu klasy monotonicznej. Dla dowolnej aproksymacji
$H^{(n)}$ izometria daje zbieżność w $L^2$ zmiennych
$I_T(H^{(n)})$. Nierówność maksymalna Dooba zapewnia
\[
\E\sup_{t\le T}|I_t(H^{(n)})-I_t(H^{(m)})|^2
\le4\,\E\int_0^T\int_E|H^{(n)}-H^{(m)}|^2
\,\dd\nu\,\dd t.
\]
Wybieramy podciąg z sumowalnymi kwadratami różnic
w normie supremum, uzyskując przez Borela--Cantellego
jednostajną granicę prawie na pewno. Ponieważ jednostajna
granica funkcji càdlàg jest càdlàg, otrzymujemy taką
wersję $I(H)$. Jednoznaczność granicy w normie $L^2$
i nierówność maksymalna zapewniają niezależność od ciągu
aproksymującego i nierozróżnialność. Martyngałowość
wynika z granic w $L^2$, a izometria z ciągłości iloczynu
skalarnego w $L^2$. Zbieżność ucp zachodzi dla całej
aproksymacji, nie tylko podciągu.
\end{proof}

\begin{wniosek}[Nawias przewidywalny i kowariacja]
\label{cor:prm-angle}
Dla $H,K$ jak w Twierdzeniu \ref{th:prm-extend} zachodzi
\[
\langle I(H),I(K)\rangle_t
=\int_0^t\int_EH_s(z)K_s(z)\,\dd\nu(z)\,\dd s.
\]
W szczególności proces $I_t(H)^2-\int_0^t\int_EH_s(z)^2
\,\dd\nu(z)\,\dd s$ jest martyngałem.
\end{wniosek}
\begin{proof}
Dla elementarnego $H$ obliczamy przyrosty kwadratu.
Po warunkowaniu względem przeszłości przyrost kwadratu
części martyngałowej ma średnią
$\int_{t_j}^{t_{j+1}}\int_EH_s(z)^2\,\dd\nu\,\dd s$;
wyrazy mieszane mają średnią zero. Dla ogólnego $H$
wybieramy aproksymację w normie $L^2$. Izometria i
nierówność Cauchy'ego--Schwarza dają zbieżność
$I_t(H^{(n)})^2\to I_t(H)^2$ w $L^1$, a
$\int_0^t\int(H^{(n)})^2\dd\nu\dd s\to
\int_0^t\int H^2\dd\nu\dd s$ w $L^1$.
Warunkowa definicja martyngału przechodzi więc do granicy.
Tożsamość z kowariacją dostajemy przez polaryzację.
\end{proof}

\section{Skoki całki, wariacja kwadratowa i interpretacja}

\begin{twierdzenie}[Wariacja kwadratowa całki skompensowanej]
\label{th:prm-bracket}
Dla przewidywalnego $H$ spełniającego założenia
Twierdzenia \ref{th:prm-extend} i procesu $M=I(H)$ zachodzi
\[
[M]_t=\int_{(0,t]\times E}H_s(z)^2\,N(\dd s,\dd z),
\qquad
\langle M\rangle_t=\int_0^t\int_EH_s(z)^2\,\dd\nu(z)\,\dd s.
\tag{11.3}\label{eq:prm-brackets}
\]
W szczególności $\Delta M_t=\int_EH_t(z)N(\{t\},\dd z)$
dla wszystkich $t$ poza jednym zbiorem zerowym.
\end{twierdzenie}
\begin{proof}
Najpierw ustalmy $t\le T$ oraz deterministyczny podział
$\pi=\{0=t_0<\cdots<t_m=t\}$; piszemy
$Q_\pi(Y)=\sum_{j=1}^m(Y_{t_j}-Y_{t_{j-1}})^2$.
Dla procesu elementarnego $H^{(n)}$, ograniczonego i
wspieranego na zbiorze o skończonej intensywności, całka
$I(H^{(n)})$ ma na $[0,t]$ skończenie wiele skoków, a między
nimi zmienia się w sposób absolutnie ciągły. Z rozdzielenia
przedziałów zawierających skoki i ciągłych odcinków trajektorii
wynika, że dla każdego ciągu podziałów o średnicy dążącej do zera
\[
 Q_\pi(I(H^{(n)}))\longrightarrow
 \int_{(0,t]\times E}(H_s^{(n)}(z))^2\,N(\dd s,\dd z)
 \quad\text{prawie na pewno}.
\]
To jest własność wariacji kwadratowej funkcji o skończonej
wariacji złożonej ze skończenie wielu skoków, a nie tylko
stwierdzenie o zbieżności procesów w supremum.

Niech teraz $H^{(n)}\to H$ w
$L^2(\Prob\otimes\dd s\otimes\nu)$ i oznaczmy
$M^{(n)}=I(H^{(n)})$, $M=I(H)$. Dla dowolnego martyngału
kwadratowo całkowalnego $D$ o $D_0=0$ ortogonalność
przyrostów daje
$\E Q_\pi(D)=\E D_t^2$. Stosując
$|a^2-b^2|\le |a-b|\,|a+b|$ i nierówność
Cauchy'ego--Schwarza do sum po podziale, otrzymujemy
\[
\begin{split}
\E|Q_\pi(M)-Q_\pi(M^{(n)})|
&\le (\E Q_\pi(M-M^{(n)}))^{1/2}
      (\E Q_\pi(M+M^{(n)}))^{1/2}\\
&\le \norm{H-H^{(n)}}_2
      (\norm H_2+\norm{H^{(n)}}_2).
\end{split}
\]
Oszacowanie jest \emph{jednostajne względem podziału}.
Ponadto formuła kompensacji i ponowne zastosowanie
Cauchy'ego--Schwarza dają
\[
\E\left|\int ((H^{(n)})^2-H^2)\,\dd N\right|
\le\norm{H-H^{(n)}}_2
  (\norm H_2+\norm{H^{(n)}}_2)\longrightarrow0.
\]
Przyjmując najpierw $n$ duże, a następnie dzieląc coraz
drobniej przedział $[0,t]$, dowodzimy więc, że
$Q_\pi(M)\to\int H^2\dd N$ w prawdopodobieństwie.
Ta granica jest wariacją kwadratową $[M]_t$.

W przypadku procesów elementarnych, rozwijając kwadrat
skompensowanych przyrostów, otrzymujemy
$M_t^2-\int_{(0,t]\times E}H_s(z)^2N(\dd s,\dd z)$
jako martyngał. Dla ogólnego $H$ przejście do granicy jest
uprawnione w $L^1$: izometria kontroluje kwadraty $M_t^{(n)}$,
a poprzednie oszacowanie kontroluje całki względem $N$.
Nawias przewidywalny wynika z Wniosku~\ref{cor:prm-angle}.
Tożsamość dla skoków uzyskujemy najpierw dla procesów
elementarnych, a następnie wzdłuż podciągu jednostajnie
zbieżnego prawie na pewno oraz przez aproksymację całki
względem $N$ na zbiorach o skończonej intensywności.
\end{proof}

\begin{uwaga}[Dlaczego nawiasy są różne?]
Dla $M_t=N_t-\lambda t$, gdzie $N$ jest procesem
Poissona, mamy $[M]_t=N_t$, ale
$\langle M\rangle_t=\lambda t$. Pierwszy proces
jest losowym \emph{zrealizowanym} sumowaniem kwadratów
skoków, drugi jego przewidywalną średnią.
Ten drugi pojawia się w izometrii $L^2$.
\end{uwaga}

\begin{zadanie}[Skoki dodatnie i ujemne, wyprowadzenie dwóch nawiasów]
Niech $N^+,N^-$ będą niezależnymi procesami Poissona z
intensywnościami $\lambda_+,\lambda_->0$. Dla
$M_t=(N_t^+-\lambda_+t)-(N_t^--\lambda_-t)$ obliczyć
$[M]_t$ i $\langle M\rangle_t$.

\emph{Rozwiązanie.} Bierzemy $E=\{-1,1\}$,
$\nu=\lambda_+\delta_1+\lambda_-\delta_{-1}$,
$H(t,z)=z$. Przy prawdopodobieństwie jeden zliczenia
$N^+,N^-$ nie skaczą w tej samej chwili. Zatem
$[M]_t=\sum_{s\le t}(\Delta M_s)^2=N_t^++N_t^-$ oraz
$\langle M\rangle_t=(\lambda_++\lambda_-)t$.
Tożsamości weryfikuje też wariancja
$\E M_t^2=t(\lambda_++\lambda_-)$.
\end{zadanie}

\section{Lokalizacja i skoki o nieskończonej aktywności}

\begin{twierdzenie}[Lokalna całkowalność kwadratowa]
\label{th:prm-local}
Jeżeli $H$ jest przewidywalna i
$\int_0^T\int_EH_t(z)^2\,\dd\nu(z)\,\dd t<\infty$
prawie na pewno, to istnieje càdlàg lokalny martyngał
$\int_0^t\int_EH_s(z)\,\widetilde N(\dd s,\dd z)$,
który po zatrzymaniu w odpowiednim ciągu czasów
jest całką z Twierdzenia \ref{th:prm-extend}.
\end{twierdzenie}
\begin{proof}
Niech $A_t=\int_0^t\int_EH_s(z)^2\,\dd\nu(z)\,\dd s$.
Jest procesem ciągłym, rosnącym i adaptowanym.
Definiujemy $\tau_n=\inf\{t:A_t>n\}\wedge T$.
Ciągłość $A$ daje $A_{\tau_n}\le n$, więc proces
$H\ind{\{t\le\tau_n\}}$ należy do przestrzeni $L^2$
z normą nie większą od $n$. Własność zatrzymywania dla
całki skompensowanej wynika najpierw z definicji dla
procesów elementarnych, a następnie przez aproksymację
w normie $L^2$. Całki na kolejnych poziomach lokalizacji
są zgodne do chwili $\tau_n$, więc sklejają się w jeden
proces, określony z dokładnością do nierozróżnialności.
Ponieważ $A_T<\infty$ prawie na pewno, dla każdej
trajektorii $\tau_n=T$ od pewnego $n$.
\end{proof}

\begin{przyklad}[Miara stabilna i dlaczego kompensacja jest potrzebna]
Niech w jednym wymiarze $\dd\nu(z)=c|z|^{-1-\alpha}\dd z$
na $\{0<|z|\le1\}$, gdzie $0<\alpha<2$ i $c>0$.
Wtedy $\int_{|z|\le1}z^2\,\dd\nu(z)=2c/(2-\alpha)$,
więc $M_t=\int_0^t\int_{|z|\le1}z\,\widetilde N(\dd s,\dd z)$
istnieje w $L^2$. Jeśli $\alpha\ge1$, to
$\int_{|z|\le1}|z|\,\dd\nu(z)=\infty$; nie możemy
zastąpić całki skompensowanej zwykłą sumą bezwzględną
skoków. Mimo to jej wariacja kwadratowa jest skończona
prawie na pewno na $[0,T]$, bo
$\E[M]_T=T\,2c/(2-\alpha)<\infty$.
\end{przyklad}

\begin{zadanie}[Czekanie na pierwszy skok i losowa funkcja podcałkowa]
Niech $N_t$ będzie jednorodnym procesem Poissona o intensywności
$\lambda>0$, a $\tau=\inf\{t>0:N_t\ge1\}$ jego pierwszą chwilą
skoku. Zdefiniujmy $H_t=\ind{\{N_{t-}=0\}}$.
Wyznaczyć $M_t=\int_0^tH_s\,\dd(N_s-\lambda s)$, jego
średnią i drugi moment.

\emph{Rozwiązanie.} Lewostronnie ciągły, adaptowany proces
$H$ jest przewidywalny. Do pierwszego skoku jest równy $1$,
a po pierwszym skoku jest równy $0$. W całce względem
$\dd N$ występuje więc dokładnie jeden skok wielkości $1$,
jeśli $\tau\le t$. W całce względem czasu otrzymujemy
długość odcinka przed pierwszym skokiem. Stąd
\[
M_t=\ind{\{\tau\le t\}}-\lambda(t\wedge\tau).
\]
Ponieważ $\Prob(\tau>s)=e^{-\lambda s}$,
$\E(t\wedge\tau)=\int_0^t\Prob(\tau>s)\,\dd s
=(1-e^{-\lambda t})/\lambda$.
Wynika stąd $\E M_t=0$. Z izometrii Poissona,
\[
\E M_t^2=\lambda\E\int_0^tH_s^2\,\dd s
=\lambda\E(t\wedge\tau)=1-e^{-\lambda t}.
\]
To przykład całki z \emph{losowym} współczynnikiem,
który można obliczyć bez znajomości rozkładu całego procesu
całkowego. Próba użycia $\ind{\{N_t=0\}}$ jako współczynnika
na całym przedziale $(0,t]$ korzystałaby z informacji
przyszłej i nie spełniałaby tego argumentu.
\end{zadanie}

\section{Dekompozycja Lévy'ego--Itô: dokładne znaczenie formuły}

\begin{twierdzenie}[Reprezentacja konstruowanego procesu Lévy'ego]
\label{th:levy-ito-made}
Niech $b\in\R^d$, $Q$ będzie symetryczną macierzą
nieujemnie określoną, a $\nu$ miarą Lévy'ego.
Na pewnej przestrzeni probabilistycznej można zbudować
niezależne: $d$-wymiarowy proces Wienera $W$ i miarę
Poissona $N$ na $(0,\infty)\times(\R^d\setminus\{0\})$
o intensywności $\dd t\,\dd\nu(z)$, tak że proces
\[
X_t=bt+Q^{1/2}W_t
+\int_0^t\int_{\|z\|\le1}z\,\widetilde N(\dd s,\dd z)
+\int_0^t\int_{\|z\|>1}z\,N(\dd s,\dd z)
\tag{11.4}\label{eq:levy-ito-made}
\]
ma wykładnik Lévy'ego--Chinczyna odpowiadający trójce
$(b,Q,\nu)$. Całka dużych skoków jest sumą skończenie
wielu wyrazów na każdym przedziale zwartym, nawet jeśli
nie ma skończonej wartości oczekiwanej.
\end{twierdzenie}
\begin{proof}
Z Twierdzenia \ref{th:prm-existence} bierzemy miarę $N$;
z konstrukcji produktowej uniezależniamy od niej $W$.
Ponieważ $\int_{\|z\|\le1}\|z\|^2\,\dd\nu(z)<\infty$,
mała całka skompensowana istnieje na każdym horyzoncie,
a z \eqref{eq:prm-isometry} jest ograniczona w $L^2$.
Ponieważ $\nu(\{\|z\|>1\})<\infty$, zliczenia dużych
skoków są lokalnie skończone. Dla $\varepsilon>0$
zastąpmy małą całkę przez skompensowaną całkę po
$\varepsilon<\|z\|\le1$. Jest to złożony proces
Poissona pomniejszony o jego liniowy średni dryf.
Jego wykładnik, z rachunku w Rozdziale \ref{ch:levy},
wynosi $t\int_{\varepsilon<\|z\|\le1}
(e^{i\langle u,z\rangle}-1-i\langle u,z\rangle)\,\dd\nu(z)$.
Izometria i nierówność Dooba zapewniają zbieżność ucp
małych całek przy $\varepsilon\downarrow0$.
Ponieważ $|e^{iv}-1-iv|\le v^2/2$, po przejściu
do granicy otrzymujemy małoskokowy składnik wykładnika.
Dodając niezależne części: $bt$, $Q^{1/2}W$ i duże skoki,
dostajemy dokładnie formułę \eqref{eq:LK}.
Własności Lévy'ego wynikają z niezależnych stacjonarnych
przyrostów składników i prawostronnej ciągłości ich
trajektorii. Uzupełnia to konstrukcyjny kierunek
Twierdzenia \ref{th:LK}; kierunek klasyfikacyjny pozostaje
przyjętym twierdzeniem pomocniczym.
\end{proof}

\begin{zadanie}[Poisson i wykładnik charakterystyczny z całki]
Dla $\nu=\lambda\delta_a$ na $\R\setminus\{0\}$,
$\lambda>0$, $a\ne0$, obliczyć
$Y_t=\int_0^t\int z\,N(\dd s,\dd z)$ oraz jego wykładnik.

\emph{Rozwiązanie.} Jedynym dopuszczalnym rozmiarem skoku
jest $a$, a liczba skoków do chwili $t$ ma rozkład
Poissona o parametrze $\lambda t$. Zatem $Y_t=aN_t$ i
$\E e^{iuY_t}=\exp(\lambda t(e^{iua}-1))$.
Jeżeli $|a|\le1$, to w trójce z obcięciem $z\ind{\{|z|\le1\}}$
parametr $b=\lambda a$; jeśli $|a|>1$, mamy $b=0$.
To dlatego ta sama suma skoków ma inny zapis współczynnika
$b$ przy innej funkcji obcięcia.
\end{zadanie}

\begin{center}
\begin{tikzpicture}[>=Latex,scale=.9]
\draw[->] (0,0)--(9,0) node[right] {$t$};
\draw[->] (0,-1.3)--(0,3.2) node[above] {$M_t$};
\draw[blue!65!black,very thick]
(0,0)--(1.4,-.27)--(1.4,1.08)--(3,-.05)--(3,1.07)--
(4.6,.75)--(4.6,2.14)--(6.2,1.82)--(6.2,1.12)--(8.5,.72);
\draw[orange!80!black,thick,dashed] (0,0)--(8.4,-1.0);
\node[anchor=west,scale=.85,blue!65!black] at (5.6,2.75)
{skoki kompensowane};
\node[anchor=west,scale=.82,orange!80!black] at (4.9,-.92)
{przewidywalny trend};
\end{tikzpicture}
\par\smallskip\small Schemat procesu skokowego: odjęcie kompensatora
zmienia zachowanie między skokami, ale nie ich rozmiary.
Rysunek jest poglądowy, nie stanowi symulacji.
\end{center}

\section{Granica teorii i przygotowanie do półmartyngałów}
Zdefiniowaliśmy całkę względem skompensowanej miary
Poissona w klasie kwadratowo całkowalnej i przez lokalizację.
Udało się również uzasadnić budowę procesu z zadanej trójki
Lévy'ego bez ukrytego założenia o całkowalności dużych
skoków. Ogólne teoretyczne pojęcie \emph{kompensatora}
losowej miary skoków, której intensywność może zależeć
od przeszłości, wymaga teorii procesów przewidywalnych
i półmartyngałów; do niego wrócimy w rozdziale następnym.

% ===========================================================================
% ROZDZIAL 12 -- Polmartyngaly i wzor Ito ze skokami
% ===========================================================================
\chapter{Półmartyngały i wzór Itô dla procesów ze skokami}
\label{ch:semimartingale-jumps}

\section{Dlaczego połączenie części ciągłej i skokowej wymaga nowego języka?}

Wcześniej skonstruowaliśmy całkę Itô względem procesu Wienera,
a w Rozdziale~\ref{ch:poisson-random-measure} --- względem
skompensowanej miary Poissona. Oba przypadki łączy rozkład
zmienności procesu na część lokalnie martyngałową i część o
skończonej wariacji. Ta obserwacja motywuje klasę
\emph{półmartyngałów}. Wprowadzimy ją na poziomie potrzebnym do
jawnych SDE ze skokami, a głębsze twierdzenia o wszystkich
półmartyngałach sformułujemy z wyraźnym zaznaczeniem zakresu dowodu.

\begin{definicja}[Półmartyngał i półmartyngał szczególny]
\label{def:semimartingale}
Niech $(\calF_t)$ będzie filtracją spełniającą warunki zwykłe,
a $X$ adaptowanym procesem càdlàg na $[0,T]$.
Proces $X$ nazywamy \emph{półmartyngałem}, jeżeli istnieją
càdlàg lokalny martyngał $M$ oraz adaptowany càdlàg proces
$A$ o skończonej wariacji na każdym przedziale $[0,t]$,
takie że $X=X_0+M+A$ i $M_0=A_0=0$.
Jeżeli można wybrać $A$ przewidywalny, to $X$ nazywamy
\emph{półmartyngałem szczególnym}. W takim przypadku
rozkład na przewidywalny $A$ i lokalny martyngał $M$
jest jednoznaczny przy powyższej normalizacji.
\end{definicja}

\begin{proof}[Uzasadnienie jednoznaczności rozkładu szczególnego]
Załóżmy, że $X=X_0+M+A=X_0+M'+A'$, gdzie $A,A'$ są
przewidywalne i mają lokalnie skończoną wariację.
Proces $L=M-M'=A'-A$ jest przewidywalnym lokalnym
martyngałem o skończonej wariacji. Ponieważ jego skoki są
przewidywalne, dla każdego czasu przewidywalnego $\tau$
warunek martyngałowy daje
$\E[\Delta L_\tau\mid\calF_{\tau-}]=0$ po lokalizacji.
Ale $\Delta L_\tau$ jest $\calF_{\tau-}$-mierzalny,
a więc $\Delta L_\tau=0$. Twierdzenie o anonsowaniu
skoków procesu przewidywalnego pozwala wywnioskować,
że $L$ nie ma skoków. Ciągły lokalny martyngał o
skończonej wariacji jest stały (jego wariacja kwadratowa
wynosi zero), a $L_0=0$. Otrzymujemy $L=0$.
\end{proof}

\begin{uwaga}[Granica powyższego argumentu]
Użyte w ostatnim dowodzie twierdzenie o anonsowaniu skoków
przewidywalnych jest standardowym wynikiem ogólnej teorii
procesów; jego pełna konstrukcja nie została tutaj
przeprowadzona. W dalszych konkretnych przykładach
jednoznaczność kompensatorów wynika bezpośrednio z
obliczeń warunkowych. Nie zastępujemy w ten sposób
pełnego dowodu ogólnego twierdzenia Dooba--Meyera.
\end{uwaga}

\begin{przyklad}[Trzy różne półmartyngały]
\begin{enumerate}
\item $X_t=W_t+\int_0^tb_s\,\dd s$: część
lokalnie martyngałowa jest ciągła, a skończona wariacja
pochodzi z dryfu.
\item $X_t=N_t-\lambda t$ jest martyngałem, mimo że
między skokami ma ujemne nachylenie; jego wariacja
całkowita jest skończona na każdym zwartym przedziale.
\item Proces gamma jest niemalejącym procesem o
skończonej wariacji, mimo iż na każdym przedziale
niezerowej długości wykonuje nieskończenie wiele skoków.
\end{enumerate}
Drugi przykład pokazuje, że proces o skończonej wariacji
\emph{nie musi} być przewidywalny ani wolny od części
martyngałowej. Sam wygląd trajektorii nie wyznacza
kanonicznej dekompozycji.
\end{przyklad}

\section{Całka względem półmartyngału: konstrukcja w dostępnej klasie}

Zaczynamy od klasy, którą potrafimy zbudować przy pomocy
już udowodnionych izometrii. Niech $X=X_0+M+A$, gdzie
$M$ jest sumą całek względem procesu Wienera i
skompensowanej miary Poissona, lokalnie kwadratowo
całkowalnych w sensie poprzednich rozdziałów, a $A$
adaptowanym procesem càdlàg o lokalnie skończonej wariacji.
Dla ograniczonego procesu przewidywalnego $H$, dla którego
oba składniki $H$ razy funkcje podcałkowe $M$ należą lokalnie
do odpowiednich przestrzeni $L^2$, definiujemy
\[
\int_0^t H_s\,\dd X_s
:=\int_0^tH_s\,\dd M_s+\int_{(0,t]}H_s\,\dd A_s.
\tag{12.1}\label{eq:semi-integral}
\]
Pierwszą całkę definiujemy przez lokalizację w $L^2$,
a drugą trajektoria po trajektorii jako całkę Lebesgue'a--Stieltjesa.
W miejscu skoku współczynnik ma wartość $H_s$, dostępną
z informacji sprzed skoku w sensie przewidywalności. Jeśli
$H_s=Y_{s-}$ dla adaptowanego procesu càdlàg $Y$, jest to
konkretnie lewostronna granica $Y$; nie każda przewidywalna
funkcja podcałkowa musi mieć postać lewostronnej granicy.

\begin{lemat}[Zatrzymywanie i skok całki]
\label{lem:semi-jumps}
W sytuacji z \eqref{eq:semi-integral}, jeśli $H$ jest
ograniczony i przewidywalny, a $\tau\le T$ czasem zatrzymania,
to
\[
\left(\int_0^\cdot H_s\,\dd X_s\right)^\tau_t
=\int_0^tH_s\ind{\{s\le\tau\}}\,\dd X_s,
\qquad
\Delta\left(\int_0^\cdot H_s\,\dd X_s\right)_t
=H_t\Delta X_t.
\]
\end{lemat}
\begin{proof}
Dla całki względem $A$ oba wzory wynikają bezpośrednio
z definicji całki Stieltjesa z atomami i z własności
miary $\dd A$. Dla elementarnego przewidywalnego $H$
są to tożsamości sum przyrostów lokalnego martyngału.
W ogólnym przypadku wykorzystujemy czasy lokalizacji z konstrukcji
całek Wienera i Poissona, aby odpowiednie całki miały skończone
drugie momenty. Aproksymacja
procesów elementarnych i nierówność maksymalna dla
całek martyngałowych dają zbieżność ucp, zatem
przechodzimy do granicy w pierwszej równości.
Jeśli $Y_n\to Y$ ucp i wszystkie procesy są càdlàg,
to wzdłuż podciągu zachodzi zbieżność jednostajna
prawie na pewno, a stąd $\Delta Y_n(t)\to\Delta Y(t)$
jednostajnie w $t$ na ustalonym zwartym przedziale.
To przenosi równość o skokach, a lokalizacja kończy dowód.
\end{proof}

\begin{uwaga}[Co zostaje poza konstrukcją?]
Dla lokalnych martyngałów, które nie są lokalnie
kwadratowo całkowalne, pełne zdefiniowanie
$H\cdot X$ wymaga ogólnej teorii całkowania względem
półmartyngałów. Jej ścisła wersja obejmuje ciągłość
operatora całkowego w topologii półmartyngałów.
Nie utożsamiamy tego rozszerzenia z samą izometrią $L^2$;
w dalszych dowodach pracujemy w jawnej klasie
procesów Wienera i skompensowanych miar Poissona.
\end{uwaga}

\section{Wariacja kwadratowa: część ciągła i rzeczywiste skoki}

\begin{twierdzenie}[Wariacja kwadratowa w klasie dyfuzji ze skokami]
\label{th:jump-qv}
Niech $X$ będzie rzeczywistym procesem postaci
\[
X_t=x+\int_0^t b_s\,\dd s+\int_0^t\sigma_s\,\dd W_s
+\int_{(0,t]\times E_0}\gamma_s(z)\,\widetilde N(\dd s,\dd z)
+\int_{(0,t]\times E_1}\Gamma_s(z)\,N(\dd s,\dd z),
\]
gdzie $E_0,E_1\subset E$ są rozłączne,
$\nu(E_1)<\infty$, a współczynniki są przewidywalne i spełniają
$\int_0^T(|b_s|+\sigma_s^2+\int_{E_0}\gamma_s(z)^2\dd\nu(z))\dd s
<\infty$ prawie na pewno.
Wówczas dla podziałów deterministycznych o średnicy dążącej
do zera sumy kwadratów przyrostów zbiegają według
prawdopodobieństwa do
\[
[X]_t=\int_0^t\sigma_s^2\,\dd s
+\int_{(0,t]\times E_0}\gamma_s(z)^2\,N(\dd s,\dd z)
+\int_{(0,t]\times E_1}\Gamma_s(z)^2\,N(\dd s,\dd z).
\tag{12.2}\label{eq:jump-qv}
\]
Zbieżność dotyczy każdej ustalonej chwili $t$; przy
silniejszych oszacowaniach po lokalizacji otrzymujemy ucp.
\end{twierdzenie}
\begin{proof}
Najpierw obcinamy zakres małych skoków tak, aby jego
intensywność była skończona. Dla procesu mającego
skończenie wiele skoków na $[0,T]$ oraz ciągłą część
Itô możemy rozbić sumę kwadratów przyrostów na
wyrazy: kwadrat ciągłej części, kwadrat skokowej
części i wyrazy mieszane. Dla ciągłej części
wariacja kwadratowa wynosi $\int\sigma^2\dd s$;
ciągła część skończonej wariacji ma wariację kwadratową
zero. Każdy ustalony skok znajdzie się dla dostatecznie
drobnego podziału w odrębnym przedziale, a ciągłość
pozostałej części oznacza, że wyrazy mieszane dążą do
zera. Otrzymujemy \eqref{eq:jump-qv} dla obciętego procesu.

Dla części o nieskończonej aktywności stosujemy
aproksymację $X^{(n)}$ przez ucięcie miary skoków.
Najpierw lokalizujemy procesy tak, aby całkowita wariacja
składnika dryfu oraz całki kwadratów współczynników
ciągłych i małoskokowych były ograniczone stałą. Ponieważ
dużych skoków jest na $[0,T]$ skończenie wiele i ich rozmiary
są skończone prawie na pewno, możemy także obciąć ich
rozmiary i liczbę; prawdopodobieństwo zmiany pierwotnej
trajektorii przez to obcięcie dąży do zera. W tym
ograniczonym modelu izometria Poissona i nierówność maksymalna
dają $X^{(n)}\to X$ ucp oraz
$\E\int\!|(\gamma^{(n)})^2-\gamma^2|\,\dd N\to0$.
Sama zbieżność ucp nie wystarczyłaby do zbieżności wariacji
kwadratowej. Aby uzasadnić przejście dla sum kwadratów
przyrostów, używamy nierówności Cauchy'ego--Schwarza
\[
\left|\sum_j(\Delta_jX)^2-\sum_j(\Delta_jX^{(n)})^2\right|
\le\left(\sum_j(\Delta_j(X-X^{(n)}))^2\right)^{1/2}
\left(\sum_j(\Delta_j(X+X^{(n)}))^2\right)^{1/2}.
\]
Dla sum martyngałowych po lokalizacji pierwszy czynnik
ma drugi moment kontrolowany przez izometrię, jednostajnie
po podziałach, a drugi ma ograniczony drugi moment
(np. przez zatrzymanie procesów i oszacowania dla części
skończonej wariacji). Zatem najpierw wybieramy $n$ duże,
a potem średnicę podziału małą. Otrzymujemy zbieżność
według prawdopodobieństwa do prawej strony
\eqref{eq:jump-qv}. Dla nieograniczonych dużych skoków
stosujemy lokalizację według liczby i rozmiaru skoków.
\end{proof}

\begin{przyklad}[Skoki nie tworzą członu $\dd t$ w chwili skoku]
Dla $M_t=N_t-\lambda t$ zachodzi $[M]_t=N_t$,
a nie $\lambda t$. Proces $\lambda t$ jest jego nawiasem
przewidywalnym. Natomiast dla niezależnego Wienera
$W$ mamy $[W+M]_t=t+N_t$.
\end{przyklad}

\section{Wzór Itô ze skokami: najpierw przypadek skończony}

\begin{lemat}[Wzór dla skończonej liczby skoków]
\label{lem:finite-jump-ito}
Niech $X=C+J$, gdzie $C$ jest ciągłym procesem Itô
na $[0,T]$, a $J_t=\sum_{k:S_k\le t}U_k$ adaptowanym procesem
o skończenie wielu skokach na $[0,T]$ (bez dodatkowej
ciągłej części skończonej wariacji).
Dla $f\in C^{1,2}([0,T]\times\R)$ zachodzi
\[
\begin{split}
f(t,X_t)-f(0,X_0)={}&\int_0^t f_t(s,X_{s-})\,\dd s
+\int_0^t f_x(s,X_{s-})\,\dd X_s\\
&+\frac12\int_0^t f_{xx}(s,X_{s-})\,\dd[C]_s\\
&+\sum_{0<s\le t}\bigl(f(s,X_s)-f(s,X_{s-})
-f_x(s,X_{s-})\Delta X_s\bigr).
\end{split}
\tag{12.3}\label{eq:ito-finite-jumps}
\]
\end{lemat}
\begin{proof}
Uporządkujmy chwile skoków $J$ w $[0,t]$:
$0<\tau_1<\cdots<\tau_m\le t$. Pomiędzy skokami
proces $X$ jest ciągłym procesem Itô z dodaną
ciągłą częścią procesu o skończonej wariacji.
Zastosowanie ciągłego wzoru Itô na kolejnych przedziałach
$[0,\tau_1)$, $(\tau_1,\tau_2)$, itd., po lokalizacji i
przejściu do granic jednostronnych, daje składniki
całkowe. Przy każdym skoku występuje dokładna różnica
$f(\tau_j,X_{\tau_j})-f(\tau_j,X_{\tau_j-})$.
W całce $\int f_x\,\dd X$ skok uwzględniony jest jako
$f_x(\tau_j,X_{\tau_j-})\Delta X_{\tau_j}$.
Ich różnica to właśnie składnik sumy w tezie.
Sumowanie skończonej liczby przedziałów i skoków
kończy dowód. Nie ma tu przejścia do granicy w liczbie skoków.
\end{proof}

\section{Wzór Itô dla kompensowanych małych skoków}

\begin{twierdzenie}[Wzór Itô dla skokowej dyfuzji]
\label{th:ito-jump-diffusion}
Niech $T>0$, $N$ będzie losową miarą Poissona na
$(0,T]\times E$ o intensywności $\dd s\,\dd\nu(z)$,
$W$ zaś procesem Wienera względem wspólnej filtracji,
niezależnym od $N$. Niech $E=E_0\sqcup E_1$ i
$\nu(E_1)<\infty$. Niech $b,\sigma$ będą rzeczywistymi
procesami przewidywalnymi, a $\gamma,\Gamma$ funkcjami
przewidywalnymi na $\Omega\times(0,T]\times E_0$ i $E_1$
odpowiednio, spełniającymi
\[
\E\int_0^T\left(|b_s|+\sigma_s^2+
\int_{E_0}\gamma_s(z)^2\,\dd\nu(z)\right)\dd s<\infty.
\]
Nie nakładamy warunku momentowego na duże skoki $\Gamma$;
zakładamy jedynie jej wartości rzeczywiste $N$-prawie wszędzie.
Niech
\[
X_t=x+\int_0^tb_s\,\dd s+\int_0^t\sigma_s\,\dd W_s
+\int_{(0,t]\times E_0}\gamma_s(z)\,\widetilde N(\dd s,\dd z)
+\int_{(0,t]\times E_1}\Gamma_s(z)\,N(\dd s,\dd z).
\]
Dla $f\in C^{1,2}([0,T]\times\R)$, którego
$f,f_t,f_x,f_{xx}$ są ograniczone, zachodzi prawie na pewno
jednocześnie dla wszystkich $t\in[0,T]$:
\[
\begin{split}
f(t,X_t)={}&f(0,x)
+\int_0^t\left(f_t+b_sf_x+\tfrac12\sigma_s^2f_{xx}\right)(s,X_{s-})\,\dd s\\
&+\int_0^t\sigma_s f_x(s,X_{s-})\,\dd W_s\\
&+\int_{(0,t]\times E_0}
\bigl[f(s,X_{s-}+\gamma_s(z))-f(s,X_{s-})\bigr]
\,\widetilde N(\dd s,\dd z)\\
&+\int_0^t\int_{E_0}
\bigl[f(s,X_{s-}+\gamma_s(z))-f(s,X_{s-})
-f_x(s,X_{s-})\gamma_s(z)\bigr] \,\dd\nu(z)\,\dd s\\
&+\int_{(0,t]\times E_1}
\bigl[f(s,X_{s-}+\Gamma_s(z))-f(s,X_{s-})\bigr]
\,N(\dd s,\dd z).
\end{split}
\tag{12.4}\label{eq:ito-jump-general}
\]
We wszystkich wyrażeniach $f_t,f_x,f_{xx}$ są oceniane
w punkcie $(s,X_{s-})$, o ile nie zapisano inaczej.
\end{twierdzenie}
\begin{proof}
\emph{Krok 1: obcięcie małych skoków.} Wybieramy rosnące
zbiory $E_0^{(n)}\uparrow E_0$ z
$\nu(E_0^{(n)})<\infty$ i
$\E\int_0^T\int_{E_0\setminus E_0^{(n)}}\gamma_s(z)^2
\,\dd\nu(z)\,\dd s\to0$ (istnieją przez $\sigma$-skończoność
$\nu$ oraz zbieżność zdominowaną).
Zastępujemy w definicji $X$ małą skompensowaną całkę
przez całkę po $E_0^{(n)}$, otrzymując $X^{(n)}$.
Z izometrii \eqref{eq:prm-isometry} i nierówności Dooba
\[
\E\sup_{t\le T}|X_t-X_t^{(n)}|^2
\le4\E\int_0^T\int_{E_0\setminus E_0^{(n)}}\gamma_s(z)^2
\,\dd\nu(z)\,\dd s\longrightarrow0.
\]

\emph{Krok 2: skończenie wiele skoków.} Na $E_0^{(n)}$
całka skompensowana ma postać sumy skoków minus
ciągły proces skończonej wariacji
$\int_0^t\int_{E_0^{(n)}}\gamma_s(z)\,\dd\nu(z)\,\dd s$.
Ten drugi proces ma skończoną wariację prawie na pewno,
ponieważ Cauchy--Schwarz i skończoność $\nu(E_0^{(n)})$
dają całkowalność bezwzględną. Wobec tego
Lemat~\ref{lem:finite-jump-ito} stosuje się do $X^{(n)}$.
Rozpisujemy wszystkie skoki i grupujemy ich wkłady:
całka $f_x\gamma\,\widetilde N$ oraz suma poprawek
$R=f(X_-+\gamma)-f(X_-)-f_x(X_-)\gamma$ dają
po przekształceniu
$[f(X_-+\gamma)-f(X_-)]\,\widetilde N+R\,\dd s\,\dd\nu$.
Duże skoki pozostają nieskompensowane. Otrzymujemy
\eqref{eq:ito-jump-general} z $E_0$ zastąpionym przez
$E_0^{(n)}$ i $X$ przez $X^{(n)}$.

\emph{Krok 3: przejście do granicy w członach ciągłych.}
Przechodzimy do podciągu, po którym $X^{(n)}\to X$
jednostajnie prawie na pewno. Na każdym przedziale
$[0,T]$ trajektoria $X$ jest ograniczona; jednostajna
ciągłość pochodnych $f$ na zwartych prostokątach
i zbieżność zdominowana prowadzą do zbieżności
zwykłych całek po czasie. Dla całki Wienera stosujemy
izometrię i ograniczenie $|f_x|\le C$:
\[
\E\int_0^T\sigma_s^2
|f_x(s,X^{(n)}_{s-})-f_x(s,X_{s-})|^2\,\dd s\to0.
\]
Wynika to z dominacji przez $4C^2\sigma_s^2$ oraz
zbieżności punktowej po podciągu, prawie wszędzie
w $\Prob\otimes\dd s$.

\emph{Krok 4: małe skoki.} Z twierdzenia o wartości
średniej mamy
$|f(s,x+u)-f(s,x)|\le C|u|$ i
$|R(s,x,u)|\le\frac12C u^2$, gdzie $C$ ogranicza
$|f_x|,|f_{xx}|$. Izometria Poissona pozwala więc
przejść do granicy w skompensowanych całkach:
dla wspólnej części $E_0^{(n)}$ różnice argumentów
zbiegają punktowo i są dominowane przez $2C|\gamma|$,
a na usuniętej części $E_0\setminus E_0^{(n)}$
norma kwadratowa dąży do zera. Podobnie poprawki
$R(s,X^{(n)}_{s-},\gamma_s(z))$ zbiegają w
$L^1(\Prob\otimes\dd s\otimes\nu)$ dzięki dominacji
przez $C\gamma^2$. Po lokalizacji pozostaje też
ciągłość składnika dużych skoków: na $[0,T]$ jest ich
skończenie wiele, a argumenty $X^{(n)}_{s-}$ zbiegają
jednostajnie wzdłuż podciągu. Wszystkie człony prawej
strony zbiegają zatem do odpowiednich członów w tezie.
Zbieżność ucp i dobór wspólnego podciągu dla skończonej
rodziny całek zapewniają równość procesów z dokładnością
do nierozróżnialności.
\end{proof}

\begin{uwaga}[Dlaczego poprawka ma rząd kwadratowy?]
Zapis $f(x+u)-f(x)-f'(x)u$ jest błędem liniowego
przybliżenia funkcji. Przy $f\in C^2$ jego wielkość jest
rzędu $u^2$; dlatego właśnie warunek
$\int_{|z|\le1}|z|^2\,\dd\nu(z)<\infty$
wystarcza do całkowania poprawek po nieskończenie wielu
małych skokach. To jest odpowiednik roli wariacji
kwadratowej ruchu Browna w ciągłym wzorze Itô.
\end{uwaga}

\begin{zadanie}[Dlaczego zwykły wzór łańcuchowy jest niepoprawny?]
Niech $N$ będzie procesem Poissona o intensywności
$\lambda>0$ i $f(x)=x^2$. Wyznaczyć $N_t^2$ przez
sumowanie po skokach oraz wskazać brakujący człon
w naiwnym wzorze $\dd f(N_t)=f'(N_{t-})\,\dd N_t$.

\emph{Rozwiązanie.} W każdej chwili skoku $\Delta N_s=1$,
więc $(N_{s-}+1)^2-N_{s-}^2=2N_{s-}+1$.
Sumując po chwilach skoków, otrzymujemy
\[
N_t^2=\int_{(0,t]}2N_{s-}\,\dd N_s+N_t.
\]
Drugi składnik jest dokładnie
$\sum_{s\le t}(\Delta N_s)^2$.
Gdybyśmy bez poprawki zastosowali klasyczną
regułę łańcuchową, zgubilibyśmy jeden przy każdym skoku.
Dla gładkiego $f$ odpowiednia poprawka jest równa
$f(N_{s-}+1)-f(N_{s-})-f'(N_{s-})$.
\end{zadanie}

\section{Iloczyny, wykładnik stochastyczny i zastosowania}

\begin{wniosek}[Reguła iloczynu z kowariacją skokową]
\label{cor:jump-product}
Dla dwóch procesów $X,Y$ klasy z
Twierdzenia~\ref{th:ito-jump-diffusion}, napędzanych
wspólnym procesem Wienera i wspólną miarą Poissona,
zachodzi
\[
X_tY_t=X_0Y_0+\int_0^tX_{s-}\,\dd Y_s
+\int_0^tY_{s-}\,\dd X_s+[X,Y]_t,
\]
gdzie $[X,Y]$ zawiera kowariację części ciągłych
oraz sumę $\sum_{s\le t}\Delta X_s\Delta Y_s$.
\end{wniosek}
\begin{proof}
Stosujemy wzór Itô do dwuwymiarowego procesu
$(X,Y)$ i funkcji $f(x,y)=xy$, po uprzednim zatrzymaniu
procesów na przedziałach, na których $|X|+|Y|$ jest
ograniczone. Poza skokami jedyny składnik drugiego
rzędu to $\dd[X,Y]^c$, a na skoku zachodzi dokładna
tożsamość
$(x+u)(y+v)-xy-yu-xv=uv$.
Sumowanie tych składników jest definicją kowariacji.
Formalną wielowymiarową wersję uzyskujemy również
przez polaryzację jednowymiarowego wzoru dla kwadratu.
\end{proof}

\begin{twierdzenie}[Wykładnik stochastyczny w przypadku skończonej aktywności]
\label{th:dolean-finite}
Niech $Z$ będzie rzeczywistym procesem postaci
$Z_t=\int_0^t a_s\,\dd s+\int_0^t\sigma_s\,\dd W_s
+\sum_{k:S_k\le t}U_k$, gdzie na $[0,T]$ suma zawiera
skończenie wiele skoków, $a,\sigma$ są przewidywalne,
$\int_0^T(|a_s|+\sigma_s^2)\dd s<\infty$ prawie na pewno, a
$U_k>-1$. Wówczas równanie
$\dd Y_t=Y_{t-}\,\dd Z_t$, $Y_0=y_0$, ma dokładnie jedno
adaptowane rozwiązanie càdlàg oraz
\[
Y_t=y_0\exp\left(Z_t-Z_0-\frac12\int_0^t\sigma_s^2\,\dd s
-\sum_{S_k\le t}U_k\right)
\prod_{S_k\le t}(1+U_k).
\tag{12.5}\label{eq:dolean-finite}
\]
\end{twierdzenie}
\begin{proof}
Między kolejnymi skokami równanie jest ciągłym liniowym
SDE. Wzór Itô zastosowany do wykładnika pokazuje,
że jego rozwiązanie ma postać
$Y_s=Y_r\exp(Z_s-Z_r-\frac12\int_r^s\sigma_u^2\dd u)$.
W chwili $S_k$ równanie całkowe daje
$\Delta Y_{S_k}=Y_{S_k-}U_k$, a więc
$Y_{S_k}=Y_{S_k-}(1+U_k)$.
Łącząc te równości dla skończonej liczby skoków,
otrzymujemy \eqref{eq:dolean-finite}. Jednoznaczność
wynika z jednoznaczności liniowego SDE między skokami
i jednoznacznego przeliczenia w każdej chwili skoku.
Dodatniość $1+U_k$ zapewnia zachowanie znaku $Y$;
gdy $U_k=-1$, rozwiązanie po skoku staje się zerem,
co wymaga osobnego uwzględnienia we wzorze logarytmicznym.
\end{proof}

\begin{zadanie}[Cena z losowymi skokami i poprawka dryfu]
Niech $\mu\in\R$, $\sigma\ge0$, $\lambda>0$, $\kappa>-1$
oraz $N$ będzie procesem Poissona o intensywności $\lambda$,
niezależnym od $W$. Rozwiązać
\[
\dd S_t=\mu S_{t-}\,\dd t+\sigma S_{t-}\,\dd W_t
+\kappa S_{t-}\,\dd(N_t-\lambda t),\qquad S_0=s_0>0.
\]
Obliczyć $\E S_t$ i wariancję, jeżeli jest skończona.

\emph{Rozwiązanie.} Między skokami efektywny dryf to
$\mu-\lambda\kappa$. Każdy skok mnoży wartość procesu
przez $1+\kappa$. Otrzymujemy
\[
S_t=s_0\exp\left((\mu-\lambda\kappa-\tfrac12\sigma^2)t
+\sigma W_t\right)(1+\kappa)^{N_t}.
\]
Niezależność daje
$\E e^{\sigma W_t-\sigma^2t/2}=1$ oraz
$\E(1+\kappa)^{N_t}=e^{\lambda t\kappa}$,
zatem $\E S_t=s_0e^{\mu t}$.
Dla drugiego momentu używamy
$\E e^{2\sigma W_t-\sigma^2t}=e^{\sigma^2t}$ oraz
$\E(1+\kappa)^{2N_t}=e^{\lambda t(2\kappa+\kappa^2)}$.
Stąd
\[
\E S_t^2=s_0^2e^{(2\mu+\sigma^2+\lambda\kappa^2)t},
\quad
\Var(S_t)=s_0^2e^{2\mu t}
\bigl(e^{(\sigma^2+\lambda\kappa^2)t}-1\bigr).
\]
To zadanie pokazuje, dlaczego kompensacja zmienia dryf,
ale nie rozmiar pojedynczego skoku.
\end{zadanie}

\begin{zadanie}[Proces Ornsteina--Uhlenbecka wymuszany skokami]
Niech $N$ będzie procesem Poissona o intensywności
$\lambda>0$, $a>0$, a $J_t=c(N_t-\lambda t)$ dla
$c\in\R$. Rozwiązać równanie
$\dd X_t=-aX_{t-}\dd t+\dd J_t$ z deterministycznym
$X_0=x$ oraz obliczyć $\E X_t$ i $\Var(X_t)$.

\emph{Rozwiązanie.} Całkowanie przez części dla
$e^{at}X_t$ (czynnik deterministyczny nie ma skoków)
daje
$X_t=xe^{-at}+c\int_0^te^{-a(t-s)}\,\dd(N_s-\lambda s)$.
Całka ma średnią zero, więc $\E X_t=xe^{-at}$.
Izometria z Rozdziału~\ref{ch:poisson-random-measure} daje
\[
\Var(X_t)=\lambda c^2\int_0^t e^{-2a(t-s)}\,\dd s
=\frac{\lambda c^2}{2a}(1-e^{-2at}).
\]
W odróżnieniu od gaussowskiego procesu
Ornsteina--Uhlenbecka rozkład $X_t$ ma tu, dla
$c\ne0$, składnik atomowy odpowiadający brakowi skoków
na $[0,t]$. Nie jest zatem ogólnie rozkładem normalnym.
\end{zadanie}

\section{Generator dyfuzji ze skokami i równanie wsteczne}

\begin{twierdzenie}[Generator w przypadku jednorodnych współczynników]
\label{th:jump-generator}
Niech $X^x$ będzie czasowo jednorodnym procesem Markowa
spełniającym równanie skokowe
z $b=b(x)$, $\sigma=\sigma(x)$, $\gamma=\gamma(x,z)$,
$\Gamma=\Gamma(x,z)$, z istnieniem rozwiązania i
odpowiednią ciągłością w warunku początkowym.
Dla $f\in C_b^2(\R)$ niech
\[
\begin{split}
Lf(x)={}&b(x)f'(x)+\frac12\sigma(x)^2f''(x)\\
&+\int_{E_0}[f(x+\gamma(x,z))-f(x)-f'(x)\gamma(x,z)]\,\dd\nu(z)\\
&+\int_{E_1}[f(x+\Gamma(x,z))-f(x)]\,\dd\nu(z).
\end{split}
\tag{12.6}\label{eq:jump-generator}
\]
Jeżeli powyższe całki są skończone i $Lf$ jest
ograniczona oraz ciągła, to
$f(X_t^x)-f(x)-\int_0^tLf(X_{s-}^x)\,\dd s$
jest lokalnym martyngałem. Ponadto, gdy można
uzasadnić zbieżność zdominowaną dla $Lf(X_s^x)$ przy
$s\downarrow0$, zachodzi
$\lim_{t\downarrow0}t^{-1}(\E f(X_t^x)-f(x))=Lf(x)$.
\end{twierdzenie}
\begin{proof}
Wstawiamy do wzoru \eqref{eq:ito-jump-general}
$f(t,x)=f(x)$. Dla dużych skoków, których liczba jest
skończona na odcinkach zwartych, formuła kompensacji
\eqref{eq:compensation} pokazuje, że
\[
\int_{(0,t]\times E_1}\Delta_zf(X_{s-})\,N(\dd s,\dd z)
-\int_0^t\int_{E_1}\Delta_zf(X_{s-})\,\dd\nu(z)\,\dd s
\]
jest lokalnym martyngałem; $\Delta_zf(x)=f(x+\Gamma(x,z))-f(x)$.
Składniki Wienera i małych skoków są lokalnymi martyngałami
z poprzednich rozdziałów. Pozostałe człony skończonej
wariacji sumują się dokładnie do $\int Lf(X_{s-})\dd s$.
Po wzięciu wartości oczekiwanych (przy zapewnionej
całkowalności) otrzymujemy
$\E f(X_t^x)-f(x)=\E\int_0^tLf(X_s^x)\dd s$.
Ciągłość stochastyczna w zerze i ciągłość $Lf$ prowadzą
do $\E Lf(X_s^x)\to Lf(x)$ przy $s\downarrow0$,
a ograniczoność $Lf$ pozwala zastosować zbieżność
zdominowaną. Dzielenie przez $t$ kończy dowód.
\end{proof}

\begin{przyklad}[Generator procesu złożonego Poissona]
Dla $X_t^x=x+\sum_{j=1}^{N_t}Y_j$, gdzie
$N$ ma intensywność $\lambda$ i znaczniki mają prawo $\rho$,
$Lf(x)=\lambda\int_{\R}(f(x+y)-f(x))\,\dd\rho(y)$.
Jeśli $Y_j=a$ deterministycznie, mamy
$Lf(x)=\lambda(f(x+a)-f(x))$ i równanie wsteczne
$\partial_tu(t,x)=\lambda[u(t,x+a)-u(t,x)]$.
Dla $u(0,x)=f(x)$ jego rozwiązanie dane przez
półgrupę ma postać
$u(t,x)=\sum_{n\ge0}e^{-\lambda t}(\lambda t)^n f(x+na)/n!$.
Sprawdzenie różniczkowania szeregu jest dozwolone dla
ograniczonej $f$, bo szereg pochodnych po $t$ zbiega
jednostajnie na zwartych przedziałach czasu.
\end{przyklad}

\begin{center}
\begin{tikzpicture}[>=Latex,scale=.95]
\draw[->] (0,0)--(9,0) node[right] {$t$};
\draw[->] (0,-.8)--(0,3.0) node[above] {$X_t$};
\draw[blue!75!black,thick] plot[smooth] coordinates
{(0,.5) (.8,.7) (1.5,.45) (2.3,.8) (3,1.0) (3.8,.65)};
\draw[blue!75!black,thick] (3.8,.65)--(3.8,1.85);
\draw[blue!75!black,thick] plot[smooth] coordinates
{(3.8,1.85) (4.5,2.02) (5.2,1.75) (5.8,2.0)};
\draw[blue!75!black,thick] (5.8,2.0)--(5.8,1.1);
\draw[blue!75!black,thick] plot[smooth] coordinates
{(5.8,1.1) (6.7,1.5) (7.5,1.2) (8.5,1.65)};
\draw[orange!80!black,dashed,thick] (3.8,.65)--(3.8,1.85);
\draw[orange!80!black,dashed,thick] (5.8,1.1)--(5.8,2.0);
\node[anchor=west,scale=.87,orange!80!black] at (6.1,2.55)
{skoki $\Delta X$};
\node[anchor=west,scale=.87,blue!70!black] at (.7,2.52)
{część ciągła};
\end{tikzpicture}
\par\smallskip\small Trajektoria skokowej dyfuzji:
między skokami zachodzi ciągła ewolucja, a w chwilach
losowych pojawiają się nieciągłości. Wzór Itô uwzględnia
obie części i dokładną poprawkę $f(X_s)-f(X_{s-})$.
\end{center}

\section{Co zostało dowiedzione, a co należy do ogólnej teorii?}

W tym rozdziale udowodniliśmy wzór Itô dla
\emph{dyfuzji z małymi skokami całkowanymi w $L^2$
i skończoną intensywnością dużych skoków}.
To wystarcza do całych rodzin procesów Lévy'ego oraz
wielu równań skokowych. Nie dowiedliśmy natomiast
ogólnego twierdzenia charakteryzującego wszystkie
półmartyngały jako dobre całkowniki, ani ogólnego
rozkładu Dooba--Meyera dla podmartyngałów klasy (D).
Podajemy ich właściwe miejsce w teorii, ale nie sugerujemy,
że wynikają z izometrii Poissona bez dalszej pracy.
W kolejnej części tomu zastosujemy przedstawiony rachunek
do równań stochastycznych ze skokami, a następnie wrócimy
do ścisłych całek po trajektoriach i mechaniki kwantowej.

% ===========================================================================
% ROZDZIAŁ 13 -- Stochastyczne równania różniczkowe ze skokami
% ===========================================================================
\chapter{Stochastyczne równania różniczkowe ze skokami}
\label{ch:jump-sde}

\section{Motywacja: dlaczego zwykłe SDE nie opisuje nagłych zmian?}

Model $\dd X_t=b(X_t)\,\dd t+\sigma(X_t)\,\dd W_t$ ma ciągłe
trajektorie. Nie opisuje zatem nagłego uszkodzenia elementu, wpłaty
kapitału ani natychmiastowego przyrostu liczby cząstek. Dodanie
skokowego składnika prowadzi do procesu càdlàg. Jego wartości
bezpośrednio przed skokiem i bezpośrednio po skoku nie są tym samym:
$X_{t-}$ zawiera informację dostępną przed skokiem, a $X_t$ ---
już wynik skoku. Dlatego poprawne równanie ma po lewej stronie
$\dd X_t$, ale w prawych współczynnikach skokowych pojawia się
$X_{t-}$.

Są tu dwie naturalne konwencje. Całka względem $N$ dodaje
rzeczywiste skoki, natomiast całka względem
$\widetilde N=N-\dd t\,\dd\nu$ odejmuje od nich średni dryf.
Pozornie kosmetyczny wybór wpływa na współczynnik dryfu, lecz
nie zmienia skoków procesu.

\begin{definicja}[Równanie skokowe i jego silne rozwiązanie]
\label{def:jump-sde}
Niech $T>0$, $E$ będzie przestrzenią polską z miarą
$\sigma$-skończoną $\nu$, $E=E_0\sqcup E_1$ będzie podziałem
na zbiory borelowskie, gdzie $\nu(E_1)<\infty$.
Na przestrzeni z filtracją spełniającą warunki zwykłe niech $W$
będzie $m$-wymiarowym procesem Wienera, a $N$ losową miarą
Poissona o intensywności $\dd t\,\dd\nu(z)$, przy czym
$(W,N)$ ma niezależne przyrosty względem wspólnej filtracji.
Niech $x\in\R^d$ oraz
\[
\begin{gathered}
b:[0,T]\times\R^d\to\R^d,\quad
\sigma:[0,T]\times\R^d\to\R^{d\times m},\\
\gamma:[0,T]\times\R^d\times E_0\to\R^d,\quad
\Gamma:[0,T]\times\R^d\times E_1\to\R^d
\end{gathered}
\]
będą funkcjami borelowskimi. \emph{Silnym rozwiązaniem} równania
\[
\begin{split}
X_t=x&+\int_0^t b(s,X_{s-})\,\dd s
 +\int_0^t\sigma(s,X_{s-})\,\dd W_s\\
&+\int_{(0,t]\times E_0}\gamma(s,X_{s-},z)\,\widetilde N(\dd s,\dd z)
 +\int_{(0,t]\times E_1}\Gamma(s,X_{s-},z)\,N(\dd s,\dd z)
\end{split}
\tag{13.1}\label{eq:jump-sde}
\]
nazywamy proces adaptowany càdlàg $X$ o $X_0=x$, dla którego
wszystkie całki istnieją we wskazanym sensie, a tożsamość zachodzi
prawie na pewno jednocześnie dla wszystkich $t\in[0,T]$.
Przez $\|\sigma\|_{\mathrm{HS}}$ oznaczamy normę Hilberta--Schmidta.
\end{definicja}

\begin{uwaga}[Dwie skale skoków]
Na $E_1$ występuje skończenie wiele punktów miary $N$ na każdym
przedziale zwartym, zatem składnik z $\Gamma$ jest zwykłą losową
sumą. Zbiór $E_0$ może mieć nieskończoną miarę, a małe skoki
trzeba sumować po kompensacji. Dla samego istnienia procesu
skokowego nie żądamy drugiego momentu $\Gamma$; w pierwszym
twierdzeniu przyjmiemy go celowo, aby podać zupełny dowód
metodą Picarda w przestrzeni $L^2$.
\end{uwaga}

\section{Przygotowanie: zamiana dużych skoków na całkę skompensowaną}

\begin{lemat}[Przepisanie równania w klasie $L^2$]
\label{lem:rewrite-jump-sde}
Załóżmy, że $\nu(E_1)<\infty$ i dla każdego $(t,x)$ zachodzi
$\int_{E_1}\|\Gamma(t,x,z)\|^2\,\dd\nu(z)<\infty$.
Oznaczmy
\[
 \beta(t,x)=b(t,x)+\int_{E_1}\Gamma(t,x,z)\,\dd\nu(z),\qquad
 h(t,x,z)=\gamma(t,x,z)\ind{E_0}(z)+\Gamma(t,x,z)\ind{E_1}(z).
\]
Jeśli obie strony mają sens całkowy, równanie
\eqref{eq:jump-sde} jest równoważne równaniu
\[
 X_t=x+\int_0^t\beta(s,X_{s-})\,\dd s
 +\int_0^t\sigma(s,X_{s-})\,\dd W_s
 +\int_{(0,t]\times E}h(s,X_{s-},z)\,\widetilde N(\dd s,\dd z).
 \tag{13.2}\label{eq:jump-sde-comp}
\]
\end{lemat}
\begin{proof}
Nierówność Cauchy'ego--Schwarza daje
\[
\int_{E_1}\|\Gamma(t,x,z)\|\,\dd\nu(z)
\le \nu(E_1)^{1/2}
\left(\int_{E_1}\|\Gamma(t,x,z)\|^2\,\dd\nu(z)\right)^{1/2}.
\]
Można więc odjąć skończony kompensator dużych skoków:
\[
\int_{(0,t]\times E_1}\Gamma\,\dd N
=\int_{(0,t]\times E_1}\Gamma\,\dd\widetilde N
 +\int_0^t\int_{E_1}\Gamma\,\dd\nu\,\dd s.
\]
Tożsamość sprawdzamy najpierw dla procesów prostych,
a następnie dla ogólnych współczynników przez lokalizację
i izometrię z rozdziału~\ref{ch:poisson-random-measure}.
Zsumowanie składników dryfu daje \eqref{eq:jump-sde-comp}.
\end{proof}

\begin{przyklad}[Różnica pomiędzy dryfem a średnim przyrostem]
Niech $N$ będzie procesem Poissona o intensywności
$\lambda>0$. Dla $X_t=x+bt+cN_t$ mamy
$\E X_t=x+(b+\lambda c)t$. Jeśli zamiast tego zapiszemy
$X_t=x+\widetilde b t+c(N_t-\lambda t)$, to oba opisy
są identyczne dokładnie przy $\widetilde b=b+\lambda c$.
Nie należy więc porównywać współczynników przy $\dd t$ bez
sprawdzenia, czy skoki zostały skompensowane.
\end{przyklad}

\section{Istnienie, jednoznaczność i dokładny dowód metodą Picarda}

\begin{twierdzenie}[Silne istnienie i jednoznaczność w klasie $L^2$]
\label{th:jump-picard}
Niech dane będą jak w Definicji~\ref{def:jump-sde}.
Załóżmy, że istnieją $L,K_0<\infty$ takie, iż dla wszystkich
$t\in[0,T]$ i $x,y\in\R^d$ zachodzą
\[
\begin{split}
&\|b(t,x)-b(t,y)\|^2+
 \|\sigma(t,x)-\sigma(t,y)\|_{\mathrm{HS}}^2\\
&\quad+\int_{E_0}\|\gamma(t,x,z)-\gamma(t,y,z)\|^2\dd\nu(z)
 +\int_{E_1}\|\Gamma(t,x,z)-\Gamma(t,y,z)\|^2\dd\nu(z)
 \le L\|x-y\|^2,
\end{split}
\tag{13.3}\label{eq:jump-lip}
\]
oraz
\[
\sup_{t\in[0,T]}\left(
\|b(t,0)\|^2+\|\sigma(t,0)\|_{\mathrm{HS}}^2
+\int_{E_0}\|\gamma(t,0,z)\|^2\dd\nu(z)
+\int_{E_1}\|\Gamma(t,0,z)\|^2\dd\nu(z)
\right)\le K_0.
\tag{13.4}\label{eq:jump-at-zero}
\]
Wtedy istnieje adaptowane silne rozwiązanie càdlàg $X$
równania \eqref{eq:jump-sde}, jedyne z dokładnością
do nierozróżnialności, i
$\E\sup_{t\le T}\|X_t\|^2<\infty$.
Dla dwóch rozwiązań napędzanych tą samą parą $(W,N)$
i z tym samym warunkiem początkowym zachodzi
$\Prob(X_t=Y_t\ \text{dla wszystkich }t\le T)=1$.
\end{twierdzenie}
\begin{proof}
\emph{Krok 1: kontrola współczynników po kompensacji.}
Niech $\beta,h$ będą z Lematu~\ref{lem:rewrite-jump-sde}.
Z \eqref{eq:jump-lip} i nierówności Cauchy'ego--Schwarza
wynika, z pewną stałą $C_1=C_1(L,\nu(E_1))$,
\[
\|\beta(t,x)-\beta(t,y)\|^2
+\|\sigma(t,x)-\sigma(t,y)\|_{\mathrm{HS}}^2
+\int_E\|h(t,x,z)-h(t,y,z)\|^2\dd\nu(z)
\le C_1\|x-y\|^2.
\]
Zastosowanie tych samych nierówności przy $y=0$ daje
oszacowanie wzrostu
$\|\beta(t,x)\|^2+\|\sigma(t,x)\|_{\mathrm{HS}}^2+
\int_E\|h(t,x,z)\|^2\dd\nu(z)
\le C_2(1+\|x\|^2)$.

\emph{Krok 2: iteracje.} Ustalmy $X^{(0)}_t=x$ i
zdefiniujmy rekurencyjnie $X^{(n+1)}$ prawą stroną
\eqref{eq:jump-sde-comp} z $X^{(n)}_{s-}$ w miejscu $X_{s-}$.
Indukcja pokazuje, że każda iteracja jest adaptowana i càdlàg,
a wszystkie całki są kwadratowo całkowalne. Istotnie,
ze wzrostu liniowego oraz z nierówności maksymalnych
ciągłych i skokowych martyngałów otrzymujemy
\[
\E\sup_{s\le t}\|X^{(n+1)}_s\|^2
\le C_3\left(1+\|x\|^2+
 \int_0^t\E\sup_{u\le r}\|X^{(n)}_u\|^2\dd r\right),
\]
gdzie $C_3$ zależy od $T$ i danych, ale nie od $n$.

\emph{Krok 3: oszacowanie różnic.} Niech
$D^{(n)}=X^{(n+1)}-X^{(n)}$ oraz
$F_n(t)=\E\sup_{s\le t}\|D^{(n)}_s\|^2$.
Dla $n\ge1$ stosujemy kolejno
$\|u+v+w\|^2\le3(\|u\|^2+\|v\|^2+\|w\|^2)$,
$\sup_{s\le t}\|\int_0^s a_r\dd r\|^2
\le t\int_0^t\|a_r\|^2\dd r$, nierówność Dooba
oraz izometrie całek Wienera i Poissona. Dostajemy
\[
 F_n(t)\le C_T\int_0^t F_{n-1}(r)\,\dd r,
 \qquad C_T=3TC_1+12C_1.
 \tag{13.5}\label{eq:jump-picard-est}
\]
Pierwsza różnica spełnia $F_0(t)\le A$ dla $0\le t\le T$,
z pewną skończoną $A$ wynikającą z \eqref{eq:jump-at-zero}.
Indukcja całkowa prowadzi do
\[
 F_n(t)\le A\frac{(C_Tt)^n}{n!}.
 \tag{13.6}\label{eq:jump-picard-factorial}
\]
Tutaj dokładnie pojawia się silnia: każde kolejne
całkowanie po odcinku $[0,t]$ zwiększa wykładnik $t$
i dzieli oszacowanie przez następny indeks.

\emph{Krok 4: granica i równanie całkowe.} Z oszacowania
\eqref{eq:jump-picard-factorial} otrzymujemy
\[
\sum_{n\ge0}\E\sup_{t\le T}\|D^{(n)}_t\|
\le \sqrt A\sum_{n\ge0}\frac{(C_TT)^{n/2}}{\sqrt{n!}}<\infty.
\]
Twierdzenie o zbieżności monotonicznej daje więc
$\sum_n\sup_{t\le T}\|D^{(n)}_t\|<\infty$ prawie na pewno.
Iteracje mają jednostajną granicę $X$, której trajektorie
są càdlàg; proces pozostaje adaptowany.
Jednocześnie iteracje zbiegają w normie
$\E\sup_{t\le T}\|\cdot\|^2$ (nierówność Minkowskiego
i zbieżny szereg pierwiastków $F_n(T)$).
Warunek Lipschitza i oszacowania maksymalne zapewniają
zbieżność każdej z trzech całek w \eqref{eq:jump-sde-comp}
w tej samej normie. Przechodząc do granicy, otrzymujemy
tożsamość procesów, a Lemat~\ref{lem:rewrite-jump-sde}
przywraca pierwotną postać \eqref{eq:jump-sde}.

\emph{Krok 5: jednoznaczność.} Dla dwóch rozwiązań $X,Y$
ze wspólnym szumem, których normy maksymalne mają skończony
drugi moment, stosujemy \eqref{eq:jump-picard-est}
bez zmiany argumentu, zastępując $D^{(n)}$ przez $X-Y$.
W przypadku dowolnych rozwiązań o lokalnie całkowalnych
charakterystykach stosujemy najpierw wspólne czasy
zatrzymania i Lemat~\ref{lem:jump-exit-audit} podany niżej,
a następnie usuwamy lokalizację. Dla rozwiązań o
skończonym momencie maksymalnym wynika stąd
$G(t):=\E\sup_{s\le t}\|X_s-Y_s\|^2
\le C_T\int_0^tG(r)\dd r$.
Lemat Gronwalla daje $G=0$. Równość w normie supremum
daje nierozróżnialność, a nie tylko równość w ustalonej
chwili.
\end{proof}

\begin{uwaga}[Co jeżeli duże skoki nie mają drugiego momentu?]
Metoda Picarda powyżej działa w $L^2$, dlatego wymaga
\eqref{eq:jump-at-zero} dla $\Gamma$. Jeżeli $\nu(E_1)<\infty$,
można zamiast tego rozwiązywać równanie między kolejnymi dużymi
skokami, a w chwili skoku aktualizować stan według
$X_s=X_{s-}+\Gamma(s,X_{s-},Z_s)$. W tej konstrukcji
nie potrzeba momentu drugiego znacznika, ale trzeba
sprawdzić nieeksplozję i jednoznaczność dynamiki między
skokami. Nie jest to już wniosek z podanej izometrii $L^2$.
\end{uwaga}

\section{Zależność od danych, lokalizacja i nieeksplozja}

\begin{twierdzenie}[Stabilność względem stanu początkowego]
\label{th:jump-stability}
Przy założeniach Twierdzenia~\ref{th:jump-picard} niech
$X^x,X^y$ będą rozwiązaniami napędzanymi tymi samymi
$W,N$, z warunkami początkowymi $x,y\in\R^d$.
Istnieje stała $C_T$ zależna od współczynników i $T$, taka że
\[
 \E\sup_{t\le T}\|X_t^x-X_t^y\|^2\le C_T\|x-y\|^2.
 \tag{13.7}\label{eq:jump-stability}
\]
\end{twierdzenie}
\begin{proof}
Odejmując oba równania, otrzymujemy początkowy wyraz
$x-y$ oraz całki z różnic współczynników. Oszacowania
z dowodu Twierdzenia~\ref{th:jump-picard} dają
$G(t)\le4\|x-y\|^2+C\int_0^tG(s)\dd s$.
Lemat Gronwalla dowodzi $G(T)\le4e^{CT}\|x-y\|^2$.
\end{proof}

\begin{twierdzenie}[Warunek lokalny i brak eksplozji]
\label{th:jump-local-explosion}
Niech współczynniki równania \eqref{eq:jump-sde-comp}
będą lokalnie Lipschitzowskie w następującym sensie:
dla każdego $R>0$ istnieje $L_R<\infty$, dla którego
odpowiednik \eqref{eq:jump-lip} zachodzi przy
$\|x\|,\|y\|\le R$. Jeżeli dodatkowo
\[
 \|\beta(t,x)\|^2+\|\sigma(t,x)\|_{\mathrm{HS}}^2
 +\int_E\|h(t,x,z)\|^2\dd\nu(z)
 \le K(1+\|x\|^2)
 \quad(t\in[0,T],\ x\in\R^d),
 \tag{13.8}\label{eq:jump-growth}
\]
to istnieje jedyne globalne silne rozwiązanie równania
\eqref{eq:jump-sde-comp} na $[0,T]$, a
$\E\sup_{t\le T}\|X_t\|^2\le C_{K,T}(1+\|x\|^2)$.
\end{twierdzenie}
\begin{proof}
Niech $\pi_R$ będzie rzutem na domkniętą kulę promienia
$R$. Współczynniki $\beta(t,\pi_Rx)$,
$\sigma(t,\pi_Rx)$, $h(t,\pi_Rx,z)$ są globalnie
Lipschitzowskie, ponieważ rzut jest $1$-Lipschitzowski,
a argumenty pozostają w kuli promienia $R$.
Twierdzenie~\ref{th:jump-picard} daje rozwiązanie $X^R$.
Dwa rozwiązania z różnymi promieniami pokrywają się
(do chwili wyjścia z mniejszej kuli) przez jednoznaczność
z zatrzymaniem. Przy skoku równość obowiązuje także
w chwili wyjścia, gdyż współczynniki skokowe oblicza się
w stanie \emph{przed} skokiem.

Z nierówności maksymalnych i \eqref{eq:jump-growth},
niezależnie od $R$, otrzymujemy
\[
 \E\sup_{t\le T}\|X^R_t\|^2
 \le C(1+\|x\|^2)+C\int_0^T
 \E\sup_{u\le s}\|X^R_u\|^2\dd s
 \le C_{K,T}(1+\|x\|^2).
\]
Dla $\tau_R=\inf\{t\le T:\|X^R_t\|\ge R\}\wedge T$
nierówność Markowa daje
$\Prob(\tau_R<T)\le C_{K,T}(1+\|x\|^2)/R^2$.
Lewa strona dąży do zera, zatem ze zgodności lokalnych
rozwiązań otrzymujemy jedno rozwiązanie na całym $[0,T]$.
Wynik o momencie przechodzi do granicy przez lemat Fatou.
\end{proof}

\begin{uwaga}[Dlaczego lokalność nie wystarcza sama w sobie?]
Dla deterministycznego równania $x'(t)=x(t)^2$, $x(0)=1$,
wyraz szumowy jest równy zero, a dryf jest lokalnie
Lipschitzowski. Rozwiązanie $x(t)=(1-t)^{-1}$ eksploduje
w chwili $1$. Jest to szczególny przypadek SDE, więc
nie wolno wyciągać globalnego istnienia wyłącznie z
lokalnej regularności współczynników.
\end{uwaga}


\section{Lokalizacja i skok opuszczający obszar: uzupełnienie dowodowe}

\begin{lemat}[Zgodność rozwiązań także w chwili skoku wyjściowego]
\label{lem:jump-exit-audit}
Niech $X,Y$ będą dwoma adaptowanymi rozwiązaniami càdlàg
równania \eqref{eq:jump-sde-comp}, napędzanymi tą samą parą
$(W,N)$, z tym samym stanem początkowym. Załóżmy, że
współczynniki spełniają warunek Lipschitza na każdej kuli
w normie $\|\cdot\|^2+\|\cdot\|_{\mathrm{HS}}^2+L^2(\nu)$.
Dla $R>\|x\|$ oznaczmy
$\tau_R=\inf\{t\ge0:\|X_t\|\vee\|Y_t\|\ge R\}\wedge T$.
Wówczas $X_{t\wedge\tau_R}=Y_{t\wedge\tau_R}$ dla
wszystkich $t\in[0,T]$ prawie na pewno.
\end{lemat}
\begin{proof}
Ponieważ $X,Y$ są càdlàg, dla $s\le\tau_R$ zachodzi
$\|X_{s-}\|,\|Y_{s-}\|\le R$, również gdy w chwili
$s=\tau_R$ następuje skok przekraczający promień $R$.
Zatrzymujemy różnicę równań w chwili $\tau_R$. W całce
Poissona po $(0,t\wedge\tau_R]$ uwzględniony jest skok
w chwili $\tau_R$; jego współczynnik jest jednak obliczany
w stanie lewostronnym i nadal spełnia lokalne oszacowanie
Lipschitza. Stosując nierówność Cauchy'ego--Schwarza,
nierówność maksymalną Dooba i izometrie obu całek,
otrzymujemy po standardowej lokalizacji całkowalności
\[
 G_R(t):=\E\sup_{u\le t}
 \|X_{u\wedge\tau_R}-Y_{u\wedge\tau_R}\|^2
 \le C_{R,T}\int_0^t G_R(s)\,\dd s.
\]
Jeżeli zatrzymane całki nie mają z góry skończonego
momentu, najpierw zatrzymujemy je po przekroczeniu poziomu
$n$ przez ich kwadratowe charakterystyki, stosujemy to samo
oszacowanie i dopuszczamy $n\to\infty$; na odcinku do
$\tau_R$ charakterystyki są skończone prawie na pewno.
Lemat Gronwalla daje $G_R=0$. W ten sposób nie gubimy
skoku w chwili wyjścia, co byłoby błędem przy
naśladowaniu argumentu dla trajektorii ciągłych.
\end{proof}

\begin{uwaga}[Które założenie daje globalność?]
Sama lokalna lipschitzowskość zapewnia rozwiązanie do chwili
eksplozji. Oszacowanie liniowego wzrostu
\eqref{eq:jump-growth} daje ograniczenie momentu niezależne
od promienia obcięcia, a zatem $\Prob(\tau_R<T)\to0$.
Oba etapy są logicznie odrębne: zgodność rozwiązań
zatrzymanych nie oznacza jeszcze, że czasy wyjścia
zbiegają do nieskończoności.
\end{uwaga}

\section{Zadania rozwiązane i przykłady interpretacyjne}

\begin{zadanie}[Skokowy proces Ornsteina--Uhlenbecka ze znacznikami]
Niech $a,\lambda>0$, $N$ będzie miarą Poissona na
$(0,\infty)\times\R$ o intensywności
$\lambda\,\dd t\,\dd\rho(z)$, gdzie $\rho$ jest
rozkładem prawdopodobieństwa o skończonym drugim momencie.
Rozwiązać
\[
 X_t=x-a\int_0^tX_s\,\dd s
 +\int_{(0,t]\times\R}z\,\widetilde N(\dd s,\dd z),
\]
a następnie wyznaczyć średnią, wariancję i granicę wariancji.

\emph{Rozwiązanie.} Mnożymy równanie przez deterministyczny
czynnik całkujący $e^{at}$. Nie ma dodatkowego członu
kowariacji, gdyż $e^{at}$ ma skończoną wariację i jest ciągły.
Stąd
\[
 X_t=xe^{-at}+\int_{(0,t]\times\R}
 e^{-a(t-s)}z\,\widetilde N(\dd s,\dd z).
\]
Z kompensacji wynika $\E X_t=xe^{-at}$, natomiast z
izometrii Poissona
\[
 \Var(X_t)=\lambda\E_\rho Z^2\int_0^te^{-2a(t-s)}\dd s
 =\frac{\lambda\E_\rho Z^2}{2a}(1-e^{-2at}).
\]
Wariancja zbiega do $\lambda\E_\rho Z^2/(2a)$, ale
nie oznacza to automatycznie gaussowskości rozkładu
stacjonarnego.
\end{zadanie}

\begin{zadanie}[Intensywność zależna od stanu przez rozrzedzanie]
Niech $\overline N$ będzie miarą Poissona na
$(0,\infty)\times(0,\Lambda]$ o intensywności
$\dd t\,\dd u$, gdzie $0<\Lambda<\infty$, a
$\lambda:\N_0\to[0,\Lambda]$ dowolną funkcją.
Zdefiniujmy proces zliczający wzorem
\[
 X_t=x+\int_{(0,t]\times(0,\Lambda]}
  \ind{\{u\le\lambda(X_{s-})\}}\,\overline N(\dd s,\dd u),
 \qquad x\in\N_0.
\]
Wyjaśnić jego konstrukcję i obliczyć prawdopodobieństwo
jednego przejścia na krótkim przedziale.

\emph{Rozwiązanie.} Na $[0,T]$ miara $\overline N$ ma
skończenie wiele punktów prawie na pewno. Porządkujemy
je chronologicznie; w chwili punktu $(s,u)$ wykonujemy skok
$X_s=X_{s-}+1$, gdy $u\le\lambda(X_{s-})$, a w przeciwnym
razie nie zmieniamy stanu. W ten sposób otrzymujemy
jedyny proces skokowy, niekorzystający z informacji przyszłej.
Na krótkim przedziale $[t,t+h]$, pod warunkiem $X_t=k$,
prawdopodobieństwo zaakceptowania dokładnie jednego punktu
wynosi $\lambda(k)h+O(h^2)$, a prawdopodobieństwo co najmniej
dwóch akceptacji wynosi $O(h^2)$, jednostajnie w $k$.
Istotnie, liczba wszystkich punktów kandydackich ma
rozkład Poissona z parametrem $\Lambda h$, więc zdarzenie
dwóch lub więcej ma prawdopodobieństwo $O(h^2)$; przy
jednym kandydacie znacznik jest jednostajny i daje
prawdopodobieństwo akceptacji $\lambda(k)/\Lambda$.
Generator to $Lf(k)=\lambda(k)(f(k+1)-f(k))$.
\end{zadanie}

\begin{przyklad}[Dlaczego rozwiązania nie muszą zachowywać porządku?]
W całce skokowej sam warunek Lipschitza nie oznacza
porównywania trajektorii. Załóżmy, że w równaniu
$\dd X_t=-2X_{t-}\,\dd N_t$ mamy dwa warunki
początkowe $x=0$ i $y=1$. Przed pierwszym skokiem
$X^0=0$, $X^1=1$, lecz po nim $X^0=0$, $X^1=-1$.
Porządek został odwrócony. Aby uzyskać zasadę porównania,
trzeba dodatkowo kontrolować monotoniczność przekształceń
$x\mapsto x+\Gamma(t,x,z)$ wywoływanych przez skoki.
\end{przyklad}

\section{Związek z generatorami i własnością Markowa}

\begin{stwierdzenie}[Własność Markowa dla jednorodnych współczynników]
\label{prop:jump-markov}
Przy hipotezach Twierdzenia~\ref{th:jump-picard}, jeśli
$b,\sigma,\gamma,\Gamma$ nie zależą jawnie od czasu,
rodzina rozwiązań $(X^x)_{x\in\R^d}$ wyznacza
jednorodny proces Markowa o półgrupie
$P_tf(x)=\E[f(X_t^x)]$ dla każdej ograniczonej funkcji
borelowskiej $f$. Co więcej, proces ma mocną własność Markowa.
\end{stwierdzenie}
\begin{proof}
Iteracje Picarda są mierzalnymi funkcjonałami stanu
początkowego oraz przyrostów pary $(W,N)$; ich granica
jest mierzalna, a jednoznaczność ścieżkowa zapewnia
zgodność rozwiązań przesuniętych w czasie. Przy ustalonej
chwili $s$ przyszłe przyrosty pary $(W,N)$ są niezależne
od $\calF_s$. Stąd dla ograniczonej $f$
$\E[f(X_{s+t}^x)\mid\calF_s]=P_tf(X_s^x)$.
Dla ograniczonego czasu zatrzymania $\tau$ zastępujemy
$s$ przez przybliżenia $\tau_n$ o skończenie wielu
wartościach, z $\tau_n\downarrow\tau$. Prawostronna
ciągłość trajektorii, stabilność w stanie początkowym
i ciągłość stochastyczna w czasie pozwalają przejść
do granicy dla ograniczonych ciągłych $f$.
Dla dowolnej ograniczonej borelowskiej $f$ uogólniamy
wynik przez twierdzenie o klasie monotonicznej.
\end{proof}

\begin{uwaga}[Od jakiej informacji zależy przyszłość?]
Mocna własność Markowa oznacza, że po czasie zatrzymania
$\tau$ dalsza ewolucja zależy od aktualnego stanu
$X_\tau$, a nie od całej przeszłej trajektorii. Nie znaczy
to, że wartości $X_t$ w różnych chwilach są niezależne.
Właśnie zależność od $X_{s-}$ pozwala modelować
sprzężenie bieżącej dynamiki z losowymi rozmiarami skoków.
\end{uwaga}

\section{Podsumowanie i przygotowanie do rozwiązań słabych}
Metoda Picarda wraz z izometrią całki skompensowanej
zapewnia silne rozwiązania skokowych SDE przy globalnym
warunku Lipschitza. Lokalizacja poszerza zakres do
współczynników lokalnie Lipschitzowskich o liniowym
wzroście. Te wyniki nie odpowiadają jednak na pytanie,
czy można skonstruować rozwiązanie, \emph{wybierając również}
przestrzeń probabilistyczną i proces napędzający, ani czy
jednoznaczność rozkładu pociąga za sobą jednoznaczność
trajektorii. Temu poświęcimy rozdział następny.

% ===========================================================================
% ROZDZIAŁ 14 -- Słabe rozwiązania, jednoznaczność co do prawa, Yamada--Watanabe
% ===========================================================================
\chapter{Rozwiązania silne i słabe oraz twierdzenie Yamada--Watanabego}
\label{ch:weak-strong-sde}

\section{Motywacja: co należy ustalić przed rozwiązaniem równania?}

Gdy rozwiązujemy klasyczne równanie różniczkowe, znamy funkcję
wymuszającą i szukamy odpowiedniej trajektorii. W równaniu
stochastycznym sytuacja jest subtelniejsza. Możemy albo z góry
ustalić przestrzeń probabilistyczną, filtrację i proces Wienera,
a następnie szukać procesu $X$ na tej przestrzeni, albo
pozwolić, by przestrzeń i szum były częścią konstruowanego
rozwiązania. Pierwszy problem nazywamy \emph{silnym}, a drugi
\emph{słabym}. Przymiotniki te nie dotyczą rozmiaru współczynników
ani siły zastosowanej normy.

Co więcej, pytanie o \emph{istnienie} rozwiązania i pytanie
o jego \emph{jednoznaczność} są niezależne. Musimy również
odróżnić dwie jednoznaczności: zgodność trajektorii przy
wspólnym szumie oraz zgodność samych praw rozkładu. To
rozróżnienie ma znaczenie nawet wtedy, gdy rozważamy wyłącznie
równania o ciągłych trajektoriach.

\section{Dokładne definicje i podstawowe implikacje}

Niech $T>0$, $d,m\in\N$, $x\in\R^d$,
$b:[0,T]\times\R^d\to\R^d$ oraz
$\sigma:[0,T]\times\R^d\to\R^{d\times m}$ będą
funkcjami borelowskimi. W tym rozdziale rozważamy
\[
 X_t=x+\int_0^t b(s,X_s)\,\dd s
 +\int_0^t\sigma(s,X_s)\,\dd W_s,
 \qquad 0\le t\le T.
 \tag{14.1}\label{eq:weak-strong-main}
\]

\begin{definicja}[Rozwiązanie słabe]
\label{def:weak-solution}
\emph{Rozwiązaniem słabym} równania \eqref{eq:weak-strong-main}
nazywamy układ
$(\Omega,\calF,\Prob,(\calF_t)_{0\le t\le T},W,X)$,
w którym $(\Omega,\calF,\Prob)$ jest przestrzenią
probabilistyczną, filtracja spełnia warunki zwykłe,
$W$ jest $m$-wymiarowym procesem Wienera \emph{względem tej
filtracji}, a $X$ jest ciągłym, adaptowanym procesem
$\R^d$-wartościowym o $X_0=x$ i
\[
 \int_0^T\bigl(\|b(s,X_s)\|+
 \|\sigma(s,X_s)\|_{\mathrm{HS}}^2\bigr)\,\dd s<\infty
 \quad\text{prawie na pewno}.
\]
Wymagamy, aby \eqref{eq:weak-strong-main} zachodziło
prawie na pewno dla wszystkich $t$ jednocześnie.
Przestrzeń probabilistyczna i proces Wienera są częścią
poszukiwanego rozwiązania.
\end{definicja}

\begin{definicja}[Rozwiązanie silne]
\label{def:strong-solution-precise}
\emph{Silnym rozwiązaniem} na zadanej przestrzeni z zadanym
procesem Wienera $W$ nazywamy rozwiązanie
\eqref{eq:weak-strong-main}, adaptowane do uzupełnionej,
prawostronnie ciągłej filtracji naturalnej procesu $W$.
Równoważnie, w powyższym przypadku deterministycznego stanu
początkowego, można wybrać proces $X$ będący nieantycypującym,
mierzalnym funkcjonałem trajektorii $W$. Dla losowego stanu
początkowego należy do danych tego funkcjonału dodać stan
początkowy, niezależny od przyszłych przyrostów szumu.
\end{definicja}

\begin{definicja}[Dwa rodzaje jednoznaczności]
\label{def:path-law-uniq}
Równanie \eqref{eq:weak-strong-main} ma
\emph{jednoznaczność ścieżkową}, jeżeli na każdej wspólnej
przestrzeni probabilistycznej z tym samym procesem Wienera
każde dwa rozwiązania $X,Y$ o tym samym stanie początkowym
są nierozróżnialne, to jest
$\Prob(X_t=Y_t\text{ dla wszystkich }t\le T)=1$.

Mówimy o \emph{jednoznaczności co do prawa}, jeżeli prawa
procesu $X$ na $C([0,T];\R^d)$ są jednakowe dla wszystkich
słabych rozwiązań tego równania o wskazanym stanie początkowym.
Nie jest to z definicji to samo, co jednoznaczność wspólnego
prawa pary $(X,W)$.
\end{definicja}

\begin{uwaga}[Nie mylimy przestrzeni praw z przestrzenią trajektorii]
Równość $\Law(X)=\Law(Y)$ jest równością miar na
$C([0,T];\R^d)$ i nie oznacza, że $X_t(\omega)=Y_t(\omega)$.
Nawet jeśli dwa procesy mają to samo prawo, mogą być
zdefiniowane na różnych przestrzeniach probabilistycznych.
Jednoznaczność ścieżkowa wymaga porównania dwóch procesów
na \emph{tej samej} przestrzeni i przy \emph{tym samym} szumie.
\end{uwaga}

\begin{przyklad}[Istnienie bez jednoznaczności]
Niech $d=m=1$, $\sigma\equiv0$,
$b(x)=2\sqrt{\max\{x,0\}}$ i $x=0$.
Dla każdego $c\ge0$ funkcja
$X_t^{(c)}=(t-c)_+^2$ spełnia
$X_t^{(c)}=\int_0^t2\sqrt{X_s^{(c)}}\,\dd s$.
Istnieje też rozwiązanie identycznie zerowe.
Są to rozwiązania silne na każdej przestrzeni z procesem
Wienera, ale nie zachodzi ani jednoznaczność ścieżkowa,
ani jednoznaczność co do prawa. Nie można więc zastąpić
hipotezy Lipschitza samą ciągłością współczynnika.
\end{przyklad}

\section{Przygotowanie do twierdzenia Yamada--Watanabego}

Dowód nie polega na formalnej manipulacji z prawami
rozkładu. Kluczowa idea jest następująca: \emph{dwa
warunkowo niezależne rozwiązania przy tej samej realizacji
szumu muszą, na mocy jednoznaczności ścieżkowej, być równe}.
Z tego wywnioskujemy, że przy ustalonym szumie rozkład
warunkowy rozwiązania skupia się w jednym punkcie.
Aby uniknąć ukrytego założenia o mierzalności, osobno
uzasadnimy wszystkie kroki tej konstrukcji.

\begin{lemat}[Niezależność przyszłych przyrostów i przyczynowość]
\label{lem:YW-causality}
Niech $(W,X)$ będzie słabym rozwiązaniem
\eqref{eq:weak-strong-main}. Wybierzmy regularny rozkład
warunkowy $K(w,\dd y)$ całej trajektorii $X$ pod warunkiem
$W=w$, gdzie $w\in C([0,T];\R^m)$ i
$y\in C([0,T];\R^d)$. Dla każdego $t\in[0,T]$ oraz
każdej ograniczonej borelowskiej funkcji $F$ zależnej
wyłącznie od $y|_{[0,t]}$ istnieje wersja funkcji
$w\mapsto\int F(y)K(w,\dd y)$ mierzalna względem
$\sigma(w_s:0\le s\le t)$ po uzupełnieniu miary Wienera.
\end{lemat}
\begin{proof}
Zmienna $F(X)$ jest $\calF_t$-mierzalna.
Przyszły proces przyrostów $(W_{t+s}-W_t)_{0\le s\le T-t}$
jest niezależny od $\calF_t$, a więc także od $F(X)$
oraz historii $W$ do chwili $t$. Zatem dla każdej ograniczonej
funkcji $G$ przeszłości $W$ i każdej ograniczonej funkcji $H$
przyszłych przyrostów mamy
\[
 \E[F(X)G(W_{\le t})H(W_{t+\cdot}-W_t)]
 =\E[F(X)G(W_{\le t})]\,\E H(W_{t+\cdot}-W_t).
\]
Twierdzenie o klasie monotonicznej rozciąga tę równość
na ograniczone funkcje całej trajektorii $W$. Oznacza to,
że $\E[F(X)\mid\sigma(W_s:0\le s\le T)]$ jest równe
$\E[F(X)\mid\sigma(W_s:0\le s\le t)]$ prawie na pewno.
Lewa strona jest $\int F(y)K(W,\dd y)$, co daje tezę.
\end{proof}

\begin{lemat}[Równanie Itô jako własność wspólnego prawa]
\label{lem:YW-transfer}
Jeśli dwie pary ciągłych procesów $(X,W)$ i
$(\widehat X,\widehat W)$ mają to samo prawo na
$C([0,T];\R^d)\times C([0,T];\R^m)$, a oba procesy
$W,\widehat W$ są procesami Wienera względem odpowiednich
filtracji zawierających historie $X,\widehat X$, to
spełnianie równania \eqref{eq:weak-strong-main} przez
pierwszą parę implikuje jego spełnianie przez drugą
(z tymi samymi współczynnikami i stanem początkowym).
\end{lemat}
\begin{proof}
Całka $\int_0^t b(s,X_s)\,\dd s$ jest borelowskim
funkcjonałem trajektorii $X$ tam, gdzie istnieje.
Dla całki stochastycznej oznaczmy
$H_s=\sigma(s,X_s)$ i najpierw obetnijmy jej wartości,
$H_s^{(K)}=H_s\ind{\{\|H_s\|\le K\}}$.
Dla ustalonego $K$ uśrednienia lewostronne
\[
 H_s^{(K,n)}=n\int_{(s-1/n)^+}^{s}H_r^{(K)}\,\dd r
\]
są ciągłymi, adaptowanymi funkcjami czasu i zbieżne do
$H^{(K)}$ w $L^2(\Prob\otimes\dd s)$ na podstawie
ciągłości przesunięć w $L^2$ oraz zbieżności zdominowanej.
W dalszym przybliżeniu zastępujemy je wartościami na
lewych końcach drobnego deterministycznego podziału.
Otrzymana całka jest skończoną sumą
$\sum_j H_{t_j}^{(K,n)}(W_{t_{j+1}}-W_{t_j})$,
a więc \emph{borelowskim funkcjonałem pary} $(X,W)$.
Izometria Itô daje zbieżność tych sum do
$\int_0^tH_s^{(K)}\,\dd W_s$ w prawdopodobieństwie.
Po dopuszczeniu $K\to\infty$ i lokalizacji przy
$\int_0^T\|H_s\|^2\,\dd s$ otrzymujemy całą całkę.
Wybierając podciągi, wyrażamy zatem całkę w każdej
wymiernej chwili jako prawie wszędzie określoną granicę
borelowskich funkcjonałów pary $(X,W)$.
Ta sama konstrukcja i te same rozkłady par pozwalają
przenieść tożsamość równania na $(\widehat X,\widehat W)$.
Ciągłość obu stron rozszerza ją ze zbioru chwil wymiernych
na wszystkie $t$.
\end{proof}

\begin{twierdzenie}[Yamada--Watanabe: wersja dla równań Browna]
\label{th:YW}
Niech $b,\sigma$ będą funkcjami borelowskimi określonymi
przed \eqref{eq:weak-strong-main}, a $x\in\R^d$ będzie
ustalony. Załóżmy, że istnieje przynajmniej jedno słabe
rozwiązanie \eqref{eq:weak-strong-main}, i że zachodzi
jednoznaczność ścieżkowa w sensie
Definicji~\ref{def:path-law-uniq}. Wówczas:
\begin{enumerate}
\item istnieje silne rozwiązanie dla każdego $m$-wymiarowego
procesu Wienera na przestrzeni z filtracją naturalną,
po jej zwykłym uzupełnieniu;
\item rozwiązanie jest nieantycypującym mierzalnym funkcjonałem
trajektorii $W$, z dokładnością do nierozróżnialności;
\item zachodzi jednoznaczność co do prawa procesu $X$,
a nawet jednoznaczność wspólnego prawa pary $(X,W)$.
\end{enumerate}
\end{twierdzenie}
\begin{proof}
\emph{Krok 1: warunkowe kopie.} Ustalmy słabe rozwiązanie
i oznaczmy jego wspólne prawo $(W,X)$ przez $\mu$.
Ponieważ przestrzenie ciągłych trajektorii są polskie,
istnieje regularny rozkład warunkowy $K(w,\dd y)$
rozwiązania przy ustalonej trajektorii szumu.
Na przestrzeni kanonicznej trójek $(w,y_1,y_2)$
wprowadzamy miarę
\[
 \mathsf Q(\dd w,\dd y_1,\dd y_2)
 =\mu_W(\dd w)K(w,\dd y_1)K(w,\dd y_2),
\]
gdzie $\mu_W$ jest miarą Wienera. Warunkowo względem
$w$ procesy $y_1,y_2$ są niezależnymi kopiami z prawa
$K(w,\cdot)$.

\emph{Krok 2: wspólny proces Wienera.} Nie możemy po prostu
założyć, że $w$ jest procesem Wienera względem filtracji
zawierającej historie obu $y_i$: trzeba to udowodnić.
Niech $\calG_t$ będzie uzupełnioną filtracją generowaną
przez $w_s,y_{1,s},y_{2,s}$ dla $s\le t$ (następnie bierzemy
jej wersję prawostronnie ciągłą). Dla ograniczonej funkcji
$F$ zależnej od tych historii jej warunkowa średnia pod
warunkiem całego $w$ jest funkcją historii $w_{\le t}$,
ponieważ warunkowe prawo dwóch kopii jest iloczynem
$K(w,\cdot)$ i zachodzi Lemat~\ref{lem:YW-causality}.
Przyszłe przyrosty $w$ są niezależne od jego przeszłości,
zatem niezależne od $\calG_t$. Przejście do filtracji
prawostronnie ciągłej uzasadnia się przez przybliżenie
chwil $t$ z góry i prawostronną ciągłość trajektorii.
Proces $w$ pozostaje więc procesem Wienera względem
wspólnej filtracji.

\emph{Krok 3: obie kopie rozwiązują to samo równanie.}
Z definicji $\mathsf Q$ każda para $(y_i,w)$ ma prawo
$\mu$. Lematem~\ref{lem:YW-transfer} przenosimy zatem
równanie Itô na każdą z par, względem wspólnej filtracji
z kroku 2. Obie kopie mają wspólny szum $w$ i wspólny
stan początkowy. Z jednoznaczności ścieżkowej wynika
$y_1=y_2$ jako trajektorie, $\mathsf Q$-prawie na pewno.

\emph{Krok 4: skupienie rozkładu warunkowego w punkcie.}
Warunkowa niezależność i poprzednia równość dają
\[
 \int K(w,\dd y_1)K(w,\dd y_2)
 \ind{\{y_1\ne y_2\}}=0
 \quad\text{dla }\mu_W\text{-prawie każdego }w.
\]
Lemat~\ref{lem:dirac-audit}, udowodniony poniżej,
pokazuje, że taka miara musi być miarą Diraca.
Wobec tego $K(w,\cdot)=\delta_{F(w)}$ prawie wszędzie,
dla pewnego mierzalnego odwzorowania
$F:C([0,T];\R^m)\to C([0,T];\R^d)$.
Lemat~\ref{lem:YW-causality} pokazuje, że ograniczenie
$F(w)|_{[0,t]}$ zależy, z dokładnością do zbiorów zerowych,
wyłącznie od $w|_{[0,t]}$. Stąd proces $F(W)$ jest
adaptowany do uzupełnionej filtracji naturalnej i,
na mocy Lematu~\ref{lem:YW-transfer}, spełnia SDE.
Otrzymaliśmy silne istnienie. To samo odwzorowanie możemy
zastosować do dowolnego innego procesu Wienera, gdyż jego
prawo jest $\mu_W$.

\emph{Krok 5: jednoznaczność prawa.} Dla dwóch dowolnych
słabych rozwiązań wybierzmy ich jądra warunkowe
$K_1(w,\cdot)$ i $K_2(w,\cdot)$ i powtórzmy kroki 1--3,
tym razem losując jedną kopię z $K_1$, a drugą z $K_2$.
Jednoznaczność ścieżkowa znów wymusza równość kopii.
Ponieważ obie są warunkowo niezależne, ich jądra są
jednocześnie jednakowymi miarami Diraca. Otrzymujemy
to samo prawo pary $(X,W)$, a po marginalizacji ---
to samo prawo $X$.
\end{proof}

\begin{uwaga}[Zakres twierdzenia]
Powyższy dowód korzysta ze standardowego twierdzenia o
istnieniu regularnych rozkładów warunkowych na przestrzeniach
polskich oraz z konstrukcji miary na iloczynie za pomocą jąder.
Nie jest to bezpośredni wniosek z elementarnego rachunku
Itô; potrzebny jest też Lemat~\ref{lem:YW-transfer}.
Analogiczna idea działa dla wielu równań napędzanych
procesami ze skokami, ale wymaga ponownej kontroli
przyczynowości i kanonicznej przestrzeni losowych miar.
\end{uwaga}


\section{Doprecyzowanie argumentu warunkowego w twierdzeniu Yamada--Watanabego}

\begin{lemat}[Kiedy warunkowy rozkład jest miarą Diraca?]
\label{lem:dirac-audit}
Niech $S$ będzie przestrzenią polską, $K$ miarą
probabilistyczną na $S$ i niech $Y_1,Y_2$ będą
niezależnymi zmiennymi o rozkładzie $K$. Jeżeli
$\Prob(Y_1=Y_2)=1$, to $K=\delta_y$ dla pewnego $y\in S$.
Jeśli ponadto $w\mapsto K(w,\cdot)$ jest jądrem
probabilistycznym przyjmującym prawie wszędzie miary
Diraca, to jego atom $F(w)$ można wybrać mierzalnie
względem uzupełnienia odpowiedniej miary bazowej.
\end{lemat}
\begin{proof}
Wybieramy przeliczalną rodzinę zbiorów borelowskich
rozdzielającą punkty $S$ (na przykład kule o wymiernych
promieniach i środkach z przeliczalnego zbioru gęstego).
Dla każdego zbioru $A$ z tej rodziny niezależność daje
$\Prob(Y_1\in A,Y_2\notin A)=K(A)(1-K(A))=0$.
Wszystkie te wartości są więc równe zeru lub jedności.
Dwie niezależne realizacje należą prawie na pewno do
identycznych elementów całej rodziny, a rodzina rozdziela
punkty, więc są równe. Gdyby $K$ nie było miarą Diraca,
istniałby zbiór borelowski o prawdopodobieństwie ściśle
między zerem a jednością: można wziąć dwie rozłączne
okolice punktów należących do nośnika $K$.
Otrzymujemy sprzeczność.

Dla mierzalności rozważamy odwzorowanie
$\delta:S\to\mathcal P(S)$, $y\mapsto\delta_y$.
Jest ono ciągłe i różnowartościowe. Jego obraz jest
nawet domknięty: jeśli $\delta_{y_n}$ zbiega słabo do
$\mu$, to ciasność zbieżnej rodziny daje wspólny
zwarty zbiór zawierający prawie wszystkie $y_n$;
podciąg $y_{n_k}\to y$ pokazuje $\mu=\delta_y$.
Każdy zbieżny podciąg punktów ma tę samą granicę,
więc $y_n\to y$. Odwrotność obrazu jest zatem
ciągła, a w szczególności borelowska. Na zbiorze pełnej miary, na którym
$K(w,\cdot)$ jest miarą Diraca, kładziemy
$F(w)=\delta^{-1}(K(w,\cdot))$; poza nim wybieramy
ustalony punkt. To daje wymagany funkcjonał.
\end{proof}

\begin{uwaga}[Dlaczego sama równość praw nie wystarcza?]
W dowodzie Twierdzenia~\ref{th:YW} wspólne prawo pary
$(X,W)$ nie jest zastępnikiem warunku adaptacji.
Najpierw potrzebujemy Lematu~\ref{lem:YW-causality}, aby
udowodnić, że warunkowe kopie nie mają informacji o
przyszłych przyrostach $W$. Dopiero wtedy wspólny proces
$W$ pozostaje procesem Wienera względem filtracji
powiększonej o historie obu kopii i wolno zastosować
jednoznaczność ścieżkową. Lemat~\ref{lem:dirac-audit}
uzasadnia ostatni etap, w którym warunkowe prawo
zostaje zastąpione mierzalnym funkcjonałem $F(W)$.
\end{uwaga}

\section{Przykłady i zastosowania jednoznaczności}

\begin{zadanie}[Silne rozwiązanie jest automatycznie słabym]
Przypuśćmy, że współczynniki $b,\sigma$ spełniają
globalny warunek Lipschitza i liniowego wzrostu.
Co dokładnie wynika o wszystkich trzech rodzajach
jednoznaczności?

\emph{Rozwiązanie.} Z twierdzenia konstrukcyjnego
rozdziału~5 istnieje silne rozwiązanie i zachodzi
jednoznaczność ścieżkowa. W szczególności rozwiązanie
silne jest słabym, bo jego filtracja i szum spełniają
Definicję~\ref{def:weak-solution}. Hipotezy
Twierdzenia~\ref{th:YW} są spełnione, a więc zachodzi
jednoznaczność co do prawa, także dla dowolnych słabych
rozwiązań konstruowanych na innych przestrzeniach.
Rozwiązanie silne nie jest z definicji ``bardziej
regularną'' trajektorią; chodzi wyłącznie o źródło
losowości i mierzalność względem niego.
\end{zadanie}

\begin{twierdzenie}[Słabe istnienie przy ograniczonym dryfie]
\label{th:bounded-borel-weak}
Niech $b:[0,T]\times\R\to\R$ będzie ograniczoną funkcją
borelowską, a $x\in\R$. Równanie
\[
 X_t=x+\int_0^t b(s,X_s)\,\dd s+W_t
 \tag{14.2}\label{eq:borel-drift}
\]
ma słabe rozwiązanie, a jego prawo na $C([0,T];\R)$
jest jednoznaczne.
\end{twierdzenie}
\begin{proof}
\emph{Istnienie.} Na przestrzeni z procesem Wienera $B$
pod miarą $\mathsf P$ połóżmy $X_t=x+B_t$ i
\[
 Z_T=\exp\left(\int_0^Tb(s,X_s)\,\dd B_s
 -\tfrac12\int_0^Tb(s,X_s)^2\,\dd s\right).
\]
Ograniczoność $b$ zapewnia warunek Novikova,
więc $\E_{\mathsf P}Z_T=1$ i możemy określić
$\dd\mathsf Q=Z_T\dd\mathsf P$.
Twierdzenie Girsanowa z rozdziału~8 mówi, że
$W_t=B_t-\int_0^tb(s,X_s)\dd s$ jest
procesem Wienera pod $\mathsf Q$.
Wobec tego $X=x+\int b(s,X_s)\dd s+W$,
co daje słabe rozwiązanie.

\emph{Jednoznaczność prawa.} Niech $\widehat X$ będzie
dowolnym słabym rozwiązaniem napędzanym przez
$\widehat W$ pod miarą $\widehat{\mathsf P}$.
Girsanow zastosowany z przeciwnym znakiem do
$\theta_s=b(s,\widehat X_s)$ daje miarę
$\widehat{\mathsf Q}$, równoważną
$\widehat{\mathsf P}$, pod którą
$\widehat X_t-x=\widehat W_t+
\int_0^tb(s,\widehat X_s)\dd s$ jest
standardowym procesem Wienera. Jednocześnie
\[
 \frac{\dd\widehat{\mathsf P}}{\dd\widehat{\mathsf Q}}
 =\exp\left(\int_0^T b(s,\widehat X_s)\,\dd\widehat X_s
 -\frac12\int_0^Tb(s,\widehat X_s)^2\,\dd s\right).
\]
Całka po prawej jest całką Itô względem Browna
$\widehat X-x$ pod $\widehat{\mathsf Q}$ i jest
mierzalnym funkcjonałem trajektorii $\widehat X$
z dokładnością do zbioru zerowego, co uzasadnia
Lemat~\ref{lem:YW-transfer}. Ponieważ pod
$\widehat{\mathsf Q}$ prawo $\widehat X$ jest zawsze
prawem procesu $x+B$, powyższy wzór wyznacza jednoznacznie
jego prawo pod $\widehat{\mathsf P}$.
\end{proof}

\begin{uwaga}[Czego twierdzenie o ograniczonym dryfie nie dowodzi?]
W Twierdzeniu~\ref{th:bounded-borel-weak} skonstruowano
\emph{słabe} rozwiązanie i wykazano jednoznaczność co do
prawa. Nie udowodniono na tej podstawie jednoznaczności
ścieżkowej ani istnienia silnego rozwiązania przy dowolnym
ograniczonym borelowskim dryfie. Takie wzmocnienia wymagają
osobnych argumentów regularizacyjnych.
\end{uwaga}

\begin{przyklad}[Dlaczego sama jednoznaczność prawa nie zastępuje ścieżkowej?]
Dwa procesy $X,Y$ mogą być niezależnymi ruchami Browna.
Wówczas $\Law(X)=\Law(Y)$, ale
$\Prob(X_t=Y_t\text{ dla każdego }t\in[0,T])=0$.
Ta obserwacja dotyczy ogólnie praw procesów. Nie jest
kontrprzykładem do Twierdzenia~\ref{th:YW}, ponieważ
procesy nie zostały przedstawione jako dwa rozwiązania
tego samego SDE napędzanego przez \emph{wspólny} proces
Wienera.
\end{przyklad}

\begin{center}
\begin{tikzpicture}[>=Latex,node distance=2cm]
\node[draw,rounded corners,align=center,minimum width=3cm,minimum height=.7cm]
 (w) {zadany szum $W$};
\node[draw,rounded corners,align=center,below left=1.3cm and 1.9cm of w,
 minimum width=3.3cm,minimum height=.9cm] (strong)
 {rozwiązanie silne\\$X=F(W)$};
\node[draw,rounded corners,align=center,below right=1.3cm and 1.9cm of w,
 minimum width=3.4cm,minimum height=.9cm] (weak)
 {rozwiązanie słabe\\$\Law(W,X)$};
\draw[->,thick,blue!70!black] (w)--(strong)
 node[midway,left,scale=.85] {konstrukcja};
\draw[->,thick,blue!70!black] (w)--(weak)
 node[midway,right,scale=.85] {wybór prawa};
\draw[->,thick,dashed] (strong)--(weak)
 node[midway,below,scale=.85] {implikuje};
\end{tikzpicture}
\par\smallskip\small
Różnica między ujęciem silnym i słabym: w pierwszym
szum jest dany przed konstrukcją rozwiązania, a w drugim
możemy wybierać go razem z przestrzenią probabilistyczną.
\end{center}

\section{Podsumowanie i kierunki dalszej teorii}

Otrzymaliśmy wnikliwe rozróżnienie istnienia silnego i słabego,
jednoznaczności ścieżkowej oraz jednoznaczności co do prawa.
Twierdzenie Yamada--Watanabego wyjaśnia, dlaczego przy
istnieniu słabego rozwiązania jednoznaczność ścieżkowa
wymusza silne istnienie: rozkład warunkowy rozwiązania
przy danej trajektorii szumu musi się skoncentrować w jednym
punkcie. Dowód wymagał regularnych rozkładów warunkowych,
argumentu przyczynowości i kontroli całek Itô przy
przenoszeniu praw par procesów.

W kolejnych rozdziałach rozwiniemy całki po trajektoriach
w języku miary Wienera i wprowadzimy fizyczne całki
Feynmana, z uwzględnieniem różnicy pomiędzy miarą
probabilistyczną i oscylacyjnym formalizmem propagatora.




% ===========================================================================
% ROZDZIAŁ 15: całki po trajektoriach, most Browna, jądra Feynmana--Kaca
% ===========================================================================
% Część końcowa: ścisłe całki funkcyjne i rachunek operatorowy.
% W całym rozdziale 16 twierdzenie spektralne, klasyczne jądra gaussowskie
% i podstawowa teoria generatorów półgrup są jawnie przyjmowanymi narzędziami.
\chapter{Całki po trajektoriach i funkcjonały miary Wienera}
\label{ch:path-integrals}
\paragraph{Odsyłacze międzytomowe.} Wykorzystanie metod próbkowania do aproksymacji całek omawia tom IV, rozdział~\ref{IV-ch:monte-carlo}.


\section{Motywacja: całkujemy po przestrzeni funkcji}

W zwykłej całce Lebesgue'a argumentem funkcji podcałkowej jest liczba
albo punkt przestrzeni euklidesowej. W całce \emph{po trajektoriach}
argumentem jest cała funkcja $w:[0,T]\to\R^d$. Przykładowo
$F(w)=\exp(-\int_0^T V(w(s))\,\dd s)$ zależy od tego, w jakich
obszarach przestrzeni przebywała trajektoria w całym przedziale czasu.
Chcemy nadać ścisły sens wyrażeniu $\int F(w)\,\dd\mu(w)$, a nie
wprowadzać nieokreślony symbol $\mathcal D w$.

W rozdziale~\ref{ch:cameron-martin} skonstruowaliśmy miarę Wienera.
Teraz użyjemy jej do obliczeń, zbudujemy \emph{miarę mostu Browna},
a następnie przejdziemy od funkcjonałów trajektorii do jądra
operatora Feynmana--Kaca. Jest to punkt wyjścia do matematycznej
interpretacji niektórych całek stosowanych w mechanice kwantowej.
Przypominamy rozróżnienie: $\int_0^T H_s\,\dd W_s$ jest zmienną
losową, natomiast $\int F(w)\,\dd\mu_W(w)$ jest liczbą, o ile
funkcjonał $F$ jest całkowalny.

\begin{definicja}[Kanoniczna przestrzeń trajektorii i miara Wienera]
\label{def:path-wiener}
Niech $d\in\N$, $T>0$ i $x\in\R^d$. Oznaczmy
$\calC_{x,T}=\{w\in C([0,T];\R^d):w(0)=x\}$, z normą
supremum i borelowskim $\sigma$-ciałem. Niech $B$ będzie
$d$-wymiarowym standardowym ruchem Browna. Miara
$\mu_{x,T}:=\Law((x+B_s)_{0\le s\le T})$ na $\calC_{x,T}$
nazywa się \emph{miarą Wienera rozpoczynającą się w $x$}.
Jeżeli $F:\calC_{x,T}\to\R$ jest funkcjonałem borelowskim
i $\int|F|\,\dd\mu_{x,T}<\infty$, to jego \emph{całką
po trajektoriach} jest $\int_{\calC_{x,T}}F(w)\,\dd\mu_{x,T}(w)$.
\end{definicja}

Definicja nie używa fikcyjnego „jednostajnego rozkładu” na
nieskończeniewymiarowej przestrzeni funkcji. Miara Wienera jest
skoncentrowana na trajektoriach ciągłych, ale prawie na pewno
nigdzie nieróżniczkowalnych. Tych własności nie można ignorować
w całkach po trajektoriach.

Dla $t>0$ oznaczmy przez
$p_t(x,y)=(2\pi t)^{-d/2}\exp(-\|y-x\|^2/(2t))$
jądro cieplne generatora $\frac12\Delta$.

\begin{twierdzenie}[Całki cylindryczne i rozkłady skończeniewymiarowe]
\label{th:path-cylinder}
Niech $0<t_1<\cdots<t_n\le T$, niech
$\Phi:(\R^d)^n\to\R$ będzie borelowska i całkowalna
względem rozkładu $(x+B_{t_1},\ldots,x+B_{t_n})$.
Wówczas
\[
\begin{split}
 &\int_{\calC_{x,T}}\Phi(w(t_1),\ldots,w(t_n))\,
       \dd\mu_{x,T}(w)\\
 &\quad=\int_{(\R^d)^n}\Phi(x_1,\ldots,x_n)
   \prod_{k=1}^n p_{t_k-t_{k-1}}(x_{k-1},x_k)
   \,\dd\lambda_{dn}(x_1,\ldots,x_n),
\end{split}
\]
przy konwencji $t_0=0$, $x_0=x$. Funkcjonały cylindryczne
wyznaczają jednoznacznie całki wszystkich ograniczonych
borelowskich funkcjonałów na $\calC_{x,T}$.
\end{twierdzenie}
\begin{proof}
Przyrosty $B_{t_k}-B_{t_{k-1}}$ są niezależne i mają gęstości
$p_{t_k-t_{k-1}}(0,\cdot)$. Odwzorowanie liniowe wysyłające
przyrosty na kolejne położenia ma wyznacznik równy $1$,
więc wspólna gęstość położeń jest iloczynem zapisanym w tezie.
Otrzymujemy wzór najpierw dla funkcji wskaźnikowych
prostokątów, następnie przez liniowość i twierdzenie o
zbieżności monotonicznej dla wszystkich nieujemnych
funkcji borelowskich, a przez rozdzielenie części dodatniej
i ujemnej dla całkowalnych.

Wreszcie borelowskie $\sigma$-ciało na przestrzeni
ciągłych trajektorii jest generowane przez odwzorowania
$w\mapsto w(q)$ dla wymiernych $q\in[0,T]$.
Istotnie, norma supremum dwóch trajektorii jest supremum
odległości w chwilach wymiernych. Klasa zbiorów, dla
których miary się zgadzają, jest klasą Dynkina zawierającą
cylindry i ich skończone przecięcia; twierdzenie o
$\pi$-$\lambda$-układach daje jednoznaczność.
\end{proof}

\section{Most Browna: prawidłowe warunkowanie względem końca trajektorii}

Zdarzenie $\{x+B_T=y\}$ ma prawdopodobieństwo zero dla
każdego ustalonego $y$. Dlatego nie możemy zdefiniować mostu
przez elementarny iloraz prawdopodobieństw.
Rozwiązaniem jest \emph{regularny rozkład warunkowy}, który
możemy w tym przypadku skonstruować wprost.

\begin{definicja}[Most Browna pomiędzy ustalonymi końcami]
\label{def:brownian-bridge-path}
Niech $x,y\in\R^d$, $T>0$, a $B$ będzie standardowym
$d$-wymiarowym ruchem Browna. Proces
\[
 \beta_s^{x,y,T}=x+\frac{s}{T}(y-x)+B_s-\frac{s}{T}B_T,
 \qquad 0\le s\le T,
\]
nazywamy \emph{mostem Browna z $x$ do $y$ w czasie $T$}.
Jego prawo na
$\calC_{x,y,T}:=\{w\in C([0,T];\R^d):w(0)=x,\ w(T)=y\}$
oznaczamy przez $\mu_{x,y,T}^{\mathrm{br}}$.
\end{definicja}

\begin{twierdzenie}[Konstrukcja gaussowska i warunkowanie]
\label{th:bridge-conditional}
Dla każdego $x,y\in\R^d$ proces $\beta^{x,y,T}$ ma
ciągłe trajektorie, średnią
$m_s=x+\frac{s}{T}(y-x)$ i macierz kowariancji
\[
\operatorname{Cov}(\beta_s^{x,y,T},\beta_r^{x,y,T})
 =\left(\min(s,r)-\frac{sr}{T}\right)I_d.
\]
Dla $0<t_1<\cdots<t_n<T$ jego rozkład wektora
$(\beta_{t_1},\ldots,\beta_{t_n})$ ma gęstość
\[
\frac{p_{t_1}(x,x_1)\,p_{t_2-t_1}(x_1,x_2)
 \cdots p_{t_n-t_{n-1}}(x_{n-1},x_n)
 p_{T-t_n}(x_n,y)}{p_T(x,y)}.
\tag{15.1}\label{eq:bridge-density}
\]
Ponadto rodzina $y\mapsto\mu_{x,y,T}^{\mathrm{br}}$
jest wersją regularnego rozkładu warunkowego miary Wienera
$\mu_{x,T}$ względem $w(T)=y$: dla nieujemnego borelowskiego
funkcjonału $F$ i nieujemnej borelowskiej funkcji $g$,
\[
 \int F(w)g(w(T))\,\dd\mu_{x,T}(w)
 =\int_{\R^d}p_T(x,y)g(y)
    \left(\int F(w)\,\dd\mu_{x,y,T}^{\mathrm{br}}(w)\right)
    \dd\lambda_d(y).
\tag{15.2}\label{eq:bridge-disintegration}
\]
\end{twierdzenie}
\begin{proof}
Wektor $(B_{t_1},\ldots,B_{t_n},B_T)$ jest gaussowski.
Dla $s,r\le T$ obliczamy
$\operatorname{Cov}(B_s-\frac sT B_T,
B_r-\frac rT B_T)=\min(s,r)-sr/T$.
Co więcej, proces $B_s-sB_T/T$ jest niezależny od $B_T$:
ma zerową kowariancję z $B_T$, a cała rodzina jest wspólnie
gaussowska. Zatem warunkowanie procesu $x+B$ przy wartości
końcowej $y$ daje proces o tej samej kowariancji, lecz z
podstawioną średnią $m_s$. To jest dokładnie prawo
$\beta^{x,y,T}$.

Gęstość wspólnego rozkładu
$(x+B_{t_1},\ldots,x+B_{t_n},x+B_T)$ jest iloczynem
jąder cieplnych. Dzieląc przez dodatnią gęstość brzegową
$p_T(x,y)$, otrzymujemy \eqref{eq:bridge-density} dla
każdego $y$ (przy takim wyborze wersji jądra).
Tożsamość \eqref{eq:bridge-disintegration} wynika więc
najpierw dla cylindrycznych $F$ i prostych $g$ z twierdzenia
Fubiniego. Klasa tych funkcjonałów generuje całe borelowskie
$\sigma$-ciało przestrzeni trajektorii; twierdzenie o klasie
monotonicznej rozszerza tożsamość na wszystkie nieujemne
borelowskie $F$ i $g$. Mierzalność względem $y$ wynika
również z jawnej reprezentacji
$\beta^{x,y,T}=\beta^{x,0,T}+(s/T)y$.
\end{proof}

\begin{wniosek}[Gęstość początkowego fragmentu mostu]
\label{cor:bridge-initial-density}
Niech $0<s<T$, $y\in\R^d$ i niech $F$ będzie
ograniczonym borelowskim funkcjonałem zależnym wyłącznie
od trajektorii do chwili $s$. Wówczas
\[
 \E_{x,y,T}^{\mathrm{br}}F
 =\frac{1}{p_T(x,y)}\E_x\left[
    F(x+B)\,p_{T-s}(x+B_s,y)\right].
\tag{15.3}\label{eq:bridge-likelihood}
\]
Dla każdego $s<T$ prawo mostu na $[0,s]$ jest więc
równoważne prawu zwykłego ruchu Browna na $[0,s]$.
Na pełnym przedziale $[0,T]$ miary te są natomiast
wzajemnie singularne.
\end{wniosek}
\begin{proof}
Po ustaleniu przeszłości do $s$ niezależny przyrost
$B_T-B_s$ ma gęstość $p_{T-s}(x+B_s,\cdot)$.
Stosujemy tę własność do wspólnej miary
$\E_x[F\ind{\{x+B_T\in A\}}]$, a następnie porównujemy
gęstości względem $\lambda_d(y)$ i używamy ciągłej wersji
jądra z \eqref{eq:bridge-disintegration}. Współczynnik
$p_{T-s}/p_T$ jest ściśle dodatni i skończony,
co daje równoważność miar na $[0,s]$.
W chwili końcowej most spełnia $w(T)=y$ z
prawdopodobieństwem $1$, a zwykła miara Wienera przypisuje
temu zdarzeniu prawdopodobieństwo $0$. To dowodzi
singularności na $[0,T]$.
\end{proof}

\begin{przyklad}[Położenie w połowie drogi]
W jednowymiarowym moście z $x$ do $y$ w czasie $T$
zmienna $\beta_{T/2}$ ma rozkład
$\Normal((x+y)/2,T/4)$, a nie rozkład jednostajny.
Średnia jest liniową interpolacją końców, natomiast
wariancja $T/4$ jest konsekwencją wymagania zakończenia
trajektorii dokładnie w $y$.
\end{przyklad}



\begin{zadanie}[Prawdopodobieństwo przejścia mostu w przód]
Niech $0\le s<t<T$ i $\beta^{x,y,T}$ będzie
$d$-wymiarowym mostem Browna. Wyznaczyć gęstość warunkową
$\beta_t$ przy zadanej wartości $\beta_s=z$.

\emph{Rozwiązanie.} Z formuły
\eqref{eq:bridge-density}, po wycałkowaniu pozostałych
współrzędnych, otrzymujemy jądro
\[
 p^{\mathrm{br}}_{s,t}(z,z')=
 \frac{p_{t-s}(z,z')p_{T-t}(z',y)}{p_{T-s}(z,y)}.
\]
Mianownik jest dodatni, a całka prawej strony względem
$z'$ wynosi $1$ dzięki półgrupowej własności jądra
cieplnego. Jądro zależy od obu chwil $s,t$, zatem most
jest \emph{niejednorodnym w czasie} procesem Markowa.
Dodatkowy czynnik $p_{T-t}(z',y)$ nadaje większą wagę
położeniom, z których łatwiej dotrzeć do końca $y$.
To wyjaśnia kierunek dryfu w następnym twierdzeniu.
\end{zadanie}

\section{Dynamika mostu: dryf, który wymusza punkt końcowy}

\begin{twierdzenie}[Stochastyczne równanie mostu Browna]
\label{th:bridge-sde}
Niech $x,y\in\R$, $T>0$, a $W$ będzie standardowym
ruchem Browna. Dla $0\le s<T$ określmy
\[
 Y_s=x+\frac{s}{T}(y-x)+(T-s)\int_0^s\frac1{T-r}\,\dd W_r.
\tag{15.4}\label{eq:bridge-adapted}
\]
Proces $Y$ ma to samo prawo co most Browna z $x$ do $y$
i przyjmuje ciągłą modyfikację na $[0,T]$ spełniającą
$Y_T=y$. Na $[0,T)$ jest jedynym silnym rozwiązaniem
\[
 \dd Y_s=\frac{y-Y_s}{T-s}\,\dd s+\dd W_s,
 \qquad Y_0=x.
\tag{15.5}\label{eq:bridge-sde}
\]
\end{twierdzenie}
\begin{proof}
Na każdym $[0,T-\varepsilon]$ funkcja
$(T-s)^{-1}$ jest ograniczona i kwadratowo całkowalna,
więc całka Itô w \eqref{eq:bridge-adapted} jest
poprawnie określona i adaptowana. Z izometrii Wienera,
dla $s\le r<T$,
\[
\begin{split}
&\Cov\left((T-s)\int_0^s\frac{\dd W_u}{T-u},
 (T-r)\int_0^r\frac{\dd W_u}{T-u}\right)\\
&\qquad=(T-s)(T-r)\int_0^s\frac{\dd u}{(T-u)^2}
 =s-\frac{sr}{T}.
\end{split}
\]
Średnia $Y_s$ jest równa $x+s(y-x)/T$.
Wspólna gaussowskość daje zatem te same rozkłady
skończeniewymiarowe, co w
Twierdzeniu~\ref{th:bridge-conditional}.
Na każdym $[0,T-\varepsilon]$ proces $Y$ ma trajektorie
ciągłe. Uzasadnijmy osobno zachowanie w chwili $T$.
Wybierzmy zbiór $D=\mathbb Q\cap[0,T)$, gęsty w $[0,T)$.
Zgodność rozkładów skończeniewymiarowych $Y$ i mostu
$\beta^{x,y,T}$ oznacza zgodność ich praw na
przeliczalnej przestrzeni współrzędnych $\mathbb R^D$.
Dla mostu z ciągłości na $[0,T]$ wynika, że dla każdego
$k\in\N$ istnieje $\delta>0$ takie, że
$|\beta_q-y|<1/k$ dla wszystkich $q\in D$,
$T-\delta<q<T$. Zdarzenie to jest borelowskie
w przestrzeni $\mathbb R^D$, zatem ten sam warunek
spełnia $Y$ prawie na pewno. Ponieważ $Y$ jest ciągły
na $[0,T)$, oszacowanie rozciąga się z wymiernych $q$
na wszystkie $s$ dostatecznie bliskie $T$. Otrzymujemy
$\lim_{s\uparrow T}Y_s=y$ prawie na pewno. Ustalając
$Y_T=y$, dostajemy ciągłe przedłużenie i prawo mostu
na całym $[0,T]$. Nie wystarczyłaby tu sama zbieżność
$Y_s\to y$ według prawdopodobieństwa.

Z reguły iloczynu, ponieważ $s\mapsto T-s$ ma
skończoną wariację, otrzymujemy
\[
\dd\left((T-s)\int_0^s\frac{\dd W_r}{T-r}\right)
 =-\left(\int_0^s\frac{\dd W_r}{T-r}\right)\dd s+\dd W_s.
\]
Po dodaniu pochodnej deterministycznej średniej
upraszcza się to do \eqref{eq:bridge-sde}.
Jednoznaczność na $[0,T-\varepsilon]$ wynika z
liniowości równania lub z mnożenia różnicy dwóch
rozwiązań przez $(T-s)^{-1}$; ponieważ
$\varepsilon>0$ jest dowolne, otrzymujemy
jednoznaczność na $[0,T)$.
\end{proof}

\begin{uwaga}[Intuicja dryfu mostu]
Współczynnik $(y-Y_s)/(T-s)$ jest tym większy, im
mniej czasu pozostało do końca i im dalej proces
znajduje się od $y$. Nie jest on ograniczony przy
$s\uparrow T$, ale równanie jest dobrze określone
na każdym przedziale $[0,T-\varepsilon]$, a w granicy
trajektoria dochodzi do wyznaczonego punktu.
To bardzo konkretny przykład \emph{niejednorodnego
w czasie procesu Markowa}.
\end{uwaga}

\section{Funkcjonały liniowe i pierwsze obliczenia całek po trajektoriach}

\begin{twierdzenie}[Całka trajektorii mostu jest gaussowska]
\label{th:bridge-integrated}
Jeżeli $d=1$, $T>0$ i $x,y\in\R$, to
\[
 A_T:=\int_0^T\beta_s^{x,y,T}\,\dd s
 \sim\Normal\!\left(\frac{T(x+y)}2,\frac{T^3}{12}\right).
\]
W szczególności dla każdego $u\in\R$,
\[
 \int_{\calC_{x,y,T}}\exp\left(iu\int_0^Tw(s)\,\dd s\right)
     \,\dd\mu_{x,y,T}^{\mathrm{br}}(w)
 =\exp\left(\frac{iuT(x+y)}2-\frac{u^2T^3}{24}\right).
 \tag{15.6}\label{eq:bridge-area-char}
\]
\end{twierdzenie}
\begin{proof}
Sumy Riemanna są liniowymi kombinacjami wspólnie
gaussowskich zmiennych, a z ciągłości trajektorii zbiegają
prawie na pewno. Aby uzasadnić również zbieżność w $L^2$,
niech $\beta^{(n)}_s$ oznacza funkcję schodkową równą
$\beta_{t_k}$ na przedziale $[t_k,t_{k+1})$.
Z nierówności Cauchy'ego--Schwarza
\[
\E\left|\int_0^T(\beta^{(n)}_s-\beta_s)\,\dd s\right|^2
\le T\int_0^T\E|\beta^{(n)}_s-\beta_s|^2\,\dd s.
\]
Dla mostu
$\E|\beta_t-\beta_s|^2\le |t-s|+
\|y-x\|^2|t-s|^2/T^2$,
bo wariancja przyrostu wynosi
$|t-s|-(t-s)^2/T$; oba składniki znikają wraz
ze średnicą podziału. Prawa strona
zmierza więc do zera wraz ze średnicą podziału. Funkcje
charakterystyczne gaussowskich sum przechodzą do granicy,
co pokazuje, że sama całka ma rozkład gaussowski.
Jej średnia jest całką funkcji $m_s$ i wynosi
$T(x+y)/2$. Z twierdzenia Fubiniego
\[
 \Var(A_T)=\int_0^T\!\int_0^T
 \left(\min(s,r)-\frac{sr}{T}\right)\dd s\,\dd r
 =\frac{T^3}{3}-\frac{T^3}{4}=\frac{T^3}{12}.
\]
Wzór \eqref{eq:bridge-area-char} jest teraz funkcją
charakterystyczną rozkładu normalnego.
\end{proof}

\begin{zadanie}[Potencjał liniowy: jawna całka funkcyjna]
Dla $\lambda\in\R$ obliczyć
$\E_{x,y,T}^{\mathrm{br}}
\exp(-\lambda\int_0^T\beta_s\,\dd s)$.

\emph{Rozwiązanie.} Z twierdzenia powyżej i wzoru na
transformatę Laplace'a rozkładu normalnego otrzymujemy
\[
 \E_{x,y,T}^{\mathrm{br}}
  \exp\!\left(-\lambda\int_0^T\beta_s\,\dd s\right)
 =\exp\left(-\frac{\lambda T(x+y)}2
           +\frac{\lambda^2T^3}{24}\right).
\]
Funkcjonał jest całkowalny pomimo tego, że potencjał
$V(z)=\lambda z$ nie jest ograniczony. Wynik pochodzi
z rachunku gaussowskiego, a nie z ogólnego twierdzenia
o całkach z dowolnym nieograniczonym potencjałem.
\end{zadanie}


\section{Nieliniowy funkcjonał: całka z kwadratu ruchu Browna}

\begin{twierdzenie}[Transformata Laplace'a energii trajektorii]
\label{th:brownian-square-laplace}
Niech $B$ będzie standardowym jednowymiarowym
ruchem Browna z $B_0=0$. Dla $\alpha\ge0$ i $T>0$,
\[
 \E\exp\left(-\alpha\int_0^T B_s^2\,\dd s\right)
 =\frac1{\sqrt{\cosh(T\sqrt{2\alpha})}}.
\tag{15.7}\label{eq:brownian-square-laplace}
\]
\end{twierdzenie}
\begin{proof}
Dla $x\in\R$ rozważmy $X_s=x+B_s$ i funkcję
$u(t,x)=\E_x\exp(-\alpha\int_0^tX_s^2\,\dd s)$.
Zamiast zakładać z góry, że $u$ jest rozwiązaniem PDE,
znajdziemy jego kandydata i zweryfikujemy go przez
wzór Itô. Szukamy funkcji $v(t,x)=A(t)e^{-q(t)x^2}$
spełniającej
$\partial_tv=\tfrac12\partial_{xx}v-\alpha x^2v$,
$v(0,x)=1$. Po podstawieniu otrzymujemy
\[
 q'(t)=\alpha-2q(t)^2,\qquad
 A'(t)=-q(t)A(t),\qquad q(0)=0,\quad A(0)=1.
\]
Dla $\alpha>0$ rozwiązaniami są
\[
 q(t)=\sqrt{\alpha/2}\tanh(t\sqrt{2\alpha}),\qquad
 A(t)=\cosh(t\sqrt{2\alpha})^{-1/2}.
\]
Zatem $0<v(t,x)\le1$ oraz $v(0,x)=1$.
Dla ustalonego $t>0$ stosujemy wzór Itô do procesu
\[
 M_s=\exp\left(-\alpha\int_0^sX_r^2\,\dd r\right)
 v(t-s,X_s),\qquad 0\le s\le t.
\]
Drift znika na mocy równania PDE. Po zatrzymaniu w
$\tau_R=\inf\{s:|X_s|\ge R\}$ część stochastyczna
jest martyngałem. Ponieważ $0\le M_s\le1$, można
usunąć zatrzymanie przez zbieżność zdominowaną,
a w chwili $t$ mamy $v(0,X_t)=1$. Dostajemy
$v(t,x)=\E_xe^{-\alpha\int_0^tX_s^2\,\dd s}$.
Podstawiając $x=0$, otrzymujemy tezę. Dla
$\alpha=0$ obie strony są równe $1$.
\end{proof}

\begin{uwaga}[Dwa rodzaje funkcjonałów]
Całka $\int_0^T B_s\,\dd s$ jest gaussowska,
ponieważ jest liniowym funkcjonałem procesu
gaussowskiego. Całka $\int_0^T B_s^2\,\dd s$
nie jest gaussowska: jest nieujemna i ma transformację
Laplace'a \eqref{eq:brownian-square-laplace}.
Obliczenie drugiej wymaga więcej niż znajomości
funkcji kowariancji i prowadzi naturalnie do
równania parabolicznego.
\end{uwaga}

\section{Feynmana--Kaca półgrupa ważona i jej jądro}

\begin{definicja}[Półgrupa z potencjałem]
\label{def:path-weighted-semigroup}
Niech $V:\R^d\to\R$ będzie ograniczoną funkcją borelowską.
Dla $t\ge0$ i ograniczonej borelowskiej funkcji
$f:\R^d\to\R$ określamy
\[
 (P_t^Vf)(x)=\E_x\left[
 \exp\left(-\int_0^tV(X_s)\,\dd s\right)f(X_t)\right],
 \quad X_s=x+B_s.
\tag{15.8}\label{eq:path-FK-semigroup}
\]
Całka względem $s$ jest zwykłą całką Lebesgue'a;
całość jest całką po prawie trajektorii $X$.
\end{definicja}

\begin{twierdzenie}[Własność półgrupowa i jądro mostowe]
\label{th:path-FK-kernel}
Przy założeniach Definicji~\ref{def:path-weighted-semigroup}:
\begin{enumerate}
\item $P_{s+t}^V=P_s^VP_t^V$ dla $s,t\ge0$,
$P_0^V=I$ oraz
$\|P_t^V f\|_\infty\le e^{t\|V\|_\infty}\|f\|_\infty$;
\item $P_t^V$ jest operatorem dodatnim;
\item dla $t>0$ funkcja
\[
 k_t^V(x,y)=p_t(x,y)\,
 \E_{x,y,t}^{\mathrm{br}}\left[
   \exp\left(-\int_0^tV(\beta_s)\,\dd s\right)\right]
 \tag{15.9}\label{eq:path-FK-kernel}
\]
jest mierzalnym jądrem całkowym, tj.
$(P_t^Vf)(x)=\int_{\R^d}k_t^V(x,y)f(y)\,\dd\lambda_d(y)$;
\item jądra spełniają tożsamość Chapmana--Kołmogorowa:
\[
 k_{s+t}^V(x,y)=\int_{\R^d}k_s^V(x,z)k_t^V(z,y)
 \,\dd\lambda_d(z)
\]
dla $s,t>0$ i prawie wszystkich $y$
(przy ustalonych $x,s,t$).
\end{enumerate}
\end{twierdzenie}
\begin{proof}
Niech $X=x+B$. Dla ustalonych $s,t\ge0$ rozdzielamy
całkę potencjału na $[0,s]$ i $[s,s+t]$. Dzięki własności
Markowa, dla funkcjonału dalszej trajektorii otrzymujemy
\[
 \E_x\left[e^{-\int_s^{s+t}V(X_r)\,\dd r}
    f(X_{s+t})\mid\calF_s\right]=(P_t^V f)(X_s).
\]
Tożsamość najpierw zachodzi dla prostych funkcjonałów
zależnych od skończenie wielu chwil, a przez aproksymację
monotoniczną dla funkcjonałów mierzalnych całej przyszłej
trajektorii. Funkcjonał $\int_s^{s+t}V(X_r)\,\dd r$
jest mierzalny: można go otrzymać jako granicę całek
funkcji prostych względem czasu, ponieważ
$(r,\omega)\mapsto V(X_r(\omega))$ jest produktowo
mierzalne. Własność wieży daje $P_{s+t}^V=P_s^VP_t^V$.
Dodatniość i oszacowanie normy wynikają z granic
$e^{-t\|V\|_\infty}\le e^{-\int_0^tV}\le e^{t\|V\|_\infty}$.

Aby otrzymać jądro, stosujemy rozkład warunkowy
\eqref{eq:bridge-disintegration} do funkcjonału
$F(w)=\exp(-\int_0^tV(w(r))\,\dd r)$
i funkcji $g(w(t))=f(w(t))$ (najpierw dla funkcji
nieujemnych, potem przez rozkład na części dodatnią
i ujemną). Dostajemy \eqref{eq:path-FK-kernel}.
Złożenie dwóch operatorów całkowych, twierdzenie
Fubiniego i własność półgrupowa dają ostatnią tezę
jako równość jąder dla prawie wszystkich $y$.
\end{proof}

\begin{uwaga}[Brak konieczności istnienia gęstości trajektorii]
Wyrażenie \eqref{eq:path-FK-kernel} nie sugeruje, że
istnieje gęstość miary Wienera względem jakiejś
„miary Lebesgue'a” na $C([0,t])$. Czynnik $p_t(x,y)$
pochodzi wyłącznie z gęstości \emph{końcowego położenia}.
Pozostały czynnik jest zwykłą wartością oczekiwaną po
moście, który ma dobrze określoną miarę probabilistyczną.
\end{uwaga}

\section{Dyskretyzacja czasu: dlaczego powstaje całka po trajektoriach?}

\begin{twierdzenie}[Skończeniewymiarowe aproksymacje całki po trajektoriach]
\label{th:path-time-slice}
Niech $T>0$, $x\in\R^d$, $V:\R^d\to\R$ będzie
\emph{ograniczona i ciągła}, a $f:\R^d\to\R$
ograniczona i borelowska. Niech $\Delta=T/n$,
$t_k=k\Delta$. Wtedy
\[
\begin{split}
 (P_T^Vf)(x)&=\lim_{n\to\infty}
 \int_{(\R^d)^n}f(x_n)
 \exp\left(-\Delta\sum_{k=0}^{n-1}V(x_k)\right)\\
 &\hspace{2cm}\cdot\prod_{k=1}^n p_\Delta(x_{k-1},x_k)
 \,\dd\lambda_{dn}(x_1,\ldots,x_n),\qquad x_0=x.
\end{split}
\tag{15.10}\label{eq:path-time-slice}
\]
Dla ustalonych $x,y$ jądro \eqref{eq:path-FK-kernel}
jest granicą analogicznych całek po
$(\R^d)^{n-1}$ z $x_n=y$ (\emph{bez} końcowego
czynnika $f(x_n)$). W tych całkach iloczyn jąder
cieplnych pozostaje bez normalizacji; dopiero po
podzieleniu go przez $p_T(x,y)$ otrzymuje się gęstość
rozkładu punktów pośrednich mostu.
\end{twierdzenie}
\begin{proof}
Z Twierdzenia~\ref{th:path-cylinder} całka po
$(\R^d)^n$ jest równa
$\E_x[e^{-\Delta\sum_{k=0}^{n-1}V(X_{k\Delta})}f(X_T)]$.
Do zbieżności wykładnika nie potrzeba ciągłości $f$:
wartość $f(X_T)$ pozostaje taka sama dla wszystkich $n$.
Dla każdej ciągłej trajektorii $w$ jej obraz $w([0,T])$
jest zwarty, a $V$ jest na tym obrazie jednostajnie ciągła.
Zatem lewostronne sumy Riemanna zbiegają do
$\int_0^TV(w(s))\,\dd s$. Wykładniki są ograniczone
przez $e^{T\|V\|_\infty}$, więc twierdzenie o zbieżności
zdominowanej prowadzi do \eqref{eq:path-time-slice}.

Dla mostu \eqref{eq:bridge-density} daje rozkład
kolejnych położeń przed chwilą $T$; ten sam argument
Riemanna i zbieżności zdominowanej obowiązuje przy
każdym ustalonym $y$. Mnożąc przez $p_T(x,y)$,
otrzymujemy jądro \eqref{eq:path-FK-kernel}.
\end{proof}

\begin{uwaga}[Dlaczego wymagana jest ciągłość $V$?]
Dla ograniczonego potencjału borelowskiego
$\int_0^TV(X_s)\,\dd s$ nadal jest określona,
a twierdzenie o jądrze pozostaje prawdziwe.
Nie wolno jednak automatycznie twierdzić, że \emph{sumy
Riemanna przy ustalonej siatce} zbiegają dla każdej
trajektorii: funkcja $s\mapsto V(w(s))$ może być
nieciągła i niecałkowalna w sensie Riemanna.
Takie uogólnienie wymaga innej aproksymacji.
\end{uwaga}

\begin{center}
\begin{tikzpicture}[x=0.9cm,y=0.8cm,>=Latex]
\draw[->] (0,0)--(10.2,0) node[right] {$s$};
\draw[->] (0,-1.2)--(0,2.7) node[above] {$w(s)$};
\draw[thick,blue!65!black] plot[smooth] coordinates
 {(0,0.1) (0.6,0.5) (1.0,-0.05) (1.5,1.05) (2.0,0.55)
 (2.4,0.95) (3.1,0.3) (3.8,1.45) (4.2,0.85)
 (5.0,1.65) (5.8,0.95) (6.5,1.85) (7.0,0.8)
 (7.8,1.4) (8.3,0.75) (9.0,1.7) (9.8,1.15)};
\foreach \x in {0,2,4,6,8,10} {\draw[gray!70] (\x,-0.14)--(\x,0.14);}
\draw[dashed,gray!65] (2,-1.08)--(2,2.25);
\draw[dashed,gray!65] (4,-1.08)--(4,2.25);
\draw[dashed,gray!65] (6,-1.08)--(6,2.25);
\draw[dashed,gray!65] (8,-1.08)--(8,2.25);
\node[font=\small,align=center] at (5,-1.72)
 {Jedna trajektoria ciągła i kolejne chwile dyskretyzacji;\\
  całka po trajektoriach uśrednia funkcjonały po wszystkich takich drogach.};
\end{tikzpicture}
\end{center}

\section{Dwa zadania łączące probabilistykę z potencjałami}

\begin{zadanie}[Potencjał stały i ubytek masy]
Niech $V(x)\equiv v\ge0$. Wyznaczyć $P_t^V1$,
jądro i wyjaśnić sens probabilistyczny.

\emph{Rozwiązanie.} Ponieważ całka $\int_0^t V(X_s)\,\dd s=vt$
nie zależy od trajektorii, mamy
$(P_t^Vf)(x)=e^{-vt}\E_xf(X_t)$ i
$k_t^V(x,y)=e^{-vt}p_t(x,y)$. W szczególności
$(P_t^V1)(x)=e^{-vt}$. Ubytek masy interpretuje się
jako przeżycie cząstki zabijanej z intensywnością $v$;
nie jest to przejściowe jądro procesu konserwatywnego,
bo całka jądra po $y$ jest mniejsza niż $1$ dla $t>0$.
\end{zadanie}

\begin{zadanie}[Liczba odwiedzin a czas przebywania]
Niech $A\subset\R^d$ będzie borelowski,
$\lambda>0$ i $V=\lambda\ind{A}$. Wykazać, że
\[
 (P_T^V1)(x)=\E_x e^{-\lambda\,\ell_T(A)},
 \qquad \ell_T(A)=\int_0^T\ind{\{X_s\in A\}}\,\dd s.
\]

\emph{Rozwiązanie.} Podstawiamy $V$ w
\eqref{eq:path-FK-semigroup}. Zmienna $\ell_T(A)$
to \emph{czas okupacji} zbioru, a nie liczba jego odwiedzin.
Dla Brownowskiej trajektorii pojedyncze przecięcie granicy
zbioru może nie mieć sensownej liczności, lecz czas
okupacji pozostaje poprawnie określoną zmienną z przedziału
$[0,T]$. Z monotoniczności funkcji wykładniczej wynika
$0\le (P_T^V1)(x)\le1$, co odpowiada prawdopodobieństwu
przeżycia przy intensywności zabijania $\lambda$ wyłącznie
podczas pobytu w $A$.
\end{zadanie}

\section{Przejście do mechaniki kwantowej}

Wszystkie obiekty tego rozdziału są ściśle określonymi
całkami Lebesgue'a względem \emph{dodatnich} miar
probabilistycznych. W szczególności
$e^{-\int_0^T V(w(s))\,\dd s}$ jest dodatnią wagą dla
rzeczywistego $V$. W fizyce kwantowej pojawia się natomiast
czynnik $e^{iS[w]/\hbar}$, gdzie $S$ jest działaniem
klasycznym. W następnym rozdziale pokażemy, skąd bierze się
formalna analogia oraz dlaczego nie wolno zastępować
miary Wienera nieokreśloną „miarą fazową” bez oddzielnego
uzasadnienia matematycznego.


% ===========================================================================
% ROZDZIAŁ 16: całki Feynmana, Wick, propagator, oscylator
% ===========================================================================
\chapter{Całki Feynmana, obrót Wicka i mechanika kwantowa}
\label{ch:feynman-path}

\section{Dwa ewolucyjne równania i dwie rodziny operatorów}

W mechanice kwantowej funkcja falowa $\psi(t,\cdot)$ zmienia się
zgodnie z równaniem Schrödingera. Dla cząstki o masie $m>0$
i stałej Plancka $\hbar>0$ rozważmy w $L^2(\R^d)$
operator
\[
 H_0=-\frac{\hbar^2}{2m}\Delta,
 \qquad H=H_0+V,
\tag{16.1}\label{eq:QM-H}
\]
gdzie $V:\R^d\to\R$ jest rzeczywistą, ograniczoną funkcją
borelowską i $V$ oznacza operator mnożenia przez tę funkcję.
Dziedziną $H_0$ jest $H^2(\R^d)$, a $H$ jest
samiosprzężony na tej samej dziedzinie: dodanie ograniczonego
operatora samosprzężonego zachowuje samosprzężoność.
Wtedy $H\ge-\|V\|_\infty I$. Z twierdzenia spektralnego
otrzymujemy dwie rodziny operatorów:
\[
 U(t)=e^{-itH/\hbar},\quad t\in\R,\qquad
 S(s)=e^{-sH/\hbar},\quad s\ge0.
\tag{16.2}\label{eq:unitary-euclidean}
\]
Pierwsza jest \emph{grupą unitarną}, druga --- \emph{półgrupą
cieplną} (niekoniecznie kontrakcją, jeśli $V$ bywa ujemny).

\begin{stwierdzenie}[W jakim sensie operatory rozwiązują równania?]
\label{prop:unitary-generator}
Niech $f\in H^2(\R^d)$. Wówczas $u(t)=U(t)f$ jest
różniczkowalne jako funkcja o wartościach w $L^2(\R^d)$
i spełnia $i\hbar u'(t)=Hu(t)$, $u(0)=f$.
Podobnie $v(s)=S(s)f$ spełnia $\hbar v'(s)=-Hv(s)$.
Dla dowolnego $f\in L^2$ operatorowe wzory
\eqref{eq:unitary-euclidean} nadal mają sens, ale
pochodne w $L^2$ nie muszą istnieć.
\end{stwierdzenie}
\begin{proof}
Dla miary spektralnej $E_H$ mamy
$U(t)f=\int_{\R}e^{-it\lambda/\hbar}\,\dd E_H(\lambda)f$.
Dla $f\in\mathcal D(H)$ ilorazy różnicowe zbieżne są
w normie $L^2$ na mocy nierówności
$|(e^{-ih\lambda/\hbar}-1)/h|\le|\lambda|/\hbar$
i twierdzenia o zbieżności zdominowanej względem
miary spektralnej $\|E_H(\dd\lambda)f\|^2$.
Dostajemy pochodną $-(i/\hbar)HU(t)f$.
Analogiczny rachunek dla $S$ jest dozwolony, ponieważ
widmo $H$ jest ograniczone od dołu. Ciągłość silna
obu rodzin pochodzi z tego samego argumentu.
\end{proof}

\section{Euklidesowy propagator jest całką względem miary Wienera}

\begin{twierdzenie}[Feynman--Kac z jawnymi stałymi fizycznymi]
\label{th:physical-FK}
Niech $V\in L^\infty(\R^d;\R)$ i $H$ będzie operatorem
z \eqref{eq:QM-H}. Dla $s>0$, $f\in L^2(\R^d)$,
po wybraniu odpowiedniego reprezentanta prawej strony,
\[
 (S(s)f)(x)=\E\left[
 f\left(x+\sqrt{\frac{\hbar}{m}}B_s\right)
 \exp\left(-\frac1\hbar\int_0^s
 V\left(x+\sqrt{\frac{\hbar}{m}}B_r\right)\,\dd r\right)
 \right]
\tag{16.3}\label{eq:physical-FK}
\]
dla prawie wszystkich $x\in\R^d$, gdzie $B$ jest
$d$-wymiarowym standardowym ruchem Browna.
Dla ograniczonej funkcji borelowskiej $f$ prawa strona
jest określona w każdym $x$.
\end{twierdzenie}
\begin{proof}
Po podzieleniu równania $\hbar\partial_s v=-Hv$ przez
$\hbar$ otrzymujemy operator generatora
$\frac{\hbar}{2m}\Delta-\frac1\hbar V$.
Proces $X_s=x+\sqrt{\hbar/m}B_s$ ma generator
$\frac{\hbar}{2m}\Delta$.
Wzór Feynmana--Kaca z rozdziału~\ref{ch:path-integrals},
z potencjałem $V/\hbar$, daje więc operator $T(s)$ po
prawej stronie \eqref{eq:physical-FK}. Wykażemy, że jest
on dokładnie $S(s)$, nie zakładając istnienia klasycznego
rozwiązania PDE.

Niech $A=(\hbar/(2m))\Delta$ w $L^2$, $B=-V/\hbar$
(będzie to ograniczony operator mnożenia). Ważona
półgrupa $T(s)$ spełnia równanie całkowe Duhamela
\[
 T(s)f=e^{sA}f+\int_0^s e^{rA}B\,T(s-r)f\,\dd r.
\tag{16.4}\label{eq:FK-duhamel}
\]
Wyprowadźmy tę tożsamość bez odwoływania się do
formalnego różniczkowania półgrupy. Dla każdej trajektorii
połóżmy $g(r)=-V(X_r)/\hbar$; symbol $B=-V/\hbar$
użyty w równaniu Duhamela oznacza natomiast operator mnożenia.
Ponieważ $g$ jest ograniczona i mierzalna, z fundamentalnego
twierdzenia rachunku całkowego dla funkcji absolutnie
ciągłych otrzymujemy
\[
 e^{\int_0^s g(r)\,\dd r}
 =1+\int_0^s g(r)e^{\int_r^s g(u)\,\dd u}\,\dd r.
\]
Mnożymy przez $f(X_s)$ i bierzemy wartość oczekiwaną.
Po warunkowaniu względem $\calF_r$ czynnik pochodzący
z odcinka $[r,s]$ ma warunkową średnią
$(T(s-r)f)(X_r)$; jest to zastosowanie własności Markowa
do całej przyszłej trajektorii. Otrzymujemy
\eqref{eq:FK-duhamel} najpierw dla ograniczonych
funkcji borelowskich $f$ należących do $L^2$.
Tonelli lub Fubini wolno stosować, ponieważ
$|g|\le\|V\|_\infty/\hbar$ i waga jest ograniczona przez
$e^{s\|V\|_\infty/\hbar}$.
Wreszcie zachodzi punktowe oszacowanie
$|T(s)f|\le e^{s\|V\|_\infty/\hbar}e^{sA}|f|$.
Półgrupa cieplna jest kontrakcją w $L^2$, a więc
$\|T(s)f\|_2\le e^{s\|V\|_\infty/\hbar}\|f\|_2$.
Daje to ciągłe rozszerzenie operatorów $T(s)$ i
równania Duhamela na całe $L^2$.

Na drugiej stronie wzoru \eqref{eq:unitary-euclidean}
operator $S(s)=e^{s(A+B)}$ spełnia dokładnie to samo
równanie Duhamela dla ograniczonego zaburzenia $B$.
Jeżeli $D(s)=T(s)-S(s)$, to
$\|D(s)f\|_2\le\|B\|\int_0^s\|D(r)f\|_2\,\dd r$.
Lemat Gronwalla daje $D(s)f=0$. Wzór obowiązuje
najpierw na gęstym zbiorze ograniczonych funkcji
całkowalnych, a potem na całym $L^2$.
\end{proof}

\begin{uwaga}[Normalizacja fizyczna]
W rozdziale~\ref{ch:path-integrals} rozważaliśmy
$H=-\frac12\Delta+V$ i czas euklidesowy $t$.
Teraz pojawiają się $\hbar$ i $m$: wariancja każdej
współrzędnej przyrostu procesu $X$ wynosi
$(\hbar/m)s$, a potencjał występuje w wykładniku
z czynnikiem $1/\hbar$. Pominięcie tych stałych
prowadziłoby do błędnej skali propagatora.
\end{uwaga}

\section{Obrót Wicka w sensie twierdzenia spektralnego}

Formalna zamiana $s=it$ bywa nazywana \emph{obrotem Wicka}.
Właściwy matematyczny argument nie polega na
podstawieniu liczby zespolonej do losowej trajektorii.
Jest nim przedłużenie operatorów
$e^{-sH/\hbar}$ na prawą półpłaszczyznę
zespolonego parametru.

\begin{twierdzenie}[Przejście od półgrupy cieplnej do propagatora]
\label{th:wick-spectral}
Niech $H$ będzie samosprzężonym operatorem w zespolonej
przestrzeni Hilberta $\calH$, a jego widmo niech będzie
ograniczone od dołu. Dla $\operatorname{Re}z>0$ określmy
$S(z)=e^{-zH/\hbar}$ przez rachunek spektralny.
Wówczas $z\mapsto S(z)$ jest holomorficzne w prawej
półpłaszczyźnie w normie operatorowej lokalnie, a dla
każdego $t\in\R$ i $f\in\calH$ zachodzi
\[
 \lim_{\varepsilon\downarrow0}
  S(\varepsilon+it)f=e^{-itH/\hbar}f
  \quad\text{w normie }\calH.
\tag{16.5}\label{eq:wick-strong}
\]
Nie wynika stąd samo przez się zbieżność punktowa jąder
ani istnienie zespolonej miary na przestrzeni trajektorii.
\end{twierdzenie}
\begin{proof}
Niech $H\ge-MI$. Dla $\varepsilon\in(0,1]$ i
$\lambda\ge-M$ mamy $|e^{-\varepsilon\lambda/\hbar}|
\le e^{M/\hbar}$. Z rachunku spektralnego
\[
\begin{split}
&\|S(\varepsilon+it)f-e^{-itH/\hbar}f\|_{\calH}^2\\
&\quad=\int_{\R}\left|e^{-\varepsilon\lambda/\hbar}-1\right|^2
 \,\dd\langle E_H(\lambda)f,f\rangle.
\end{split}
\]
Całkowana funkcja dąży punktowo do zera i jest
jednostajnie ograniczona przez $(1+e^{M/\hbar})^2$.
Zbieżność zdominowana daje \eqref{eq:wick-strong}.

Holomorficzność wynika z tego, że na każdym zwartym
podzbiorze prawej półpłaszczyzny
$\lambda^k e^{-z\lambda/\hbar}$ jest jednostajnie
ograniczona dla $\lambda\ge-M$, dla każdego $k\ge0$.
Można więc różniczkować rachunek spektralny w normie
operatorowej. Tożsamość na brzegu jest jednak
\emph{zbieżnością silną operatorów}, nie całką po
trajektoriach z nową miarą probabilistyczną.
\end{proof}

\section{Cząstka swobodna: obliczenie propagatora}

\begin{twierdzenie}[Jądro cząstki swobodnej]
\label{th:free-schrodinger-kernel}
Niech $m,\hbar>0$, $H_0=-\hbar^2\Delta/(2m)$ w
$L^2(\R^d)$. Dla $t>0$ operator
$U_0(t)=e^{-itH_0/\hbar}$ ma, na funkcjach Schwartza,
oscylacyjną reprezentację
\[
 (U_0(t)f)(x)=\int_{\R^d}^{\mathrm{osc}}
 K_0(t;x,y)f(y)\,\dd\lambda_d(y),
\quad
 K_0(t;x,y)=\left(\frac{m}{2\pi i\hbar t}\right)^{d/2}
 \exp\left(\frac{im\|x-y\|^2}{2\hbar t}\right).
\tag{16.6}\label{eq:free-propagator}
\]
Potęga pierwiastka jest określona przez granicę z
$z=\varepsilon+it$, $\varepsilon>0$. Wyrażenie jest
granicą regularizowanych całek gaussowskich; nie jest
jądrem przejścia miary probabilistycznej.
\end{twierdzenie}
\begin{proof}
Stosujemy transformatę Fouriera z konwencją
$\widehat f(\xi)=\int e^{-i\langle x,\xi\rangle}
 f(x)\,\dd\lambda_d(x)$. Wtedy
\[
 \widehat{U_0(t)f}(\xi)
 =e^{-it\hbar\|\xi\|^2/(2m)}\widehat f(\xi).
\]
Dla parametru $z$ z $\operatorname{Re}z>0$
transformata odwrotna mnożnika
$e^{-z\hbar\|\xi\|^2/(2m)}$ jest klasyczną gęstością
Gaussa
\[
 q_z(x,y)=\left(\frac{m}{2\pi\hbar z}\right)^{d/2}
 \exp\left(-\frac{m\|x-y\|^2}{2\hbar z}\right).
\]
Stosujemy $q_{\varepsilon+it}$ do $f$ i dopuszczamy
$\varepsilon\downarrow0$. Twierdzenie~\ref{th:wick-spectral}
gwarantuje zbieżność w $L^2$. Dla funkcji Schwartza
można również przejść do granicy w odwróconej
transformacie Fouriera w sensie rozkładów temperowanych.
Zastępując $z$ przez $it$ w jądrach, otrzymujemy
\eqref{eq:free-propagator} ze wskazanym wyborem gałęzi.
Ponieważ $f\in\calS(\R^d)\subset L^1(\R^d)$,
sam zapis $\int K_0(t;x,y)f(y)\,\dd\lambda_d(y)$
może być odczytany jako zwykła całka absolutnie zbieżna
względem $y$. Dopisek $\mathrm{osc}$ odnosi się do
charakteru jądra i do rozszerzenia reprezentacji na
szersze klasy funkcji, dla których całka bezwzględna
nie musi istnieć.
\end{proof}

\begin{uwaga}[Dlaczego zwykła miara zespolona zawodzi już w przykładzie swobodnym?]
Dla $t>0$ moduł jądra \eqref{eq:free-propagator}
nie zależy od $y$:
$|K_0(t;x,y)|=(m/(2\pi\hbar t))^{d/2}$.
Dlatego $\int_{\R^d}|K_0(t;x,y)|\,\dd\lambda_d(y)=\infty$.
Nie można więc potraktować tego jądra jako gęstości
\emph{skończonej zespolonej miary przejścia}.
Tym bardziej nie wynika z niego istnienie przeliczalnie
addytywnej miary fazowej o skończonej wariacji na
całej przestrzeni trajektorii.
\end{uwaga}

\section{Dyskretyzacja czasu: od działania klasycznego do iloczynów operatorów}

Dla gładkiej trajektorii $x:[0,T]\to\R^d$ klasyczne
działanie cząstki w potencjale $V$ jest równe
$S[x]=\int_0^T(\frac m2\|\dot x(s)\|^2-V(x(s)))\,\dd s$.
Typowa trajektoria Brownowska nie ma jednak pochodnej:
ten wzór \emph{nie jest} funkcjonałem określonym prawie
wszędzie względem miary Wienera. Przybliżamy działanie
dla łamanych o wierzchołkach $x_0,\ldots,x_n$.

\begin{twierdzenie}[Iloczyny Trottera dla ograniczonego potencjału]
\label{th:QM-trotter}
Niech $H_0$ będzie jak w Twierdzeniu~\ref{th:free-schrodinger-kernel},
a $V\in L^\infty(\R^d;\R)$. Dla każdego $t\in\R$
i $f\in L^2(\R^d)$ zachodzi
\[
 e^{-it(H_0+V)/\hbar}f
 =\lim_{n\to\infty}
 \left(e^{-itV/(n\hbar)}e^{-itH_0/(n\hbar)}\right)^n f
 \quad\text{w }L^2(\R^d).
\tag{16.7}\label{eq:trotter-quantum}
\]
Jeżeli dodatkowo $V$ jest gładkie i wszystkie jego
pochodne są ograniczone, iloczyny można na gładkich
szybko malejących funkcjach zapisywać przez iterowane
całki oscylacyjne o dyskretnym działaniu
\[
 S_n(x_0,\ldots,x_n)
 =\sum_{k=0}^{n-1}\left(
    \frac{m\|x_{k+1}-x_k\|^2}{2\Delta}
    -\Delta V(x_k)\right),\quad\Delta=t/n,
\tag{16.8}\label{eq:discrete-action}
\]
przy czym $t>0$, $x_0=x$ i $x_n=y$ są końcami
trajektorii. Zbieżność w \eqref{eq:trotter-quantum}
nie oznacza istnienia całki Lebesgue'a z
$e^{iS[x]/\hbar}$ po nieskończonej przestrzeni trajektorii.
\end{twierdzenie}
\begin{proof}
Podajemy argument ograniczonego zaburzenia.
Niech $A=-iH_0/\hbar$ i $B=-iV/\hbar$; wtedy $A$
jest generatorem grupy unitarnej, a $B$ jest
ograniczonym operatorem antysamiosprzężonym.
Niech $U(t)=e^{t(A+B)}$ i $U_A(t)=e^{tA}$.
Wzór Duhamela dla ograniczonego zaburzenia daje
\[
 U(h)f=U_A(h)f+\int_0^h U_A(h-s)BU(s)f\,\dd s.
\]
Dla $h\downarrow0$ z ciągłości silnej $U,U_A$
wynika, że całka po prawej jest równa
$hBf+o_f(h)$. Ponadto
$e^{hB}U_A(h)f=U_A(h)f+hBf+o_f(h)$,
gdyż $B$ jest ograniczony. Mamy zatem
$h^{-1}(e^{hB}U_A(h)-U(h))f\to0$ dla każdego $f$.
Normy operatorów
$h^{-1}(e^{hB}U_A(h)-U(h))$ są jednakowo
ograniczone dla małych dodatnich $h$ na mocy
wzoru Duhamela i nierówności
$\|e^{hB}-I\|\le h\|B\|e^{h\|B\|}$.
Zbieżność silna jest wobec tego jednostajna na
zwartych zbiorach wektorów. Dla ustalonego $t>0$
zbiór $\{U(s)f:0\le s\le t\}$ jest zwarty w $L^2$.

Korzystamy z tożsamości teleskopowej, przy
$h=t/n$ i $Q_h=e^{hB}U_A(h)$:
\[
 Q_h^n f-U(t)f
 =\sum_{k=0}^{n-1} Q_h^{n-1-k}
   (Q_h-U(h))\,U(kh)f.
\]
Każdy czynnik $Q_h$ i $U(h)$ jest unitarny, dlatego
norma prawej strony jest najwyżej
$n\sup_{0\le s\le t}\|(Q_h-U(h))U(s)f\|=n\,o_f(h)=o(1)$.
Dla $t<0$ postępujemy tak samo po zamianie znaku
parametru albo stosujemy sprzężenia. Dowodzi to
\eqref{eq:trotter-quantum}.

Dla gładkiego ograniczonego $V$ każdy czynnik
$e^{-itV/(n\hbar)}$ stojący po lewej stronie czynnika
swobodnego mnoży przez fazę
$e^{-i\Delta V(x)/\hbar}$, natomiast
$e^{-itH_0/(n\hbar)}$ ma jądro z
\eqref{eq:free-propagator}. Ich iterowane składanie daje
sumę faz kinetycznych i potencjalnych
$S_n/\hbar$. Czynnik normalizacyjny
$(m/(2\pi i\hbar\Delta))^{nd/2}$ to produkt
prefaktorów jąder swobodnych. Wszystkie całki
odczytujemy jako całki oscylacyjne lub granice
regularizowane; wcześniejszy argument zapewnia
zbieżność \emph{operatorów}, a nie miar fazowych.
\end{proof}

\begin{przyklad}[Potencjał stały]
Jeżeli $V(x)\equiv v_0\in\R$, to
$[H_0,V]=0$ i $U(t)=e^{-itv_0/\hbar}U_0(t)$.
Propagator zmienia się tylko o wspólną fazę:
$K_V(t;x,y)=e^{-itv_0/\hbar}K_0(t;x,y)$.
Jest to dokładna wersja formalnego wzoru
$e^{iS[x]/\hbar}$, w którym do działania
kinetycznego dodano wyraz $-v_0T$.
\end{przyklad}

\section{Oscylator harmoniczny: przykład, który można policzyć dokładnie}

Potencjał $V(x)=\frac12 m\omega^2x^2$, $m,\omega>0$,
nie jest ograniczony. Nie podpada więc pod hipotezy
Twierdzenia~\ref{th:QM-trotter}, ale operator
$H_{\mathrm{osc}}=-\frac{\hbar^2}{2m}\frac{\dd^2}{\dd x^2}
+\frac12m\omega^2x^2$ jest samosprzężony i
ograniczony od dołu. Możemy nadal rozważyć jego
półgrupę cieplną i propagator.

\begin{twierdzenie}[Jądro Mehlera w czasie euklidesowym]
\label{th:Mehler-euclidean}
Dla $s>0$, $x,y\in\R$ jądro operatora
$e^{-sH_{\mathrm{osc}}/\hbar}$ jest równe
\[
\begin{split}
 K_s^{\mathrm E}(x,y)
  ={}&\sqrt{\frac{m\omega}{2\pi\hbar\sinh(\omega s)}}\\
 &\cdot\exp\left[-\frac{m\omega}{2\hbar\sinh(\omega s)}
 \bigl((x^2+y^2)\cosh(\omega s)-2xy\bigr)\right].
\end{split}
\tag{16.9}\label{eq:Mehler-E}
\]
Dla ograniczonej ciągłej $f$ prawa strona określa
naturalne rozszerzenie półgrupy, nawet gdy $f\notin L^2$.
W tym sensie jądro można zapisać jako
\[
\begin{split}
 (e^{-sH_{\mathrm{osc}}/\hbar}f)(x)
 =\E_x\Big[&f(X_s)\exp\Big(
  -\frac{m\omega^2}{2\hbar}\int_0^sX_r^2\,\dd r
 \Big)\Big],\\
 &X_r=x+\sqrt{\hbar/m}B_r.
\end{split}
\tag{16.10}\label{eq:Mehler-FK}
\]
\end{twierdzenie}
\begin{proof}
Najpierw sprawdzamy wzór \eqref{eq:Mehler-E}
bez odwoływania się do całek formalnych.
Szukamy dodatniego jądra w postaci
$K_s(x,y)=A_s\exp[-a_s(x^2+y^2)+b_sxy]$.
Podstawiając do równania
\[
 \partial_sK=\frac{\hbar}{2m}\partial_{xx}K
    -\frac{m\omega^2}{2\hbar}x^2K
\]
i porównując współczynniki wyrazów $x^2$, $xy$
oraz stałego, dostajemy układ
\[
 a_s'= -\frac{2\hbar}{m}a_s^2+
       \frac{m\omega^2}{2\hbar},\qquad
 b_s'=-\frac{2\hbar}{m}a_sb_s,\qquad
 \frac{A_s'}{A_s}=-\frac{\hbar}{m}a_s,
\tag{16.11}\label{eq:Mehler-odes}
\]
oraz warunek zgodności
$b_s^2=4a_s^2-(m\omega/\hbar)^2$.
Rozwiązanie spełniające asymptotykę jądra Gaussa
przy $s\downarrow0$ jest równe
\[
 a_s=\frac{m\omega}{2\hbar}\coth(\omega s),\quad
 b_s=\frac{m\omega}{\hbar\sinh(\omega s)},\quad
 A_s=\sqrt{\frac{m\omega}{2\pi\hbar\sinh(\omega s)}}.
\]
Rzeczywiście $\coth' r=-1/\sinh^2r$ oraz
$\coth^2r-1=1/\sinh^2r$ weryfikują
\eqref{eq:Mehler-odes}. Gdy $s\downarrow0$,
$\sinh(\omega s)\sim\omega s$ i
$\cosh(\omega s)\sim1$, więc $K_s(x,y)$ jest
asymptotycznie jądrem
$\sqrt{m/(2\pi\hbar s)}
\exp[-m(x-y)^2/(2\hbar s)]$.
Z własności przybliżonej jedności wynika
$\int K_s(x,y)f(y)\,\dd y\to f(x)$
dla ciągłych funkcji o zwartym nośniku.

Sprawdźmy własność półgrupową, której nie wolno
wnioskować tylko z podstawienia do PDE. Przy oznaczeniach
$a_s,b_s,A_s$ wprowadzonych powyżej z całki Gaussa wynika
\[
\begin{split}
&\int_{\R}K_s(x,z)K_t(z,y)\,\dd z\\
&=A_sA_t\sqrt{\frac{\pi}{a_s+a_t}}
\exp\!\left[-\left(a_s-\frac{b_s^2}{4(a_s+a_t)}\right)x^2
-\left(a_t-\frac{b_t^2}{4(a_s+a_t)}\right)y^2
+\frac{b_sb_t}{2(a_s+a_t)}xy\right].
\end{split}
\]
Wstawiając wzory na $a_s,b_s,A_s$ i korzystając z
$\sinh(\omega(s+t))=\sinh(\omega s)\cosh(\omega t)
+\cosh(\omega s)\sinh(\omega t)$, sprawdzamy, że
prefaktor wynosi $A_{s+t}$, oba współczynniki przy
$x^2,y^2$ są równe $-a_{s+t}$, a współczynnik przy
$xy$ jest równy $b_{s+t}$. Otrzymujemy więc dokładnie
$K_{s+t}(x,y)$.
Dla $s>0$ całkowanie Gaussa daje też
$\int K_s(x,y)\,\dd y\le1$; symetria jądra i
oszacowanie Schura zapewniają
$\|K_sf\|_2\le\|f\|_2$.
Własność przybliżonej jedności daje silną ciągłość
w zerze na gęstym $C_c^\infty$, a oszacowanie normy
rozszerza ją na całe $L^2$. Samą identyfikację
generatora z domknięciem operatora oscylatora
uzasadnimy poniżej przez jednoznaczność równania
w $L^2$; nie zakładamy jej na podstawie samej
zgodności formalnych pochodnych.

Do reprezentacji probabilistycznej nie potrzebujemy
ogólnego twierdzenia o monotonicznej zbieżności form.
Dla $f\in C_c^\infty(\R)$, $f\ge0$, połóżmy
$u(s,x)=\int_{\R}K_s^{\mathrm E}(x,y)f(y)\,\dd y$.
Obliczając całkę Gaussa, dostajemy
\[
 \int_{\R} K_s^{\mathrm E}(x,y)\,\dd y
 =\frac{\exp[-(m\omega/(2\hbar))x^2\tanh(\omega s)]}
        {\sqrt{\cosh(\omega s)}}\le1.
\]
Zatem $0\le u(s,x)\le\|f\|_\infty$.
Funkcja $u$ spełnia klasyczne równanie
$\partial_su=(\hbar/(2m))u_{xx}-(m\omega^2/(2\hbar))x^2u$
przy $s>0$ oraz $u(s,x)\to f(x)$ lokalnie jednostajnie,
gdy $s\downarrow0$.

Niech $X_r=x+\sqrt{\hbar/m}B_r$. Dla
$0<\varepsilon<s$ i $R>|x|$ stosujemy wzór Itô do
\[
 Z_r=\exp\left(-\frac{m\omega^2}{2\hbar}
 \int_0^rX_v^2\,\dd v\right)u(s-r,X_r),
 \qquad 0\le r\le(s-\varepsilon)\wedge\tau_R,
\]
gdzie $\tau_R=\inf\{r\ge0:|X_r|\ge R\}$.
Dryf znika dzięki równaniu PDE, a po zatrzymaniu
współczynnik stochastyczny jest ograniczony na
odpowiednim zwartym zbiorze czasu i przestrzeni,
więc $Z$ jest martyngałem. Mamy $0\le Z\le\|f\|_\infty$.
Zatem $\E Z_{(s-\varepsilon)\wedge\tau_R}=u(s,x)$.
Dopuszczamy najpierw $R\to\infty$: ciągłość Brownowskich
trajektorii daje $\tau_R\to\infty$ prawie na pewno,
a zbieżność zdominowana uzasadnia przejście do granicy.
Następnie $\varepsilon\downarrow0$; ciągłość $X$ i
lokalnie jednostajna zbieżność $u(\varepsilon,\cdot)\to f$
dają właśnie \eqref{eq:Mehler-FK}.
Dla $f$ dowolnej ograniczonej ciągłej stosujemy obcięcia
i zbieżność zdominowaną, a dla funkcji zmieniających znak
rozkład na części dodatnią i ujemną.

Aby zidentyfikować otrzymaną półgrupę z
$e^{-sH_{\mathrm{osc}}/\hbar}$ w $L^2$, wykorzystujemy
jednoznaczność rozwiązań równania cieplnego z
potencjałem nieujemnym w klasie $L^2$: dla różnicy
dwóch rozwiązań z zerową daną początkową iloczyn skalarny
z rozwiązaniem i całkowanie przez części dają
\[
 \frac{\dd}{\dd s}\|v(s)\|_2^2
 =-\frac{\hbar}{m}\|\partial_xv(s)\|_2^2
  -\frac{m\omega^2}{\hbar}\|xv(s)\|_2^2\le0.
\]
Wyjaśnijmy dokładniej przejście poza funkcje o zwartym
nośniku. Wybierzmy $\chi\in C_c^\infty(\R)$, $0\le\chi\le1$,
równą $1$ na $[-1,1]$, i połóżmy $\chi_R(x)=\chi(x/R)$.
Mnożąc równanie dla różnicy $v$ przez $\chi_R^2\bar v$,
całkując i biorąc części rzeczywiste, otrzymujemy po
nierówności Cauchy'ego--Schwarza
\[
\frac{\dd}{\dd s}\int\chi_R^2|v|^2\,\dd x
\le \frac{C}{R^2}\int_{\R}|v|^2\,\dd x.
\]
Istotnie, wyraz mieszany zawiera
$2\chi_R\chi_R'\bar v\,\partial_xv$ i można go
wchłonąć w ujemny wyraz
$-\frac{\hbar}{m}\int\chi_R^2|\partial_xv|^2$,
pozostawiając $C\int|\chi_R'|^2|v|^2$;
składnik potencjału jest niedodatni. Na przedziale
$[\delta,s]$ rachunek jest klasyczny dla rozwiązań
gładkich; w przestrzeni formowej otrzymuje się go przez
aproksymację w normie formy. Dopuszczamy następnie
$R\to\infty$ i $\delta\downarrow0$, korzystając z
ciągłości w $L^2$ i ograniczenia normy na zwartych
przedziałach czasu. Dostajemy
$\|v(s)\|_2^2\le\|v(0)\|_2^2=0$.
Półgrupa spektralna i operator z jądrem $K_s^{\mathrm E}$
dają rozwiązania z tymi samymi danymi, więc są identyczne.
Nie wykorzystaliśmy tutaj nieudowodnionej jednoznaczności
klasycznej dla dowolnych nieograniczonych potencjałów.
\end{proof}

\begin{wniosek}[Propagator oscylatora w czasie rzeczywistym]
\label{cor:Mehler-real}
Niech $0<t<\pi/\omega$. Wtedy jądro propagatora
$e^{-itH_{\mathrm{osc}}/\hbar}$ ma oscylacyjną postać
\[
\begin{split}
 K_t^{\mathrm R}(x,y)
  ={}&\sqrt{\frac{m\omega}{2\pi i\hbar\sin(\omega t)}}\\
 &\cdot\exp\left[\frac{im\omega}{2\hbar\sin(\omega t)}
 \bigl((x^2+y^2)\cos(\omega t)-2xy\bigr)\right].
\end{split}
\tag{16.12}\label{eq:Mehler-R}
\]
Dla dłuższych czasów konieczny jest wybór fazy
zgodnej z ciągłością grupy unitarnej; w chwilach
$\sin(\omega t)=0$ propagator nie ma zwykłego jądra
funkcyjnego tej postaci.
\end{wniosek}
\begin{proof}
Stosujemy Twierdzenie~\ref{th:wick-spectral} do
$H_{\mathrm{osc}}$ i dla $z=\varepsilon+it$ korzystamy
z analitycznej kontynuacji wzoru \eqref{eq:Mehler-E}.
Ponieważ $\sinh(i\omega t)=i\sin(\omega t)$ i
$\cosh(i\omega t)=\cos(\omega t)$, otrzymujemy
\eqref{eq:Mehler-R} jako granicę oscylacyjnych jąder
przy $\varepsilon\downarrow0$. W przedziale
$0<t<\pi/\omega$ gałąź pierwiastka jest wybrana
jednoznacznie przez zbieżność od $t=0+$.
Granica operatorów jest silna w $L^2$, co uzasadnia
propagator, bez wymagania istnienia miary zespolonej
na przestrzeni trajektorii.
\end{proof}

\begin{zadanie}[Przypadek zerowej częstości]
Sprawdzić, że gdy $\omega\downarrow0$, jądro
\eqref{eq:Mehler-E} zbiega do jądra cząstki swobodnej
w czasie euklidesowym.

\emph{Rozwiązanie.} Korzystamy z
$\sinh(\omega s)\sim\omega s$,
$\cosh(\omega s)=1+O(\omega^2s^2)$.
Prefaktor dąży do $\sqrt{m/(2\pi\hbar s)}$,
a wykładnik do $-m(x-y)^2/(2\hbar s)$.
Zatem $K_s^{\mathrm E}(x,y)\to q_s(x,y)$.
W tym rachunku zanik potencjału prowadzi dokładnie do
półgrupy cieplnej cząstki swobodnej.
\end{zadanie}


\section{Ślad półgrupy oscylatora i funkcja podziału}

\begin{zadanie}[Gaussowski rachunek śladu]
Dla $H_{\mathrm{osc}}$ z
Twierdzenia~\ref{th:Mehler-euclidean} obliczyć
$\int_{\R}K_s^{\mathrm E}(x,x)\,\dd x$ i porównać
z sumą po poziomach energetycznych.

\emph{Rozwiązanie.} Podstawiając $y=x$ do
\eqref{eq:Mehler-E} i korzystając z
$(\cosh r-1)/\sinh r=\tanh(r/2)$, mamy
\[
 K_s^{\mathrm E}(x,x)=
 \sqrt{\frac{m\omega}{2\pi\hbar\sinh(\omega s)}}
 \exp\left(-\frac{m\omega}{\hbar}
     \tanh(\omega s/2)x^2\right).
\]
Całka Gaussa daje
\[
 \int_{\R}K_s^{\mathrm E}(x,x)\,\dd x
 =\frac1{2\sinh(\omega s/2)}.
\]
Aby nie zakładać śladowości bez uzasadnienia,
zauważmy, że dodatnie jądro jest symetryczne i spełnia
własność półgrupową. Wobec tego
\[
\int_{\R^2}|K_{s/2}^{\mathrm E}(x,y)|^2\,
\dd\lambda_2(x,y)=\int_{\R}K_s^{\mathrm E}(x,x)\,
\dd\lambda_1(x)<\infty.
\]
Operator o jądrze $K_{s/2}^{\mathrm E}$ jest zatem
Hilberta--Schmidta, a jego kwadrat, równy
$e^{-sH_{\mathrm{osc}}/\hbar}$, jest dodatnim
operatorem śladowym. To uzasadnia utożsamienie
całki przekątnej jądra z jego śladem.
Znane widmo oscylatora,
$E_n=\hbar\omega(n+\frac12)$, $n\ge0$,
daje niezależne sprawdzenie:
\[
 \operatorname{Tr}(e^{-sH_{\mathrm{osc}}/\hbar})
 =\sum_{n=0}^\infty e^{-s\omega(n+1/2)}
 =\frac1{2\sinh(\omega s/2)}.
\]
W ten sposób jądro otrzymane z całki po trajektoriach
koduje wszystkie poziomy energetyczne. Śladowość
uzasadniliśmy powyżej; wzór na widmo oscylatora
pozostaje przywołanym wynikiem analizy operatorowej,
nie jest wyprowadzany z samego rachunku całki Gaussa.
W języku mechaniki statystycznej parametr $s$ odpowiada
odwrotnej temperaturze przez $s=\beta\hbar$;
funkcja podziału wynosi wtedy
$Z(\beta)=\operatorname{Tr}(e^{-\beta H_{\mathrm{osc}}})$.
\end{zadanie}

\section{Co jest dowiedzione, a co pozostaje formalną analogią?}

Nasze rozróżnienie nie jest jedynie słowne. W czasie
euklidesowym wzór \eqref{eq:physical-FK} jest zwykłą
wartością oczekiwaną po przestrzeni trajektorii, o ile
potencjał spełnia przyjęte założenia. W czasie rzeczywistym
propagator istnieje jako grupa unitarna z twierdzenia
spektralnego. Dla cząstki swobodnej i oscylatora
harmonicznego znamy nawet jawne jądra oscylacyjne.
Natomiast formalny zapis
\[
 \int_{x(0)=x}^{x(T)=y}
  \exp\left(\frac{i}{\hbar}S[x]\right)\,\mathcal D x
\]
\emph{nie został} zinterpretowany jako całka Lebesgue'a
względem przeliczalnie addytywnej miary zespolonej na
przestrzeni wszystkich trajektorii.
Wynik \eqref{eq:trotter-quantum} daje ścisły sens
aproksymacjom \emph{operatorowym} przez skończone iloczyny,
a Twierdzenie~\ref{th:wick-spectral} daje ścisły sens
przejściu między dodatnią półgrupą a unitarną ewolucją.
Nie należy utożsamiać tych faktów z istnieniem
ogólnej miary Feynmana.

\begin{center}
\begin{tikzpicture}[>=Latex, node distance=8mm and 12mm]
\node[draw,rounded corners,align=center,text width=5.5cm,
 minimum height=1.05cm] (a)
 {\textbf{Czas euklidesowy $s>0$}\\
  $e^{-sH/\hbar}$, miara Wienera, dodatnie wagi};
\node[draw,rounded corners,align=center,text width=5.5cm,
 minimum height=1.05cm,right=of a] (b)
 {\textbf{Czas rzeczywisty $t\in\R$}\\
  $e^{-itH/\hbar}$, ewolucja unitarna};
\draw[->,thick,blue!70!black] (a)--(b);
\node[draw,rounded corners,align=center,text width=5.5cm,
 minimum height=0.9cm,below=14mm of a] (c)
 {Całka po trajektoriach\\
  jako wartość oczekiwana};
\node[draw,rounded corners,align=center,text width=5.5cm,
 minimum height=0.9cm,below=14mm of b] (d)
 {Oscylacyjne całki skończeniewymiarowe\\
  i granice iloczynów operatorów};
\draw[->,thick] (a)--(c);
\draw[->,thick] (b)--(d);
\draw[->,densely dashed] (c)--(d);
\end{tikzpicture}
\par\smallskip\small
Strzałka górna oznacza silną granicę spektralną; strzałka
przerywana przypomina, że przejście do całek oscylacyjnych
nie jest zwykłą zmianą miary na przestrzeni trajektorii.
\end{center}

\paragraph{Zakończenie tomu III.}
Droga od izometrii Wienera, przez całkę Itô, martyngały,
SDE, procesy Lévy'ego i półmartyngały, prowadzi do
probabilistycznych reprezentacji półgrup. Całki po
trajektoriach są ścisłym narzędziem analizy; ich związek
z całkami Feynmana i mechaniką kwantową jest potężny,
ale wymaga kontrolowania samosprzężoności operatorów,
regularności potencjałów i rodzaju używanej zbieżności.

\end{document}
